\chapter{Relationen}

\FormulaDefDelta[Relation (Menge geordneter Paare)]{F \subseteq A \times B}{
\DeltaRow{Mengen}{A\dsep B\dsep F}
}

\section{Begriff der totalen Relation}

\subsection{Axiom der Totalität}

\FormulaAxiomDelta[Totalität]
{x\in A \vdash \exists y\,(x,y)\in F}
{
\DeltaRow{Mengen}{x\dsep y\dsep A\dsep F}
}
\subsection{Definition der totalen Relation}

\FormulaDefDeltaK[Begriff der totalen Relation]{\TotRel{R,A,B}}{Totale Relation}{
  \DeltaRow{Mengen}{A \dsep B \dsep R \dsep x \dsep y}
  \DeltaRow{\textbf{Axiome}}{}
  \DeltaRow{Menge geordneter Paare}
           {R \subseteq A \times B}
           [\FormulaRefAuto{F \subseteq A \times B}]
  \DeltaRow{Totalität}
           {x \in A \vdash \exists y\,(x,y)\in R}
           [\FormulaRefAuto{x\in A \vdash \exists y\,(x,y)\in F}]
  \DeltaRow{\textbf{Neue Symbole}}{}
  \DeltaPrem{Totale Relationen}{}
}

\subsection{Beispiele}

\subsubsection{Mitgliedschaftsrelation}

% ------------------------------------------------------------
% 0) Definition: Mitgliedschaftsrelation der Familie
% ------------------------------------------------------------
\FormulaDefDelta[Mitgliedschaftsrelation]{%
  \MemRel(M) \coloneqq \{\, (X,a)\in M\times \bigcup M \mid a\in X \,\}
}{
  \DeltaRow{Mengen}{M \dsep X \dsep a}
}

% ------------------------------------------------------------
% 1) Lemma: \MemRel(M) ist Teilmenge des Produkts
% ------------------------------------------------------------
\FormulaThmDelta[Relation]
{\MemRel(M) \subseteq M\times \bigcup M}
{
  \DeltaRow{Mengen}{M \dsep X \dsep a}
}
\begin{tabproof}
  \proofstep{}{\MemRel(M)=\{\, (X,a)\in M\times \bigcup M \mid a\in X \,\}}{\rA}
  \proofstep{}{\MemRel(M) \subseteq M\times \bigcup M}{%
    \FormulaRefAuto{\{\,x\in A \mid P(x)\,\}\subseteq A}{}}
\end{tabproof}

% ------------------------------------------------------------
% 2) Lemma: Aus (X,a)\in \MemRel(M) folgt X\in M, a\in \bigcup M und a\in X
% ------------------------------------------------------------
\FormulaThmDelta
{(X,a)\in \MemRel(M) \vdash X\in M \land a\in \bigcup M \land a\in X}
{
  \DeltaRow{Mengen}{M \dsep X \dsep a}
}
\begin{tabproof}
  \proofstep{1}{(X,a)\in \MemRel(M)}{\rA}
  \proofstep{1}{(X,a)\in M\times \bigcup M \land a\in X}{%
    \FormulaRefAuto{x \in \{u \in A \mid P(u)\} \eqvdash x \in A \land P(x)}{1}}
  \proofstep{1}{(X,a)\in M\times \bigcup M}{\rAEa{2}}
  \proofstep{1}{a\in X}{\rAEb{2}}
  \proofstep{1}{X\in M \land a\in \bigcup M}{%
    \FormulaRefAuto{(a,b)\in A\times B \eqvdash a\in A \land b\in B}{3}}
  \proofstep{1}{X\in M}{\rAEa{5}}
  \proofstep{1}{a\in \bigcup M}{\rAEb{5}}
  \proofstep{1}{X\in M \land a\in \bigcup M \land a\in X}{\FormulaRefAuto{P\dsep Q\dsep R\vdash P\land Q\land R}{6,7,4}}
\end{tabproof}

\FormulaThmDelta
{(X,a)\in \MemRel(M) \vdash X\in M}
{
  \DeltaRow{Mengen}{M \dsep X \dsep a}
}
\begin{tabproof}
  \proofstep{1}{(X,a)\in \MemRel(M)}{\rA}
  \proofstep{1}{X\in M \land a\in \bigcup M \land a\in X}{\FormulaRefAuto{(X,a)\in \MemRel(M) \vdash X\in M \land a\in \bigcup M \land a\in X}{1}}
  \proofstep{1}{X\in M}{\rAEa{2}}
\end{tabproof}

% ------------------------------------------------------------
% 3) Lemma: Aus X\in M und a\in X folgt (X,a)\in \MemRel(M)
% ------------------------------------------------------------
\FormulaThmDelta
{X\in M \dsep a\in X \vdash (X,a)\in \MemRel(M)}
{
  \DeltaRow{Mengen}{M \dsep X \dsep a}
}
\begin{tabproof}
  \proofstep{1}{X\in M}{\rA}
  \proofstep{2}{a\in X}{\rA}
  \proofstep{1,2}{a\in \bigcup M}{%
    \FormulaRefAuto{B\in A\dsep x\in B\vdash x \in \bigcup A}{1,2}}
  \proofstep{1,2}{(X,a)\in M\times \bigcup M}{%
    \FormulaRefAuto{(a,b)\in A\times B \eqvdash a\in A \land b\in B}{1,3}}
  \proofstep{}{\{\, (X,a)\in M\times \bigcup M \mid a\in X \,\}=\MemRel(M)}{\FormulaRefAuto{a=b\vdash b=a}{\FormulaRefAuto{\MemRel(M) \coloneqq \{\, (X,a)\in M\times \bigcup M \mid a\in X \,\}}}}
  \proofstep{1,2}{(X,a)\in \{\, (X,a)\in M\times \bigcup M \mid a\in X \,\}}{%
    \FormulaRefAuto{x\in A \dsep P(x) \vdash x\in \{\,x\in A \mid P(x)\,\}}{4,2}}
  \proofstep{1,2}{(X,a)\in \MemRel(M)}{%
    \rIE{5,6}}
\end{tabproof}

% ============================================================
% (1) Totalität von \MemRel(M) (als eigenes Lemma)
% ============================================================

\FormulaThmDelta[Totalität von $\MemRel(M)$]{%
  X\in M \vdash \exists a\,\bigl((X,a)\in \MemRel(M)\bigr)
}{
  \DeltaRow{Mengen}{M \dsep X \dsep a}
  \DeltaPrem{\makecell[l]{Paarweise disjunkte\\ nichtleere Familie}}{\DisjFam(M)}
}
\begin{tabproof}
  \proofstep{1}{X\in M}{\rA}

  \proofstep{1}{X\neq\varnothing}{%
    \FormulaRefAuto{X\in A \vdash X\neq\varnothing}{1}}

  \proofstep{1}{\exists a\,(a\in X)}{%
    \FormulaRefAuto{A\neq\varnothing \eqvdash \exists x\,(x\in A)}{2}}

  \proofstep{4}{a\in X}{\rA}

  \proofstep{1,4}{(X,a)\in \MemRel(M)}{%
    \FormulaRefAuto{X\in M \dsep a\in X \vdash (X,a)\in \MemRel(M)}{1,4}}

  \proofstep{1,4}{\exists a\,\bigl((X,a)\in \MemRel(M)\bigr)}{\rEI{5}}

  \proofstep{1}{\exists a\,\bigl((X,a)\in \MemRel(M)\bigr)}{\rEE{3,4,6}}

\end{tabproof}


% ============================================================
% (2) TotRel(\MemRel(M),M,⋃M) als "Definition-Style"-Theorem (ohne tabproof)
% ============================================================

\FormulaThmDelta[Totale Relation]{\TotRel{\MemRel(M),M,\bigcup M}}{
  \DeltaRow{Mengen}{M \dsep X \dsep a}
  \DeltaPrem{\makecell[l]{Paarweise disjunkte\\ nichtleere Familie}}{\DisjFam(M)}
  \DeltaRow{\textbf{Begründung}}{}
  \DeltaRow{Menge geordneter Paare}{\MemRel(M) \subseteq M\times \bigcup M}
    [\FormulaRefAuto{\MemRel(M) \subseteq M\times \bigcup M}]
  \DeltaRow{Totalität}{X\in M \vdash \exists a\,\bigl((X,a)\in \MemRel(M)\bigr)}
    [\FormulaRefAuto{X\in M \vdash \exists a\,\bigl((X,a)\in \MemRel(M)\bigr)}]
}{}

\subsubsection{Strikte Obermengenrelation in Potenzmengenfamilien}

\FormulaDefDelta[Strikte Obermengenrelation auf einer Potenzmengenfamilie]{%
  B\subseteq \powerset(M)\vdash
  \StrictSupRel{B}\coloneqq
  \{(x,y)\in B\times B\mid x\subset y\}%
}{%
  \DeltaRow{Mengen}{M\dsep B\dsep x\dsep y}
  \DeltaRow{Relationen}{\StrictSupRel{B}}
}

\FormulaThmDelta[Einfuehrung in die strikte Obermengenrelation]{%
  B\subseteq \powerset(M)\dsep x\in B\dsep y\in B\dsep x\subset y
  \vdash (x,y)\in \StrictSupRel{B}%
}{%
  \DeltaRow{Mengen}{M\dsep B\dsep x\dsep y}
  \DeltaRow{Relationen}{\StrictSupRel{B}}
}
\begin{tabproof}
  \proofstep{1}{B\subseteq \powerset(M)}{\rA}
  \proofstep{2}{x\in B}{\rA}
  \proofstep{3}{y\in B}{\rA}
  \proofstep{4}{x\subset y}{\rA}
  \proofstep{2,3}{(x,y)\in B\times B}{%
    \FormulaRefAuto{a\in A\dsep b\in B\vdash (a,b)\in A\times B}{2,3}}
  \proofstep{2,3,4}{(x,y)\in \{(u,v)\in B\times B\mid u\subset v\}}{%
    \FormulaRefAuto{x \in A\dsep P(x)\vdash x \in \{x \in A \mid P(x)\}}{5,4}}
  \proofstep{1,2,3,4}{(x,y)\in \StrictSupRel{B}}{%
    \rIE{\FormulaRefAuto{B\subseteq \powerset(M)\vdash \StrictSupRel{B}\coloneqq \{(x,y)\in B\times B\mid x\subset y\}}{1},6}}
\end{tabproof}

\FormulaThmDelta[Elimination der strikten Obermengenrelation]{%
  B\subseteq \powerset(M)\dsep (x,y)\in \StrictSupRel{B}
  \vdash x\in B\land y\in B\land x\subset y%
}{%
  \DeltaRow{Mengen}{M\dsep B\dsep x\dsep y}
  \DeltaRow{Relationen}{\StrictSupRel{B}}
}
\begin{tabproof}
  \proofstep{1}{B\subseteq \powerset(M)}{\rA}
  \proofstep{2}{(x,y)\in \StrictSupRel{B}}{\rA}
  \proofstep{1}{\StrictSupRel{B}=\{(u,v)\in B\times B\mid u\subset v\}}{%
    \FormulaRefAuto{B\subseteq \powerset(M)\vdash \StrictSupRel{B}\coloneqq \{(x,y)\in B\times B\mid x\subset y\}}{1}}
  \proofstep{1,2}{(x,y)\in \{(u,v)\in B\times B\mid u\subset v\}}{\rIE{3,2}}
  \proofstep{1,2}{(x,y)\in B\times B\land x\subset y}{%
    \FormulaRefAuto{x\in\{u\in A\mid P(u)\}\eqvdash x\in A\land P(x)}{4}}
  \proofstep{1,2}{(x,y)\in B\times B}{\rAEa{5}}
  \proofstep{1,2}{x\subset y}{\rAEb{5}}
  \proofstep{1,2}{x\in B}{\FormulaRefAuto{(a,b)\in A\times B\vdash a\in A}{6}}
  \proofstep{1,2}{y\in B}{\FormulaRefAuto{(a,b)\in A\times B\vdash b\in B}{6}}
  \proofstep{1,2}{x\in B\land y\in B\land x\subset y}{\rAI{\rAI{8,9},7}}
\end{tabproof}

\FormulaThmDelta[Typisierung der strikten Obermengenrelation]{%
  B\subseteq \powerset(M)\vdash \StrictSupRel{B}\subseteq B\times B%
}{%
  \DeltaRow{Mengen}{M\dsep B\dsep z\dsep x\dsep y}
  \DeltaRow{Relationen}{\StrictSupRel{B}}
}
\begin{tabproof}
  \proofstep{1}{B\subseteq \powerset(M)}{\rA}
  \proofstep{2}{z\in \StrictSupRel{B}}{\rA}
  \proofstep{1}{\StrictSupRel{B}=\{(x,y)\in B\times B\mid x\subset y\}}{%
    \FormulaRefAuto{B\subseteq \powerset(M)\vdash \StrictSupRel{B}\coloneqq \{(x,y)\in B\times B\mid x\subset y\}}{1}}
  \proofstep{1,2}{z\in \{(x,y)\in B\times B\mid x\subset y\}}{\rIE{3,2}}
  \proofstep{1,2}{z\in B\times B}{%
    \FormulaRefAuto{x \in \{x \in A \mid P(x)\}\vdash x\in A}{4}}
  \proofstep{1}{\StrictSupRel{B}\subseteq B\times B}{%
    \FormulaRefAuto{A \subseteq B \coloneq \forall x\,(x\in A \rightarrow x\in B)}{\rUI{\rRI{2,5}}}}
\end{tabproof}

\FormulaThmDelta[Totalitaet der strikten Obermengenrelation]{%
  B\subseteq \powerset(M)\dsep
  \forall x\in B\,\exists y\in B\,\bigl(x\subset y\bigr)
  \vdash \TotRel{\StrictSupRel{B},B,B}%
}{%
  \DeltaRow{Mengen}{M\dsep B\dsep x\dsep y}
  \DeltaRow{Relationen}{\StrictSupRel{B}}
}
\begin{tabproof}
  \proofstep{1}{B\subseteq \powerset(M)}{\rA}
  \proofstep{2}{\forall x\in B\,\exists y\in B\,\bigl(x\subset y\bigr)}{\rA}
  \proofstep{1}{\StrictSupRel{B}\subseteq B\times B}{%
    \FormulaRefAuto{B\subseteq \powerset(M)\vdash \StrictSupRel{B}\subseteq B\times B}{1}}

  \proofstep{4}{x\in B}{\rA}
  \proofstep{2,4}{\exists y\in B\,\bigl(x\subset y\bigr)}{\rRE{4,\rUE{2}}}
  \proofstep{6}{y\in B\land x\subset y}{\rA}
  \proofstep{6}{y\in B}{\rAEa{6}}
  \proofstep{6}{x\subset y}{\rAEb{6}}
  \proofstep{1,4,6}{(x,y)\in \StrictSupRel{B}}{%
    \FormulaRefAuto{B\subseteq \powerset(M)\dsep x\in B\dsep y\in B\dsep x\subset y \vdash (x,y)\in \StrictSupRel{B}}{1,4,7,8}}
  \proofstep{1,4,6}{\exists y\,(x,y)\in \StrictSupRel{B}}{\rEI{9}}
  \proofstep{1,2,4}{\exists y\,(x,y)\in \StrictSupRel{B}}{\rEE{5,6,10}}
  \proofstep{1,2}{x\in B\rightarrow \exists y\,(x,y)\in \StrictSupRel{B}}{\rRI{4,11}}
  \proofstep{1,2}{\TotRel{\StrictSupRel{B},B,B}}{%
    \FormulaRefAuto{Totale Relation}{3,12}}
\end{tabproof}

\section{Begriff der Äquivalenzrelation}

\subsection{Axiome einer Äquivalenzrelation}

\FormulaAxiomDelta[Reflexivität]{%
  x\in A \vdash (x,x)\in R%
}{
  \DeltaRow{Mengen}{A \dsep R \dsep x}
}

\FormulaAxiomDelta[Symmetrie]{%
  (x,y)\in R \vdash (y,x)\in R%
}{
  \DeltaRow{Mengen}{A \dsep R \dsep x \dsep y}
}

\FormulaAxiomDelta[Transitivität]{%
  (x,y)\in R \dsep (y,z)\in R \vdash (x,z)\in R%
}{
  \DeltaRow{Mengen}{A \dsep R \dsep x \dsep y \dsep z}
}

\subsection{Definition der Äquivalenzrelation}

\FormulaDefDeltaK[Begriff der Äquivalenzrelation]{\EqRel{R,A}}{Äquivalenzrelation}{
  \DeltaRow{Mengen}{A \dsep R \dsep x \dsep y \dsep z}
  \DeltaRow{\textbf{Axiome}}{}
  \DeltaRow{Relation auf \(A\)}
           {R \subseteq A \times A}
           [\FormulaRefAuto{F \subseteq A \times B}]
  \DeltaRow{Reflexivität}
           {x\in A \vdash (x,x)\in R}
           [\FormulaRefAuto{x\in A \vdash (x,x)\in R}]
  \DeltaRow{Symmetrie}
           {(x,y)\in R \vdash (y,x)\in R}
           [\FormulaRefAuto{(x,y)\in R \vdash (y,x)\in R}]
  \DeltaRow{Transitivität}
           {(x,y)\in R \dsep (y,z)\in R \vdash (x,z)\in R}
           [\FormulaRefAuto{(x,y)\in R \dsep (y,z)\in R \vdash (x,z)\in R}]
  \DeltaRow{\textbf{Neue Symbole}}{}
  \DeltaPrem{Äquivalenzrelationen}{}
}

\section{Begriff der Ordnung}

\subsection{Axiome einer partiellen Ordnung}

\FormulaAxiomDelta[Antisymmetrie]{%
  (x,y)\in R \dsep (y,x)\in R \vdash x=y%
}{
  \DeltaRow{Mengen}{A \dsep R \dsep x \dsep y}
}

\subsection{Definition der partiellen Ordnung}

\FormulaDefDeltaK[Begriff der partiellen Ordnung]{\PartOrd{R,A}}{Partielle Ordnung}{
  \DeltaRow{Mengen}{A \dsep R \dsep x \dsep y \dsep z}
  \DeltaRow{\textbf{Axiome}}{}
  \DeltaRow{Relation auf \(A\)}
           {R \subseteq A \times A}
           [\FormulaRefAuto{F \subseteq A \times B}]
  \DeltaRow{Reflexivität}
           {x\in A \vdash (x,x)\in R}
           [\FormulaRefAuto{x\in A \vdash (x,x)\in R}]
  \DeltaRow{Transitivität}
           {(x,y)\in R \dsep (y,z)\in R \vdash (x,z)\in R}
           [\FormulaRefAuto{(x,y)\in R \dsep (y,z)\in R \vdash (x,z)\in R}]
  \DeltaRow{Antisymmetrie}
           {(x,y)\in R \dsep (y,x)\in R \vdash x=y}
           [\FormulaRefAuto{(x,y)\in R \dsep (y,x)\in R \vdash x=y}]
  \DeltaRow{\textbf{Neue Symbole}}{}
  \DeltaPrem{Partielle Ordnungen}{}
}

\subsection{Axiom einer totalen Ordnung}

\FormulaAxiomDelta[Vergleichbarkeit]{%
  x\in A \dsep y\in A \vdash (x,y)\in R \lor (y,x)\in R%
}{
  \DeltaRow{Mengen}{A \dsep R \dsep x \dsep y}
}

\subsection{Definition der totalen Ordnung}

\FormulaDefDeltaK[Begriff der totalen Ordnung]{\TotOrd{R,A}}{Totale Ordnung}{
  \DeltaRow{Mengen}{A \dsep R \dsep x \dsep y \dsep z}
  \DeltaRow{\textbf{Axiome}}{}
  \DeltaRow{Partielle Ordnung}
           {\PartOrd{R,A}}
           [\FormulaRefAuto{Partielle Ordnung}]
  \DeltaRow{Vergleichbarkeit}
           {x\dsep y\in A \vdash (x,y)\in R \lor (y,x)\in R}
           [\FormulaRefAuto{x\in A \dsep y\in A \vdash (x,y)\in R \lor (y,x)\in R}]
  \DeltaRow{\textbf{Neue Symbole}}{}
  \DeltaPrem{Totale Ordnungen}{}
}

% ============================================================
% Kleinste / groesste Elemente, Minimum / Maximum
% (auf Basis einer partiellen Ordnung)
% ============================================================

% ============================================================
% Minimum / Maximum (als eindeutige Terme via \iota)
% Kleinstes/Groesstes Element als Spezialfall T=A
% ============================================================

\subsection{Minimum und Maximum}

% ------------------------------------------------------------
% (1) Eindeutigkeit des Minimum-Kriteriums
% ------------------------------------------------------------
\FormulaThmDelta[Eindeutigkeit eines Minimums]{%
  \exists m\in T\,\forall x\in T\,\bigl((m,x)\in R\bigr)
  \vdash \exists! m\in T\,\forall x\in T\,\bigl((m,x)\in R\bigr)%
}{
  \DeltaRow{Mengen}{A \dsep R \dsep T \dsep m \dsep n \dsep x}
  \DeltaPrem{Partielle Ordnung}{\PartOrd{R,A}}
}
\begin{tabproof}
  \proofstep{1}{\exists m\in T\,\forall x\in T\,\bigl((m,x)\in R\bigr)}{\rA}

  \proofstep{2}{m\in T \land \forall x\in T\,\bigl((m,x)\in R\bigr)}{\rA}
  \proofstep{3}{n\in T \land \forall x\in T\,\bigl((n,x)\in R\bigr)}{\rA}

  \proofstep{2}{m\in T}{\rAEa{2}}
  \proofstep{2}{\forall x\in T\,\bigl((m,x)\in R\bigr)}{\rAEb{2}}
  \proofstep{3}{n\in T}{\rAEa{3}}
  \proofstep{3}{\forall x\in T\,\bigl((n,x)\in R\bigr)}{\rAEb{3}}

  \proofstep{2,3}{(m,n)\in R}{\rRE{6,\rUE{5}}}
  \proofstep{2,3}{(n,m)\in R}{\rRE{4,\rUE{7}}}

  \proofstep{2,3}{m=n}{%
    \FormulaRefAuto{(x,y)\in R \dsep (y,x)\in R \vdash x=y}{8,9}}

  \proofstep{1}{\exists! m\in T\,\forall x\in T\,\bigl((m,x)\in R\bigr)}{%
    \UEI{1,2,3,10}}
\end{tabproof}


% ------------------------------------------------------------
% (2) Definition: Minimum als \iota-Term (partielle Definition)
% ------------------------------------------------------------
\FormulaDefDelta[Minimum]{%
  \begin{aligned}
    &\exists m\in T\,\forall x\in T\,\bigl((m,x)\in R\bigr)\\
    &\vdash \Min{R,A}(T)\coloneqq
      \iota m\in T\,\forall x\in T\,\bigl((m,x)\in R\bigr)
  \end{aligned}%
}{
  \DeltaRow{Mengen}{A \dsep R \dsep T \dsep m \dsep x}
  \DeltaPrem{Partielle Ordnung}{\PartOrd{R,A}}
}

% ------------------------------------------------------------
% (3) Eindeutigkeit des Maximum-Kriteriums (kurz mit  \exists m\in T ...)
% ------------------------------------------------------------
\FormulaThmDelta[Eindeutigkeit eines Maximums]{%
  \exists m\in T\,\forall x\in T\,\bigl((x,m)\in R\bigr)
  \vdash \exists! m\in T\,\forall x\in T\,\bigl((x,m)\in R\bigr)%
}{
  \DeltaRow{Mengen}{A \dsep R \dsep T \dsep m \dsep n \dsep x}
  \DeltaPrem{Partielle Ordnung}{\PartOrd{R,A}}
}
\begin{tabproof}
  \proofstep{1}{\exists m\in T\,\forall x\in T\,\bigl((x,m)\in R\bigr)}{\rA}

  \proofstep{2}{m\in T \land \forall x\in T\,\bigl((x,m)\in R\bigr)}{\rA}
  \proofstep{3}{n\in T \land \forall x\in T\,\bigl((x,n)\in R\bigr)}{\rA}

  \proofstep{2}{m\in T}{\rAEa{2}}
  \proofstep{2}{\forall x\in T\,\bigl((x,m)\in R\bigr)}{\rAEb{2}}
  \proofstep{3}{n\in T}{\rAEa{3}}
  \proofstep{3}{\forall x\in T\,\bigl((x,n)\in R\bigr)}{\rAEb{3}}

  \proofstep{2,3}{(m,n)\in R}{\rRE{4,\rUE{7}}}
  \proofstep{2,3}{(n,m)\in R}{\rRE{6,\rUE{5}}}

  \proofstep{2,3}{m=n}{%
    \FormulaRefAuto{(x,y)\in R \dsep (y,x)\in R \vdash x=y}{8,9}}

  \proofstep{1}{\exists! m\in T\,\forall x\in T\,\bigl((x,m)\in R\bigr)}{%
    \UEI{1,2,3,10}}
\end{tabproof}

% ------------------------------------------------------------
% (4) Definition: Maximum als \iota-Term (partielle Definition)
% ------------------------------------------------------------
\FormulaDefDelta[Maximum]{%
  \begin{aligned}
    &\exists m\in T\,\forall x\in T\,\bigl((x,m)\in R\bigr)\\
    &\vdash \Max{R,A}(T)\coloneqq
      \iota m\in T\,\forall x\in T\,\bigl((x,m)\in R\bigr)
  \end{aligned}%
}{
  \DeltaRow{Mengen}{A \dsep R \dsep T \dsep m \dsep x}
  \DeltaPrem{Partielle Ordnung}{\PartOrd{R,A}}
}



% ============================================================
% Schranken
% ============================================================

\subsection{Schranken}

\subsubsection{Untere und obere Schrankenmengen}

\FormulaDefDelta[Untere Schrankenmenge]{%
  T\subseteq A \vdash
  \LB_{R,A}(T)\coloneqq
  \{\, a\in A \mid \forall x\in T\,\bigl((a,x)\in R\bigr)\,\}%
}{
  \DeltaRow{Mengen}{A \dsep R \dsep T \dsep a \dsep x}
  \DeltaPrem{Partielle Ordnung}{\PartOrd{R,A}}
}

\FormulaDefDelta[Obere Schrankenmenge]{%
  T\subseteq A \vdash
  \UB_{R,A}(T)\coloneqq
  \{\, a\in A \mid \forall x\in T\,\bigl((x,a)\in R\bigr)\,\}%
}{
  \DeltaRow{Mengen}{A \dsep R \dsep T \dsep a \dsep x}
  \DeltaPrem{Partielle Ordnung}{\PartOrd{R,A}}
}

\subsubsection{Infimum und Supremum}

\FormulaDefDelta[Infimum]{%
  \begin{aligned}
    &T\subseteq A \dsep \exists m\in \LB_{R,A}(T)\,
      \forall a\in \LB_{R,A}(T)\,\bigl((a,m)\in R\bigr)\\
    &\vdash \inf_{R,A}(T)\coloneqq \Max{R,A}\bigl(\LB_{R,A}(T)\bigr)
  \end{aligned}%
}{
  \DeltaRow{Mengen}{A \dsep R \dsep T \dsep m \dsep a}
  \DeltaPrem{Partielle Ordnung}{\PartOrd{R,A}}
}

\FormulaDefDelta[Supremum]{%
  \begin{aligned}
    &T\subseteq A \dsep \exists m\in \UB_{R,A}(T)\,
      \forall a\in \UB_{R,A}(T)\,\bigl((m,a)\in R\bigr)\\
    &\vdash \sup_{R,A}(T)\coloneqq \Min{R,A}\bigl(\UB_{R,A}(T)\bigr)
  \end{aligned}%
}{
  \DeltaRow{Mengen}{A \dsep R \dsep T \dsep m \dsep a}
  \DeltaPrem{Partielle Ordnung}{\PartOrd{R,A}}
}


\subsection{Restriktion einer partiellen Ordnung}

% ============================================================
% Restriktion einer partiellen Ordnung auf T
% ============================================================

% ------------------------------------------------------------
% (0) Lemma: Relation auf T
% ------------------------------------------------------------
\FormulaThmDelta[Relation auf $T$]{%
  R\cap (T\times T)\subseteq T\times T%
}{
  \DeltaRow{Mengen}{R \dsep T}
}
\begin{tabproof}
  \proofstep{}{\;R\cap (T\times T)\subseteq T\times T}{%
    \FormulaRefAuto{A \cap B \subseteq B}}
\end{tabproof}

% ------------------------------------------------------------
% (1) Lemma: Reflexivität der Restriktion
% ------------------------------------------------------------
\FormulaThmDelta[Reflexivität der Restriktion]{%
  T\subseteq A \dsep x\in T
  \vdash (x,x)\in R\cap (T\times T)%
}{
  \DeltaRow{Mengen}{A \dsep R \dsep T \dsep x}
  \DeltaPrem{Teilmenge}{T\subseteq A}
  \DeltaPrem{Partielle Ordnung}{\PartOrd{R,A}}
}
\begin{tabproof}
  \proofstep{1}{T\subseteq A}{\rA}
  \proofstep{2}{x\in T}{\rA}

  \proofstep{1,2}{x\in A}{\FormulaRefAuto{A\subseteq B, x\in A \vdash x\in B}{1,2}}

  \proofstep{1,2}{(x,x)\in R}{%
    \FormulaRefAuto{x\in A \vdash (x,x)\in R}{3}}

  \proofstep{2}{(x,x)\in T\times T}{%
    \FormulaRefAuto{a\in A \dsep b\in B \vdash (a,b)\in A\times B}{2,2}}

  \proofstep{1,2}{(x,x)\in R\cap (T\times T)}{%
    \FormulaRefAuto{x\in A \dsep x\in B \vdash x\in A\cap B}{4,5}}
\end{tabproof}

% ------------------------------------------------------------
% (2) Lemma: Transitivität der Restriktion
% ------------------------------------------------------------
\FormulaThmDelta[Transitivität der Restriktion]{%
  (x,y)\in R\cap (T\times T) \dsep (y,z)\in R\cap (T\times T)
  \vdash (x,z)\in R\cap (T\times T)%
}{
  \DeltaRow{Mengen}{A \dsep R \dsep T \dsep x \dsep y \dsep z}
  \DeltaPrem{Partielle Ordnung}{\PartOrd{R,A}}
}
\begin{tabproof}
  \proofstep{1}{(x,y)\in R\cap (T\times T)}{\rA}
  \proofstep{2}{(y,z)\in R\cap (T\times T)}{\rA}

  \proofstep{1}{(x,y)\in R}{\FormulaRefAuto{x\in A\cap B \vdash x\in A}{1}}
  \proofstep{2}{(y,z)\in R}{\FormulaRefAuto{x\in A\cap B \vdash x\in A}{2}}

  \proofstep{1,2}{(x,z)\in R}{%
    \FormulaRefAuto{(x,y)\in R \dsep (y,z)\in R \vdash (x,z)\in R}{3,4}}

  \proofstep{1}{(x,y)\in T\times T}{\FormulaRefAuto{x\in A\cap B \vdash x\in B}{1}}
  \proofstep{2}{(y,z)\in T\times T}{\FormulaRefAuto{x\in A\cap B \vdash x\in B}{2}}

  \proofstep{1}{x\in T}{\FormulaRefAuto{(a,b)\in A\times B \vdash a\in A}{6}}
  \proofstep{2}{z\in T}{\FormulaRefAuto{(a,b)\in A\times B \vdash b\in B}{7}}

  \proofstep{1,2}{(x,z)\in T\times T}{%
    \FormulaRefAuto{a\in A \dsep b\in B \vdash (a,b)\in A\times B}{8,9}}

  \proofstep{1,2}{(x,z)\in R\cap (T\times T)}{%
    \FormulaRefAuto{x\in A \dsep x\in B \vdash x\in A\cap B}{5,10}}
\end{tabproof}


% ------------------------------------------------------------
% (3) Lemma: Antisymmetrie der Restriktion
% ------------------------------------------------------------
\FormulaThmDelta[Antisymmetrie der Restriktion]{%
  (x,y)\in R\cap (T\times T) \dsep (y,x)\in R\cap (T\times T)
  \vdash x=y%
}{
  \DeltaRow{Mengen}{A \dsep R \dsep T \dsep x \dsep y}
  \DeltaPrem{Partielle Ordnung}{\PartOrd{R,A}}
}
\begin{tabproof}
  \proofstep{1}{(x,y)\in R\cap (T\times T)}{\rA}
  \proofstep{2}{(y,x)\in R\cap (T\times T)}{\rA}

  \proofstep{1}{(x,y)\in R}{\FormulaRefAuto{x\in A\cap B \vdash x\in A}{1}}
  \proofstep{2}{(y,x)\in R}{\FormulaRefAuto{x\in A\cap B \vdash x\in A}{2}}

  \proofstep{1,2}{x=y}{%
    \FormulaRefAuto{(x,y)\in R \dsep (y,x)\in R \vdash x=y}{3,4}}
\end{tabproof}


% ------------------------------------------------------------
% (4) Hauptsatz: Restriktion ist wieder partielle Ordnung
% ------------------------------------------------------------
\FormulaThmDelta[Restriktion einer partiellen Ordnung]{%
  T\subseteq A \dsep \PartOrd{R,A}\vdash \PartOrd{R\cap (T\times T),T}%
}{
  \DeltaRow{Mengen}{A \dsep R \dsep T \dsep x \dsep y \dsep z}
  \DeltaPrem{Teilmenge}{T\subseteq A}
  \DeltaPrem{Partielle Ordnung}{\PartOrd{R,A}}
  \DeltaRow{\textbf{Begründung}}{}
  \DeltaRow{Relation auf \(T\)}{R\cap (T\times T)\subseteq T\times T}
    [\FormulaRefAuto{R\cap (T\times T)\subseteq T\times T}]
  \DeltaRow{Reflexivität}{x\in T \vdash (x,x)\in R\cap (T\times T)}
    [\FormulaRefAuto{T\subseteq A \dsep x\in T \vdash (x,x)\in R\cap (T\times T)}]
  \DeltaRow{Transitivität}{%
    \begin{aligned}
      &(x,y)\dsep (y,z)\in R\cap (T\times T)\\
      &\vdash (x,z)\in R\cap (T\times T)
    \end{aligned}}
    [\FormulaRefAuto{(x,y)\in R\cap (T\times T) \dsep (y,z)\in R\cap (T\times T) \vdash (x,z)\in R\cap (T\times T)}]
  \DeltaRow{Antisymmetrie}{%
    \begin{aligned}
      &(x,y)\dsep (y,x)\in R\cap (T\times T)\\
      &\vdash x=y
    \end{aligned}}
    [\FormulaRefAuto{(x,y)\in R\cap (T\times T) \dsep (y,x)\in R\cap (T\times T) \vdash x=y}]
}{}



% ============================================================
% Wohlordnung
% (totale Ordnung + jedes nichtleere T ⊆ A besitzt ein Minimum)
% ============================================================

\subsection{Wohlordnung}
\subsubsection{Axiom einer Wohlordnung}

\FormulaAxiomDelta[Wohlordnungsprinzip]{%
  T\subseteq A \dsep T\neq\varnothing
  \vdash \exists m\in T\,\forall x\in T\,\bigl((m,x)\in R\bigr)%
}{
  \DeltaRow{Mengen}{A \dsep R \dsep T \dsep m \dsep x}
}

\subsubsection{Definition der Wohlordnung}

\FormulaDefDeltaK[Begriff der Wohlordnung]{\WellOrd{R,A}}{Wohlordnung}{
  \DeltaRow{Mengen}{A \dsep R \dsep T \dsep m \dsep x \dsep y \dsep z}
  \DeltaRow{\textbf{Axiome}}{}
  \DeltaRow{Totale Ordnung}
           {\TotOrd{R,A}}
           [\FormulaRefAuto{Totale Ordnung}]
  \DeltaRow{Wohlordnungsprinzip}{%
    \begin{aligned}
      &T\subseteq A \dsep T\neq\varnothing\\
      &\vdash \exists m\in T\,\forall x\in T\,\bigl((m,x)\in R\bigr)
    \end{aligned}}
           [\FormulaRefAuto{T\subseteq A \dsep T\neq\varnothing \vdash \exists m\in T\,\forall x\in T\,\bigl((m,x)\in R\bigr)}]
  \DeltaRow{\textbf{Neue Symbole}}{}
  \DeltaPrem{Wohlordnungen}{}
}

\FormulaDefDeltaK[Begriff der Wohlordenbarkeit]{%
  \WellOrdable{A}\coloneqq \exists R\,\WellOrd{R,A}%
}{WellOrdable}{
  \DeltaRow{Mengen}{A \dsep R}
  \DeltaRow{Wohlordnung}
           {\WellOrd{R,A}}
           [\FormulaRefAuto{Wohlordnung}]
  \DeltaRow{\textbf{Neue Symbole}}{}
  \DeltaPrem{Wohlordenbare Mengen}{}
}

\FormulaDefDeltaK[Wohlordnungssatz]{%
  \WOP\coloneqq \forall A\,\WellOrdable{A}%
}{WOP}{
  \DeltaRow{Mengen}{A}
  \DeltaRow{Wohlordenbarkeit}
           {\WellOrdable{A}}
           [\FormulaRefAuto{WellOrdable}]
}




\section{Begriff der Funktion}

\subsection{Axiom der funktionalen Eindeutigkeit}

\FormulaAxiomDelta[Funktionale Eindeutigkeit]
{(x,y)\in F\dsep (x,z)\in F \vdash y=z}
{
\DeltaRow{Mengen}{x\dsep y\dsep z\dsep F}
}

\subsection{Definition der Funktion}


\FormulaDefDeltaK[Begriff der Funktion]{F\colon A\to B}{Funktion}{
  \DeltaRow{Mengen}{A\dsep B\dsep F\dsep x\dsep y\dsep z}
  \DeltaRow{\textbf{Axiome}}{}
  %
  % — Existenz als Menge geordneter Paare —
  \DeltaRow{\makecell[l]{Menge geordneter\\Paare}}
           {F \subseteq A \times B}
           [\FormulaRefAuto{F \subseteq A \times B}]
  %
  % — Funktionale Eindeutigkeit —
  \DeltaRow{\makecell[l]{Funktionale \\Eindeutigkeit}}
           {(x,y)\in F\dsep (x,z)\in F \vdash y=z}
           [\FormulaRefAuto{(x,y)\in F\dsep (x,z)\in F \vdash y=z}]
  %
  % — Aussonderung (Schema) —
  \DeltaRow{Totalität}
           {x\in A \vdash \exists y\,(x,y)\in F}
           [\FormulaRefAuto{x\in A \vdash \exists y\,(x,y)\in F}]
%
    \DeltaRow{\textbf{Neue Symbole}}{}
    \DeltaPrem{Funktionen}{ }
}

\section{Grundlegende Eigenschaften}

\FormulaThmDelta{F\colon A\to B \vdash \TotRel{F,A,B}}{
\DeltaRow{Mengen}{A\dsep B\dsep F}
}
\begin{tabproof}
    \proofstep{1}{F\colon A\to B}{\rA}
    \proofstep{1}{F \subseteq A \times B}{\FormulaRefAuto{R \subseteq A \times B}{1}}
    \proofstep{1}{x \in A \vdash \exists y\,\bigl((x,y)\in F\bigr)}{\FormulaRefAuto{x\in A \vdash \exists y\,(x,y)\in F}{1}}
    \proofstep{1}{\TotRel{F,A,B}}{\FormulaRefAuto{Totale Relation}{2,3}}
\end{tabproof}


\FormulaThmDelta{(x,y)\in F\vdash (x,y)\in A\times B}{
\DeltaRow{Mengen}{x\dsep y\dsep A\dsep B}
\DeltaPrem{Funktionen}{F\colon A\to B}
}
\begin{tabproof}
  \proofstep{1}{(x,y)\in F}{\rA}
  \proofstep{}{F\subseteq A\times B}{\FormulaRefAuto{F \subseteq A \times B}}
  \proofstep{1}{(x,y)\in A\times B}{\FormulaRefAuto{A\subseteq B,\, x\in A \vdash x\in B}{2,1}}
\end{tabproof}

\FormulaThmDelta{(x,y)\in F\vdash x\in A}{
\DeltaRow{Mengen}{x\dsep y\dsep A\dsep B}
\DeltaPrem{Funktionen}{F\colon A\to B}
}
\begin{tabproof}
  \proofstep{1}{(x,y)\in F}{\rA}
  \proofstep{1}{(x,y)\in A\times B}{\FormulaRefAuto{(x,y)\in F\vdash (x,y)\in A\times B}{1}}
  \proofstep{1}{x\in A}{\FormulaRefAuto{(a,b)\in A\times B\vdash a\in A}{2}}
\end{tabproof}

\FormulaThmDelta{(x,y)\in F\vdash y\in B}{
\DeltaRow{Mengen}{x\dsep y\dsep A\dsep B}
\DeltaPrem{Funktionen}{F\colon A\to B}
}
\begin{tabproof}
  \proofstep{1}{(x,y)\in F}{\rA}
  \proofstep{1}{(x,y)\in A\times B}{\FormulaRefAuto{(x,y)\in F\vdash (x,y)\in A\times B}{1}}
  \proofstep{1}{y\in B}{\FormulaRefAuto{(a,b)\in A\times B\vdash b\in B}{2}}
\end{tabproof}

\section{Funktionswert}

\FormulaThmDelta[Eindeutigkeit des Funktionswertes]{x\in A\vdash \exists! y\;(x,y)\in F}{
\DeltaRow{Mengen}{x\dsep A}
\DeltaPrem{Funktionen}{F\colon A\to B}
}
\begin{tabproof}
  \proofstep{1}{x\in A}{\rA}
  \proofstep{1}{\exists y(x,y)\in F}{\FormulaRefAuto{x\in A \vdash \exists y\,(x,y)\in F}{1}}
  \proofstep{2}{(x,y)\in F}{\rA}
  \proofstep{3}{(x,z)\in F}{\rA}
  \proofstep{2,3}{y=z}{\FormulaRefAuto{(x,y)\in F\dsep (x,z)\in F \vdash y=z}{2,3}}
  \proofstep{1}{\exists! y(x,y)\in F}{\UEI{1,2,3,4}}
\end{tabproof}
%%begin novalidate
\FormulaDefDelta[Funktionswert]{x\in A\vdash (F(x) \coloneqq \iota y\,\bigl((x,y)\in F\bigr)}{
\DeltaRow{Mengen}{x\dsep A}
\DeltaPrem{Funktionen}{F\colon A\to B}
}
%%end novalidate

\FormulaThmDelta{x\in A\vdash (x,F(x))\in F}{\DeltaRow{Mengen}{x\dsep A}
\DeltaPrem{Funktionen}{F\colon A\to B}
}
\begin{tabproof}
    \proofstep{1}{x\in A}{\rA}
    \proofstep{1}{(x,F(x))\in F}{\FormulaRefAuto{x\in A\vdash (F(x) \coloneqq \iota y\,\bigl((x,y)\in F\bigr)}{1}}
\end{tabproof}

\subsection{Eigenschaften des Funktionswertes}

\FormulaThmDelta[Ersetzung im Funktionsargument]{%
  x = y \dsep a = F(y) \vdash a = F(x)
}{
  \DeltaRow{Mengen}{A \dsep B \dsep x \dsep y \dsep a}
  \DeltaPrem{Funktionen}{F\colon A\to B}
}
\begin{tabproof}
  \proofstep{1}{x = y}{\rA}
  \proofstep{2}{a = F(y)}{\rA}
  \proofstep{1}{y = x}{\FormulaRefAuto{a = b \vdash b = a}{1}}
  \proofstep{1,2}{a = F(x)}{\rIE{3,2}}
\end{tabproof}



\FormulaThmDelta{x\in A\dsep y\in A\dsep x=y\vdash F(x)=F(y)}
{
\DeltaRow{Mengen}{x\dsep y\dsep A}
\DeltaPrem{Funktionen}{F\colon A\to B}
}
\begin{tabproof}
    \proofstep{1}{x\in A}{\rA}
    \proofstep{2}{y\in A}{\rA}
    \proofstep{3}{x=y}{\rA}
    \proofstep{1}{(x,F(x))\in F}{\FormulaRefAuto{x\in A\vdash (F(x) \coloneqq \iota y\,\bigl((x,y)\in F\bigr)}{1}}
    \proofstep{2}{(y,F(y))\in F}{\FormulaRefAuto{x\in A\vdash (F(x) \coloneqq \iota y\,\bigl((x,y)\in F\bigr)}{2}}
    \proofstep{1,3}{(y,F(x))\in F}{\rIE{3,4}}
    \proofstep{1,2,3}{F(x)=F(y)}{\FormulaRefAuto{(x,y)\in F\dsep (x,z)\in F \vdash y=z}{6,5}}
\end{tabproof}

\FormulaThmDelta{%
  x_1\in A \dsep x_2\in A \dsep F(x_1)\neq F(x_2) \vdash x_1\neq x_2%
}{
  \DeltaRow{Mengen}{A\dsep B\dsep x_1\dsep x_2}
  \DeltaPrem{Funktionen}{F\colon A\to B}
}
\begin{tabproof}
  \proofstep{1}{x_1\in A}{\rA}
  \proofstep{2}{x_2\in A}{\rA}
  \proofstep{3}{F(x_1)\neq F(x_2)}{\rA}

  \proofstep{4}{x_1=x_2}{\rA}
  \proofstep{1,2,4}{F(x_1)=F(x_2)}{%
    \FormulaRefAuto{x\in A\dsep y\in A\dsep x=y\vdash F(x)=F(y)}{1,2,4}}
  \proofstep{1,2,3,4}{\bot}{\rBI{3,5}}
  \proofstep{1,2,3}{x_1\neq x_2}{\rCI{4,6}}
\end{tabproof}

\FormulaThmDelta{x\in A\dsep F(x)=y\vdash (x,y)\in F}
{
\DeltaRow{Mengen}{x\dsep y\dsep A}
\DeltaPrem{Funktionen}{F\colon A\to B}
}
\begin{tabproof}
  \proofstep{1}{x\in A}{\rA}
  \proofstep{2}{F(x)=y}{\rA}
  \proofstep{1}{(x,F(x))\in F}{\FormulaRefAuto{x\in A\vdash (F(x) \coloneqq \iota y\,\bigl((x,y)\in F\bigr)}{1}}
  \proofstep{1,2}{(x,y)\in F}{\rIE{2,3}}
\end{tabproof}

\FormulaThmDelta{x\in A\dsep y=F(x)\vdash (x,y)\in F}
{
\DeltaRow{Mengen}{x\dsep y\dsep A}
\DeltaPrem{Funktionen}{F\colon A\to B}
}
\begin{tabproof}
  \proofstep{1}{x\in A}{\rA}
  \proofstep{2}{y=F(x)}{\rA}
  \proofstep{2}{F(x)=y}{\FormulaRefAuto{a=b\vdash b=a}{2}}
  \proofstep{1,2}{(x,y)\in F}{\FormulaRefAuto{x\in A\dsep F(x)=y\vdash (x,y)\in F}{1,3}}
\end{tabproof}

\FormulaThmDelta{x\in A\dsep (x,y)\in F\vdash y=F(x)}{
\DeltaRow{Mengen}{x\dsep y\dsep A}
\DeltaPrem{Funktionen}{F\colon A\to B}
}
\begin{tabproof}
  \proofstep{1}{(x,y)\in F}{\rA}
  \proofstep{2}{x\in A}{\rA}
  \proofstep{2}{(x,F(x))\in F}{\FormulaRefAuto{x\in A\vdash (F(x) \coloneqq \iota y\,\bigl((x,y)\in F\bigr)}{2}}
  \proofstep{1}{y=F(x)}{\FormulaRefAuto{(x,y)\in F\dsep (x,z)\in F \vdash y=z}{1,3}}
\end{tabproof}

\FormulaThmDelta{x\in A\dsep (x,y)\in F\vdash F(x)=y}{
\DeltaRow{Mengen}{x\dsep y\dsep A}
\DeltaPrem{Funktionen}{F\colon A\to B}
}
\begin{tabproof}
  \proofstep{1}{(x,y)\in F}{\rA}
  \proofstep{2}{x\in A}{\rA}
  \proofstep{2}{(x,F(x))\in F}{\FormulaRefAuto{x\in A\vdash (F(x) \coloneqq \iota y\,\bigl((x,y)\in F\bigr)}{2}}
  \proofstep{1,2}{F(x)=y}{\FormulaRefAuto{(x,y)\in F\dsep (x,z)\in F \vdash y=z}{3,1}}
\end{tabproof}


\FormulaThmDelta{F\colon A\to B\dsep x\in A\vdash F(x)\in B}{
\DeltaRow{Mengen}{A\dsep B}
}
\begin{tabproof}
  \proofstep{1}{F\colon A\to B}{\rA}
  \proofstep{2}{x\in A}{\rA}
  \proofstep{1}{F\subseteq A\times B}{\FormulaRefAuto{F \subseteq A \times B}{1}}
  \proofstep{2}{(x,F(x))\in F}{\FormulaRefAuto{x\in A\vdash (F(x) \coloneqq \iota y\,\bigl((x,y)\in F\bigr)}{2}}
  \proofstep{1,2}{(x,F(x))\in A\times B}{\FormulaRefAuto{ A\subseteq B,\, x\in A \vdash x\in B }{3,4}}
  \proofstep{1,2}{F(x)\in B}{\FormulaRefAuto{(a,b)\in A\times B\vdash b\in B}{5}}
\end{tabproof}

\FormulaThmDelta{(x,y)\in F\vdash y=F(x)}{
\DeltaRow{Mengen}{x\dsep y\dsep A\dsep B}
\DeltaPrem{Funktionen}{F\colon A\to B}
}
\begin{tabproof}
  \proofstep{1}{(x,y)\in F}{\rA}
  \proofstep{1}{x\in A}{\FormulaRefAuto{(x,y)\in F\vdash x\in A}{1}}
  \proofstep{1}{(x,F(x))\in F}{\FormulaRefAuto{x\in A\vdash (F(x) \coloneqq \iota y\,\bigl((x,y)\in F\bigr)}{2}}
  \proofstep{1}{y=F(x)}{\FormulaRefAuto{(x,y)\in F\dsep (x,z)\in F \vdash y=z}{1,3}}
\end{tabproof}

\FormulaThmDeltaR{%
  (x,y)\in F\vdash F(x)=y
}{%
  F\colon A\to B\dsep (x,y)\in F\vdash F(x)=y
}{%
  \DeltaRow{Mengen}{x\dsep y\dsep A\dsep B}
  \DeltaPrem{Funktionen}{F\colon A\to B}
}
\begin{tabproof}
  \proofstep{1}{(x,y)\in F}{\rA}
  \proofstep{1}{y=F(x)}{%
    \FormulaRefAuto{(x,y)\in F\vdash y=F(x)}{1}}
  \proofstep{1}{F(x)=y}{%
    \FormulaRefAuto{a=b\vdash b=a}{2}}
\end{tabproof}

\FormulaThmDelta{(x,y)\in F\vdash F(x)=y}{
\DeltaRow{Mengen}{x\dsep y\dsep A\dsep B}
\DeltaPrem{Funktionen}{F\colon A\to B}
}
\begin{tabproof}
  \proofstep{1}{(x,y)\in F}{\rA}
  \proofstep{1}{x\in A}{\FormulaRefAuto{(x,y)\in F\vdash x\in A}{1}}
  \proofstep{1}{(x,F(x))\in F}{\FormulaRefAuto{x\in A\vdash (F(x) \coloneqq \iota y\,\bigl((x,y)\in F\bigr)}{2}}
  \proofstep{1}{F(x)=y}{\FormulaRefAuto{(x,y)\in F\dsep (x,z)\in F \vdash y=z}{3,1}}
\end{tabproof}

\FormulaThmDelta{F(x)=G(x)\dsep (x,y)\in F\vdash (x,y)\in G}{
\DeltaRow{Mengen}{x\dsep A\dsep B}
\DeltaPrem{Funktionen}{F,G\colon A\to B}
}
\begin{tabproof}
  \proofstep{1}{F(x)=G(x)}{\rA}
  \proofstep{2}{(x,y)\in F}{\rA}
  \proofstep{2}{y=F(x)}{\FormulaRefAuto{(x,y)\in F\vdash F(x)=y}{2}}
  \proofstep{1,2}{y=G(x)}{\FormulaRefAuto{a=b, b=c \vdash a=c}{3,1}}
  \proofstep{2}{x\in A}{\FormulaRefAuto{(x,y)\in F\vdash x\in A}{2}}
  \proofstep{1,2}{(x,y)\in G}{\FormulaRefAuto{x\in A\dsep y=F(x)\vdash (x,y)\in F}{5,4}}
\end{tabproof}

\FormulaThmDelta{%
  \{x\in A\mid F(x)\in \{a\}\}=\{x\in A\mid F(x)=a\}%
}{
  \DeltaRow{Mengen}{A\dsep a\dsep x}
  \DeltaPrem{Funktionen}{F}
}
\begin{tabproof}
  \proofstep{1}{x\in \{x\in A\mid F(x)\in \{a\}\}}{\rA}
  \proofstep{1}{x\in A\land F(x)\in \{a\}}{%
    \FormulaRefAuto{x \in \{u \in A \mid P(u)\} \eqvdash x \in A \land P(x)}{1}}
  \proofstep{1}{x\in A}{\rAEa{2}}
  \proofstep{1}{F(x)\in \{a\}}{\rAEb{2}}
  \proofstep{1}{F(x)=a}{\FormulaRefAuto{x \in \{a\} \eqvdash x = a}{4}}
  \proofstep{1}{x\in \{x\in A\mid F(x)=a\}}{%
    \FormulaRefAuto{x \in A\dsep P(x)\vdash x \in \{x \in A \mid P(x)\}}{3,5}}
  \proofstep{}{\{x\in A\mid F(x)\in \{a\}\}\subseteq \{x\in A\mid F(x)=a\}}{%
    \FormulaRefAuto{A \subseteq B \coloneq \forall x\,(x\in A \rightarrow x\in B)}{\rUI{\rRI{1,6}}}}

  \proofstep{8}{x\in \{x\in A\mid F(x)=a\}}{\rA}
  \proofstep{8}{x\in A\land F(x)=a}{%
    \FormulaRefAuto{x \in \{u \in A \mid P(u)\} \eqvdash x \in A \land P(x)}{8}}
  \proofstep{8}{x\in A}{\rAEa{9}}
  \proofstep{8}{F(x)=a}{\rAEb{9}}
  \proofstep{8}{F(x)\in \{a\}}{\FormulaRefAuto{x \in \{a\} \eqvdash x = a}{11}}
  \proofstep{8}{x\in \{x\in A\mid F(x)\in \{a\}\}}{%
    \FormulaRefAuto{x \in A\dsep P(x)\vdash x \in \{x \in A \mid P(x)\}}{10,12}}
  \proofstep{}{\{x\in A\mid F(x)=a\}\subseteq \{x\in A\mid F(x)\in \{a\}\}}{%
    \FormulaRefAuto{A \subseteq B \coloneq \forall x\,(x\in A \rightarrow x\in B)}{\rUI{\rRI{8,13}}}}

  \proofstep{}{\{x\in A\mid F(x)\in \{a\}\}=\{x\in A\mid F(x)=a\}}{%
    \FormulaRefAuto{A \subseteq B, B \subseteq A \vdash A = B}{7,14}}
\end{tabproof}


\subsection{Zur Gleichheit von Funktionen}

\FormulaThmDelta{\forall x\in A (F(x)=G(x)) \vdash \forall (x,y)((x,y)\in F\rightarrow (x,y)\in G)}{
\DeltaRow{Mengen}{A\dsep B}
\DeltaPrem{Funktionen}{F,G\colon A\to B}
}
\begin{tabproof}
  \proofstep{1}{\forall x\in A (F(x)=G(x))}{\rA}
  \proofstep{2}{(x,y)\in F}{\rA}
  \proofstep{2}{x\in A}{\FormulaRefAuto{(x,y)\in F\vdash x\in A}{2}}
  \proofstep{1,2}{F(x)=G(x)}{\rRE{\rUE{1},3}}
  \proofstep{1,2}{(x,y)\in G}{\FormulaRefAuto{F(x)=G(x)\dsep (x,y)\in F\vdash (x,y)\in G}{4,2}}
  \proofstep{1}{(x,y)\in F\rightarrow (x,y)\in G}{\rRI{2,5}}
  \proofstep{1}{\forall (x,y)((x,y)\in F\rightarrow (x,y)\in G)}{\rUI{6}}
\end{tabproof}



\FormulaThmDelta{\forall x\in A (F(x)=G(x)) \vdash \forall (x,y)((x,y)\in G\rightarrow (x,y)\in F)}{
\DeltaRow{Mengen}{A\dsep B}
\DeltaPrem{Funktionen}{F,G\colon A\to B}
}
\begin{tabproof}
  \proofstep{1}{\forall x\in A (F(x)=G(x))}{\rA}
  \proofstep{2}{x\in A}{\rA}
  \proofstep{1,2}{F(x)=G(x)}{\rRE{\rUE{1},2}}
  \proofstep{1,2}{G(x)=F(x)}{\FormulaRefAuto{a=b\vdash b=a}{3}}
  \proofstep{1}{\forall x\in A(G(x)=F(x))}{\rUI{\rRI{2,4}}}
  \proofstep{1}{\forall (x,y)((x,y)\in G\rightarrow (x,y)\in F)}{\FormulaRefAuto{\forall x\in A (F(x)=G(x)) \vdash \forall (x,y)((x,y)\in F\rightarrow (x,y)\in G)}{5}}
\end{tabproof}

\FormulaThmDelta{\forall x\in A (F(x)=G(x)) \eqvdash \forall (x,y)((x,y)\in F\leftrightarrow (x,y)\in G)}{
\DeltaRow{Mengen}{A\dsep B}
\DeltaPrem{Funktionen}{F,G\colon A\to B}
}
\begin{tabproofsplit}
\proofpart{\(\vdash\)}
  \proofstep{1}{\forall x\in A (F(x)=G(x))}{\rA}
  \proofstep{1}{\forall (x,y)((x,y)\in F\rightarrow (x,y)\in G)}{\FormulaRefAuto{\forall x\in A (F(x)=G(x)) \vdash \forall (x,y)((x,y)\in F\rightarrow (x,y)\in G)}{1}}
  \proofstep{1}{\forall (x,y)((x,y)\in G\rightarrow (x,y)\in F)}{\FormulaRefAuto{\forall x\in A (F(x)=G(x)) \vdash \forall (x,y)((x,y)\in G\rightarrow (x,y)\in F)}{1}}
  \proofstep{1}{\forall (x,y)((x,y)\in F\leftrightarrow (x,y)\in G)}{\FormulaRefAuto{\forall x (P(x) \leftrightarrow Q(x)) \eqvdash \forall x (P(x) \rightarrow Q(x)) \land \forall x (Q(x) \rightarrow P(x))}{2,3}}
\closeproofpart
\proofpart{\(\dashv\)}
  \proofstep{1}{\forall (x,y)((x,y)\in F\leftrightarrow (x,y)\in G)}{\rA}
  \proofstep{2}{x\in A}{\rA}
  \proofstep{2}{(x,F(x))\in F}{\FormulaRefAuto{x\in A\vdash (F(x) \coloneqq \iota y\,\bigl((x,y)\in F\bigr)}{2}}
  \proofstep{1}{(x,F(x))\in F\leftrightarrow (x,F(x))\in G}{\rUI{1}}
  \proofstep{1,2}{(x,F(x))\in G}{\FormulaRefAuto{P \leftrightarrow Q, P \vdash Q}{4,3}}
  \proofstep{2}{(x,G(x))\in G}{\FormulaRefAuto{x\in A\vdash (F(x) \coloneqq \iota y\,\bigl((x,y)\in F\bigr)}{2}}
  \proofstep{1,2}{F(x)=G(x)}{\FormulaRefAuto{(x,y)\in F\dsep (x,z)\in F \vdash y=z}{5,6}}
  \proofstep{1,2}{\forall x\in A(F(x)=G(x))}{\rUI{\rRI{2,7}}}
\closeproofpart
\end{tabproofsplit}

\FormulaThmDelta{\forall x\in A (F(x)=G(x)) \eqvdash F=G}{
\DeltaRow{Mengen}{A\dsep B}
\DeltaPrem{Funktionen}{F,G\colon A\to B}
}
\begin{tabproofwide}
  \proofstepwide{\forall x\in A (F(x)=G(x))}{\leftrightarrow}{\forall (x,y)((x,y)\in F\leftrightarrow (x,y)\in G)}%
    {\FormulaRefAuto{\forall x\in A (F(x)=G(x)) \eqvdash \forall (x,y)((x,y)\in F\leftrightarrow (x,y)\in G)}}
  \proofstepwide{}{\leftrightarrow}{F=G}%
    {\FormulaRefAuto{\forall x\, (x \in A \leftrightarrow x \in B) \eqvdash A = B}{1}}
   \proofstepwide{\forall x\in A (F(x)=G(x))}{\leftrightarrow}{F=G}%
    {\rChain{1,2}}
\end{tabproofwide}

\FormulaThmDelta{x\in A\dsep F=G\vdash F(x)=G(x)}{
\DeltaRow{Mengen}{A\dsep B}
\DeltaPrem{Funktionen}{F,G\colon A\to B}
}
\begin{tabproof}
    \proofstep{1}{x\in A}{\rA}
    \proofstep{2}{F=G}{\rA}
    \proofstep{2}{\forall x\in A (F(x)=G(x))}{\FormulaRefAuto{\forall x\in A (F(x)=G(x)) \eqvdash F=G}{2}}
    \proofstep{1,2}{\forall x\in A (F(x)=G(x))}{\rRE{1,\rUI{3}}}
\end{tabproof}


\section{Axiom der Ersetzung}

\FormulaAxiomDelta[Ersetzung]{
\exists B\;\forall y\;\bigl( y\in B\;\leftrightarrow\; \exists x\in A\;y=F(x) \bigr)
}{
\DeltaRow{Mengen}{A}
\DeltaPrem{Funktionen}{F\colon A\to B}
}

\section{Die Bildmenge}

\subsection{Bildmengen von Teilmengen}

%%begin novalidate
\FormulaDefDelta[Bildmenge einer Teilmenge]
{%
  C\subseteq A\dsep F\colon A\to B\vdash
  F[C]\coloneqq \iota D\Bigl(\forall y\,\bigl(y\in D \leftrightarrow \exists x\in C\,y=F(x)\bigr)\Bigr)
}
{
  \DeltaRow{Mengen}{A\dsep B\dsep C}
  \DeltaRow{Funktionen}{F}
}
%%end novalidate

\FormulaThmDelta[Charakterisierung der Bildmenge einer Teilmenge]{%
  C\subseteq A\dsep F\colon A\to B\vdash y\in F[C]\leftrightarrow \exists x\in C\,y=F(x)
}
{
  \DeltaRow{Mengen}{A\dsep B\dsep C\dsep x\dsep y}
  \DeltaRow{Funktionen}{F}
}
\begin{tabproof}
  \proofstep{1}{C\subseteq A}{\rA}
  \proofstep{2}{F\colon A\to B}{\rA}
  \proofstep{1,2}{\forall y\,\bigl(y\in F[C]\leftrightarrow \exists x\in C\,y=F(x)\bigr)}{%
    \FormulaRefAuto{C\subseteq A\dsep F\colon A\to B\vdash F[C]\coloneqq \iota D\Bigl(\forall y\,\bigl(y\in D \leftrightarrow \exists x\in C\,y=F(x)\bigr)\Bigr)}{1,2}}
  \proofstep{1,2}{y\in F[C]\leftrightarrow \exists x\in C\,y=F(x)}{\rUE{3}}
\end{tabproof}

\FormulaThmDelta[Aus einem Bildelement einer Teilmenge folgt ein Urbild in der Teilmenge]{%
  C\subseteq A\dsep F\colon A\to B\dsep y\in F[C]\vdash \exists x\in C\,y=F(x)
}
{
  \DeltaRow{Mengen}{A\dsep B\dsep C\dsep x\dsep y}
  \DeltaRow{Funktionen}{F}
}
\begin{tabproof}
  \proofstep{1}{C\subseteq A}{\rA}
  \proofstep{2}{F\colon A\to B}{\rA}
  \proofstep{3}{y\in F[C]}{\rA}
  \proofstep{1,2}{y\in F[C]\leftrightarrow \exists x\in C\,y=F(x)}{%
    \FormulaRefAuto{C\subseteq A\dsep F\colon A\to B\vdash y\in F[C]\leftrightarrow \exists x\in C\,y=F(x)}{1,2}}
  \proofstep{1,2,3}{\exists x\in C\,y=F(x)}{%
    \FormulaRefAuto{P\leftrightarrow Q, P \vdash Q}{4,3}}
\end{tabproof}

\FormulaThmDelta[Ein Urbild in einer Teilmenge liefert ein Bildelement der Teilmenge]{%
  C\subseteq A\dsep F\colon A\to B\dsep \exists x\in C\,y=F(x)\vdash y\in F[C]
}
{
  \DeltaRow{Mengen}{A\dsep B\dsep C\dsep x\dsep y}
  \DeltaRow{Funktionen}{F}
}
\begin{tabproof}
  \proofstep{1}{C\subseteq A}{\rA}
  \proofstep{2}{F\colon A\to B}{\rA}
  \proofstep{3}{\exists x\in C\,y=F(x)}{\rA}
  \proofstep{1,2}{y\in F[C]\leftrightarrow \exists x\in C\,y=F(x)}{%
    \FormulaRefAuto{C\subseteq A\dsep F\colon A\to B\vdash y\in F[C]\leftrightarrow \exists x\in C\,y=F(x)}{1,2}}
  \proofstep{1,2,3}{y\in F[C]}{%
    \FormulaRefAuto{P\leftrightarrow Q, Q \vdash P}{4,3}}
\end{tabproof}

\FormulaThmDelta[Elemente einer Teilmenge liegen im Bild der Teilmenge]{%
  C\subseteq A\dsep F\colon A\to B\dsep x\in C\vdash F(x)\in F[C]
}
{
  \DeltaRow{Mengen}{A\dsep B\dsep C\dsep x\dsep z}
  \DeltaRow{Funktionen}{F}
}
\begin{tabproof}
  \proofstep{1}{C\subseteq A}{\rA}
  \proofstep{2}{F\colon A\to B}{\rA}
  \proofstep{3}{x\in C}{\rA}
  \proofstep{}{F(x)=F(x)}{\rII}
  \proofstep{3}{\exists z\in C\,F(x)=F(z)}{\rEI{\rAI{3,4}}}
  \proofstep{1,2,3}{F(x)\in F[C]}{%
    \FormulaRefAuto{C\subseteq A\dsep F\colon A\to B\dsep \exists x\in C\,y=F(x)\vdash y\in F[C]}{1,2,5}}
\end{tabproof}

\paragraph{Spezialfall \(C=A\)}

\FormulaThmDelta{F\colon A\to B\vdash y\in F[A]\leftrightarrow \exists x\in A\,y=F(x)}
{
  \DeltaRow{Mengen}{A\dsep B\dsep x\dsep y}
  \DeltaRow{Funktionen}{F}
}
\begin{tabproof}
  \proofstep{}{A\subseteq A}{\FormulaRefAuto{A\subseteq A}}
  \proofstep{2}{F\colon A\to B}{\rA}
  \proofstep{1,2}{y\in F[A]\leftrightarrow \exists x\in A\,y=F(x)}{%
    \FormulaRefAuto{C\subseteq A\dsep F\colon A\to B\vdash y\in F[C]\leftrightarrow \exists x\in C\,y=F(x)}{1,2}}
\end{tabproof}

\FormulaThmDelta[Aus einem Bildelement folgt ein Urbild]{%
  F\colon A\to B\dsep y\in F[A]\vdash \exists x\in A\,y=F(x)
}{
  \DeltaRow{Mengen}{A\dsep B\dsep x\dsep y}
  \DeltaRow{Funktionen}{F}
}
\begin{tabproof}
  \proofstep{}{A\subseteq A}{\FormulaRefAuto{A\subseteq A}}
  \proofstep{2}{F\colon A\to B}{\rA}
  \proofstep{3}{y\in F[A]}{\rA}
  \proofstep{1,2,3}{\exists x\in A\,y=F(x)}{%
    \FormulaRefAuto{C\subseteq A\dsep F\colon A\to B\dsep y\in F[C]\vdash \exists x\in C\,y=F(x)}{1,2,3}}
\end{tabproof}

\FormulaThmDelta[Ein Urbild liefert ein Bildelement]{%
  F\colon A\to B\dsep \exists x\in A\,y=F(x)\vdash y\in F[A]
}{
  \DeltaRow{Mengen}{A\dsep B\dsep x\dsep y}
  \DeltaRow{Funktionen}{F}
}
\begin{tabproof}
  \proofstep{}{A\subseteq A}{\FormulaRefAuto{A\subseteq A}}
  \proofstep{2}{F\colon A\to B}{\rA}
  \proofstep{3}{\exists x\in A\,y=F(x)}{\rA}
  \proofstep{1,2,3}{y\in F[A]}{%
    \FormulaRefAuto{C\subseteq A\dsep F\colon A\to B\dsep \exists x\in C\,y=F(x)\vdash y\in F[C]}{1,2,3}}
\end{tabproof}

\subsection{Rechenregeln für Bildmengen}

\FormulaThmDelta[Bild einer Einermenge]{%
  x\in A \dsep F\colon A\to B \vdash F[\{x\}]=\{F(x)\}
}{%
  \DeltaRow{Mengen}{A\dsep B\dsep x\dsep y\dsep z}
  \DeltaRow{Funktionen}{F}
}
\begin{tabproofsplitwide}
  \proofpartwideind[Teilmengenrichtung \(F[\{x\}]\subseteq \{F(x)\}\)]{%
    x\in A \dsep F\colon A\to B \vdash F[\{x\}]\subseteq \{F(x)\}
  }[%
    x\in A \dsep F\colon A\to B \vdash F[\{x\}]\subseteq \{F(x)\}
  ]
    \proofstepwidestar[1]{x\in A}{\rA}
    \proofstepwidestar[2]{F\colon A\to B}{\rA}
    \proofstepwidestar[3]{y\in F[\{x\}]}{\rA}
    \proofstepwidestar[1]{\{x\}\subseteq A}{\FormulaRefAuto{a\in A \vdash \{a\}\subseteq A}{1}}
    \proofstepwidestar[1,2,3]{\exists z\in \{x\}\,y=F(z)}{%
      \FormulaRefAuto{C\subseteq A\dsep F\colon A\to B\dsep y\in F[C]\vdash \exists x\in C\,y=F(x)}{4,2,3}}

    \proofstepwidestar[6]{z\in \{x\}\land y=F(z)}{\rA}
    \proofstepwidestar[6]{z\in \{x\}}{\rAEa{6}}
    \proofstepwidestar[6]{y=F(z)}{\rAEb{6}}
    \proofstepwidestar[6]{z=x}{\FormulaRefAuto{x\in \{a\}\eqvdash x=a}{7}}
    \proofstepwidestar[6]{x=z}{\FormulaRefAuto{a=b\vdash b=a}{9}}
    \proofstepwidestar[6]{y=F(x)}{%
      \FormulaRefAuto{x=y \dsep a=F(y) \vdash a=F(x)}{10,8}}
    \proofstepwidestar[6]{y\in \{F(x)\}}{\FormulaRefAuto{x\in \{a\}\eqvdash x=a}{11}}

    \proofstepwidestar[1,2,3]{y\in \{F(x)\}}{\rEE{5,6,12}}
    \proofstepwidestar[1,2]{\forall y\,(y\in F[\{x\}]\rightarrow y\in \{F(x)\})}{\rUI{\rRI{3,13}}}
    \proofstepwidestar[1,2]{F[\{x\}]\subseteq \{F(x)\}}{%
      \FormulaRefAuto{A\subseteq B \coloneq \forall x\,(x\in A \rightarrow x\in B)}{14}}
  \closeproofpartwideind

  \proofpartwideind[Teilmengenrichtung \(\{F(x)\}\subseteq F[\{x\}]\)]{%
    x\in A \dsep F\colon A\to B \vdash \{F(x)\}\subseteq F[\{x\}]
  }[%
    x\in A \dsep F\colon A\to B \vdash \{F(x)\}\subseteq F[\{x\}]
  ]
    \proofstepwidestar[1]{x\in A}{\rA}
    \proofstepwidestar[2]{F\colon A\to B}{\rA}
    \proofstepwidestar[3]{y\in \{F(x)\}}{\rA}
    \proofstepwidestar[3]{y=F(x)}{\FormulaRefAuto{x\in \{a\}\eqvdash x=a}{3}}
    \proofstepwidestar[]{x=x}{\rII}
    \proofstepwidestar[]{x\in \{x\}}{\FormulaRefAuto{x\in \{a\}\eqvdash x=a}{5}}
    \proofstepwidestar[3]{x\in \{x\}\land y=F(x)}{\rAI{6,4}}
    \proofstepwidestar[3]{\exists z\in \{x\}\,y=F(z)}{\rEI{7}}
    \proofstepwidestar[1]{\{x\}\subseteq A}{\FormulaRefAuto{a\in A \vdash \{a\}\subseteq A}{1}}
    \proofstepwidestar[1,2,3]{y\in F[\{x\}]}{%
      \FormulaRefAuto{C\subseteq A\dsep F\colon A\to B\dsep \exists x\in C\,y=F(x)\vdash y\in F[C]}{9,2,8}}
    \proofstepwidestar[1,2]{\forall y\,(y\in \{F(x)\}\rightarrow y\in F[\{x\}])}{\rUI{\rRI{3,10}}}
    \proofstepwidestar[1,2]{\{F(x)\}\subseteq F[\{x\}]}{%
      \FormulaRefAuto{A\subseteq B \coloneq \forall x\,(x\in A \rightarrow x\in B)}{11}}
  \closeproofpartwideind

  \proofpartwide{Zusammenfassung}
    \proofstepwidestar[1]{x\in A}{\rA}
    \proofstepwidestar[2]{F\colon A\to B}{\rA}
    \proofstepwidestar[1,2]{F[\{x\}]\subseteq \{F(x)\}}{%
      \FormulaRefAuto{x\in A \dsep F\colon A\to B \vdash F[\{x\}]\subseteq \{F(x)\}}{1,2}}
    \proofstepwidestar[1,2]{\{F(x)\}\subseteq F[\{x\}]}{%
      \FormulaRefAuto{x\in A \dsep F\colon A\to B \vdash \{F(x)\}\subseteq F[\{x\}]}{1,2}}
    \proofstepwidestar[1,2]{F[\{x\}]=\{F(x)\}}{%
      \FormulaRefAuto{A\subseteq B\dsep B\subseteq A \vdash A=B}{3,4}}
  \closeproofpartwide
\end{tabproofsplitwide}

\FormulaThmDelta[Monotonie der Bildmenge]{%
  F\colon M\to N \dsep A\subseteq B \dsep B\subseteq M \vdash F[A]\subseteq F[B]
}{%
  \DeltaRow{Mengen}{A\dsep B\dsep M\dsep N\dsep x\dsep y}
  \DeltaRow{Funktionen}{F}
}
\begin{tabproof}
  \proofstep{1}{F\colon M\to N}{\rA}
  \proofstep{2}{A\subseteq B}{\rA}
  \proofstep{3}{B\subseteq M}{\rA}
  \proofstep{4}{y\in F[A]}{\rA}
  \proofstep{2,3}{A\subseteq M}{%
    \FormulaRefAuto{A\subseteq B,\, B\subseteq C \vdash A\subseteq C}{2,3}}
  \proofstep{1,2,3,4}{\exists x\in A\,y=F(x)}{%
    \FormulaRefAuto{C\subseteq A\dsep F\colon A\to B\dsep y\in F[C]\vdash \exists x\in C\,y=F(x)}{5,1,4}}
  \proofstep{7}{x\in A\land y=F(x)}{\rA}
  \proofstep{7}{x\in A}{\rAEa{7}}
  \proofstep{7}{y=F(x)}{\rAEb{7}}
  \proofstep{2,8}{x\in B}{\FormulaRefAuto{A\subseteq B\dsep x\in A\vdash x\in B}{2,8}}
  \proofstep{2,7}{x\in B\land y=F(x)}{\rAI{10,9}}
  \proofstep{2,7}{\exists x\in B\,y=F(x)}{\rEI{11}}
  \proofstep{1,2,3,7}{y\in F[B]}{%
    \FormulaRefAuto{C\subseteq A\dsep F\colon A\to B\dsep \exists x\in C\,y=F(x)\vdash y\in F[C]}{3,1,12}}
  \proofstep{1,2,3,4}{y\in F[B]}{\rEE{6,7,13}}
  \proofstep{1,2,3}{\forall y\,(y\in F[A]\rightarrow y\in F[B])}{\rUI{\rRI{4,14}}}
  \proofstep{1,2,3}{F[A]\subseteq F[B]}{%
    \FormulaRefAuto{A\subseteq B \coloneq \forall x\,(x\in A \rightarrow x\in B)}{15}}
\end{tabproof}

\FormulaThmDelta[Bild eines Elements aus einer Vereinigung]{%
  F\colon M\to N \dsep A\subseteq M \dsep B\subseteq M \dsep x\in A\cup B \vdash F(x)\in F[A]\cup F[B]
}{%
  \DeltaRow{Mengen}{A\dsep B\dsep M\dsep N\dsep x}
  \DeltaRow{Funktionen}{F}
}
\begin{tabproof}
  \proofstep{1}{F\colon M\to N}{\rA}
  \proofstep{2}{A\subseteq M}{\rA}
  \proofstep{3}{B\subseteq M}{\rA}
  \proofstep{4}{x\in A\cup B}{\rA}
  \proofstep{4}{x\in A\lor x\in B}{%
    \FormulaRefAuto{x\in A\cup B\eqvdash x\in A\lor x\in B}{4}}

  \proofstep{6}{x\in A}{\rA}
  \proofstep{1,2,6}{F(x)\in F[A]}{%
    \FormulaRefAuto{C\subseteq A\dsep F\colon A\to B\dsep x\in C\vdash F(x)\in F[C]}{2,1,6}}
  \proofstep{1,2,6}{F(x)\in F[A]\cup F[B]}{%
    \FormulaRefAuto{y\in A\vdash y\in A\cup B}{7}}

  \proofstep{9}{x\in B}{\rA}
  \proofstep{1,3,9}{F(x)\in F[B]}{%
    \FormulaRefAuto{C\subseteq A\dsep F\colon A\to B\dsep x\in C\vdash F(x)\in F[C]}{3,1,9}}
  \proofstep{1,3,9}{F(x)\in F[A]\cup F[B]}{%
    \FormulaRefAuto{y\in B\vdash y\in A\cup B}{10}}

  \proofstep{1,2,3,4}{F(x)\in F[A]\cup F[B]}{\rOE{5,6,8,9,11}}
\end{tabproof}

\FormulaThmDelta[Bild einer Vereinigung]{%
  F\colon M\to N \dsep A\subseteq M \dsep B\subseteq M \vdash F[A\cup B]=F[A]\cup F[B]
}{%
  \DeltaRow{Mengen}{A\dsep B\dsep M\dsep N\dsep x\dsep y}
  \DeltaRow{Funktionen}{F}
}
\begin{tabproofsplitwide}
  \proofpartwideind[Teilmengenrichtung \(F[A\cup B]\subseteq F[A]\cup F[B]\)]{%
    F\colon M\to N \dsep A\subseteq M \dsep B\subseteq M \vdash F[A\cup B]\subseteq F[A]\cup F[B]
  }[%
    F\colon M\to N \dsep A\subseteq M \dsep B\subseteq M \vdash F[A\cup B]\subseteq F[A]\cup F[B]
  ]
    \proofstepwidestar[1]{F\colon M\to N}{\rA}
    \proofstepwidestar[2]{A\subseteq M}{\rA}
    \proofstepwidestar[3]{B\subseteq M}{\rA}
    \proofstepwidestar[4]{y\in F[A\cup B]}{\rA}
    \proofstepwidestar[2,3]{A\cup B\subseteq M}{%
      \FormulaRefAuto{A\subseteq M\dsep B\subseteq M \vdash A\cup B\subseteq M}{2,3}}
    \proofstepwidestar[1,2,3,4]{\exists x\in A\cup B\,y=F(x)}{%
      \FormulaRefAuto{C\subseteq A\dsep F\colon A\to B\dsep y\in F[C]\vdash \exists x\in C\,y=F(x)}{5,1,4}}
    \proofstepwidestar[7]{x\in A\cup B\land y=F(x)}{\rA}
    \proofstepwidestar[7]{x\in A\cup B}{\rAEa{7}}
    \proofstepwidestar[7]{y=F(x)}{\rAEb{7}}
    \proofstepwidestar[1,2,3,8]{F(x)\in F[A]\cup F[B]}{%
      \FormulaRefAuto{F\colon M\to N \dsep A\subseteq M \dsep B\subseteq M \dsep x\in A\cup B \vdash F(x)\in F[A]\cup F[B]}{1,2,3,8}}
    \proofstepwidestar[1,2,3,7]{y\in F[A]\cup F[B]}{\rIE{9,10}}
    \proofstepwidestar[1,2,3,4]{y\in F[A]\cup F[B]}{\rEE{6,7,11}}
    \proofstepwidestar[1,2,3]{\forall y\,\bigl(y\in F[A\cup B]\rightarrow y\in F[A]\cup F[B]\bigr)}{\rUI{\rRI{4,12}}}
    \proofstepwidestar[1,2,3]{F[A\cup B]\subseteq F[A]\cup F[B]}{%
      \FormulaRefAuto{A\subseteq B \coloneq \forall x\,(x\in A \rightarrow x\in B)}{13}}
  \closeproofpartwideind

  \proofpartwideind[Teilmengenrichtung \(F[A]\cup F[B]\subseteq F[A\cup B]\)]{%
    F\colon M\to N \dsep A\subseteq M \dsep B\subseteq M \vdash F[A]\cup F[B]\subseteq F[A\cup B]
  }[%
    F\colon M\to N \dsep A\subseteq M \dsep B\subseteq M \vdash F[A]\cup F[B]\subseteq F[A\cup B]
  ]
    \proofstepwidestar[1]{F\colon M\to N}{\rA}
    \proofstepwidestar[2]{A\subseteq M}{\rA}
    \proofstepwidestar[3]{B\subseteq M}{\rA}
    \proofstepwidestar[2]{A\subseteq A\cup B}{\FormulaRefAuto{A\subseteq A\cup B}{2}}
    \proofstepwidestar[3]{B\subseteq A\cup B}{\FormulaRefAuto{B\subseteq A\cup B}{3}}
    \proofstepwidestar[2,3]{A\cup B\subseteq M}{%
      \FormulaRefAuto{A\subseteq M\dsep B\subseteq M \vdash A\cup B\subseteq M}{2,3}}
    \proofstepwidestar[1,2,3]{F[A]\subseteq F[A\cup B]}{%
      \FormulaRefAuto{F\colon M\to N \dsep A\subseteq B \dsep B\subseteq M \vdash F[A]\subseteq F[B]}{1,4,6}}
    \proofstepwidestar[1,2,3]{F[B]\subseteq F[A\cup B]}{%
      \FormulaRefAuto{F\colon M\to N \dsep A\subseteq B \dsep B\subseteq M \vdash F[A]\subseteq F[B]}{1,5,6}}
    \proofstepwidestar[1,2,3]{F[A]\cup F[B]\subseteq F[A\cup B]}{%
      \FormulaRefAuto{X\subseteq Z\dsep Y\subseteq Z \vdash X\cup Y\subseteq Z}{7,8}}
  \closeproofpartwideind

  \proofpartwide{Zusammenfassung}
    \proofstepwidestar[1]{F\colon M\to N}{\rA}
    \proofstepwidestar[2]{A\subseteq M}{\rA}
    \proofstepwidestar[3]{B\subseteq M}{\rA}
    \proofstepwidestar[1,2,3]{F[A\cup B]\subseteq F[A]\cup F[B]}{%
      \FormulaRefAuto{F\colon M\to N \dsep A\subseteq M \dsep B\subseteq M \vdash F[A\cup B]\subseteq F[A]\cup F[B]}{1,2,3}}
    \proofstepwidestar[1,2,3]{F[A]\cup F[B]\subseteq F[A\cup B]}{%
      \FormulaRefAuto{F\colon M\to N \dsep A\subseteq M \dsep B\subseteq M \vdash F[A]\cup F[B]\subseteq F[A\cup B]}{1,2,3}}
    \proofstepwidestar[1,2,3]{F[A\cup B]=F[A]\cup F[B]}{%
      \FormulaRefAuto{A\subseteq B\dsep B\subseteq A \vdash A=B}{4,5}}
  \closeproofpartwide
\end{tabproofsplitwide}

\FormulaThmDelta[Bild eines Durchschnitts liegt im Durchschnitt der Bilder]{%
  F\colon M\to N \dsep A\subseteq M \dsep B\subseteq M \vdash F[A\cap B]\subseteq F[A]\cap F[B]
}{%
  \DeltaRow{Mengen}{A\dsep B\dsep M\dsep N}
  \DeltaRow{Funktionen}{F}
}
\begin{tabproof}
  \proofstep{1}{F\colon M\to N}{\rA}
  \proofstep{2}{A\subseteq M}{\rA}
  \proofstep{3}{B\subseteq M}{\rA}
  \proofstep{}{A\cap B\subseteq A}{\FormulaRefAuto{A\cap B\subseteq A}}
  \proofstep{}{A\cap B\subseteq B}{\FormulaRefAuto{A\cap B\subseteq B}}
  \proofstep{1,2}{F[A\cap B]\subseteq F[A]}{%
    \FormulaRefAuto{F\colon M\to N \dsep A\subseteq B \dsep B\subseteq M \vdash F[A]\subseteq F[B]}{1,4,2}}
  \proofstep{1,3}{F[A\cap B]\subseteq F[B]}{%
    \FormulaRefAuto{F\colon M\to N \dsep A\subseteq B \dsep B\subseteq M \vdash F[A]\subseteq F[B]}{1,5,3}}
  \proofstep{1,2,3}{F[A\cap B]\subseteq F[A]\cap F[B]}{%
    \FormulaRefAuto{A\subseteq B \dsep A\subseteq C \vdash A\subseteq B\cap C}{6,7}}
\end{tabproof}

\FormulaThmDelta[Das Bild der leeren Menge ist leer]{%
  F\colon A\to B \vdash F[\varnothing]=\varnothing
}{%
  \DeltaRow{Mengen}{A\dsep B\dsep y}
  \DeltaRow{Funktionen}{F}
}
\begin{tabproofsplitwide}
  \proofpartwideind[Teilmengenrichtung \(F[\varnothing]\subseteq \varnothing\)]{%
    F\colon A\to B \vdash F[\varnothing]\subseteq \varnothing
  }[%
    F\colon A\to B \vdash F[\varnothing]\subseteq \varnothing
  ]
    \proofstepwidestar[1]{F\colon A\to B}{\rA}
    \proofstepwidestar[2]{y\in F[\varnothing]}{\rA}
    \proofstepwidestar[]{\varnothing\subseteq A}{\FormulaRefAuto{\varnothing\subseteq A}}
    \proofstepwidestar[1,2]{\exists x\in \varnothing\,y=F(x)}{%
      \FormulaRefAuto{C\subseteq A\dsep F\colon A\to B\dsep y\in F[C]\vdash \exists x\in C\,y=F(x)}{3,1,2}}
    \proofstepwidestar[1,2]{y\in \varnothing}{%
      \FormulaRefAuto{\exists x\in \varnothing\,P(x)\vdash Q}{4}}
    \proofstepwidestar[1]{\forall y\,(y\in F[\varnothing]\rightarrow y\in \varnothing)}{%
      \rUI{\rRI{2,5}}}
    \proofstepwidestar[1]{F[\varnothing]\subseteq \varnothing}{%
      \FormulaRefAuto{A\subseteq B \coloneq \forall x\,(x\in A \rightarrow x\in B)}{6}}
  \closeproofpartwideind

  \proofpartwideind[Teilmengenrichtung \(\varnothing\subseteq F[\varnothing]\)]{%
    F\colon A\to B \vdash \varnothing\subseteq F[\varnothing]
  }[%
    F\colon A\to B \vdash \varnothing\subseteq F[\varnothing]
  ]
    \proofstepwidestar[1]{F\colon A\to B}{\rA}
    \proofstepwidestar[1]{\varnothing\subseteq F[\varnothing]}{\FormulaRefAuto{\varnothing\subseteq A}}
  \closeproofpartwideind

  \proofpartwide{Zusammenfassung}
    \proofstepwidestar[1]{F\colon A\to B}{\rA}
    \proofstepwidestar[1]{F[\varnothing]\subseteq \varnothing}{%
      \FormulaRefAuto{F\colon A\to B \vdash F[\varnothing]\subseteq \varnothing}{1}}
    \proofstepwidestar[1]{\varnothing\subseteq F[\varnothing]}{%
      \FormulaRefAuto{F\colon A\to B \vdash \varnothing\subseteq F[\varnothing]}{1}}
    \proofstepwidestar[1]{F[\varnothing]=\varnothing}{%
      \FormulaRefAuto{A\subseteq B\dsep B\subseteq A \vdash A=B}{2,3}}
  \closeproofpartwide
\end{tabproofsplitwide}

\subsection{Bildmengen im Kontext des Wertebereichs}

\FormulaThmDelta{F\subseteq A\times F[A]}
{
  \DeltaRow{Mengen}{A\dsep B\dsep x\dsep y\dsep z}
  \DeltaPrem{Funktionen}{F\colon A\to B}
}
\begin{tabproof}
  \proofstep{1}{F\colon A\to B}{\rA}
  \proofstep{2}{(x,y)\in F}{\rA}
  \proofstep{2}{x\in A}{\FormulaRefAuto{(x,y)\in F\vdash x\in A}{2}}
  \proofstep{2}{y=F(x)}{\FormulaRefAuto{x\in A\dsep (x,y)\in F\vdash F(x)=y}{3,2}}
  \proofstep{}{F(x)=F(x)}{\rII}
  \proofstep{3}{\exists z\in A\,F(x)=F(z)}{\rEI{\rAI{3,5}}}
  \proofstep{1,3}{F(x)\in F[A]}{%
    \FormulaRefAuto{F\colon A\to B\dsep \exists x\in A\,y=F(x)\vdash y\in F[A]}{1,6}}
  \proofstep{1,2,3}{y\in F[A]}{\rIE{4,7}}
  \proofstep{1,2,3}{(x,y)\in A\times F[A]}{%
    \FormulaRefAuto{a\in A\dsep b\in B\vdash (a,b)\in A\times B}{3,8}}
  \proofstep{1}{F\subseteq A\times F[A]}{%
    \FormulaRefAuto{A \subseteq B \coloneq \forall x\,(x\in A \rightarrow x\in B)}{\rUI{\rRI{2,9}}}}
\end{tabproof}

\FormulaThmDelta[Das Bild einer Teilmenge liegt im Zielbereich]{%
  C\subseteq A \dsep F\colon A\to B \vdash F[C]\subseteq B
}{%
  \DeltaRow{Mengen}{A\dsep B\dsep C\dsep x\dsep y}
  \DeltaRow{Funktionen}{F}
}
\begin{tabproof}
  \proofstep{1}{C\subseteq A}{\rA}
  \proofstep{2}{F\colon A\to B}{\rA}
  \proofstep{3}{y\in F[C]}{\rA}
  \proofstep{1,2,3}{\exists x\in C\,y=F(x)}{%
    \FormulaRefAuto{C\subseteq A\dsep F\colon A\to B\dsep y\in F[C]\vdash \exists x\in C\,y=F(x)}{1,2,3}}
  \proofstep{5}{x\in C\land y=F(x)}{\rA}
  \proofstep{5}{x\in C}{\rAEa{5}}
  \proofstep{5}{y=F(x)}{\rAEb{5}}
  \proofstep{1,6}{x\in A}{%
    \FormulaRefAuto{A\subseteq B\dsep x\in A\vdash x\in B}{1,6}}
  \proofstep{1,2,5}{F(x)\in B}{%
    \FormulaRefAuto{F\colon A\to B\dsep x\in A\vdash F(x)\in B}{2,8}}
  \proofstep{1,2,5}{y\in B}{\rIE{7,9}}
  \proofstep{1,2,3}{y\in B}{\rEE{4,5,10}}
  \proofstep{1,2}{\forall y\,(y\in F[C]\rightarrow y\in B)}{\rUI{\rRI{3,11}}}
  \proofstep{1,2}{F[C]\subseteq B}{%
    \FormulaRefAuto{A\subseteq B \coloneq \forall x\,(x\in A \rightarrow x\in B)}{12}}
\end{tabproof}

\FormulaThmDelta[Bilder liegen im Zielbereich]{%
  F\colon A\to B \vdash F[A]\subseteq B
}{
  \DeltaRow{Mengen}{A\dsep B}
  \DeltaRow{Funktionen}{F}
}
\begin{tabproof}
  \proofstep{1}{F\colon A\to B}{\rA}
  \proofstep{}{A\subseteq A}{\FormulaRefAuto{A\subseteq A}}
  \proofstep{1}{F[A]\subseteq B}{%
    \FormulaRefAuto{C\subseteq A \dsep F\colon A\to B \vdash F[C]\subseteq B}{2,1}}
\end{tabproof}

\FormulaThmDelta[Bild einer bijektiven Funktion ist die Zielmenge]{%
  F\colon A\bij B \vdash F[A]=B
}{%
  \DeltaRow{Mengen}{A\dsep B\dsep x\dsep y}
  \DeltaRow{Funktionen}{F}
}
\begin{tabproofsplitwide}
  \proofpartwideind[Teilmengenrichtung \(F[A]\subseteq B\)]{%
    F\colon A\bij B \vdash F[A]\subseteq B%
  }[%
    F\colon A\bij B \vdash F[A]\subseteq B%
  ]
    \proofstepwidestar[1]{F\colon A\bij B}{\rA}
    \proofstepwidestar[1]{F\colon A\to B}{\FormulaRefAuto{Bijektive Funktion}{1}}
    \proofstepwidestar[1]{F[A]\subseteq B}{%
      \FormulaRefAuto{F\colon A\to B \vdash F[A]\subseteq B}{2}}
  \closeproofpartwideind

  \proofpartwideind[Teilmengenrichtung \(B\subseteq F[A]\)]{%
    F\colon A\bij B \vdash B\subseteq F[A]%
  }[%
    F\colon A\bij B \vdash B\subseteq F[A]%
  ]
    \proofstepwidestar[1]{F\colon A\bij B}{\rA}
    \proofstepwidestar[1]{F\colon A\sur B}{\FormulaRefAuto{Bijektive Funktion}{1}}
    \proofstepwidestar[1]{F\colon A\to B}{\FormulaRefAuto{Bijektive Funktion}{1}}

    \proofstepwidestar[4]{y\in B}{\rA}
    \proofstepwidestar[1,4]{\exists x\in A\,F(x)=y}{%
      \FormulaRefAuto{y\in B\vdash \exists x\in A\; F(x)=y}{4}}

    \proofstepwidestar[6]{x\in A\land F(x)=y}{\rA}
    \proofstepwidestar[6]{x\in A}{\rAEa{6}}
    \proofstepwidestar[6]{F(x)=y}{\rAEb{6}}
    \proofstepwidestar[6]{y=F(x)}{\FormulaRefAuto{a=b \vdash b=a}{8}}
    \proofstepwidestar[6]{x\in A\land y=F(x)}{\rAI{7,9}}
    \proofstepwidestar[6]{\exists x\in A\,y=F(x)}{\rEI{10}}
    \proofstepwidestar[1,4,6]{y\in F[A]}{%
      \FormulaRefAuto{F\colon A\to B\dsep \exists x\in A\,y=F(x)\vdash y\in F[A]}{3,11}}
    \proofstepwidestar[1,4]{y\in F[A]}{\rEE{5,6,12}}

    \proofstepwidestar[1]{\forall y\,\bigl(y\in B\rightarrow y\in F[A]\bigr)}{\rUI{\rRI{4,13}}}
    \proofstepwidestar[1]{B\subseteq F[A]}{%
      \FormulaRefAuto{A\subseteq B \coloneq \forall x\,(x\in A \rightarrow x\in B)}{14}}
  \closeproofpartwideind

  \proofpartwide{Zusammenfassung}
    \proofstepwidestar[1]{F\colon A\bij B}{\rA}
    \proofstepwidestar[1]{F[A]\subseteq B}{%
      \FormulaRefAuto{F\colon A\bij B \vdash F[A]\subseteq B}{1}}
    \proofstepwidestar[1]{B\subseteq F[A]}{%
      \FormulaRefAuto{F\colon A\bij B \vdash B\subseteq F[A]}{1}}
    \proofstepwidestar[1]{F[A]=B}{%
      \FormulaRefAuto{A\subseteq B\dsep B\subseteq A \vdash A=B}{2,3}}
  \closeproofpartwide
\end{tabproofsplitwide}

\section{Urbild}
%%begin novalidate
\FormulaDefDelta[Urbildmenge]{%
C\subseteq B \vdash F^{-1}[C] \coloneqq \{\,x\in A \mid F(x)\in C\,\}%
}{
  \DeltaRow{Mengen}{A\dsep B\dsep C\dsep x}
  \DeltaPrem{Funktionen}{F\colon A\to B}
}
%%end novalidate

\FormulaThmDelta{%
C\subseteq B \vdash \bigl(x\in F^{-1}[C] \leftrightarrow (x\in A \land F(x)\in C)\bigr)
}{
  \DeltaRow{Mengen}{A\dsep B\dsep C\dsep x}
  \DeltaPrem{Funktionen}{F\colon A\to B}
}
\begin{tabproof}
  \proofstep{1}{C\subseteq B}{\rA}

  \proofstep{1}{F^{-1}[C]=\{\,u\in A \mid F(u)\in C\,\}}{%
    \FormulaRefAuto{C\subseteq B \vdash F^{-1}[C] \coloneqq \{\,u\in A \mid F(u)\in C\,\}}{1}}

  \proofstep{}{x\in \{\,u\in A \mid F(u)\in C\,\}\leftrightarrow (x\in A \land F(x)\in C)}{%
    \FormulaRefAuto{x\in \{u\in A \mid P(u)\}\eqvdash x\in A \land P(x)}}

  \proofstep{1}{x\in F^{-1}[C]\leftrightarrow (x\in A \land F(x)\in C)}{\rIE{2,3}}
\end{tabproof}

\FormulaThmDelta[Das Urbild der leeren Menge ist leer]{%
  F\colon A\to B \vdash F^{-1}[\varnothing]=\varnothing
}{%
  \DeltaRow{Mengen}{A\dsep B\dsep x}
  \DeltaRow{Funktionen}{F}
}
\begin{tabproof}
  \proofstep{}{ \varnothing\subseteq B }{\FormulaRefAuto{\varnothing\subseteq A}}
  \proofstep{2}{x\in F^{-1}[\varnothing]}{\rA}
  \proofstep{1}{x\in F^{-1}[\varnothing]\leftrightarrow (x\in A \land F(x)\in \varnothing)}{%
    \FormulaRefAuto{C\subseteq B \vdash \bigl(x\in F^{-1}[C] \leftrightarrow (x\in A \land F(x)\in C)\bigr)}{1}}
  \proofstep{1,2}{x\in A \land F(x)\in \varnothing}{%
    \FormulaRefAuto{P \leftrightarrow Q, P \vdash Q}{3,2}}
  \proofstep{1,2}{F(x)\in \varnothing}{\rAEb{4}}
  \proofstep{}{F(x)\notin \varnothing}{\FormulaRefAuto{x \not\in \varnothing}}
  \proofstep{1,2}{x\in \varnothing}{\FormulaRefAuto{P\dsep \neg P \vdash Q}{5,6}}
  \proofstep{}{ \forall x\,(x\in F^{-1}[\varnothing]\rightarrow x\in \varnothing)}{\rUI{\rRI{2,7}}}
  \proofstep{}{F^{-1}[\varnothing]\subseteq \varnothing}{%
    \FormulaRefAuto{A\subseteq B \coloneq \forall x\,(x\in A \rightarrow x\in B)}{8}}
  \proofstep{}{F^{-1}[\varnothing]=\varnothing}{%
    \FormulaRefAuto{A\subseteq\varnothing\vdash A=\varnothing}{9}}
\end{tabproof}

\FormulaThmDelta[Urbild einer Einermenge]{%
  y\in B \vdash F^{-1}[\{y\}] = \{\,x\in A \mid F(x)=y\,\}
}{%
  \DeltaRow{Mengen}{A\dsep B\dsep x\dsep y}
  \DeltaRow{Funktionen}{F}
}
\begin{tabproofsplitwide}
  \proofpartwideind[Teilmengenrichtung \(F^{-1}[\{y\}]\subseteq \{\,x\in A \mid F(x)=y\,\}\)]{%
    y\in B \vdash F^{-1}[\{y\}]\subseteq \{\,x\in A \mid F(x)=y\,\}
  }[%
    y\in B \vdash F^{-1}[\{y\}]\subseteq \{\,x\in A \mid F(x)=y\,\}
  ]
    \proofstepwidestar[1]{y\in B}{\rA}
    \proofstepwidestar[1]{\{y\}\subseteq B}{\FormulaRefAuto{x\in A \vdash \{x\}\subseteq A}{1}}
    \proofstepwidestar[3]{x\in F^{-1}[\{y\}]}{\rA}
    \proofstepwidestar[2]{x\in F^{-1}[\{y\}]\leftrightarrow (x\in A \land F(x)\in \{y\})}{%
      \FormulaRefAuto{C\subseteq B \vdash \bigl(x\in F^{-1}[C] \leftrightarrow (x\in A \land F(x)\in C)\bigr)}{2}}
    \proofstepwidestar[2,3]{x\in A \land F(x)\in \{y\}}{%
      \FormulaRefAuto{P \leftrightarrow Q, P \vdash Q}{4,3}}
    \proofstepwidestar[2,3]{x\in A}{\rAEa{5}}
    \proofstepwidestar[2,3]{F(x)\in \{y\}}{\rAEb{5}}
    \proofstepwidestar[2,3]{F(x)=y}{\FormulaRefAuto{x\in \{a\}\eqvdash x=a}{7}}
    \proofstepwidestar[2,3]{x\in A \land F(x)=y}{\rAI{6,8}}
    \proofstepwidestar[2,3]{x\in \{\,x\in A \mid F(x)=y\,\}}{%
      \FormulaRefAuto{x\in \{u\in A \mid P(u)\}\eqvdash x\in A \land P(x)}{9}}
    \proofstepwidestar[1]{\forall x\,\bigl(x\in F^{-1}[\{y\}]\rightarrow x\in \{\,x\in A \mid F(x)=y\,\}\bigr)}{\rUI{\rRI{3,10}}}
    \proofstepwidestar[1]{F^{-1}[\{y\}]\subseteq \{\,x\in A \mid F(x)=y\,\}}{%
      \FormulaRefAuto{A\subseteq B \coloneq \forall x\,(x\in A \rightarrow x\in B)}{11}}
  \closeproofpartwideind

  \proofpartwideind[Teilmengenrichtung \(\{\,x\in A \mid F(x)=y\,\}\subseteq F^{-1}[\{y\}]\)]{%
    y\in B \vdash \{\,x\in A \mid F(x)=y\,\}\subseteq F^{-1}[\{y\}]
  }[%
    y\in B \vdash \{\,x\in A \mid F(x)=y\,\}\subseteq F^{-1}[\{y\}]
  ]
    \proofstepwidestar[1]{y\in B}{\rA}
    \proofstepwidestar[1]{\{y\}\subseteq B}{\FormulaRefAuto{x\in A \vdash \{x\}\subseteq A}{1}}
    \proofstepwidestar[3]{x\in \{\,x\in A \mid F(x)=y\,\}}{\rA}
    \proofstepwidestar[3]{x\in A \land F(x)=y}{%
      \FormulaRefAuto{x\in \{u\in A \mid P(u)\}\eqvdash x\in A \land P(x)}{3}}
    \proofstepwidestar[3]{x\in A}{\rAEa{4}}
    \proofstepwidestar[3]{F(x)=y}{\rAEb{4}}
    \proofstepwidestar[3]{F(x)\in \{y\}}{\FormulaRefAuto{x\in \{a\}\eqvdash x=a}{6}}
    \proofstepwidestar[3]{x\in A \land F(x)\in \{y\}}{\rAI{5,7}}
    \proofstepwidestar[2]{x\in F^{-1}[\{y\}]\leftrightarrow (x\in A \land F(x)\in \{y\})}{%
      \FormulaRefAuto{C\subseteq B \vdash \bigl(x\in F^{-1}[C] \leftrightarrow (x\in A \land F(x)\in C)\bigr)}{2}}
    \proofstepwidestar[2,3]{x\in F^{-1}[\{y\}]}{\FormulaRefAuto{P \leftrightarrow Q, Q \vdash P}{9,8}}
    \proofstepwidestar[1]{\forall x\,\bigl(x\in \{\,x\in A \mid F(x)=y\,\}\rightarrow x\in F^{-1}[\{y\}]\bigr)}{\rUI{\rRI{3,10}}}
    \proofstepwidestar[1]{\{\,x\in A \mid F(x)=y\,\}\subseteq F^{-1}[\{y\}]}{%
      \FormulaRefAuto{A\subseteq B \coloneq \forall x\,(x\in A \rightarrow x\in B)}{11}}
  \closeproofpartwideind

  \proofpartwide{Zusammenfassung}
    \proofstepwidestar[1]{y\in B}{\rA}
    \proofstepwidestar[1]{F^{-1}[\{y\}]\subseteq \{\,x\in A \mid F(x)=y\,\}}{%
      \FormulaRefAuto{y\in B \vdash F^{-1}[\{y\}]\subseteq \{\,x\in A \mid F(x)=y\,\}}{1}}
    \proofstepwidestar[1]{\{\,x\in A \mid F(x)=y\,\}\subseteq F^{-1}[\{y\}]}{%
      \FormulaRefAuto{y\in B \vdash \{\,x\in A \mid F(x)=y\,\}\subseteq F^{-1}[\{y\}]}{1}}
    \proofstepwidestar[1]{F^{-1}[\{y\}] = \{\,x\in A \mid F(x)=y\,\}}{%
      \FormulaRefAuto{A\subseteq B\dsep B\subseteq A \vdash A=B}{2,3}}
  \closeproofpartwide
\end{tabproofsplitwide}

\subsection{Urbilder liegen im Definitionsbereich}


\FormulaThmDelta[Urbilder liegen im Definitionsbereich]{%
  F\colon A\to B\dsep C\subseteq B \vdash F^{-1}[C]\subseteq A
}{
  \DeltaRow{Mengen}{A\dsep B\dsep C\dsep x}
  \DeltaRow{Funktionen}{F}
}
\begin{tabproof}
  \proofstep{1}{F\colon A\to B}{\rA}
  \proofstep{2}{C\subseteq B}{\rA}
  \proofstep{3}{x\in F^{-1}[C]}{\rA}
  \proofstep{1,2}{x\in F^{-1}[C]\leftrightarrow (x\in A \land F(x)\in C)}{%
    \FormulaRefAuto{C\subseteq B \vdash \bigl(x\in F^{-1}[C] \leftrightarrow (x\in A \land F(x)\in C)\bigr)}{2}}
  \proofstep{1,2,3}{x\in A \land F(x)\in C}{%
    \FormulaRefAuto{P \leftrightarrow Q, P \vdash Q}{4,3}}
  \proofstep{1,2,3}{x\in A}{\rAEa{5}}
  \proofstep{1,2}{\forall x\,\bigl(x\in F^{-1}[C]\rightarrow x\in A\bigr)}{\rUI{\rRI{3,6}}}
  \proofstep{1,2}{F^{-1}[C]\subseteq A}{%
    \FormulaRefAuto{A\subseteq B \coloneq \forall x\,(x\in A \rightarrow x\in B)}{7}}
\end{tabproof}

\subsection{Rechenregeln für Urbilder}

\FormulaThmDelta[Monotonie des Urbilds]{%
  F\colon A\to B \dsep C\subseteq D \dsep D\subseteq B \vdash F^{-1}[C]\subseteq F^{-1}[D]
}{%
  \DeltaRow{Mengen}{A\dsep B\dsep C\dsep D\dsep x}
  \DeltaRow{Funktionen}{F}
}
\begin{tabproof}
  \proofstep{1}{C\subseteq D}{\rA}
  \proofstep{2}{D\subseteq B}{\rA}
  \proofstep{1,2}{C\subseteq B}{\FormulaRefAuto{A\subseteq B\dsep B\subseteq C \vdash A\subseteq C}{1,2}}

  \proofstep{4}{x\in F^{-1}[C]}{\rA}
  \proofstep{1,2}{x\in F^{-1}[C]\leftrightarrow (x\in A \land F(x)\in C)}{%
    \FormulaRefAuto{C\subseteq B \vdash \bigl(x\in F^{-1}[C] \leftrightarrow (x\in A \land F(x)\in C)\bigr)}{3}}
  \proofstep{1,2,4}{x\in A \land F(x)\in C}{%
    \FormulaRefAuto{P \leftrightarrow Q, P \vdash Q}{5,4}}
  \proofstep{1,2,4}{x\in A}{\rAEa{6}}
  \proofstep{1,2,4}{F(x)\in C}{\rAEb{6}}
  \proofstep{1,2,4}{F(x)\in D}{\FormulaRefAuto{A\subseteq B\dsep x\in A\vdash x\in B}{1,8}}
  \proofstep{1,2,4}{x\in A \land F(x)\in D}{\rAI{7,9}}
  \proofstep{2}{x\in F^{-1}[D]\leftrightarrow (x\in A \land F(x)\in D)}{%
    \FormulaRefAuto{C\subseteq B \vdash \bigl(x\in F^{-1}[C] \leftrightarrow (x\in A \land F(x)\in C)\bigr)}{2}}
  \proofstep{1,2,4}{x\in F^{-1}[D]}{\FormulaRefAuto{P \leftrightarrow Q, Q \vdash P}{11,10}}

  \proofstep{1,2}{\forall x\,\bigl(x\in F^{-1}[C]\rightarrow x\in F^{-1}[D]\bigr)}{\rUI{\rRI{4,12}}}
  \proofstep{1,2}{F^{-1}[C]\subseteq F^{-1}[D]}{%
    \FormulaRefAuto{A\subseteq B \coloneq \forall x\,(x\in A \rightarrow x\in B)}{13}}
\end{tabproof}

\FormulaThmDelta[Element eines Urbilds einer Vereinigung]{%
  F\colon A\to B \dsep C\subseteq B \dsep D\subseteq B \dsep x\in F^{-1}[C\cup D] \vdash x\in F^{-1}[C]\cup F^{-1}[D]
}{%
  \DeltaRow{Mengen}{A\dsep B\dsep C\dsep D\dsep x}
  \DeltaRow{Funktionen}{F}
}
\begin{tabproof}
  \proofstep{1}{F\colon A\to B}{\rA}
  \proofstep{2}{C\subseteq B}{\rA}
  \proofstep{3}{D\subseteq B}{\rA}
  \proofstep{4}{x\in F^{-1}[C\cup D]}{\rA}
  \proofstep{2,3}{C\cup D\subseteq B}{\FormulaRefAuto{A\subseteq M\dsep B\subseteq M \vdash A\cup B\subseteq M}{2,3}}
  \proofstep{5}{x\in F^{-1}[C\cup D]\leftrightarrow (x\in A \land F(x)\in C\cup D)}{%
    \FormulaRefAuto{C\subseteq B \vdash \bigl(x\in F^{-1}[C] \leftrightarrow (x\in A \land F(x)\in C)\bigr)}{5}}
  \proofstep{4,5}{x\in A \land F(x)\in C\cup D}{%
    \FormulaRefAuto{P \leftrightarrow Q, P \vdash Q}{6,4}}
  \proofstep{4,5}{x\in A}{\rAEa{7}}
  \proofstep{4,5}{F(x)\in C\cup D}{\rAEb{7}}
  \proofstep{4,5}{F(x)\in C\lor F(x)\in D}{%
    \FormulaRefAuto{x\in A\cup B\eqvdash x\in A\lor x\in B}{9}}

  \proofstep{11}{F(x)\in C}{\rA}
  \proofstep{4,5,11}{x\in A \land F(x)\in C}{\rAI{8,11}}
  \proofstep{2}{x\in F^{-1}[C]\leftrightarrow (x\in A \land F(x)\in C)}{%
    \FormulaRefAuto{C\subseteq B \vdash \bigl(x\in F^{-1}[C] \leftrightarrow (x\in A \land F(x)\in C)\bigr)}{2}}
  \proofstep{2,4,5,11}{x\in F^{-1}[C]}{\FormulaRefAuto{P \leftrightarrow Q, Q \vdash P}{13,12}}
  \proofstep{2,4,5,11}{x\in F^{-1}[C]\cup F^{-1}[D]}{%
    \FormulaRefAuto{y\in A\vdash y\in A\cup B}{14}}

  \proofstep{16}{F(x)\in D}{\rA}
  \proofstep{4,5,16}{x\in A \land F(x)\in D}{\rAI{8,16}}
  \proofstep{3}{x\in F^{-1}[D]\leftrightarrow (x\in A \land F(x)\in D)}{%
    \FormulaRefAuto{C\subseteq B \vdash \bigl(x\in F^{-1}[C] \leftrightarrow (x\in A \land F(x)\in C)\bigr)}{3}}
  \proofstep{3,4,5,16}{x\in F^{-1}[D]}{\FormulaRefAuto{P \leftrightarrow Q, Q \vdash P}{18,17}}
  \proofstep{3,4,5,16}{x\in F^{-1}[C]\cup F^{-1}[D]}{%
    \FormulaRefAuto{y\in B\vdash y\in A\cup B}{19}}

  \proofstep{1,2,3,4}{x\in F^{-1}[C]\cup F^{-1}[D]}{\rOE{10,11,15,16,20}}
\end{tabproof}

\FormulaThmDelta[Urbild einer Vereinigung]{%
  F\colon A\to B \dsep C\subseteq B \dsep D\subseteq B \vdash F^{-1}[C\cup D]=F^{-1}[C]\cup F^{-1}[D]
}{%
  \DeltaRow{Mengen}{A\dsep B\dsep C\dsep D\dsep x}
  \DeltaRow{Funktionen}{F}
}
\begin{tabproofsplitwide}

  \proofpartwideind[Teilmengenrichtung \(F^{-1}[C\cup D]\subseteq F^{-1}[C]\cup F^{-1}[D]\)]{%
    F\colon A\to B \dsep C\subseteq B \dsep D\subseteq B \vdash F^{-1}[C\cup D]\subseteq F^{-1}[C]\cup F^{-1}[D]
  }[%
    F\colon A\to B \dsep C\subseteq B \dsep D\subseteq B \vdash F^{-1}[C\cup D]\subseteq F^{-1}[C]\cup F^{-1}[D]
  ]
    \proofstepwidestar[1]{F\colon A\to B}{\rA}
    \proofstepwidestar[2]{C\subseteq B}{\rA}
    \proofstepwidestar[3]{D\subseteq B}{\rA}
    \proofstepwidestar[4]{x\in F^{-1}[C\cup D]}{\rA}
    \proofstepwidestar[1,2,3,4]{x\in F^{-1}[C]\cup F^{-1}[D]}{%
      \FormulaRefAuto{F\colon A\to B \dsep C\subseteq B \dsep D\subseteq B \dsep x\in F^{-1}[C\cup D] \vdash x\in F^{-1}[C]\cup F^{-1}[D]}{1,2,3,4}}
    \proofstepwidestar[1,2,3]{\forall x\,\bigl(x\in F^{-1}[C\cup D]\rightarrow x\in F^{-1}[C]\cup F^{-1}[D]\bigr)}{\rUI{\rRI{4,5}}}
    \proofstepwidestar[1,2,3]{F^{-1}[C\cup D]\subseteq F^{-1}[C]\cup F^{-1}[D]}{%
      \FormulaRefAuto{A\subseteq B \coloneq \forall x\,(x\in A \rightarrow x\in B)}{6}}
  \closeproofpartwideind

\proofpartwideind[Teilmengenrichtung \(F^{-1}[C]\cup F^{-1}[D]\subseteq F^{-1}[C\cup D]\)]{%
  F\colon A\to B \dsep C\subseteq B \dsep D\subseteq B \vdash F^{-1}[C]\cup F^{-1}[D]\subseteq F^{-1}[C\cup D]
}[%
  F\colon A\to B \dsep C\subseteq B \dsep D\subseteq B \vdash F^{-1}[C]\cup F^{-1}[D]\subseteq F^{-1}[C\cup D]
]
  \proofstepwidestar[1]{F\colon A\to B}{\rA}
  \proofstepwidestar[2]{C\subseteq B}{\rA}
  \proofstepwidestar[3]{D\subseteq B}{\rA}
  \proofstepwidestar[2]{C\subseteq C\cup D}{\FormulaRefAuto{A\subseteq A\cup B}{2}}
  \proofstepwidestar[3]{D\subseteq C\cup D}{\FormulaRefAuto{B\subseteq A\cup B}{3}}
  \proofstepwidestar[2,3]{C\cup D\subseteq B}{%
    \FormulaRefAuto{A\subseteq M\dsep B\subseteq M \vdash A\cup B\subseteq M}{2,3}}
  \proofstepwidestar[1,2,3]{F^{-1}[C]\subseteq F^{-1}[C\cup D]}{%
    \FormulaRefAuto{F\colon A\to B \dsep C\subseteq D \dsep D\subseteq B \vdash F^{-1}[C]\subseteq F^{-1}[D]}{1,4,6}}
  \proofstepwidestar[1,2,3]{F^{-1}[D]\subseteq F^{-1}[C\cup D]}{%
    \FormulaRefAuto{F\colon A\to B \dsep C\subseteq D \dsep D\subseteq B \vdash F^{-1}[C]\subseteq F^{-1}[D]}{1,5,6}}
  \proofstepwidestar[1,2,3]{F^{-1}[C]\cup F^{-1}[D]\subseteq F^{-1}[C\cup D]}{%
    \FormulaRefAuto{X\subseteq Z\dsep Y\subseteq Z \vdash X\cup Y\subseteq Z}{7,8}}
\closeproofpartwideind

  \proofpartwide{Zusammenfassung}
    \proofstepwidestar[1]{F\colon A\to B}{\rA}
    \proofstepwidestar[2]{C\subseteq B}{\rA}
    \proofstepwidestar[3]{D\subseteq B}{\rA}
    \proofstepwidestar[1,2,3]{F^{-1}[C\cup D]\subseteq F^{-1}[C]\cup F^{-1}[D]}{%
      \FormulaRefAuto{F\colon A\to B \dsep C\subseteq B \dsep D\subseteq B \vdash F^{-1}[C\cup D]\subseteq F^{-1}[C]\cup F^{-1}[D]}{1,2,3}}
    \proofstepwidestar[1,2,3]{F^{-1}[C]\cup F^{-1}[D]\subseteq F^{-1}[C\cup D]}{%
      \FormulaRefAuto{F\colon A\to B \dsep C\subseteq B \dsep D\subseteq B \vdash F^{-1}[C]\cup F^{-1}[D]\subseteq F^{-1}[C\cup D]}{1,2,3}}
    \proofstepwidestar[1,2,3]{F^{-1}[C\cup D]=F^{-1}[C]\cup F^{-1}[D]}{%
      \FormulaRefAuto{A\subseteq B\dsep B\subseteq A \vdash A=B}{4,5}}
  \closeproofpartwide

\end{tabproofsplitwide}

\FormulaThmDelta[Urbild eines Durchschnitts liegt im Durchschnitt der Urbilder]{%
  F\colon A\to B \dsep C\subseteq B \dsep D\subseteq B \vdash F^{-1}[C\cap D]\subseteq F^{-1}[C]\cap F^{-1}[D]
}{%
  \DeltaRow{Mengen}{A\dsep B\dsep C\dsep D\dsep x}
  \DeltaRow{Funktionen}{F}
}
\begin{tabproof}
  \proofstep{1}{F\colon A\to B}{\rA}
  \proofstep{2}{C\subseteq B}{\rA}
  \proofstep{3}{D\subseteq B}{\rA}
  \proofstep{}{C\cap D\subseteq C}{\FormulaRefAuto{A\cap B\subseteq A}}
  \proofstep{}{C\cap D\subseteq D}{\FormulaRefAuto{A\cap B\subseteq B}}
  \proofstep{1,2}{F^{-1}[C\cap D]\subseteq F^{-1}[C]}{%
    \FormulaRefAuto{F\colon A\to B \dsep C\subseteq D \dsep D\subseteq B \vdash F^{-1}[C]\subseteq F^{-1}[D]}{1,4,2}}
  \proofstep{1,3}{F^{-1}[C\cap D]\subseteq F^{-1}[D]}{%
    \FormulaRefAuto{F\colon A\to B \dsep C\subseteq D \dsep D\subseteq B \vdash F^{-1}[C]\subseteq F^{-1}[D]}{1,5,3}}

  \proofstep{8}{x\in F^{-1}[C\cap D]}{\rA}
  \proofstep{6,8}{x\in F^{-1}[C]}{\FormulaRefAuto{A\subseteq B\dsep x\in A\vdash x\in B}{6,8}}
  \proofstep{7,8}{x\in F^{-1}[D]}{\FormulaRefAuto{A\subseteq B\dsep x\in A\vdash x\in B}{7,8}}
  \proofstep{6,7,8}{x\in F^{-1}[C]\land x\in F^{-1}[D]}{\rAI{9,10}}
  \proofstep{6,7,8}{x\in F^{-1}[C]\cap F^{-1}[D]}{%
    \FormulaRefAuto{x\in A\cap B\eqvdash x\in A\land x\in B}{11}}
  \proofstep{1,2,3}{\forall x\,\bigl(x\in F^{-1}[C\cap D]\rightarrow x\in F^{-1}[C]\cap F^{-1}[D]\bigr)}{\rUI{\rRI{8,12}}}
  \proofstep{1,2,3}{F^{-1}[C\cap D]\subseteq F^{-1}[C]\cap F^{-1}[D]}{%
    \FormulaRefAuto{A\subseteq B \coloneq \forall x\,(x\in A \rightarrow x\in B)}{13}}
\end{tabproof}

\FormulaThmDelta[Gemeinsame Urbildwerte liegen im Urbild des Durchschnitts]{%
  F\colon A\to B \dsep C\subseteq B \dsep D\subseteq B \dsep x\in F^{-1}[C] \dsep x\in F^{-1}[D] \vdash x\in F^{-1}[C\cap D]
}{%
  \DeltaRow{Mengen}{A\dsep B\dsep C\dsep D\dsep x}
  \DeltaRow{Funktionen}{F}
}
\begin{tabproof}
  \proofstep{1}{C\subseteq B}{\rA}
  \proofstep{2}{D\subseteq B}{\rA}
  \proofstep{3}{x\in F^{-1}[C]}{\rA}
  \proofstep{4}{x\in F^{-1}[D]}{\rA}

  \proofstep{1}{x\in F^{-1}[C]\leftrightarrow (x\in A \land F(x)\in C)}{%
    \FormulaRefAuto{C\subseteq B \vdash \bigl(x\in F^{-1}[C] \leftrightarrow (x\in A \land F(x)\in C)\bigr)}{1}}
  \proofstep{1,3}{x\in A \land F(x)\in C}{%
    \FormulaRefAuto{P \leftrightarrow Q, P \vdash Q}{5,3}}
  \proofstep{1,3}{x\in A}{\rAEa{6}}
  \proofstep{1,3}{F(x)\in C}{\rAEb{6}}

  \proofstep{2}{x\in F^{-1}[D]\leftrightarrow (x\in A \land F(x)\in D)}{%
    \FormulaRefAuto{C\subseteq B \vdash \bigl(x\in F^{-1}[C] \leftrightarrow (x\in A \land F(x)\in C)\bigr)}{2}}
  \proofstep{2,4}{x\in A \land F(x)\in D}{%
    \FormulaRefAuto{P \leftrightarrow Q, P \vdash Q}{9,4}}
  \proofstep{2,4}{F(x)\in D}{\rAEb{10}}

  \proofstep{1,2,3,4}{F(x)\in C \land F(x)\in D}{\rAI{8,11}}
  \proofstep{1,2,3,4}{F(x)\in C\cap D}{%
    \FormulaRefAuto{x\in A\cap B\eqvdash x\in A\land x\in B}{12}}
  \proofstep{2}{C\cap D\subseteq B}{%
    \FormulaRefAuto{B\subseteq M \vdash A\cap B\subseteq M}{2}}
  \proofstep{1,2,3,4}{x\in A \land F(x)\in C\cap D}{\rAI{7,13}}
  \proofstep{2}{x\in F^{-1}[C\cap D]\leftrightarrow (x\in A \land F(x)\in C\cap D)}{%
    \FormulaRefAuto{C\subseteq B \vdash \bigl(x\in F^{-1}[C] \leftrightarrow (x\in A \land F(x)\in C)\bigr)}{14}}
  \proofstep{1,2,3,4}{x\in F^{-1}[C\cap D]}{\FormulaRefAuto{P \leftrightarrow Q, Q \vdash P}{15,14}}
\end{tabproof}

\FormulaThmDelta[Urbild eines Durchschnitts]{%
  F\colon A\to B \dsep C\subseteq B \dsep D\subseteq B \vdash F^{-1}[C\cap D]=F^{-1}[C]\cap F^{-1}[D]
}{%
  \DeltaRow{Mengen}{A\dsep B\dsep C\dsep D\dsep x}
  \DeltaRow{Funktionen}{F}
}
\begin{tabproofsplitwide}

  \proofpartwideind[Teilmengenrichtung \(F^{-1}[C]\cap F^{-1}[D]\subseteq F^{-1}[C\cap D]\)]{%
    F\colon A\to B \dsep C\subseteq B \dsep D\subseteq B \vdash F^{-1}[C]\cap F^{-1}[D]\subseteq F^{-1}[C\cap D]
  }[%
    F\colon A\to B \dsep C\subseteq B \dsep D\subseteq B \vdash F^{-1}[C]\cap F^{-1}[D]\subseteq F^{-1}[C\cap D]
  ]
    \proofstepwidestar[1]{F\colon A\to B}{\rA}
    \proofstepwidestar[2]{C\subseteq B}{\rA}
    \proofstepwidestar[3]{D\subseteq B}{\rA}
    \proofstepwidestar[4]{x\in F^{-1}[C]\cap F^{-1}[D]}{\rA}
    \proofstepwidestar[4]{x\in F^{-1}[C]\land x\in F^{-1}[D]}{%
      \FormulaRefAuto{x\in A\cap B\eqvdash x\in A\land x\in B}{4}}
    \proofstepwidestar[4]{x\in F^{-1}[C]}{\rAEa{5}}
    \proofstepwidestar[4]{x\in F^{-1}[D]}{\rAEb{5}}
    \proofstepwidestar[1,2,3,6,7]{x\in F^{-1}[C\cap D]}{%
      \FormulaRefAuto{F\colon A\to B \dsep C\subseteq B \dsep D\subseteq B \dsep x\in F^{-1}[C] \dsep x\in F^{-1}[D] \vdash x\in F^{-1}[C\cap D]}{1,2,3,6,7}}
    \proofstepwidestar[1,2,3]{\forall x\,\bigl(x\in F^{-1}[C]\cap F^{-1}[D]\rightarrow x\in F^{-1}[C\cap D]\bigr)}{\rUI{\rRI{4,8}}}
    \proofstepwidestar[1,2,3]{F^{-1}[C]\cap F^{-1}[D]\subseteq F^{-1}[C\cap D]}{%
      \FormulaRefAuto{A\subseteq B \coloneq \forall x\,(x\in A \rightarrow x\in B)}{9}}
  \closeproofpartwideind

  \proofpartwide{Zusammenfassung}
    \proofstepwidestar[1]{F\colon A\to B}{\rA}
    \proofstepwidestar[2]{C\subseteq B}{\rA}
    \proofstepwidestar[3]{D\subseteq B}{\rA}
    \proofstepwidestar[1,2,3]{F^{-1}[C\cap D]\subseteq F^{-1}[C]\cap F^{-1}[D]}{%
      \FormulaRefAuto{F\colon A\to B \dsep C\subseteq B \dsep D\subseteq B \vdash F^{-1}[C\cap D]\subseteq F^{-1}[C]\cap F^{-1}[D]}{1,2,3}}
    \proofstepwidestar[1,2,3]{F^{-1}[C]\cap F^{-1}[D]\subseteq F^{-1}[C\cap D]}{%
      \FormulaRefAuto{F\colon A\to B \dsep C\subseteq B \dsep D\subseteq B \vdash F^{-1}[C]\cap F^{-1}[D]\subseteq F^{-1}[C\cap D]}{1,2,3}}
    \proofstepwidestar[1,2,3]{F^{-1}[C\cap D]=F^{-1}[C]\cap F^{-1}[D]}{%
      \FormulaRefAuto{A\subseteq B\dsep B\subseteq A \vdash A=B}{4,5}}
  \closeproofpartwide

\end{tabproofsplitwide}

\section{Graph einer Funktion}

\FormulaDefDelta[Graph einer Funktion]{\Graph(F)\coloneq\{(x,y) \in A \times B \mid y=F(x)\}}%
{
\DeltaRow{Mengen}{x\dsep y\dsep A\dsep B}
\DeltaRow{Funktionensymbol}{F}
}

\FormulaThmDelta{\Graph(F)\subseteq A\times B}%
{
\DeltaRow{Mengen}{x\dsep y\dsep A\dsep B}
\DeltaRow{Funktionensymbol}{F}
}
\begin{tabproof}
    \proofstep{1}{\Graph(F)\coloneq\{(x,y) \in A \times B \mid y=F(x)\}}{\FormulaRefAuto{\Graph(F)\coloneq\{(x,y) \in A \times B \mid y=F(x)\}}}
    \proofstep{1}{\{(x,y) \in A \times B \mid y=F(x)\}\subseteq A\times B}{\FormulaRefAuto{\{ x \in A \mid P(x)\} \subseteq A}}
    \proofstep{1}{\Graph(F)\subseteq A\times B}{\rIE{1,2}}
\end{tabproof}

\FormulaThmDelta{(x,y)\in \Graph(F)\vdash y=F(x)}%
{
\DeltaRow{Mengen}{x\dsep y\dsep A\dsep B}
\DeltaRow{Funktionensymbol}{F}
}
\begin{tabproof}
  \proofstep{1}{(x,y)\in \Graph(F)}{\rA}
  \proofstep{}{\Graph(F)=\{(x,y) \in A \times B \mid y=F(x)\}}{\FormulaRefAuto{\Graph(F)\coloneq\{(x,y) \in A \times B \mid y=F(x)\}}}
  \proofstep{1}{y=F(x)}{\FormulaRefAuto{x \in A\dsep A=\{x \in B \mid P(x)\}\vdash P(x)}{1,2}}
\end{tabproof}

\FormulaThmDelta[Funktionale Eindeutigkeit von \(\Graph(F)\)]{(x,y)\in \Graph(F)\dsep (x,z)\in \Graph(F)\vdash y=z}%
{
\DeltaRow{Mengen}{x\dsep y\dsep z}
\DeltaRow{Funktionensymbol}{F}
}
\begin{tabproof}
  \proofstep{1}{(x,y)\in \Graph(F)}{\rA}
  \proofstep{2}{(x,z)\in \Graph(F)}{\rA}
  \proofstep{1}{y = F(x)}{\FormulaRefAuto{(x,y)\in Graph(F)\vdash y=F(x)}{1}}
  \proofstep{2}{z = F(x)}{\FormulaRefAuto{(x,y)\in Graph(F)\vdash y=F(x)}{2}}
  \proofstep{5}{y = z}{\FormulaRefAuto{a = b,\, c = b \vdash a = c}{3,4}}
\end{tabproof}

\FormulaThmDelta
{\forall u\in A\,F(u)\in B\dsep x\in A \vdash \exists y\,(x,y)\in\Graph(F)}
{
\DeltaRow{Mengen}{u\dsep x\dsep y}
\DeltaRow{Funktionensymbol}{F}
}
\begin{tabproof}
  \proofstep{1}{\forall u\in A\,F(u)\in B}{\rA}
  \proofstep{2}{x\in A}{\rA}
  \proofstep{1,2}{F(x)\in B}{\rRE{\rUE{1},2}}
  \proofstep{1,2}{(x,F(x))\in A\times B}{\FormulaRefAuto{a\in A\dsep b\in B\vdash (a,b)\in A\times B}{2,3}}
  \proofstep{}{F(x)=F(x)}{\rII}
  \proofstep{1,2}{(x,F(x))\in\{(x,y)\in A \times B\mid y=F(x)\}}{\FormulaRefAuto{x \in A\dsep P(x)\vdash x \in \{x \in A \mid P(x)\}}{4,5}}
  \proofstep{1,2}{(x,F(x))\in\Graph(F)}{\rIE{\FormulaRefAuto{\Graph(F)\coloneq\{(x,y) \in A \times B \mid y=F(x)\}},6}}
  \proofstep{1,2}{\exists y\,(x,y)\in\Graph(F)}{\rEI{7}}
\end{tabproof}

\FormulaThmDelta[Graph eines Funktionensymbols ist eine Funktion]
{\forall u\in A\,F(u)\in B \vdash \Graph(F)\colon A\to B}
{
  \DeltaRow{Mengen}{A\dsep B}
  \DeltaRow{Funktionensymbol}{F}
}
\begin{tabproof}
  \proofstep{1}{\forall u\in A\,F(u)\in B}{\rA}
  \proofstep{}{\Graph(F) \subseteq A\times B }%
    {\FormulaRefAuto{\Graph(F)\subseteq A\times B}}
  \proofstep{}{\forall (x,y),(x,z)\in \Graph(F)\, y=z }%
    {\FormulaRefAuto{(x,y)\in \Graph(F)\dsep (x,z)\in \Graph(F)\vdash y=z}}
  \proofstep{1}{\forall x\in A \exists y\,(x,y)\in \Graph(F)}%
    {\FormulaRefAuto{\forall u\in A\,F(u)\in B\dsep x\in A \vdash \exists y\,(x,y)\in\Graph(F)}{1}}
  \proofstep{1}{\Graph(F)\colon A\to B}%
    {\FormulaRefAuto{Funktion}{1}}
\end{tabproof}

\FormulaThmDelta{%
  \forall u\in A\,F(u)\in B \vdash \exists G\colon A\to B%
}{%
  \DeltaDecl{Mengen}{A\dsep B\dsep u}%
  \DeltaDecl{Funktionen}{F\dsep G}%
}
\begin{tabproof}
  \proofstep{1}{\forall u\in A\,F(u)\in B}{\rA}
  \proofstep{1}{\Graph(F)\colon A\to B}{%
    \FormulaRefAuto{\forall u\in A\,F(u)\in B \vdash \Graph (F)\colon A\to B}{1}}
  \proofstep{1}{\exists G\colon A\to B}{\rEI{2}}
\end{tabproof}

\FormulaThmDelta[Funktionswerte von \(\Graph(F)\)]
{\forall u\in A\,F(u)\in B \vdash \forall x\in A\,\Graph(F)(x) = F(x)}
{
  \DeltaRow{Mengen}{A\dsep B\dsep x}
  \DeltaRow{Funktionensymbol}{F}
}
\begin{tabproof}
  \proofstep{1}{\forall u\in A\,F(u)\in B}{\rA}
  \proofstep{2}{x\in A}{\rA}
  \proofstep{1}{\Graph(F)\colon A\to B}{\FormulaRefAuto{\forall u\in A\,F(u)\in B \vdash \Graph(F)\colon A\to B}{1}}
  \proofstep{1,2}{(x,\Graph(F)(x))\in\Graph(F)}{\FormulaRefAuto{x\in A\vdash (x,F(x))\in F}{3,2}}
  \proofstep{1,2}{\Graph(F)(x) = F(x)}{\FormulaRefAuto{(x,y)\in \Graph(F)\vdash y=F(x)}{4}}
  \proofstep{1}{\forall x\in A\,\Graph(F)(x) = F(x)}{\rUI{\rRI{2,5}}}
\end{tabproof}

\FormulaThmDelta{%
  \forall u\in A\,F(u)\in B \vdash \exists G\colon A\to B\,\forall x\in A\,G(x) = F(x) %
}{%
  \DeltaDecl{Mengen}{A\dsep B\dsep u\dsep x}%
  \DeltaDecl{Funktionen}{F\dsep G}%
}
\begin{tabproof}
  \proofstep{1}{\forall u\in A\,F(u)\in B}{\rA}

  \proofstep{1}{\exists G\colon A\to B}{%
    \FormulaRefAuto{%
      \forall u\in A\,F(u)\in B \vdash \exists G\colon A\to B%
    }{1}}

  \proofstep{1}{\forall x\in A\,\Graph(F)(x)=F(x)}{%
    \FormulaRefAuto{%
      \forall u\in A\,F(u)\in B \vdash \forall x\in A\,\Graph(F)(x)=F(x)%
    }{1}}

  \proofstep{1}{\exists G\colon A\to B\,\forall x\in A\,G(x)=F(x)}{%
    \rEI{2,\rIE{3}}}
\end{tabproof}

\FormulaThmDelta[Graph einer Funktion ist die Funktion]{%
  F\colon A\to B \vdash \Graph(F)=F%
}{%
  \DeltaDecl{Mengen}{A\dsep B\dsep x}%
  \DeltaDecl{Funktionen}{F}%
}
\begin{tabproof}
  \proofstep{1}{F\colon A\to B}{\rA}

  \proofstep{2}{x\in A}{\rA}
  \proofstep{1,2}{F(x)\in B}{%
    \FormulaRefAuto{F\colon A\to B\dsep x\in A\vdash F(x)\in B}{1,2}}
  \proofstep{1}{\forall x\in A\,F(x)\in B}{\rUI{\rRI{2,3}}}

  \proofstep{1}{\Graph(F)\colon A\to B}{\FormulaRefAuto{\forall u\in A\,F(u)\in B \vdash \Graph(F)\colon A\to B}{4}}


  \proofstep{1}{\forall x\in A\,\bigl(\Graph(F)(x)=F(x)\bigr)}{\FormulaRefAuto{\forall u\in A\,F(u)\in B \vdash \forall x\in A\,\Graph(F)(x) = F(x)}{4}}

  \proofstep{1}{\Graph(F)=F}{%
    \FormulaRefAuto{\forall x\in A (F(x)=G(x)) \eqvdash F=G}{1,5,6}}
\end{tabproof}


\FormulaThmDeltaR[Funktion]{\exists! F\colon A\to B\,\forall x\in A\,F(x) = t(x)}{x\in A\vdash t(x)\coloneq y\dsep x\in A\vdash t(x)\in B\exists! F\colon A\to B\,\forall x\in A\,F(x) = t(x)}{
  \DeltaRow{Mengen}{A\dsep B\dsep F\dsep x}
  \DeltaRow{Funktionensymbol}{t}
  %
  \DeltaRow{\textbf{Bedingungen}}{}
  \DeltaRow{Definition einer Funktionsvorschrift}{x\in A\vdash t(x)\coloneq y}
  \DeltaRow{Typisierungsaxiom}{x\in A\vdash t(x)\in B}
  \DeltaRow{\textbf{Folgerung}}{}
}
\begin{tabproof}
  \proofstep{1}{\forall u\in A\,t(u)\in B}{\rA}
  \proofstep{1}{\exists F\colon A\to B\,\forall x\in A\,F(x) = t(x)}{\FormulaRefAuto{\forall u\in A\,t(u)\in B \vdash \exists F\colon A\to B\,\forall x\in A\,F(x) = t(x)}{1}}
  \proofstep{3}{F\colon A\to B\land \forall x\in A\, F(x)=t(x)}{\rA}
  \proofstep{4}{G\colon A\to B\land \forall x\in A\, G(x)=t(x)}{\rA}
  \proofstep{3}{F\colon A\to B}{\rAEa{3}}
  \proofstep{3}{\forall x\in A\, F(x)=t(x)}{\rAEb{3}}
  \proofstep{4}{G\colon A\to B}{\rAEa{4}}
  \proofstep{4}{\forall x\in A\, G(x)=t(x)}{\rAEb{4}}
  \proofstep{3,4}{\forall x\in A\, F(x)=G(x)}{\FormulaRefAuto{\forall x\, F(x)\rightarrow P(x)=Q(x)\dsep \forall x\, F(x)\rightarrow P(x)=R(x)\vdash \forall x\, F(x)\rightarrow Q(x)=R(x)}{6,8}}
  \proofstep{3,4}{F=G}{\FormulaRefAuto{\forall x\in A (F(x)=G(x)) \eqvdash F=G}{9}}
  \proofstep{3,4}{\exists! F\colon A\to B\,\forall x\in A\,F(x) = t(x)}{\UEI{2,3,4,10}}
\end{tabproof}



\section{Eigenschaften von Funktionen}
\subsection{Injektivität}

\FormulaDefDeltaK[Begriff der injektiven Funktion]{F\colon A\rightarrowtail B}{Injektive Funktion}{
  \DeltaRow{Mengen}{A\dsep B\dsep x\dsep y}
  \DeltaPrem{Funktionen}{F\colon A\to B}[\FormulaRefAuto{Funktion}]
  %
  \DeltaRow{\textbf{Axiome}}{}
  %
  % — Injektivität —
  \DeltaRow{Injektivität}
           {F\colon A\inj B\dsep x,y\in A\dsep F(x)=F(y)\vdash x=y}
           [\FormulaRefAutoFwd{F\colon A\inj B\dsep x\in A\dsep y\in A\dsep F(x)=F(y)\vdash x=y}]
  \DeltaRow{\textbf{Neue Symbole}}{}
  \DeltaRow{Injektive Funktionen}{}
}

\FormulaAxiomDeltaK[Injektivität]%
{F\colon A\inj B\dsep x\in A\dsep y\in A\dsep F(x)=F(y)\vdash x=y}%
{Injektivität}{
\DeltaRow{Mengen}{x\dsep y\dsep A\dsep B}
}

\FormulaThmDelta{F\colon A\inj B\dsep x\in A\vdash F(x)\in B}{
\DeltaRow{Mengen}{A\dsep B}
}
\begin{tabproof}
  \proofstep{1}{F\colon A\inj B}{\rA}
  \proofstep{2}{x\in A}{\rA}
  \proofstep{1}{F\colon A\to B}{\FormulaRefAuto{Injektive Funktion}{1}}
  \proofstep{1,2}{F(x)\in B}{%
    \FormulaRefAuto{F\colon A\to B\dsep x\in A\vdash F(x)\in B}{3,2}}
\end{tabproof}

\FormulaThmDelta[Elemente einer Teilmenge liegen im Bild der Teilmenge unter injektiven Funktionen]{%
  C\subseteq A\dsep F\colon A\inj B\dsep x\in C\vdash F(x)\in F[C]
}{
\DeltaRow{Mengen}{A\dsep B\dsep C\dsep x}
\DeltaRow{Funktionen}{F}
}
\begin{tabproof}
  \proofstep{1}{C\subseteq A}{\rA}
  \proofstep{2}{F\colon A\inj B}{\rA}
  \proofstep{3}{x\in C}{\rA}
  \proofstep{1,2,3}{F(x)\in F[C]}{%
    \FormulaRefAuto{C\subseteq A\dsep F\colon A\to B\dsep x\in C\vdash F(x)\in F[C]}{%
      1,\FormulaRefAuto{Injektive Funktion}{2},3}}
\end{tabproof}


\FormulaThmDelta%
{F\colon A\inj B\dsep (x,z_1)\in F\dsep (y,z_2)\in F\dsep F(x)=F(y)\vdash x=y}%
{
\DeltaRow{Mengen}{x\dsep y\dsep z_1\dsep z_2\dsep A\dsep B}
\DeltaRow{Funktionen}{F}
}
\begin{tabproof}
  \proofstep{1}{F\colon A\inj B}{\rA}
  \proofstep{2}{(x,z_1)\in F}{\rA}
  \proofstep{3}{(y,z_2)\in F}{\rA}
  \proofstep{4}{F(x)=F(y)}{\rA}
  \proofstep{2}{x\in A}{\FormulaRefAuto{(x,y)\in F\vdash x\in A}{2}}
  \proofstep{3}{y\in A}{\FormulaRefAuto{(x,y)\in F\vdash x\in A}{3}}
  \proofstep{1,2,3,4}{x=y}{%
    \FormulaRefAuto{F\colon A\inj B\dsep x\in A\dsep y\in A\dsep F(x)=F(y)\vdash x=y}{1,5,6,4}}
\end{tabproof}

\FormulaThmDelta{%
  F\colon M\inj N \dsep x\in M \dsep z\in M \dsep y=F(x) \dsep y=F(z) \vdash x=z
}{%
  \DeltaRow{Mengen}{M\dsep N\dsep x\dsep y\dsep z}
  \DeltaRow{Funktionen}{F}
}
\begin{tabproof}
  \proofstep{1}{F\colon M\inj N}{\rA}
  \proofstep{2}{x\in M}{\rA}
  \proofstep{3}{z\in M}{\rA}
  \proofstep{4}{y=F(x)}{\rA}
  \proofstep{5}{y=F(z)}{\rA}
  \proofstep{4,5}{F(x)=F(z)}{\FormulaRefAuto{a = b,\, a = c \vdash b = c}{4,5}}
  \proofstep{1,2,3,6}{x=z}{%
    \FormulaRefAuto{F\colon A\inj B\dsep x\in A\dsep y\in A\dsep F(x)=F(y)\vdash x=y}{1,2,3,6}}
\end{tabproof}

\subsubsection{Bildmengen unter injektiven Funktionen}

\FormulaThmDelta{%
  F\colon A\inj B \dsep Y\in\powerset(A)\dsep x\in A\dsep F(x)\in F[Y] \vdash x\in Y
}{%
  \DeltaRow{Mengen}{A\dsep B\dsep Y\dsep x\dsep z}
  \DeltaRow{Funktionen}{F}
}
\begin{tabproof}
  \proofstep{1}{F\colon A\inj B}{\rA}
  \proofstep{2}{Y\in\powerset(A)}{\rA}
  \proofstep{3}{x\in A}{\rA}
  \proofstep{4}{F(x)\in F[Y]}{\rA}
  \proofstep{1}{F\colon A\to B}{\FormulaRefAuto{Injektive Funktion}{1}}
  \proofstep{2}{Y\subseteq A}{\FormulaRefAuto{x \in \powerset(A)\eqvdash x \subseteq A}{2}}
  \proofstep{1,2,4}{\exists z\in Y\,F(x)=F(z)}{%
    \FormulaRefAuto{C\subseteq A\dsep F\colon A\to B\dsep y\in F[C]\vdash \exists x\in C\,y=F(x)}{6,5,4}}

  \proofstep{8}{z\in Y\land F(x)=F(z)}{\rA}
  \proofstep{8}{z\in Y}{\rAEa{8}}
  \proofstep{8}{F(x)=F(z)}{\rAEb{8}}
  \proofstep{2,8}{z\in A}{\FormulaRefAuto{A\subseteq B\dsep x\in A \vdash x\in B}{6,9}}
  \proofstep{1,2,3,8}{x=z}{\FormulaRefAuto{Injektivität}{1,3,11,10}}
  \proofstep{1,2,3,8}{z=x}{\FormulaRefAuto{a=b\vdash b=a}{12}}
  \proofstep{1,2,3,8}{x\in Y}{\rIE{13,9}}

  \proofstep{1,2,3,4}{x\in Y}{\rEE{7,8,14}}
\end{tabproof}

\FormulaThmDelta[Injektive Funktionen sind auf Potenzmengen injektiv]{%
  F\colon A\inj B \dsep X\in\powerset(A)\dsep Y\in\powerset(A)\dsep F[X]=F[Y] \vdash X=Y
}{%
  \DeltaRow{Mengen}{A\dsep B\dsep X\dsep Y\dsep x}
  \DeltaRow{Funktionen}{F}
}
\begin{tabproofsplitwide}
  \proofpartwideindR[Teilmengenrichtung \(X\subseteq Y\)]{%
    F\colon A\inj B \dsep X\in\powerset(A)\dsep Y\in\powerset(A)\dsep F[X]=F[Y] \vdash X\subseteq Y
  }{%
    F\colon A\inj B \dsep X\in\powerset(A)\dsep Y\in\powerset(A)\dsep F[X]=F[Y] \vdash X\subseteq Y
  }
    \proofstepwidestar[1]{F\colon A\inj B}{\rA}
    \proofstepwidestar[2]{X\in\powerset(A)}{\rA}
    \proofstepwidestar[3]{Y\in\powerset(A)}{\rA}
    \proofstepwidestar[4]{F[X]=F[Y]}{\rA}
    \proofstepwidestar[2]{X\subseteq A}{\FormulaRefAuto{x \in \powerset(A)\eqvdash x \subseteq A}{2}}
    \proofstepwidestar[1]{F\colon A\to B}{\FormulaRefAuto{Injektive Funktion}{1}}

    \proofstepwidestar[7]{x\in X}{\rA}
    \proofstepwidestar[2,7]{x\in A}{\FormulaRefAuto{A\subseteq B\dsep x\in A \vdash x\in B}{5,7}}
    \proofstepwidestar[1,2,7]{F(x)\in F[X]}{%
      \FormulaRefAuto{C\subseteq A\dsep F\colon A\to B\dsep x\in C\vdash F(x)\in F[C]}{5,6,7}}
    \proofstepwidestar[1,2,3,4,7]{F(x)\in F[Y]}{\rIE{4,9}}
    \proofstepwidestar[1,2,3,4,7]{x\in Y}{%
      \FormulaRefAuto{F\colon A\inj B \dsep Y\in\powerset(A)\dsep x\in A\dsep F(x)\in F[Y] \vdash x\in Y}{1,3,8,10}}

    \proofstepwidestar[1,2,3,4]{x\in X\rightarrow x\in Y}{\rRI{7,11}}
    \proofstepwidestar[1,2,3,4]{\forall x\,(x\in X\rightarrow x\in Y)}{\rUI{\rRI{7,12}}}
    \proofstepwidestar[1,2,3,4]{X\subseteq Y}{%
      \FormulaRefAuto{A\subseteq B \coloneq \forall x\,(x\in A \rightarrow x\in B)}{13}}
  \closeproofpartwideindR

  \proofpartwideindR[Teilmengenrichtung \(Y\subseteq X\)]{%
    F\colon A\inj B \dsep X\in\powerset(A)\dsep Y\in\powerset(A)\dsep F[X]=F[Y] \vdash Y\subseteq X
  }{%
    F\colon A\inj B \dsep X\in\powerset(A)\dsep Y\in\powerset(A)\dsep F[X]=F[Y] \vdash Y\subseteq X
  }
    \proofstepwidestar[1]{F\colon A\inj B}{\rA}
    \proofstepwidestar[2]{X\in\powerset(A)}{\rA}
    \proofstepwidestar[3]{Y\in\powerset(A)}{\rA}
    \proofstepwidestar[4]{F[X]=F[Y]}{\rA}
    \proofstepwidestar[4]{F[Y]=F[X]}{\FormulaRefAuto{a=b\vdash b=a}{4}}
    \proofstepwidestar[1,2,3,4]{Y\subseteq X}{%
      \FormulaRefAuto{F\colon A\inj B \dsep X\in\powerset(A)\dsep Y\in\powerset(A)\dsep F[X]=F[Y] \vdash X\subseteq Y}{1,3,2,5}}
  \closeproofpartwideindR

  \proofpartwide{Schluss}
    \proofstepwidestar[1]{F\colon A\inj B}{\rA}
    \proofstepwidestar[2]{X\in\powerset(A)}{\rA}
    \proofstepwidestar[3]{Y\in\powerset(A)}{\rA}
    \proofstepwidestar[4]{F[X]=F[Y]}{\rA}
    \proofstepwidestar[1,2,3,4]{X\subseteq Y}{%
      \FormulaRefAuto{F\colon A\inj B \dsep X\in\powerset(A)\dsep Y\in\powerset(A)\dsep F[X]=F[Y] \vdash X\subseteq Y}{1,2,3,4}}
    \proofstepwidestar[1,2,3,4]{Y\subseteq X}{%
      \FormulaRefAuto{F\colon A\inj B \dsep X\in\powerset(A)\dsep Y\in\powerset(A)\dsep F[X]=F[Y] \vdash Y\subseteq X}{1,2,3,4}}
    \proofstepwidestar[1,2,3,4]{X=Y}{%
      \FormulaRefAuto{A \subseteq B, B \subseteq A \vdash A = B}{5,6}}
  \closeproofpartwide
\end{tabproofsplitwide}

\FormulaThmDelta[Zeugen gemeinsamer Bildwerte liegen im Durchschnitt]{%
  F\colon M\inj N \dsep A\subseteq M \dsep B\subseteq M \dsep x\in A \dsep z\in B \dsep y=F(x) \dsep y=F(z) \vdash x\in A\cap B
}{%
  \DeltaRow{Mengen}{A\dsep B\dsep M\dsep N\dsep x\dsep y\dsep z}
  \DeltaRow{Funktionen}{F}
}
\begin{tabproof}
  \proofstep{1}{F\colon M\inj N}{\rA}
  \proofstep{2}{A\subseteq M}{\rA}
  \proofstep{3}{B\subseteq M}{\rA}
  \proofstep{4}{x\in A}{\rA}
  \proofstep{5}{z\in B}{\rA}
  \proofstep{6}{y=F(x)}{\rA}
  \proofstep{7}{y=F(z)}{\rA}
  \proofstep{2,4}{x\in M}{\FormulaRefAuto{A\subseteq B\dsep x\in A\vdash x\in B}{2,4}}
  \proofstep{3,5}{z\in M}{\FormulaRefAuto{A\subseteq B\dsep x\in A\vdash x\in B}{3,5}}
  \proofstep{1,6,7,8,9}{x=z}{%
    \FormulaRefAuto{F\colon M\inj N \dsep x\in M \dsep z\in M \dsep y=F(x) \dsep y=F(z) \vdash x=z}{1,8,9,6,7}}
  \proofstep{5,10}{x\in B}{\rIE{10,5}}
  \proofstep{4,11}{x\in A\land x\in B}{\rAI{4,11}}
  \proofstep{1,2,3,4,5,6,7}{x\in A\cap B}{%
    \FormulaRefAuto{x\in A\cap B\eqvdash x\in A\land x\in B}{12}}
\end{tabproof}


\FormulaThmDelta[Gemeinsame Bildwerte liegen im Bild des Durchschnitts]{%
  F\colon M\inj N \dsep A\subseteq M \dsep B\subseteq M \dsep y\in F[A] \dsep y\in F[B] \vdash y\in F[A\cap B]
}{%
  \DeltaRow{Mengen}{A\dsep B\dsep M\dsep N\dsep x\dsep y\dsep z}
  \DeltaRow{Funktionen}{F}
}
\begin{tabproof}
  \proofstep{1}{F\colon M\inj N}{\rA}
  \proofstep{2}{A\subseteq M}{\rA}
  \proofstep{3}{B\subseteq M}{\rA}
  \proofstep{4}{y\in F[A]}{\rA}
  \proofstep{5}{y\in F[B]}{\rA}
  \proofstep{1}{F\colon M\to N}{\FormulaRefAuto{Injektive Funktion}{1}}
  \proofstep{2,6,4}{\exists x\in A\,y=F(x)}{%
    \FormulaRefAuto{C\subseteq A\dsep F\colon A\to B\dsep y\in F[C]\vdash \exists x\in C\,y=F(x)}{2,6,4}}
  \proofstep{3,6,5}{\exists z\in B\,y=F(z)}{%
    \FormulaRefAuto{C\subseteq A\dsep F\colon A\to B\dsep y\in F[C]\vdash \exists x\in C\,y=F(x)}{3,6,5}}
  \proofstep{9}{x\in A\land y=F(x)}{\rA}
  \proofstep{9}{x\in A}{\rAEa{9}}
  \proofstep{9}{y=F(x)}{\rAEb{9}}
  \proofstep{12}{z\in B\land y=F(z)}{\rA}
  \proofstep{12}{z\in B}{\rAEa{12}}
  \proofstep{12}{y=F(z)}{\rAEb{12}}
  \proofstep{1,2,3,10,13,11,14}{x\in A\cap B}{%
    \FormulaRefAuto{F\colon M\inj N \dsep A\subseteq M \dsep B\subseteq M \dsep x\in A \dsep z\in B \dsep y=F(x) \dsep y=F(z) \vdash x\in A\cap B}{1,2,3,10,13,11,14}}
  \proofstep{1,2,3,9,12}{x\in A\cap B\land y=F(x)}{\rAI{15,11}}
  \proofstep{1,2,3,9,12}{\exists x\in A\cap B\,y=F(x)}{\rEI{16}}
  \proofstep{}{A\cap B\subseteq A}{\FormulaRefAuto{A\cap B\subseteq A}}
  \proofstep{2}{A\cap B\subseteq M}{\FormulaRefAuto{A\subseteq B,\, B\subseteq C \vdash A\subseteq C}{18,2}}
  \proofstep{1,2,3,9,12}{y\in F[A\cap B]}{%
    \FormulaRefAuto{C\subseteq A\dsep F\colon A\to B\dsep \exists x\in C\,y=F(x)\vdash y\in F[C]}{19,6,17}}
  \proofstep{1,2,3,5,9}{y\in F[A\cap B]}{\rEE{8,12,20}}
  \proofstep{1,2,3,4,5}{y\in F[A\cap B]}{\rEE{7,9,21}}
\end{tabproof}


\FormulaThmDelta[Bild eines Durchschnitts bei injektiven Funktionen]{%
  F\colon M\inj N \dsep A\subseteq M \dsep B\subseteq M \vdash F[A\cap B]=F[A]\cap F[B]
}{%
  \DeltaRow{Mengen}{A\dsep B\dsep M\dsep N\dsep x\dsep y\dsep z}
  \DeltaRow{Funktionen}{F}
}
\begin{tabproofsplitwide}

  \proofpartwideind[Teilmengenrichtung \(F[A\cap B]\subseteq F[A]\cap F[B]\)]{%
    F\colon M\inj N \dsep A\subseteq M \dsep B\subseteq M \vdash F[A\cap B]\subseteq F[A]\cap F[B]
  }[%
    F\colon M\inj N \dsep A\subseteq M \dsep B\subseteq M \vdash F[A\cap B]\subseteq F[A]\cap F[B]
  ]
    \proofstepwidestar[1]{F\colon M\inj N}{\rA}
    \proofstepwidestar[2]{A\subseteq M}{\rA}
    \proofstepwidestar[3]{B\subseteq M}{\rA}
    \proofstepwidestar[1]{F\colon M\to N}{\FormulaRefAuto{Injektive Funktion}{1}}
    \proofstepwidestar[1,2,3]{F[A\cap B]\subseteq F[A]\cap F[B]}{%
      \FormulaRefAuto{F\colon M\to N \dsep A\subseteq M \dsep B\subseteq M \vdash F[A\cap B]\subseteq F[A]\cap F[B]}{4,2,3}}
  \closeproofpartwideind

   \proofpartwideind[Teilmengenrichtung \(F[A]\cap F[B]\subseteq F[A\cap B]\)]{%
    F\colon M\inj N \dsep A\subseteq M \dsep B\subseteq M \vdash F[A]\cap F[B]\subseteq F[A\cap B]
  }[%
    F\colon M\inj N \dsep A\subseteq M \dsep B\subseteq M \vdash F[A]\cap F[B]\subseteq F[A\cap B]
  ]
    \proofstepwidestar[1]{F\colon M\inj N}{\rA}
    \proofstepwidestar[2]{A\subseteq M}{\rA}
    \proofstepwidestar[3]{B\subseteq M}{\rA}
    \proofstepwidestar[4]{y\in F[A]\cap F[B]}{\rA}
    \proofstepwidestar[4]{y\in F[A]\land y\in F[B]}{%
      \FormulaRefAuto{y\in A\cap B\eqvdash y\in A\land y\in B}{4}}
    \proofstepwidestar[4]{y\in F[A]}{\rAEa{5}}
    \proofstepwidestar[4]{y\in F[B]}{\rAEb{5}}
    \proofstepwidestar[1,2,3,6,7]{y\in F[A\cap B]}{%
      \FormulaRefAuto{F\colon M\inj N \dsep A\subseteq M \dsep B\subseteq M \dsep y\in F[A] \dsep y\in F[B] \vdash y\in F[A\cap B]}{1,2,3,6,7}}
    \proofstepwidestar[1]{\forall y\,\bigl(y\in F[A]\cap F[B]\rightarrow y\in F[A\cap B]\bigr)}{\rUI{\rRI{4,8}}}
    \proofstepwidestar[1]{F[A]\cap F[B]\subseteq F[A\cap B]}{%
      \FormulaRefAuto{A\subseteq B \coloneq \forall x\,(x\in A \rightarrow x\in B)}{9}}
  \closeproofpartwideind

  \proofpartwide{Zusammenfassung}
    \proofstepwidestar[1]{F\colon M\inj N}{\rA}
    \proofstepwidestar[2]{A\subseteq M}{\rA}
    \proofstepwidestar[3]{B\subseteq M}{\rA}
    \proofstepwidestar[1,2,3]{F[A\cap B]\subseteq F[A]\cap F[B]}{%
      \FormulaRefAuto{F\colon M\inj N \dsep A\subseteq M \dsep B\subseteq M \vdash F[A\cap B]\subseteq F[A]\cap F[B]}{1,2,3}}
    \proofstepwidestar[1,2,3]{F[A]\cap F[B]\subseteq F[A\cap B]}{%
      \FormulaRefAuto{F\colon M\inj N \dsep A\subseteq M \dsep B\subseteq M \vdash F[A]\cap F[B]\subseteq F[A\cap B]}{1,2,3}}
    \proofstepwidestar[1,2,3]{F[A\cap B]=F[A]\cap F[B]}{%
      \FormulaRefAuto{A\subseteq B\dsep B\subseteq A \vdash A=B}{4,5}}
  \closeproofpartwide

\end{tabproofsplitwide}


\subsubsection{Urbildmengen unter injektiven Funktionen}

\FormulaThmDelta[Urbilder liegen im Definitionsbereich injektiver Funktionen]{%
  C\subseteq B\dsep F\colon A\inj B \vdash F^{-1}[C]\subseteq A
}{
  \DeltaRow{Mengen}{A\dsep B\dsep C}
  \DeltaRow{Funktionen}{F}
}
\begin{tabproof}
  \proofstep{1}{C\subseteq B}{\rA}
  \proofstep{2}{F\colon A\inj B}{\rA}
  \proofstep{2}{F\colon A\to B}{\FormulaRefAuto{Injektive Funktion}{2}}
  \proofstep{1,2}{F^{-1}[C]\subseteq A}{%
    \FormulaRefAuto{F\colon A\to B\dsep C\subseteq B \vdash F^{-1}[C]\subseteq A}{3,1}}
\end{tabproof}

\subsection{Surjektivität}

\FormulaDefDeltaK[Begriff der surjektiven Funktion]{F\colon A\twoheadrightarrow B}{Surjektive Funktion}{
  \DeltaRow{Mengen}{A\dsep B\dsep x\dsep y}
  \DeltaPrem{Funktionen}{F\colon A\to B}[\FormulaRefAuto{Funktion}]
%
  \DeltaRow{\textbf{Axiome}}{}
  %
  % — Surjekvität —
  \DeltaRow{Surjektivität}
           {y\in B\vdash \exists x\in A\; F(x)=y}
           [\FormulaRefAutoFwd{y\in B\vdash \exists x\in A\; F(x)=y}]
  \DeltaRow{\textbf{Neue Symbole}}{}
  \DeltaPrem{Surjektive Funktionen}{ }
}

\FormulaAxiomDeltaK[Surjektivität]%
{y\in B\vdash \exists x\in A\; F(x)=y}{Surjektivität}
{
\DeltaRow{Mengen}{y\dsep A\dsep B}
\DeltaPrem{Funktionen}{F\colon A\to B}
}

\FormulaThmDelta{F\colon A\sur B\dsep x\in A\vdash F(x)\in B}{
\DeltaRow{Mengen}{A\dsep B}
}
\begin{tabproof}
  \proofstep{1}{F\colon A\sur B}{\rA}
  \proofstep{2}{x\in A}{\rA}
  \proofstep{1}{F\colon A\to B}{\FormulaRefAuto{Surjektive Funktion}{1}}
  \proofstep{1,2}{F(x)\in B}{%
    \FormulaRefAuto{F\colon A\to B\dsep x\in A\vdash F(x)\in B}{3,2}}
\end{tabproof}

\FormulaThmDelta[Elemente einer Teilmenge liegen im Bild der Teilmenge unter surjektiven Funktionen]{%
  C\subseteq A\dsep F\colon A\sur B\dsep x\in C\vdash F(x)\in F[C]
}{
\DeltaRow{Mengen}{A\dsep B\dsep C\dsep x}
\DeltaRow{Funktionen}{F}
}
\begin{tabproof}
  \proofstep{1}{C\subseteq A}{\rA}
  \proofstep{2}{F\colon A\sur B}{\rA}
  \proofstep{3}{x\in C}{\rA}
  \proofstep{1,2,3}{F(x)\in F[C]}{%
    \FormulaRefAuto{C\subseteq A\dsep F\colon A\to B\dsep x\in C\vdash F(x)\in F[C]}{%
      1,\FormulaRefAuto{Surjektive Funktion}{2},3}}
\end{tabproof}

\FormulaThmDeltaR[Nichtsurjektivität liefert ausgelassenes Element]{%
  \begin{aligned}[t]
    &F\colon M\to M\dsep \neg(F\colon M\sur M)\vdash{}\\
    &\exists m_0\in M\,\forall x\in M\,\bigl(F(x)\neq m_0\bigr)
  \end{aligned}
}{%
  F\colon M\to M\dsep \neg(F\colon M\sur M)\vdash
  \exists m_0\in M\,\forall x\in M\,\bigl(F(x)\neq m_0\bigr)
}{%
  \DeltaDecl{Mengen}{M\dsep m_0\dsep x\dsep y}%
  \DeltaDecl{Funktionen}{F}%
}
\begin{tabproof}
  \proofstep{1}{F\colon M\to M}{\rA}
  \proofstep{2}{\neg(F\colon M\sur M)}{\rA}

  \proofstep{3}{\forall y\in M\,\exists x\in M\,\bigl(F(x)=y\bigr)}{\rA}

  \proofstep{1,3}{F\colon M\sur M}{\FormulaRefAuto{Surjektive Funktion}{1,3}}

  \proofstep{1,2,3}{\bot}{\rBI{2,4}}

  \proofstep{1,2}{%
    \neg\forall y\in M\,\exists x\in M\,\bigl(F(x)=y\bigr)
  }{\rCI{3,5}}

  \proofstep{1,2}{%
    \exists m_0\in M\,\forall x\in M\,\bigl(F(x)\neq m_0\bigr)
  }{%
    \FormulaRefAuto{%
      \exists a\in A\,\forall b\in B\,(\neg F(a,b))
      \eqvdash
      \neg\forall a\in A\,\exists b\in B\,(F(a,b))
    }[6]%
  }
\end{tabproof}

\FormulaThmDelta{F\colon A\to B \vdash F\colon A\twoheadrightarrow F[A]}{
  \DeltaRow{Mengen}{A\dsep B\dsep x\dsep y}
  \DeltaRow{Funktionen}{F}
}
\begin{tabproof}
  \proofstep{1}{F\colon A\to B}{\rA}
  \proofstep{1}{F\subseteq A\times F[A]}{\FormulaRefAuto{F\subseteq A\times F[A]}{1}}
  \proofstep{1}{\forall x,y\,((x,y)\in F\dsep (x,z)\in F\rightarrow y=z)}{%
    \FormulaRefAuto{(x,y)\in F\dsep (x,z)\in F \vdash y=z}{1}}
  \proofstep{1}{\forall x\,(x\in A\rightarrow \exists y\,(x,y)\in F)}{%
    \FormulaRefAuto{x\in A \vdash \exists y\,(x,y)\in F}{1}}

  \proofstep{5}{y\in F[A]}{\rA}
  \proofstep{1,5}{\exists x\in A\,y=F(x)}{%
    \FormulaRefAuto{F\colon A\to B\dsep y\in F[A] \vdash \exists x\in A\,y=F(x)}{1,5}}

  \proofstep{7}{x\in A\land y=F(x)}{\rA}
  \proofstep{7}{x\in A}{\rAEa{7}}
  \proofstep{7}{y=F(x)}{\rAEb{7}}
  \proofstep{7}{F(x)=y}{\FormulaRefAuto{a=b\vdash b=a}{9}}
  \proofstep{7}{x\in A\land F(x)=y}{\rAI{8,10}}
  \proofstep{7}{\exists x\in A\,F(x)=y}{\rEI{11}}

  \proofstep{1,5}{\exists x\in A\,F(x)=y}{\rEE{6,7,12}}
  \proofstep{1}{\forall y\,(y\in F[A]\rightarrow \exists x\in A\,F(x)=y)}{\rUI{\rRI{5,13}}}

  \proofstep{1}{F\colon A\twoheadrightarrow F[A]}{\FormulaRefAuto{Surjektive Funktion}{2,3,4,14}}
\end{tabproof}

% ------------------------------------------------------------
% Hilfslemma: der breite Äquivalenz-Block (Diagonal-Schritt)
% ------------------------------------------------------------
\FormulaThmDelta{%
  a\in A \dsep F(a)=\{x\in A\mid x\notin F(x)\}\vdash a\in F(a)\leftrightarrow a\notin F(a)%
}{
  \DeltaRow{Mengen}{A\dsep a\dsep x}
  \DeltaPrem{Funktionen}{F\colon A\to \powerset(A)}
}
\begin{tabproofwide}
  \proofstepwidestar[1]{a\in A}{\rA}
  \proofstepwidestar[2]{F(a)=\{x\in A\mid x\notin F(x)\}}{\rA}

  \proofstepwide[1,2]{a\in F(a)}{\leftrightarrow}{%
    \begin{aligned}[t]
      a\in \{x\in A\mid {}\\
      x\notin F(x)\}
    \end{aligned}
  }{\rUE{\FormulaRefAuto{\forall x\, (x \in A \leftrightarrow x \in B) \eqvdash A = B}{2}}}

  \proofstepwide[1,2]{}{\leftrightarrow}{a\in A \land a\notin F(a)}{%
    \FormulaRefAuto{x \in \{u \in A \mid P(u)\} \eqvdash x \in A \land P(x)}{3}}

  \proofstepwide[1,2]{}{\leftrightarrow}{a\notin F(a)}{%
    \FormulaRefAuto{Q\vdash (P\leftrightarrow Q\land R)\leftrightarrow (P\leftrightarrow R)}{1}}

  \proofstepwide[1,2]{a\in F(a)}{\leftrightarrow}{a\notin F(a)}{\rChain{3,5}}
\end{tabproofwide}


% ------------------------------------------------------------
% Hauptsatz: Cantor (keine Surjektion A -> P(A))
% ------------------------------------------------------------
\FormulaThmDelta{%
  F\colon A\to \powerset(A) \vdash \neg(F\colon A\sur \powerset(A))%
}{
  \DeltaRow{Mengen}{A\dsep F}
}
\begin{tabproofwide}
  \proofstepwidestar[1]{F\colon A\sur \powerset(A)}{\rA}
  \proofstepwidestar[1]{\forall y\in \powerset(A)\,\exists x\in A\,(F(x)=y)}{%
    \FormulaRefAuto{y\in B\vdash \exists x\in A\; F(x)=y}{1}}
  \proofstepwidestar[1]{\exists x\in A\,\bigl(F(x)=\{x\in A\mid x\notin F(x)\}\bigr)}{\rUE{2}}

  \proofstepwidestar[4]{a\in A\land F(a)=\{x\in A\mid x\notin F(x)\}}{\rA}
  \proofstepwidestar[4]{a\in A}{\rAEa{4}}
  \proofstepwidestar[4]{F(a)=\{x\in A\mid x\notin F(x)\}}{\rAEb{4}}

  \proofstepwide[4]{a\in F(a)}{\leftrightarrow}{a\notin F(a)}{%
    \FormulaRefAuto{a\in A \dsep F(a)=\{x\in A\mid x\notin F(x)\}\vdash a\in F(a)\leftrightarrow a\notin F(a)}{5,6}}

  \proofstepwidestar[]{\neg (a\in F(a)\leftrightarrow a\notin F(a))}{%
    \FormulaRefAuto{\neg(P\leftrightarrow\neg P)}}

  \proofstepwidestar[4]{\bot}{\rBI{7,8}}
  \proofstepwidestar[1]{\bot}{\rEE{3,4,9}}
  \proofstepwidestar[]{\neg(F\colon A\sur \powerset(A))}{\rCI{1,10}}
\end{tabproofwide}


% ------------------------------------------------------------
% Surjektion aus Injektion (bei Nichtleere)
% (Einfügen nach \subsection{Surjektivität}, vor \subsection{Bijektivität})
% ------------------------------------------------------------

\subsubsection{Surjektion aus Injektion bei Nichtleere}


\FormulaDefDeltaR{%
  \begin{gathered}[t]
    F\colon A\inj B \vdash
    \Gsurjfrominj{F}{a_0} \coloneqq\\
    \{\, (b,a)\in B\times A \mid\\
    (b\in F[A]\land F(a)=b)\lor\\
    (b\notin F[A]\land a=a_0)\,\}
  \end{gathered}
}{%
  F\colon A\inj B \vdash
  \Gsurjfrominj{F}{a_0} \coloneqq
  \{\, (b,a)\in B\times A \mid (b\in F[A]\land F(a)=b)\lor (b\notin F[A]\land a=a_0)\,\}
}{%
  \DeltaRow{Mengen}{A\dsep B\dsep a_0}
  \DeltaRow{Funktionen}{F}
}

\FormulaThmDeltaR{%
  F\colon A\inj B \vdash \Gsurjfrominj{F}{a_0}\subseteq B\times A
}{%
  F\colon A\inj B \vdash \Gsurjfrominj{F}{a_0}\subseteq B\times A
}{%
  \DeltaRow{Mengen}{A\dsep B\dsep a_0\dsep a\dsep b}
  \DeltaRow{Funktionen}{F}
}
\begin{tabproof}
  \proofstep{1}{F\colon A\inj B}{\rA}
  \proofstep{}{%
    \begin{gathered}[t]
      \{\, (b,a)\in B\times A \mid\\
      (b\in F[A]\land F(a)=b)\lor\\
      (b\notin F[A]\land a=a_0)\,\}\\
      \subseteq B\times A
    \end{gathered}
  }{\FormulaRefAuto{\{\,x\in A \mid P(x)\,\}\subseteq A}}
  \proofstep{1}{\Gsurjfrominj{F}{a_0}\subseteq B\times A}{%
    \rIE{\FormulaRefAuto{F\colon A\inj B \vdash \Gsurjfrominj{F}{a_0} \coloneqq
    \{\, (b,a)\in B\times A \mid (b\in F[A]\land F(a)=b)\lor (b\notin F[A]\land a=a_0)\,\}}{1},2}}
\end{tabproof}

\FormulaThmDelta{%
  F\colon A\inj B \dsep (b,a)\in \Gsurjfrominj{F}{a_0}\vdash b\in B%
}{
  \DeltaRow{Mengen}{A\dsep B\dsep a_0\dsep a\dsep b}
  \DeltaRow{Funktionen}{F}
}
\begin{tabproof}
  \proofstep{1}{F\colon A\inj B}{\rA}
  \proofstep{2}{(b,a)\in \Gsurjfrominj{F}{a_0}}{\rA}
  \proofstep{1,2}{%
    \begin{aligned}[t]
      (b,a)\in \{\, (b,a)\in B\times A \mid {}&
        (b\in F[A]\land F(a)=b)\\
        &\lor (b\notin F[A]\land a=a_0)\,\}
    \end{aligned}
  }{%
    \rIE{\FormulaRefAuto{F\colon A\inj B \vdash \Gsurjfrominj{F}{a_0} \coloneqq
      \{\, (b,a)\in B\times A \mid (b\in F[A]\land F(a)=b)\lor (b\notin F[A]\land a=a_0)\,\}}{1},2}}
  \proofstep{1,2}{%
    \begin{aligned}[t]
      (b,a)\in B\times A \land \bigl({}&
        (b\in F[A]\land F(a)=b)\\
        &\lor (b\notin F[A]\land a=a_0)\bigr)
    \end{aligned}
  }{%
    \FormulaRefAuto{x \in \{u \in A \mid P(u)\} \eqvdash x \in A \land P(x)}{3}}
  \proofstep{1,2}{(b,a)\in B\times A}{\rAEa{4}}
  \proofstep{1,2}{b\in B\land a\in A}{\FormulaRefAuto{(a,b)\in A\times B\eqvdash a\in A\land b\in B}{5}}
  \proofstep{1,2}{b\in B}{\rAEa{6}}
\end{tabproof}

\FormulaThmDelta{%
  F\colon A\inj B \dsep (b,a)\in \Gsurjfrominj{F}{a_0}\vdash a\in A%
}{
  \DeltaRow{Mengen}{A\dsep B\dsep a_0\dsep a\dsep b}
  \DeltaRow{Funktionen}{F}
}
\begin{tabproof}
  \proofstep{1}{F\colon A\inj B}{\rA}
  \proofstep{2}{(b,a)\in \Gsurjfrominj{F}{a_0}}{\rA}
  \proofstep{1,2}{%
    \begin{aligned}[t]
      (b,a)\in \{\, (b,a)\in B\times A \mid {}&
        (b\in F[A]\land F(a)=b)\\
        &\lor (b\notin F[A]\land a=a_0)\,\}
    \end{aligned}
  }{%
    \rIE{\FormulaRefAuto{F\colon A\inj B \vdash \Gsurjfrominj{F}{a_0} \coloneqq
      \{\, (b,a)\in B\times A \mid (b\in F[A]\land F(a)=b)\lor (b\notin F[A]\land a=a_0)\,\}}{1},2}}
  \proofstep{1,2}{%
    \begin{aligned}[t]
      (b,a)\in B\times A \land \bigl({}&
        (b\in F[A]\land F(a)=b)\\
        &\lor (b\notin F[A]\land a=a_0)\bigr)
    \end{aligned}
  }{%
    \FormulaRefAuto{x \in \{u \in A \mid P(u)\} \eqvdash x \in A \land P(x)}{3}}
  \proofstep{1,2}{(b,a)\in B\times A}{\rAEa{4}}
  \proofstep{1,2}{b\in B\land a\in A}{\FormulaRefAuto{(a,b)\in A\times B\eqvdash a\in A\land b\in B}{5}}
  \proofstep{1,2}{a\in A}{\rAEb{6}}
\end{tabproof}

\FormulaThmDelta{%
  F\colon A\inj B \dsep (b,a)\in \Gsurjfrominj{F}{a_0}\vdash
  (b\in F[A]\land F(a)=b)\lor (b\notin F[A]\land a=a_0)%
}{
  \DeltaRow{Mengen}{A\dsep B\dsep a_0\dsep a\dsep b}
  \DeltaRow{Funktionen}{F}
}
\begin{tabproof}
  \proofstep{1}{F\colon A\inj B}{\rA}
  \proofstep{2}{(b,a)\in \Gsurjfrominj{F}{a_0}}{\rA}
  \proofstep{1,2}{%
    \begin{aligned}[t]
      (b,a)\in \{\, (b,a)\in B\times A \mid {}&
        (b\in F[A]\land F(a)=b)\\
        &\lor (b\notin F[A]\land a=a_0)\,\}
    \end{aligned}
  }{%
    \rIE{\FormulaRefAuto{F\colon A\inj B \vdash \Gsurjfrominj{F}{a_0} \coloneqq
      \{\, (b,a)\in B\times A \mid (b\in F[A]\land F(a)=b)\lor (b\notin F[A]\land a=a_0)\,\}}{1},2}}
  \proofstep{1,2}{%
    \begin{aligned}[t]
      (b,a)\in B\times A \land \bigl({}&
        (b\in F[A]\land F(a)=b)\\
        &\lor (b\notin F[A]\land a=a_0)\bigr)
    \end{aligned}
  }{%
    \FormulaRefAuto{x \in \{u \in A \mid P(u)\} \eqvdash x \in A \land P(x)}{3}}
  \proofstep{1,2}{(b\in F[A]\land F(a)=b)\lor (b\notin F[A]\land a=a_0)}{\rAEb{4}}
\end{tabproof}

\FormulaThmDeltaR{%
  \begin{gathered}[t]
    F\colon A\inj B \dsep b\in B \dsep a\in A \dsep\\
    \bigl((b\in F[A]\land F(a)=b)\lor\\
    (b\notin F[A]\land a=a_0)\bigr)\\
    \vdash (b,a)\in \Gsurjfrominj{F}{a_0}
  \end{gathered}
}{%
  F\colon A\inj B \dsep b\in B \dsep a\in A \dsep
  \bigl((b\in F[A]\land F(a)=b)\lor (b\notin F[A]\land a=a_0)\bigr)
  \vdash (b,a)\in \Gsurjfrominj{F}{a_0}
}{%
  \DeltaRow{Mengen}{A\dsep B\dsep a_0\dsep a\dsep b}
  \DeltaRow{Funktionen}{F}
}
\begin{tabproof}
  \proofstep{1}{F\colon A\inj B}{\rA}
  \proofstep{2}{b\in B}{\rA}
  \proofstep{3}{a\in A}{\rA}
  \proofstep{4}{%
    \begin{gathered}[t]
      (b\in F[A]\land F(a)=b)\lor\\
      (b\notin F[A]\land a=a_0)
    \end{gathered}
  }{\rA}
  \proofstep{2,3}{(b,a)\in B\times A}{\FormulaRefAuto{a\in A\dsep b\in B\vdash (a,b)\in A\times B}{2,3}}
  \proofstep{2,3,4}{%
    \begin{gathered}[t]
      (b,a)\in B\times A \land\\
      \bigl((b\in F[A]\land F(a)=b)\lor\\
      (b\notin F[A]\land a=a_0)\bigr)
    \end{gathered}
  }{\rAI{5,4}}
  \proofstep{2,3,4}{%
    \begin{gathered}[t]
      (b,a)\in \{\, (b,a)\in B\times A \mid\\
      (b\in F[A]\land F(a)=b)\lor\\
      (b\notin F[A]\land a=a_0)\,\}
    \end{gathered}
  }{%
    \FormulaRefAuto{x \in \{u \in A \mid P(u)\} \eqvdash x \in A \land P(x)}{6}}
  \proofstep{1,2,3,4}{(b,a)\in \Gsurjfrominj{F}{a_0}}{%
    \rIE{\FormulaRefAuto{F\colon A\inj B \vdash \Gsurjfrominj{F}{a_0} \coloneqq
    \{\, (b,a)\in B\times A \mid (b\in F[A]\land F(a)=b)\lor (b\notin F[A]\land a=a_0)\,\}}{1},7}}
\end{tabproof}

\FormulaThmDelta{%
  F\colon A\inj B \dsep b\in F[A] \dsep (b,a)\in \Gsurjfrominj{F}{a_0}\vdash F(a)=b%
}{
  \DeltaRow{Mengen}{A\dsep B\dsep a_0\dsep a\dsep b}
  \DeltaRow{Funktionen}{F}
}
\begin{tabproof}
  \proofstep{1}{F\colon A\inj B}{\rA}
  \proofstep{2}{b\in F[A]}{\rA}
  \proofstep{3}{(b,a)\in \Gsurjfrominj{F}{a_0}}{\rA}
  \proofstep{1,3}{(b\in F[A]\land F(a)=b)\lor (b\notin F[A]\land a=a_0)}{%
    \FormulaRefAuto{F\colon A\inj B \dsep (b,a)\in \Gsurjfrominj{F}{a_0}\vdash (b\in F[A]\land F(a)=b)\lor (b\notin F[A]\land a=a_0)}{1,3}}
  \proofstep{5}{b\notin F[A]\land a=a_0}{\rA}
  \proofstep{5}{b\notin F[A]}{\rAEa{5}}
  \proofstep{2,5}{\bot}{\rBI{2,6}}
  \proofstep{2}{\neg(b\notin F[A]\land a=a_0)}{\rCI{5,7}}
  \proofstep{1,2,3}{(b\in F[A]\land F(a)=b)}{%
    \FormulaRefAuto{P \lor Q, \neg Q \vdash P}{4,8}}
  \proofstep{1,2,3}{F(a)=b}{\rAEb{9}}
\end{tabproof}

\FormulaThmDelta{%
  F\colon A\inj B \dsep b\notin F[A] \dsep (b,a)\in \Gsurjfrominj{F}{a_0}\vdash a=a_0%
}{
  \DeltaRow{Mengen}{A\dsep B\dsep a_0\dsep a\dsep b}
  \DeltaRow{Funktionen}{F}
}
\begin{tabproof}
  \proofstep{1}{F\colon A\inj B}{\rA}
  \proofstep{2}{b\notin F[A]}{\rA}
  \proofstep{3}{(b,a)\in \Gsurjfrominj{F}{a_0}}{\rA}
  \proofstep{1,3}{(b\in F[A]\land F(a)=b)\lor (b\notin F[A]\land a=a_0)}{%
    \FormulaRefAuto{F\colon A\inj B \dsep (b,a)\in \Gsurjfrominj{F}{a_0}\vdash (b\in F[A]\land F(a)=b)\lor (b\notin F[A]\land a=a_0)}{1,3}}
  \proofstep{5}{b\in F[A]\land F(a)=b}{\rA}
  \proofstep{5}{b\in F[A]}{\rAEa{5}}
  \proofstep{2,5}{\bot}{\rBI{6,2}}
  \proofstep{2}{\neg(b\in F[A]\land F(a)=b)}{\rCI{5,7}}
  \proofstep{1,2,3}{(b\notin F[A]\land a=a_0)}{%
    \FormulaRefAuto{P \lor Q, \neg P \vdash Q}{4,8}}
  \proofstep{1,2,3}{a=a_0}{\rAEb{9}}
\end{tabproof}

\FormulaThmDelta{%
  F\colon A\inj B \dsep (b,a)\in \Gsurjfrominj{F}{a_0} \dsep (b,c)\in \Gsurjfrominj{F}{a_0}\vdash a=c%
}{
  \DeltaRow{Mengen}{A\dsep B\dsep a_0\dsep a\dsep b\dsep c}
  \DeltaRow{Funktionen}{F}
}
\begin{tabproof}
  \proofstep{1}{F\colon A\inj B}{\rA}
  \proofstep{2}{(b,a)\in \Gsurjfrominj{F}{a_0}}{\rA}
  \proofstep{3}{(b,c)\in \Gsurjfrominj{F}{a_0}}{\rA}
  \proofstep{1,2}{a\in A}{%
    \FormulaRefAuto{F\colon A\inj B \dsep (b,a)\in \Gsurjfrominj{F}{a_0}\vdash a\in A}{1,2}}
  \proofstep{1,3}{c\in A}{%
    \FormulaRefAuto{F\colon A\inj B \dsep (b,a)\in \Gsurjfrominj{F}{a_0}\vdash a\in A}{1,3}}
  \proofstep{}{b\in F[A]\lor b\notin F[A]}{\FormulaRefAuto{P \lor \neg P}}
  \proofstep{7}{b\in F[A]}{\rA}
  \proofstep{1,2,7}{F(a)=b}{\FormulaRefAuto{F\colon A\inj B \dsep b\in F[A] \dsep (b,a)\in \Gsurjfrominj{F}{a_0}\vdash F(a)=b}{1,7,2}}
  \proofstep{1,3,7}{F(c)=b}{\FormulaRefAuto{F\colon A\inj B \dsep b\in F[A] \dsep (b,a)\in \Gsurjfrominj{F}{a_0}\vdash F(a)=b}{1,7,3}}
  \proofstep{1,2,3,7}{F(a)=F(c)}{\FormulaRefAuto{a=b,\, c=b \vdash a=c}{8,9}}
  \proofstep{1,2,3,7}{a=c}{%
    \FormulaRefAuto{F\colon A\inj B\dsep x\in A\dsep y\in A\dsep F(x)=F(y)\vdash x=y}{1,4,5,10}}
  \proofstep{12}{b\notin F[A]}{\rA}
  \proofstep{1,2,12}{a=a_0}{\FormulaRefAuto{F\colon A\inj B \dsep b\notin F[A] \dsep (b,a)\in \Gsurjfrominj{F}{a_0}\vdash a=a_0}{1,12,2}}
  \proofstep{1,3,12}{c=a_0}{\FormulaRefAuto{F\colon A\inj B \dsep b\notin F[A] \dsep (b,a)\in \Gsurjfrominj{F}{a_0}\vdash a=a_0}{1,12,3}}
  \proofstep{1,2,3,12}{a=c}{\FormulaRefAuto{a=b,\, c=b \vdash a=c}{13,14}}
  \proofstep{1,2,3}{a=c}{\rOE{6,7,11,12,15}}
\end{tabproof}

\FormulaThmDeltaR{%
  \begin{gathered}[t]
    F\colon A\inj B \dsep a_0\in A \dsep b\in B \vdash{}\\
    \exists a\in A\,\bigl((b,a)\in \Gsurjfrominj{F}{a_0}\bigr)
  \end{gathered}
}{%
  F\colon A\inj B \dsep a_0\in A \dsep b\in B \vdash \exists a\in A\,\bigl((b,a)\in \Gsurjfrominj{F}{a_0}\bigr)
}{%
  \DeltaRow{Mengen}{A\dsep B\dsep a_0\dsep a\dsep b\dsep x}
  \DeltaRow{Funktionen}{F}
}
\begin{tabproofsplitwide}
  \proofpartwideindR[Fall \(b\in F[A]\)]{%
    F\colon A\inj B \dsep a_0\in A \dsep b\in B \dsep b\in F[A]\vdash
    \exists a\in A\,\bigl((b,a)\in \Gsurjfrominj{F}{a_0}\bigr)
  }{%
    F\colon A\inj B \dsep a_0\in A \dsep b\in B \dsep b\in F[A]\vdash
    \exists a\in A\,\bigl((b,a)\in \Gsurjfrominj{F}{a_0}\bigr)
  }
    \proofstepwidestar[1]{F\colon A\inj B}{\rA}
    \proofstepwidestar[2]{a_0\in A}{\rA}
    \proofstepwidestar[3]{b\in B}{\rA}
    \proofstepwidestar[4]{b\in F[A]}{\rA}
    \proofstepwidestar[1]{F\colon A\to B}{\FormulaRefAuto{Injektive Funktion}{1}}
    \proofstepwidestar[1,4]{\exists x\in A\,b=F(x)}{%
      \FormulaRefAuto{F\colon A\to B\dsep y\in F[A]\vdash \exists x\in A\,y=F(x)}{5,4}}

    \proofstepwidestar[7]{x\in A\land b=F(x)}{\rA}
    \proofstepwidestar[7]{x\in A}{\rAEa{7}}
    \proofstepwidestar[7]{b=F(x)}{\rAEb{7}}
    \proofstepwidestar[7]{F(x)=b}{\FormulaRefAuto{a=b\vdash b=a}{9}}
    \proofstepwidestar[4,7]{b\in F[A]\land F(x)=b}{\rAI{4,10}}
    \proofstepwidestar[4,7]{%
      \begin{gathered}[t]
        (b\in F[A]\land F(x)=b)\lor\\
        (b\notin F[A]\land x=a_0)
      \end{gathered}
    }{\rOIa{11}}
    \proofstepwidestar[1,3,7]{(b,x)\in \Gsurjfrominj{F}{a_0}}{%
      \FormulaRefAuto{F\colon A\inj B \dsep b\in B \dsep a\in A \dsep \bigl((b\in F[A]\land F(a)=b)\lor (b\notin F[A]\land a=a_0)\bigr)\vdash (b,a)\in \Gsurjfrominj{F}{a_0}}{1,3,8,12}}
    \proofstepwidestar[3,4,7]{\exists a\in A\,((b,a)\in \Gsurjfrominj{F}{a_0})}{\rEI{\rAI{8,13}}}

    \proofstepwidestar[1,3,4]{\exists a\in A\,((b,a)\in \Gsurjfrominj{F}{a_0})}{\rEE{6,7,14}}
  \closeproofpartwideindR

  \proofpartwideindR[Fall \(b\notin F[A]\)]{%
    F\colon A\inj B \dsep a_0\in A \dsep b\in B \dsep b\notin F[A]\vdash
    \exists a\in A\,\bigl((b,a)\in \Gsurjfrominj{F}{a_0}\bigr)
  }{%
    F\colon A\inj B \dsep a_0\in A \dsep b\in B \dsep b\notin F[A]\vdash
    \exists a\in A\,\bigl((b,a)\in \Gsurjfrominj{F}{a_0}\bigr)
  }
    \proofstepwidestar[1]{F\colon A\inj B}{\rA}
    \proofstepwidestar[2]{a_0\in A}{\rA}
    \proofstepwidestar[3]{b\in B}{\rA}
    \proofstepwidestar[4]{b\notin F[A]}{\rA}
    \proofstepwidestar[]{a_0=a_0}{\rII}
    \proofstepwidestar[2,4]{b\notin F[A]\land a_0=a_0}{\rAI{4,5}}
    \proofstepwidestar[2,4]{%
      \begin{gathered}[t]
        (b\in F[A]\land F(a_0)=b)\lor\\
        (b\notin F[A]\land a_0=a_0)
      \end{gathered}
    }{\rOIb{6}}
    \proofstepwidestar[1,2,3,4]{(b,a_0)\in \Gsurjfrominj{F}{a_0}}{%
      \FormulaRefAuto{F\colon A\inj B \dsep b\in B \dsep a\in A \dsep \bigl((b\in F[A]\land F(a)=b)\lor (b\notin F[A]\land a=a_0)\bigr)\vdash (b,a)\in \Gsurjfrominj{F}{a_0}}{1,3,2,7}}
    \proofstepwidestar[2,4]{\exists a\in A\,((b,a)\in \Gsurjfrominj{F}{a_0})}{\rEI{\rAI{2,8}}}
  \closeproofpartwideindR

  \proofpartwide{Schluss}
    \proofstepwidestar[1]{F\colon A\inj B}{\rA}
    \proofstepwidestar[2]{a_0\in A}{\rA}
    \proofstepwidestar[3]{b\in B}{\rA}
    \proofstepwidestar[]{b\in F[A]\lor b\notin F[A]}{\FormulaRefAuto{P \lor \neg P}}

    \proofstepwidestar[5]{b\in F[A]}{\rA}
    \proofstepwidestar[1,2,3,5]{\exists a\in A\,((b,a)\in \Gsurjfrominj{F}{a_0})}{%
      \FormulaRefAuto{F\colon A\inj B \dsep a_0\in A \dsep b\in B \dsep b\in F[A]\vdash \exists a\in A\,\bigl((b,a)\in \Gsurjfrominj{F}{a_0}\bigr)}{1,2,3,5}}

    \proofstepwidestar[7]{b\notin F[A]}{\rA}
    \proofstepwidestar[1,2,3,7]{\exists a\in A\,((b,a)\in \Gsurjfrominj{F}{a_0})}{%
      \FormulaRefAuto{F\colon A\inj B \dsep a_0\in A \dsep b\in B \dsep b\notin F[A]\vdash \exists a\in A\,\bigl((b,a)\in \Gsurjfrominj{F}{a_0}\bigr)}{1,2,3,7}}

    \proofstepwidestar[1,2,3]{\exists a\in A\,((b,a)\in \Gsurjfrominj{F}{a_0})}{\rOE{4,5,6,7,8}}
  \closeproofpartwide
\end{tabproofsplitwide}

\FormulaThmDelta{%
  F\colon A\inj B \dsep a_0\in A \vdash \Gsurjfrominj{F}{a_0}\colon B\to A%
}{
  \DeltaRow{Mengen}{A\dsep B\dsep a_0\dsep a\dsep b\dsep c}
  \DeltaRow{Funktionen}{F}
}
\begin{tabproof}
  \proofstep{1}{F\colon A\inj B}{\rA}
  \proofstep{2}{a_0\in A}{\rA}
  \proofstep{1}{\Gsurjfrominj{F}{a_0}\subseteq B\times A}{%
    \FormulaRefAuto{F\colon A\inj B \vdash \Gsurjfrominj{F}{a_0}\subseteq B\times A}{1}}
  \proofstep{4}{b\in B}{\rA}
  \proofstep{1,2,4}{\exists a\in A\,((b,a)\in \Gsurjfrominj{F}{a_0})}{%
    \FormulaRefAuto{F\colon A\inj B \dsep a_0\in A \dsep b\in B \vdash \exists a\in A\,\bigl((b,a)\in \Gsurjfrominj{F}{a_0}\bigr)}{1,2,4}}
  \proofstep{1,2}{b\in B\rightarrow \exists a\in A\,((b,a)\in \Gsurjfrominj{F}{a_0})}{\rRI{4,5}}
  \proofstep{7}{(b,a)\in \Gsurjfrominj{F}{a_0}\land (b,c)\in \Gsurjfrominj{F}{a_0}}{\rA}
  \proofstep{7}{(b,a)\in \Gsurjfrominj{F}{a_0}}{\rAEa{7}}
  \proofstep{7}{(b,c)\in \Gsurjfrominj{F}{a_0}}{\rAEb{7}}
  \proofstep{1,7}{a=c}{%
    \FormulaRefAuto{F\colon A\inj B \dsep (b,a)\in \Gsurjfrominj{F}{a_0} \dsep (b,c)\in \Gsurjfrominj{F}{a_0}\vdash a=c}{1,8,9}}
  \proofstep{1}{((b,a)\in \Gsurjfrominj{F}{a_0}\land (b,c)\in \Gsurjfrominj{F}{a_0})\rightarrow a=c}{\rRI{7,10}}
  \proofstep{1,2}{\Gsurjfrominj{F}{a_0}\colon B\to A}{\FormulaRefAuto{Funktion}{3,6,11}}
\end{tabproof}

\FormulaThmDelta{%
  F\colon A\inj B \dsep a_0\in A \vdash \Gsurjfrominj{F}{a_0}\colon B\sur A%
}{
  \DeltaRow{Mengen}{A\dsep B\dsep a_0\dsep b\dsep y\dsep x}
  \DeltaRow{Funktionen}{F}
}
\begin{tabproof}
  \proofstep{1}{F\colon A\inj B}{\rA}
  \proofstep{2}{a_0\in A}{\rA}
  \proofstep{1}{F\colon A\to B}{\FormulaRefAuto{Injektive Funktion}{1}}
  \proofstep{1,2}{\Gsurjfrominj{F}{a_0}\colon B\to A}{%
    \FormulaRefAuto{F\colon A\inj B \dsep a_0\in A \vdash \Gsurjfrominj{F}{a_0}\colon B\to A}{1,2}}
  \proofstep{5}{y\in A}{\rA}
  \proofstep{1,5}{(y,F(y))\in F}{\FormulaRefAuto{x\in A\vdash (x,F(x))\in F}{5}}
  \proofstep{1,5}{F(y)\in B}{\FormulaRefAuto{(x,y)\in F\vdash y\in B}{6}}
  \proofstep{}{F(y)=F(y)}{\rII}
  \proofstep{5}{y\in A\land F(y)=F(y)}{\rAI{5,8}}
  \proofstep{5}{\exists x\in A\,F(y)=F(x)}{\rEI{9}}
  \proofstep{1,5}{F(y)\in F[A]}{%
    \FormulaRefAuto{F\colon A\to B\dsep \exists x\in A\,y=F(x)\vdash y\in F[A]}{3,10}}
  \proofstep{1,5}{F(y)\in F[A]\land F(y)=F(y)}{\rAI{11,8}}
  \proofstep{1,5}{%
    \begin{aligned}[t]
      &(F(y)\in F[A]\land F(y)=F(y))\\
      &\lor (F(y)\notin F[A]\land y=a_0)
    \end{aligned}
  }{\rOIa{12}}
  \proofstep{1,5}{(F(y),y)\in \Gsurjfrominj{F}{a_0}}{%
    \FormulaRefAuto{F\colon A\inj B \dsep b\in B \dsep a\in A \dsep \bigl((b\in F[A]\land F(a)=b)\lor (b\notin F[A]\land a=a_0)\bigr)
    \vdash (b,a)\in \Gsurjfrominj{F}{a_0}}{1,7,5,13}}
  \proofstep{1,5}{\Gsurjfrominj{F}{a_0}(F(y))=y}{%
    \FormulaRefAuto{x\in A\dsep (x,y)\in F\vdash F(x)=y}{7,14}}
  \proofstep{1,5}{\exists b\in B\,\bigl(\Gsurjfrominj{F}{a_0}(b)=y\bigr)}{\rEI{\rAI{7,15}}}
  \proofstep{1}{y\in A\rightarrow \exists b\in B\,\bigl(\Gsurjfrominj{F}{a_0}(b)=y\bigr)}{\rRI{5,16}}
  \proofstep{1,2}{\Gsurjfrominj{F}{a_0}\colon B\sur A}{\FormulaRefAuto{Surjektive Funktion}{4,17}}
\end{tabproof}

\FormulaThmDelta{%
  F\colon A\inj B \dsep A\neq \varnothing \vdash \exists G\colon B\sur A%
}{
  \DeltaRow{Mengen}{A\dsep B}
  \DeltaRow{Funktionen}{F\dsep G}
}
\begin{tabproof}
  \proofstep{1}{F\colon A\inj B}{\rA}
  \proofstep{2}{A\neq \varnothing}{\rA}
  \proofstep{2}{\exists a_0\,(a_0\in A)}{\FormulaRefAuto{A\neq\varnothing \eqvdash \exists x\,(x \in A)}{2}}
  \proofstep{4}{a_0\in A}{\rA}
  \proofstep{1,4}{\Gsurjfrominj{F}{a_0}\colon B\sur A}{\FormulaRefAuto{F\colon A\inj B \dsep a_0\in A \vdash \Gsurjfrominj{F}{a_0}\colon B\sur A}{1,4}}
  \proofstep{1,4}{\exists G\colon B\sur A}{\rEI{5}}
  \proofstep{1,2}{\exists G\colon B\sur A}{\rEE{3,4,6}}
\end{tabproof}

\FormulaThmDelta[Linksinverses Verhalten von \(\Gsurjfrominj{F}{a_0}\) an \(y\)]{%
  F\colon A\inj B \dsep a_0\in A \dsep y\in A \vdash \Gsurjfrominj{F}{a_0}(F(y))=y
}{%
  \DeltaRow{Mengen}{A\dsep B\dsep a_0\dsep y\dsep x}
  \DeltaRow{Funktionen}{F}
}
\begin{tabproof}
  \proofstep{1}{F\colon A\inj B}{\rA}
  \proofstep{2}{a_0\in A}{\rA}
  \proofstep{3}{y\in A}{\rA}
  \proofstep{1}{F\colon A\to B}{\FormulaRefAuto{Injektive Funktion}{1}}
  \proofstep{1,2}{\Gsurjfrominj{F}{a_0}\colon B\to A}{%
    \FormulaRefAuto{F\colon A\inj B \dsep a_0\in A \vdash \Gsurjfrominj{F}{a_0}\colon B\to A}{1,2}}
  \proofstep{1,3}{F(y)\in B}{%
    \FormulaRefAuto{F\colon A\to B\dsep x\in A\vdash F(x)\in B}{4,3}}
  \proofstep{}{F(y)=F(y)}{\rII}
  \proofstep{3}{y\in A\land F(y)=F(y)}{\rAI{3,7}}
  \proofstep{3}{\exists x\in A\,F(y)=F(x)}{\rEI{8}}
  \proofstep{1,3}{F(y)\in F[A]}{%
    \FormulaRefAuto{F\colon A\to B\dsep \exists x\in A\,z=F(x)\vdash z\in F[A]}{4,9}}
  \proofstep{1,3}{%
    \begin{gathered}[t]
      (F(y)\in F[A]\land F(y)=F(y))\lor\\
      (F(y)\notin F[A]\land y=a_0)
    \end{gathered}
  }{\rOIa{\rAI{10,7}}}
  \proofstep{1,3}{(F(y),y)\in \Gsurjfrominj{F}{a_0}}{%
    \FormulaRefAuto{F\colon A\inj B \dsep b\in B \dsep a\in A \dsep \bigl((b\in F[A]\land F(a)=b)\lor (b\notin F[A]\land a=a_0)\bigr)\vdash (b,a)\in \Gsurjfrominj{F}{a_0}}{1,6,3,11}}
  \proofstep{1,2,3}{\Gsurjfrominj{F}{a_0}(F(y))=y}{%
    \FormulaRefAuto{x\in A\dsep (x,y)\in F\vdash F(x)=y}{5,6,12}}
\end{tabproof}

\FormulaThmDelta[Dualer Hilfssatz bei Nichtleere]{%
  F\colon A\inj B \dsep A\neq \varnothing \vdash \exists H\colon B\sur A\,\forall y\in A\,\bigl(H(F(y))=y\bigr)
}{%
  \DeltaRow{Mengen}{A\dsep B\dsep a_0\dsep y}
  \DeltaRow{Funktionen}{F\dsep H}
}
\begin{tabproof}
  \proofstep{1}{F\colon A\inj B}{\rA}
  \proofstep{2}{A\neq \varnothing}{\rA}
  \proofstep{2}{\exists a_0\,(a_0\in A)}{\FormulaRefAuto{A\neq\varnothing \eqvdash \exists x\,(x \in A)}{2}}
  \proofstep{4}{a_0\in A}{\rA}
  \proofstep{1,4}{\Gsurjfrominj{F}{a_0}\colon B\sur A}{%
    \FormulaRefAuto{F\colon A\inj B \dsep a_0\in A \vdash \Gsurjfrominj{F}{a_0}\colon B\sur A}{1,4}}
  \proofstep{6}{y\in A}{\rA}
  \proofstep{1,4,6}{\Gsurjfrominj{F}{a_0}(F(y))=y}{%
    \FormulaRefAuto{F\colon A\inj B \dsep a_0\in A \dsep y\in A \vdash \Gsurjfrominj{F}{a_0}(F(y))=y}{1,4,6}}
  \proofstep{1,4}{\forall y\in A\,\bigl(\Gsurjfrominj{F}{a_0}(F(y))=y\bigr)}{\rUI{\rRI{6,7}}}
  \proofstep{1,4}{%
    \Gsurjfrominj{F}{a_0}\colon B\sur A \land
    \forall y\in A\,\bigl(\Gsurjfrominj{F}{a_0}(F(y))=y\bigr)
  }{\rAI{5,8}}
  \proofstep{1,4}{\exists H\colon B\sur A\,\forall y\in A\,\bigl(H(F(y))=y\bigr)}{\rEI{9}}
  \proofstep{1,2}{\exists H\colon B\sur A\,\forall y\in A\,\bigl(H(F(y))=y\bigr)}{\rEE{3,4,10}}
\end{tabproof}

\subsubsection{Urbildmengen unter surjektiven Funktionen}

\FormulaThmDelta[Urbilder liegen im Definitionsbereich surjektiver Funktionen]{%
  C\subseteq B\dsep F\colon A\sur B \vdash F^{-1}[C]\subseteq A
}{
  \DeltaRow{Mengen}{A\dsep B\dsep C}
  \DeltaRow{Funktionen}{F}
}
\begin{tabproof}
  \proofstep{1}{C\subseteq B}{\rA}
  \proofstep{2}{F\colon A\sur B}{\rA}
  \proofstep{2}{F\colon A\to B}{\FormulaRefAuto{Surjektive Funktion}{2}}
  \proofstep{1,2}{F^{-1}[C]\subseteq A}{%
    \FormulaRefAuto{F\colon A\to B\dsep C\subseteq B \vdash F^{-1}[C]\subseteq A}{3,1}}
\end{tabproof}

\subsection{Bijektivität}

\FormulaDefDeltaK[Begriff der bijektiven Funktion]{F\colon A\bij B}{Bijektive Funktion}{
  \DeltaRow{Mengen}{A\dsep B\dsep x\dsep y}
  \DeltaPrem{Funktionen}{F\colon A\to B}[\FormulaRefAuto{Funktion}]
%
  \DeltaRow{\textbf{Axiome}}{}
  %
  % — Injektivität —
  \DeltaRow{Injektivität}
           {F\colon A\inj B\dsep x,y\in A\dsep F(x)=F(y)\vdash x=y}
           [\FormulaRefAuto{F\colon A\inj B\dsep x\in A\dsep y\in A\dsep F(x)=F(y)\vdash x=y}]
  %
  % — Surjekvität —
  \DeltaRow{Surjektivität}
           {y\in B\vdash \exists x\in A\; F(x)=y}
           [\FormulaRefAuto{y\in B\vdash \exists x\in A\; F(x)=y}]
  \DeltaRow{\textbf{Neue Symbole}}{}
  \DeltaPrem{Bijektive Funktionen}{ }
}

\FormulaThmDelta{F\colon A\bij B\dsep x\in A\vdash F(x)\in B}{
\DeltaRow{Mengen}{A\dsep B}
}
\begin{tabproof}
  \proofstep{1}{F\colon A\bij B}{\rA}
  \proofstep{2}{x\in A}{\rA}
  \proofstep{1}{F\colon A\to B}{\FormulaRefAuto{Bijektive Funktion}{1}}
  \proofstep{1,2}{F(x)\in B}{%
    \FormulaRefAuto{F\colon A\to B\dsep x\in A\vdash F(x)\in B}{3,2}}
\end{tabproof}

\FormulaThmDelta[Elemente einer Teilmenge liegen im Bild der Teilmenge unter bijektiven Funktionen]{%
  C\subseteq A\dsep F\colon A\bij B\dsep x\in C\vdash F(x)\in F[C]
}{
\DeltaRow{Mengen}{A\dsep B\dsep C\dsep x}
\DeltaRow{Funktionen}{F}
}
\begin{tabproof}
  \proofstep{1}{C\subseteq A}{\rA}
  \proofstep{2}{F\colon A\bij B}{\rA}
  \proofstep{3}{x\in C}{\rA}
  \proofstep{1,2,3}{F(x)\in F[C]}{%
    \FormulaRefAuto{C\subseteq A\dsep F\colon A\to B\dsep x\in C\vdash F(x)\in F[C]}{%
      1,\FormulaRefAuto{Bijektive Funktion}{2},3}}
\end{tabproof}

\FormulaThmDelta{F\colon A\inj B\dsep F\colon A\sur B \vdash F\colon A\bij B}{
\DeltaRow{Mengen}{A\dsep B}
\DeltaPrem{Funktionen}{F\colon A\to B}
}
\begin{tabproof}
    \proofstep{1}{F\colon A\inj B}{\rA}
    \proofstep{2}{F\colon A\sur B}{\rA}
    \proofstep{1}{\forall x,y\in A\,(F(x)=F(y)\rightarrow x=y)}{%
      \FormulaRefAuto{F\colon A\inj B\dsep x\in A\dsep y\in A\dsep F(x)=F(y)\vdash x=y}{1}}
    \proofstep{1}{\forall y\in B\,\exists x\in A\,(F(x)=y)}{\FormulaRefAuto{y\in B\vdash \exists x\in A\; F(x)=y}{2}}
    \proofstep{1,2}{F\colon A\bij B}{\FormulaRefAuto{Bijektive Funktion}{3,4}}
\end{tabproof}

\FormulaThmDelta{F\colon A\bij B \vdash F\colon A\inj B}{
\DeltaRow{Mengen}{A\dsep B}
}
\begin{tabproof}
    \proofstep{1}{F\colon A\bij B}{\rA}
    \proofstep{1}{F\colon A\to B}{\FormulaRefAuto{Bijektive Funktion}{1}}
    \proofstep{1}{\forall x,y\in A\, (F(x)=F(y)\rightarrow x=y)}{\FormulaRefAuto{Bijektive Funktion}{1}}
    \proofstep{1}{F\colon A\inj B}{\FormulaRefAuto{Injektive Funktion}{2,3}}
\end{tabproof}

\FormulaThmDelta{F\colon A\bij B \vdash F\colon A\sur B}{
\DeltaRow{Mengen}{A\dsep B}
}
\begin{tabproof}
    \proofstep{1}{F\colon A\bij B}{\rA}
    \proofstep{1}{F\colon A\to B}{\FormulaRefAuto{Bijektive Funktion}{1}}
    \proofstep{1}{\forall y\in B\exists x\in A\; F(x)=y}{\FormulaRefAuto{Bijektive Funktion}{1}}
    \proofstep{1}{F\colon A\sur B}{\FormulaRefAuto{Surjektive Funktion}{2,3}}
\end{tabproof}

\FormulaThmDelta{F\colon A\inj B\land F\colon A\sur B \eqvdash F\colon A\bij B }{
\DeltaRow{Mengen}{A\dsep B}
}
\begin{tabproofsplit}
    \proofpart{\(\vdash\)}
    \proofstep{1}{F\colon A\inj B\land F\colon A\sur B}{\rA}
    \proofstep{1}{F\colon A\inj B}{\rAEa{1}}
    \proofstep{1}{F\colon A\sur B}{\rAEb{2}}
    \proofstep{1}{F\colon A\bij B}{\FormulaRefAuto{F\colon A\inj B\dsep F\colon A\sur B \vdash F\colon A\bij B}{2,3}}
    \closeproofpart
    \proofpart{\(\dashv\)}
    \proofstep{1}{F\colon A\bij B}{\rA}
    \proofstep{1}{F\colon A\inj B}{\FormulaRefAuto{F\colon A\bij B \vdash F\colon A\inj B}{1}}
    \proofstep{1}{F\colon A\sur B}{\FormulaRefAuto{F\colon A\bij B \vdash F\colon A\sur B}{1}}
    \proofstep{1}{F\colon A\inj B\land F\colon A\sur B}{\rAI{2,3}}
    \closeproofpart
\end{tabproofsplit}

\FormulaThmDelta[Eindeutigkeit des Urbildes bei Bijektivität]{%
F\colon A\bij B \dsep y\in B \vdash \exists! x\in A\; F(x)=y
}{
  \DeltaRow{Mengen}{A \dsep B \dsep F\dsep x \dsep y \dsep z}
}
\begin{tabproof}
  % Annahmen
  \proofstep{1}{F\colon A\bij B}{\rA}
  \proofstep{2}{y\in B}{\rA}

  % Existenz eines Urbildes aus Surjektivität
  \proofstep{1,2}{\exists x\in A\,F(x)=y}{%
    \FormulaRefAuto{y\in B\vdash \exists x\in A\; F(x)=y}{1,2}}

  % Uniqueness-Teil: angenommen zwei Urbilder x,z
  \proofstep{3}{x\in A \land F(x)=y}{\rA}
  \proofstep{4}{z\in A \land F(z)=y}{\rA}

  % Aus den Konjunktionen die Einzelteile holen
  \proofstep{3}{x\in A}{\rAEa{3}}
  \proofstep{3}{F(x)=y}{\rAEb{3}}
  \proofstep{4}{z\in A}{\rAEa{4}}
  \proofstep{4}{F(z)=y}{\rAEb{4}}

  % F(x)=y und F(z)=y ⇒ F(x)=F(z) (über Gleichheitseigenschaften)
  \proofstep{4}{y=F(z)}{\FormulaRefAuto{a=b \vdash b=a}{9}}
  \proofstep{3,4}{F(x)=F(z)}{\FormulaRefAuto{a = b,\, b = c \vdash a = c}{7,10}}

  % Injektivität verwenden: aus x,z∈A und F(x)=F(z) folgt x=z
  \proofstep{1,3,4}{x=z}{%
    \FormulaRefAuto{F\colon A\inj B\dsep x\in A\dsep y\in A\dsep F(x)=F(y)\vdash x=y}{1,6,8,11}}

  % ∃!-Einführung: aus Existenz (3) und Eindeutigkeit (Subbeweis 3–11)
  \proofstep{1,2}{\exists! x\in A\,F(x)=y}{\UEI{3,4,5,12}}
\end{tabproof}

\subsubsection{Urbildmengen unter bijektiven Funktionen}

\FormulaThmDelta[Urbilder liegen im Definitionsbereich bijektiver Funktionen]{%
  C\subseteq B\dsep F\colon A\bij B \vdash F^{-1}[C]\subseteq A
}{
  \DeltaRow{Mengen}{A\dsep B\dsep C}
  \DeltaRow{Funktionen}{F}
}
\begin{tabproof}
  \proofstep{1}{C\subseteq B}{\rA}
  \proofstep{2}{F\colon A\bij B}{\rA}
  \proofstep{2}{F\colon A\to B}{\FormulaRefAuto{Bijektive Funktion}{2}}
  \proofstep{1,2}{F^{-1}[C]\subseteq A}{%
    \FormulaRefAuto{F\colon A\to B\dsep C\subseteq B \vdash F^{-1}[C]\subseteq A}{3,1}}
\end{tabproof}

\section{Beispiele von Funktionen}

\subsection{Die Identitätsfunktion}

\subsubsection{Definition der Identitätsfunktion}

\FormulaDefDelta[Funktionsvorschrift für die Identitätsfunktion auf \(A\)]
{x\in A \vdash t_{\Id_A}(x) \coloneqq x}
{
\DeltaRow{Mengen}{A\dsep x}
\DeltaRow{Funktionensymbol}{t_{\Id_A}}
}


% — Typisierungsaxiom / Wertebereich —
\FormulaThmDelta[Typisierungsaxiom]
{x\in A \vdash t_{\Id_A}(x)\in A}
{
  \DeltaRow{Mengen}{A\dsep x}
\DeltaRow{Funktionensymbol}{t_{\Id_A}}
}
\begin{tabproof}
    \proofstep{1}{x\in A}{\rA}
    \proofstep{1}{t_{\Id_A}(x)=x}{\FormulaRefAuto{x\in A \vdash t_{\Id_A}(x) \coloneqq x}{1}}
    \proofstep{1}{t_{\Id_A}(x)\in A}{\rIE{2,1}}
\end{tabproof}


\FormulaThmDelta[Identität als Funktion]{\exists! F\colon A\to A\,\forall x\in A\,F(x) = t_{\Id_A}(x)}{
  \DeltaRow{Mengen}{A\dsep x}
  \DeltaRow{Funktionensymbol}{t_{\Id_A}}
}
\begin{tabproof}
    \proofstep{}{\forall x\in A\,t_{\Id_A}(x)= x}{\FormulaRefAuto{x\in A \vdash t_{\Id_A}(x) \coloneqq x}}
    \proofstep{}{\forall x\in A\,t_{\Id_A}(x)\in A}{\FormulaRefAuto{x\in A \vdash t_{\Id_A}(x)\in A}}
    \proofstep{}{\exists! F\colon A\to A\,\forall x\in A\,F(x) = t_{\Id_A}(x)}{\FormulaRefAuto{x\in A\vdash t(x)\coloneq y\dsep x\in A\vdash t(x)\in B\exists! F\colon A\to B\,\forall x\in A\,F(x) = t(x)}{1,2}}
\end{tabproof}

\FormulaDefDelta[Identitätsfunktion auf \(A\)]
{\Id_A\coloneqq\iota F\left(F\colon A\to A \land \forall x\in A\, F(x)=t_{\Id_A}(x)\right)}
{
\DeltaRow{Mengen}{A}
\DeltaRow{Funktionensymbol}{t_{\Id_A}}
}

\FormulaThmDeltaK[Identitätsfunktion]
{\Id_A\colon A\to A}{IdA}
{
\DeltaRow{Mengen}{A}
}
\begin{tabproof}
    \proofstep{}{\Id_A\colon A\to A}{\rAEa{\FormulaRefAuto{\Id_A\coloneqq\iota F\left(F\colon A\to A \land \forall x\in A\, F(x)=t_{\Id_A}(x)\right)}}}
\end{tabproof}


\FormulaThmDelta
{\TotRel{\Id_A,A,A}}
{
\DeltaRow{Mengen}{A}
}
\begin{tabproof}
    \proofstep{}{\Id_A\colon A\to A}{\FormulaRefAuto{IdA}}
    \proofstep{}{\TotRel{\Id_A,A,A}}{\FormulaRefAuto{F\colon A\to B\vdash\TotRel{F,A,B}}{1}}
\end{tabproof}


\FormulaThmDelta[Identitätsfunktion auf \(A\)]
{x\in A\vdash \Id_A(x)=x}
{
\DeltaRow{Mengen}{A\dsep x}
}
\begin{tabproof}
    \proofstep{1}{x\in A}{\rA}
    \proofstep{}{\forall x\in A\, \Id_A(x)=t_{\Id_A}(x)}{\rAEb{\FormulaRefAuto{\Id_A\coloneqq\iota F\left(F\colon A\to A \land \forall x\in A\, F(x)=t_{\Id_A}(x)\right)}}}
    \proofstep{1}{\Id_A(x)=t_{\Id_A}(x)}{\rRE{1,\rUE{2}}}
    \proofstep{1}{t_{\Id_A}(x)=x}{\FormulaRefAuto{x\in A \vdash t_{\Id_A}(x) \coloneqq x}{1}}
    \proofstep{1}{\Id_A(x)=x}{\FormulaRefAuto{a = b,\, b = c \vdash a = c}{3,4}}
\end{tabproof}

\FormulaThmDelta{%
  x\in A\vdash x=\Id_A(x)%
}{
  \DeltaRow{Mengen}{A\dsep x}
}
\begin{tabproof}
  \proofstep{1}{x\in A}{\rA}
  \proofstep{1}{\Id_A(x)=x}{%
    \FormulaRefAuto{x\in A\vdash \Id_A(x)=x}}
  \proofstep{1}{x=\Id_A(x)}{%
    \FormulaRefAuto{a=b\vdash b=a}}
\end{tabproof}

\subsubsection{Die Bijektionseigenschaften}

\paragraph{Injektivität}
\FormulaThmDelta[Injektivität von \(\Id_A\)]{%
x\in A \dsep y\in A \dsep \Id_A(x)=\Id_A(y) \vdash x=y
}{
  \DeltaRow{Mengen}{x \dsep y \dsep A}
}
\begin{tabproof}
  \proofstep{1}{\Id_A(x)=\Id_A(y)}{\rA}
  \proofstep{2}{x\in A}{\rA}
  \proofstep{3}{y\in A}{\rA}
  \proofstep{2}{\Id_A(x)=x}{\FormulaRefAuto{x\in A\vdash \Id_A(x)=x}{2}}
  \proofstep{3}{\Id_A(y)=y}{\FormulaRefAuto{x\in A\vdash \Id_A(x)=x}{3}}
  \proofstep{1,2,3}{x=y}{\FormulaRefAuto{a = b\dsep a = c\dsep b=d\vdash c = d}{1,4,5}}
\end{tabproof}

\paragraph{Surjektivität}
\FormulaThmDelta[Surjektivität von \(\Id_A\)]{%
x\in A \vdash \exists y\in A\; \Id_A(y)=x
}{
  \DeltaRow{Mengen}{x \dsep y \dsep A}
}
\begin{tabproof}
  \proofstep{1}{x\in A}{\rA}
  \proofstep{1}{\Id_A(x)=x}{\FormulaRefAuto{x\in A\vdash \Id_A(x)=x}{1}}
  \proofstep{1}{\exists y\in A\; \Id_A(y)=x}{\rEI{\rAI{1,2}}}
\end{tabproof}

\paragraph{Bijektivität}
\FormulaThmDelta[\(\Id_A\) als bijektive Funktion]{%
\Id_A\colon A\bij A
}{
  \DeltaRow{Mengen}{A}
}
\begin{tabproof}
    \proofstep{}{\Id_A\colon A\to A}{\FormulaRefAuto{IdA}}
    \proofstep{}{\forall x,y\in A\,(\Id_A(x)=\Id_A(y) \rightarrow x=y)}{\FormulaRefAuto{x\in A \dsep y\in A \dsep \Id_A(x)=\Id_A(y) \vdash x=y}}
    \proofstep{}{\forall x\in A\exists y\in A\,\Id_A(y)=x}{\FormulaRefAuto{x\in A \vdash \exists y\in A\; \Id_A(y)=x}}
    \proofstep{}{\Id_A\colon A\bij A}{\FormulaRefAuto{Bijektive Funktion}{1,2,3}}
\end{tabproof}

\subsection{Die Projektionen auf die Komponenten}

\subsubsection{Projektion auf die erste Komponente}

\FormulaDefDelta[Funktionsvorschrift für die Projektion auf die erste Komponente]
{(x,y)\in A\times B \vdash t_{\pi_1}(x,y) \coloneqq x}
{
  \DeltaRow{Mengen}{A\dsep B\dsep x\dsep y}
  \DeltaRow{Funktionensymbol}{t_{\pi_1}}
}

% — Typisierungsaxiom / Wertebereich —
\FormulaThmDelta[Typisierungsaxiom für \(\pi_1\)]
{(x,y)\in A\times B \vdash t_{\pi_1}(x,y)\in A}
{
  \DeltaRow{Mengen}{A\dsep B\dsep x\dsep y}
  \DeltaRow{Funktionensymbol}{t_{\pi_1}}
}
\begin{tabproof}
  \proofstep{1}{(x,y)\in A\times B}{\rA}
  \proofstep{1}{x\in A \land y\in B}{%
    \FormulaRefAuto{(a,b)\in A\times B\eqvdash a\in A\land b\in B}{1}}
  \proofstep{1}{x\in A}{\rAEa{2}}
  \proofstep{1}{t_{\pi_1}(x,y)=x}{%
    \FormulaRefAuto{(x,y)\in A\times B \vdash t_{\pi_1}(x,y) \coloneqq x}{1}}
  \proofstep{1}{t_{\pi_1}(x,y)\in A}{\rIE{4,3}}
\end{tabproof}



\FormulaThmDelta[Projektion auf die erste Komponente als Funktion]{%
\exists! F\colon A\times B\to A\,\forall x\in A\,\forall y\in B\,F(x,y) = t_{\pi_1}(x,y)%
}{
  \DeltaRow{Mengen}{A\dsep B\dsep x\dsep y}
  \DeltaRow{Funktionensymbol}{t_{\pi_1}}
}
\begin{tabproof}
  \proofstep{}{\forall (x,y)\in A\times B\,t_{\pi_1}(x,y)= x}{\FormulaRefAuto{(x,y)\in A\times B \vdash t_{\pi_1}(x,y) \coloneqq x}}

  \proofstep{}{\forall (x,y)\in A\times B\,t_{\pi_1}(x,y)\in A}{\FormulaRefAuto{(x,y)\in A\times B \vdash t_{\pi_1}(x,y)\in A}}
    
  \proofstep{}{\exists! F\colon A\times B\to A\,\forall x\in A\,\forall y\in B\,F(x,y) = t_{\pi_1}(x,y)}{%
    \FormulaRefAuto{x\in A\vdash t(x)\coloneq y\dsep x\in A\vdash t(x)\in B\exists! F\colon A\to B\,\forall x\in A\,F(x) = t(x)}{1,2}}
\end{tabproof}

\FormulaDefDelta[Projektion auf die erste Komponente]
{\pi_1\coloneqq\iota F\left(F\colon A\times B\to A \land \forall x\in A\,\forall y\in B\, F(x,y)=t_{\pi_1}(x,y)\right)}
{
  \DeltaRow{Mengen}{A\dsep B}
  \DeltaRow{Funktionensymbol}{t_{\pi_1}}
}

\FormulaThmDelta[Projektion auf die erste Komponente]
{\pi_1\colon A\times B\to A}
{
  \DeltaRow{Mengen}{A\dsep B}
}
\begin{tabproof}
  \proofstep{}{\pi_1\colon A\times B\to A}{%
    \rAEa{\FormulaRefAuto{\pi_1\coloneqq\iota F\left(F\colon A\times B\to A \land \forall x\in A\,\forall y\in B\, F(x,y)=t_{\pi_1}(x,y)\right)}}}
\end{tabproof}

\FormulaThmDelta
{\TotRel{\pi_1,A\times B,A}}
{
  \DeltaRow{Mengen}{A\dsep B}
}
\begin{tabproof}
    \proofstep{}{\pi_1\colon A\times B\to A}{\FormulaRefAuto{\pi_1\colon A\times B\to A}}
    \proofstep{}{\TotRel{\pi_1,A\times B,A}}{\FormulaRefAuto{F\colon A\to B\vdash\TotRel{F,A,B}}{1}}
\end{tabproof}

\FormulaThmDelta[Grundgleichung der Projektion auf die erste Komponente]
{(x,y)\in A\times B \vdash \pi_1(x,y)=x}
{
  \DeltaRow{Mengen}{A\dsep B\dsep x\dsep y}
}
\begin{tabproof}
  \proofstep{1}{(x,y)\in A\times B}{\rA}
  \proofstep{}{%
    \forall x\in A\,\forall y\in B\, \pi_1(x,y)=t_{\pi_1}(x,y)}{%
    \rAEb{\FormulaRefAuto{\pi_1\coloneqq\iota F\left(F\colon A\times B\to A \land \forall x\in A\,\forall y\in B\, F(x,y)=t_{\pi_1}(x,y)\right)}}}
  \proofstep{1}{\pi_1(x,y)=t_{\pi_1}(x,y)}{\rRE{1,\rUE{2}}}
  \proofstep{1}{t_{\pi_1}(x,y)=x}{%
    \FormulaRefAuto{(x,y)\in A\times B \vdash t_{\pi_1}(x,y) \coloneqq x}{1}}
  \proofstep{1}{\pi_1(x,y)=x}{\FormulaRefAuto{a = b,\, b = c \vdash a = c}{3,4}}
\end{tabproof}

\FormulaThmDelta
{x\in A\dsep y\in B \vdash \pi_1(x,y)=x}
{
  \DeltaRow{Mengen}{A\dsep B\dsep x\dsep y}
  \DeltaRow{Funktionensymbol}{t_{\pi_1}}
}
\begin{tabproof}
  \proofstep{1}{x\in A}{\rA}
  \proofstep{2}{y\in B}{\rA}
  \proofstep{1,2}{(x,y)\in A\times B}{%
    \FormulaRefAuto{(a,b)\in A\times B\eqvdash a\in A\land b\in B}{1,2}}
  \proofstep{1,2}{\pi_1(x,y)=x}{%
    \FormulaRefAuto{(x,y)\in A\times B\vdash \pi_1(x,y)=x}{3}}
\end{tabproof}

\paragraph{Surjektivität}

\FormulaThmDelta[Surjektivität von \(\pi_1\)]{%
x\in A \dsep b_0\in B \vdash \exists z\in A\times B\; \pi_1(z)=x
}{
  \DeltaRow{Mengen}{A \dsep B \dsep x \dsep b_0 \dsep z}
}
\begin{tabproof}
  \proofstep{1}{x\in A}{\rA}
  \proofstep{2}{b_0\in B}{\rA}
  % Paarbildung
  \proofstep{1,2}{(x,b_0)\in A\times B}{%
    \FormulaRefAuto{(a,b)\in A\times B\eqvdash a\in A\land b\in B}{1,2}}
  % Grundgleichung der Projektion: (x,b0) in A×B ⇒ π1(x,b0)=x
  \proofstep{1,2}{\pi_1(x,b_0)=x}{%
    \FormulaRefAuto{(x,y)\in A\times B \vdash \pi_1(x,y)=x}{3}}
  % Existenzzeuge z = (x,b0) für Surjektivität
  \proofstep{1,2}{\exists z\in A\times B\; \pi_1(z)=x}{%
    \rEI{\rAI{3,4}}}
\end{tabproof}

\subsubsection{Projektion auf die zweite Komponente}

\FormulaDefDelta[Funktionsvorschrift für die Projektion auf die zweite Komponente]
{(x,y)\in A\times B \vdash t_{\pi_2}(x,y) \coloneqq y}
{
  \DeltaRow{Mengen}{A\dsep B\dsep x\dsep y}
  \DeltaRow{Funktionensymbol}{t_{\pi_2}}
}

% — Typisierungsaxiom / Wertebereich —
\FormulaThmDelta[Typisierungsaxiom für \(\pi_2\)]
{(x,y)\in A\times B \vdash t_{\pi_2}(x,y)\in B}
{
  \DeltaRow{Mengen}{A\dsep B\dsep x\dsep y}
  \DeltaRow{Funktionensymbol}{t_{\pi_2}}
}
\begin{tabproof}
  \proofstep{1}{(x,y)\in A\times B}{\rA}
  \proofstep{1}{x\in A \land y\in B}{%
    \FormulaRefAuto{(a,b)\in A\times B\eqvdash a\in A\land b\in B}{1}}
  \proofstep{1}{y\in B}{\rAEb{2}}
  \proofstep{1}{t_{\pi_2}(x,y)=y}{%
    \FormulaRefAuto{(x,y)\in A\times B \vdash t_{\pi_2}(x,y) \coloneqq y}{1}}
  \proofstep{1}{t_{\pi_2}(x,y)\in B}{\rIE{4,3}}
\end{tabproof}

\FormulaThmDelta[Projektion auf die zweite Komponente als Funktion]{%
\exists! F\colon A\times B\to B\,\forall x\in A\,\forall y\in B\,F(x,y) = t_{\pi_2}(x,y)%
}{
  \DeltaRow{Mengen}{A\dsep B\dsep x\dsep y}
  \DeltaRow{Funktionensymbol}{t_{\pi_2}}
}
\begin{tabproof}
  \proofstep{}{\forall (x,y)\in A\times B\,t_{\pi_2}(x,y)= y}{%
    \FormulaRefAuto{(x,y)\in A\times B \vdash t_{\pi_2}(x,y) \coloneqq y}}

  \proofstep{}{\forall (x,y)\in A\times B\,t_{\pi_2}(x,y)\in B}{%
    \FormulaRefAuto{(x,y)\in A\times B \vdash t_{\pi_2}(x,y)\in B}}
    
  \proofstep{}{\exists! F\colon A\times B\to B\,\forall x\in A\,\forall y\in B\,F(x,y) = t_{\pi_2}(x,y)}{%
    \FormulaRefAuto{x\in A\vdash t(x)\coloneq y\dsep x\in A\vdash t(x)\in B\exists! F\colon A\to B\,\forall x\in A\,F(x) = t(x)}{1,2}}
\end{tabproof}

\FormulaDefDelta[Projektion auf die zweite Komponente]
{\pi_2\coloneqq\iota F\left(F\colon A\times B\to B \land \forall x\in A\,\forall y\in B\, F(x,y)=t_{\pi_2}(x,y)\right)}
{
  \DeltaRow{Mengen}{A\dsep B}
  \DeltaRow{Funktionensymbol}{t_{\pi_2}}
}

\FormulaThmDelta[Projektion auf die zweite Komponente]
{\pi_2\colon A\times B\to B}
{
  \DeltaRow{Mengen}{A\dsep B}
}
\begin{tabproof}
  \proofstep{}{\pi_2\colon A\times B\to B}{%
    \rAEa{\FormulaRefAuto{\pi_2\coloneqq\iota F\left(F\colon A\times B\to B \land \forall x\in A\,\forall y\in B\, F(x,y)=t_{\pi_2}(x,y)\right)}}}
\end{tabproof}

\FormulaThmDelta
{\TotRel{\pi_2,A\times B,B}}
{
  \DeltaRow{Mengen}{A\dsep B}
}
\begin{tabproof}
    \proofstep{}{\pi_2\colon A\times B\to B}{\FormulaRefAuto{\pi_2\colon A\times B\to B}}
    \proofstep{}{\TotRel{\pi_2,A\times B,B}}{\FormulaRefAuto{F\colon A\to B\vdash\TotRel{F,A,B}}{1}}
\end{tabproof}

\FormulaThmDelta[Grundgleichung der Projektion auf die zweite Komponente]
{(x,y)\in A\times B \vdash \pi_2(x,y)=y}
{
  \DeltaRow{Mengen}{A\dsep B\dsep x\dsep y}
}
\begin{tabproof}
  \proofstep{1}{(x,y)\in A\times B}{\rA}
  \proofstep{}{%
    \forall x\in A\,\forall y\in B\, \pi_2(x,y)=t_{\pi_2}(x,y)}{%
    \rAEb{\FormulaRefAuto{\pi_2\coloneqq\iota F\left(F\colon A\times B\to B \land \forall x\in A\,\forall y\in B\, F(x,y)=t_{\pi_2}(x,y)\right)}}}
  \proofstep{1}{\pi_2(x,y)=t_{\pi_2}(x,y)}{\rRE{1,\rUE{2}}}
  \proofstep{1}{t_{\pi_2}(x,y)=y}{%
    \FormulaRefAuto{(x,y)\in A\times B \vdash t_{\pi_2}(x,y) \coloneqq y}{1}}
  \proofstep{1}{\pi_2(x,y)=y}{\FormulaRefAuto{a = b,\, b = c \vdash a = c}{3,4}}
\end{tabproof}

\FormulaThmDelta
{x\in A\dsep y\in B\,\pi_2(x,y)=y}
{
  \DeltaRow{Mengen}{A\dsep B\dsep x\dsep y}
  \DeltaRow{Funktionensymbol}{t_{\pi_2}}
}
\begin{tabproof}
  \proofstep{1}{x\in A}{\rA}
  \proofstep{2}{y\in B}{\rA}
  \proofstep{1,2}{(x,y)\in A\times B}{%
    \FormulaRefAuto{(a,b)\in A\times B\eqvdash a\in A\land b\in B}{1,2}}
  \proofstep{1,2}{\pi_2(x,y)=y}{%
    \FormulaRefAuto{(x,y)\in A\times B \vdash \pi_2(x,y)=y}{3}}
\end{tabproof}

\paragraph{Surjektivität}

\FormulaThmDelta[Surjektivität von \(\pi_2\)]{%
y\in B \dsep a_0\in A \vdash \exists z\in A\times B\; \pi_2(z)=y
}{
  \DeltaRow{Mengen}{A \dsep B \dsep y \dsep a_0 \dsep z}
}
\begin{tabproof}
  \proofstep{1}{y\in B}{\rA}
  \proofstep{2}{a_0\in A}{\rA}
  % Paarbildung
  \proofstep{1,2}{(a_0,y)\in A\times B}{%
    \FormulaRefAuto{a\in A \dsep b\in B \vdash (a,b)\in A\times B}{2,1}}
  % Grundgleichung der Projektion: (a0,y) in A×B ⇒ π2(a0,y)=y
  \proofstep{1,2}{\pi_2(a_0,y)=y}{%
    \FormulaRefAuto{(x,y)\in A\times B \vdash \pi_2(x,y)=y}{3}}
  % Existenzzeuge z = (a0,y) für Surjektivität
  \proofstep{1,2}{\exists z\in A\times B\; \pi_2(z)=y}{%
    \rEI{\rAI{3,4}}}
\end{tabproof}

\subsubsection{Eigenschaften von Projektionsabbildungen}

\FormulaThmDelta[Rekonstruktion aus Projektionen]{%
  (x,y)\in A\times B \vdash (x,y) = \bigl(\pi_1(x,y),\pi_2(x,y)\bigr)
}{
  \DeltaRow{Mengen}{A\dsep B\dsep x\dsep y}
  \DeltaRow{Projektion auf erste Komponente}{\pi_1:A\times B\to A}
  \DeltaRow{Projektion auf zweite Komponente}{\pi_2:A\times B\to B}
}
\begin{tabproofwide}
  \proofstepwidestar[1]{(x,y)\in A\times B}{\rA}
  \proofstepwidestar[1]{x=\pi_1(x,y)}{%
    \FormulaRefAuto{a=b\vdash b=a}{\FormulaRefAuto{(x,y)\in A\times B \vdash \pi_1(x,y)=x}{1}}}
  \proofstepwidestar[1]{y=\pi_2(x,y)}{%
    \FormulaRefAuto{a=b\vdash b=a}{\FormulaRefAuto{(x,y)\in A\times B \vdash \pi_2(x,y)=y}{1}}}
  \proofstepwide[1]{(x,y)}{=}{(x,y)}{\rII}
  \proofstepwide[1]{}{=}{(\pi_1(x,y),y)}{\rIE{2,4}}
  \proofstepwide[1]{}{=}{(\pi_1(x,y),\pi_2(x,y))}{\rIE{3,5}}
  \proofstepwide[1]{(x,y)}{=}{(\pi_1(x,y),\pi_2(x,y))}{\rChain{4,6}}
\end{tabproofwide}

\FormulaThmDelta[Element des Produkts als Paar]{%
  z\in A\times B \vdash z = \bigl(\pi_1(z),\pi_2(z)\bigr)
}{
  \DeltaRow{Mengen}{A\dsep B\dsep z}
  \DeltaRow{Projektion auf erste Komponente}{\pi_1:A\times B\to A}
  \DeltaRow{Projektion auf zweite Komponente}{\pi_2:A\times B\to B}
}
\begin{tabproof}
  \proofstep{1}{z\in A\times B}{\rA}
  
  \proofstep{1}{\exists x\in A\,\exists y\in B\, z=(x,y)}{%
    \FormulaRefAuto{x\in A\times B\eqvdash \exists a \in A\, \exists b \in B\,x=(a,b)}{1}}
  
  \proofstep{3}{x\in A\land y\in B\land z=(x,y)}{\rA}
  
  \proofstep{3}{x\in A}{\rAEa{3}}
  
  \proofstep{3}{y\in B}{\FormulaRefAuto{P\land Q\land R\vdash Q}{3}}
  
  \proofstep{3}{z=(x,y)}{\rAEb{3}}
  
  \proofstep{3}{(x,y)\in A\times B}{%
    \FormulaRefAuto{(a,b)\in A\times B\eqvdash a\in A\land b\in B}{4,5}}
  
  \proofstep{3}{(x,y)=\bigl(\pi_1(x,y),\pi_2(x,y)\bigr)}{%
    \FormulaRefAuto{(x,y)\in A\times B \vdash (x,y)=\bigl(\pi_1(x,y),\pi_2(x,y)\bigr)}{7}}

  \proofstep{3}{z=\bigl(\pi_1(z),\pi_2(z)\bigr)}{%
    \rIE{6,8}}
  
  \proofstep{1}{z=\bigl(\pi_1(z),\pi_2(z)\bigr)}{%
    \rEE{2,3,9}}
\end{tabproof}

\subsection{Die Inklusionsabbildung}

\subsubsection{Definition der Inklusionsabbildung}

\FormulaDefDelta[Funktionsvorschrift für die Inklusionsabbildung von \(A\) nach \(B\)]
{A\subseteq B \dsep x\in A \vdash t_{\iota_{A,B}}(x) \coloneqq x}
{
  \DeltaRow{Mengen}{A\dsep B\dsep x}
  \DeltaRow{Funktionensymbol}{t_{\iota_{A,B}}}
}

% — Typisierungsaxiom / Wertebereich —
\FormulaThmDelta[Typisierungsaxiom für \(t_{\iota_{A,B}}\)]
{A\subseteq B \dsep x\in A \vdash t_{\iota_{A,B}}(x)\in B}
{
  \DeltaRow{Mengen}{A\dsep B\dsep x}
  \DeltaRow{Funktionensymbol}{t_{\iota_{A,B}}}
}
\begin{tabproof}
  \proofstep{1}{A\subseteq B}{\rA}
  \proofstep{2}{x\in A}{\rA}
  \proofstep{1,2}{x\in B}{%
    \FormulaRefAuto{A\subseteq B, x\in A \vdash x\in B}{1,2}}
  \proofstep{1,2}{t_{\iota_{A,B}}(x)=x}{%
    \FormulaRefAuto{A\subseteq B \dsep x\in A \vdash t_{\iota_{A,B}}(x) \coloneqq x}{1,2}}
  \proofstep{1,2}{t_{\iota_{A,B}}(x)\in B}{\rIE{4,3}}
\end{tabproof}

\FormulaThmDelta[Inklusionsabbildung als Funktion]{%
A\subseteq B\vdash \exists! F\colon A\to B\,\forall x\in A\,F(x) = t_{\iota_{A,B}}(x)%
}{
  \DeltaRow{Mengen}{A\dsep B\dsep x}
  \DeltaRow{Funktionensymbol}{t_{\iota_{A,B}}}
}
\begin{tabproof}
  \proofstep{1}{A\subseteq B}{\rA}
  \proofstep{1}{\forall x\in A\,t_{\iota_{A,B}}(x)=x}{\FormulaRefAuto{A\subseteq B \dsep x\in A \vdash t_{\iota_{A,B}}(x) \coloneqq x}{1}}
  \proofstep{1}{\forall x\in A\,t_{\iota_{A,B}}(x)\in B}{\FormulaRefAuto{A\subseteq B \dsep x\in A \vdash t_{\iota_{A,B}}(x)\in B}{1}}
  \proofstep{1}{\exists! F\colon A\to B\,\forall x\in A\,F(x) = t_{\iota_{A,B}}(x)}{%
    \FormulaRefAuto{x\in A\vdash t(x)\coloneq y\dsep x\in A\vdash t(x)\in B\exists! F\colon A\to B\,\forall x\in A\,F(x) = t(x)}{2,3}}
\end{tabproof}

\FormulaDefDelta[Inklusionsabbildung von \(A\) nach \(B\)]
{A\subseteq B\vdash \iota_{A,B}\coloneqq\iota F\left(F\colon A\to B \land \forall x\in A\, F(x)=t_{\iota_{A,B}}(x)\right)}
{
  \DeltaRow{Mengen}{A\dsep B}
  \DeltaRow{Funktionensymbol}{t_{\iota_{A,B}}}
}

\FormulaThmDelta[Inklusionsabbildung von \(A\) nach \(B\)]
{A\subseteq B\vdash \iota_{A,B}\colon A\to B}
{
  \DeltaRow{Mengen}{A\dsep B}
}
\begin{tabproof}
  \proofstep{1}{A\subseteq B}{\rA}
  \proofstep{1}{\iota_{A,B}\colon A\to B}{%
    \rAEa{\FormulaRefAuto{A\subseteq B\vdash \iota_{A,B}\coloneqq\iota F\left(F\colon A\to B \land \forall x\in A\, F(x)=t_{\iota_{A,B}}(x)\right)}{1}}}
\end{tabproof}

\FormulaThmDelta
{A\subseteq B\vdash \TotRel{\iota_{A,B},A,B}}
{
  \DeltaRow{Mengen}{A\dsep B}
}
\begin{tabproof}
    \proofstep{1}{A\subseteq B}{\rA}
    \proofstep{1}{\iota_{A,B}\colon A\to B}{\FormulaRefAuto{A\subseteq B\vdash \iota_{A,B}\colon A\to B}{1}}
    \proofstep{1}{\TotRel{\iota_{A,B},A,B}}{\FormulaRefAuto{F\colon A\to B\vdash\TotRel{F,A,B}}{2}}
\end{tabproof}

\FormulaThmDelta[Grundgleichung der Inklusionsabbildung]
{A\subseteq B \dsep x\in A \vdash \iota_{A,B}(x)=x}
{
  \DeltaRow{Mengen}{A\dsep B\dsep x}
}
\begin{tabproof}
  \proofstep{1}{A\subseteq B}{\rA}
  \proofstep{2}{x\in A}{\rA}
  \proofstep{1}{%
    \forall x\in A\, \iota_{A,B}(x)=t_{\iota_{A,B}}(x)}{%
    \rAEb{\FormulaRefAuto{A\subseteq B\vdash \iota_{A,B}\coloneqq\iota F\left(F\colon A\to B \land \forall x\in A\, F(x)=t_{\iota_{A,B}}(x)\right)}{1}}}
  \proofstep{2}{\iota_{A,B}(x)=t_{\iota_{A,B}}(x)}{\rRE{2,\rUE{3}}}
  \proofstep{1,2}{t_{\iota_{A,B}}(x)=x}{%
    \FormulaRefAuto{A\subseteq B \dsep x\in A \vdash t_{\iota_{A,B}}(x) \coloneqq x}{1,2}}
  \proofstep{1,2}{\iota_{A,B}(x)=x}{\FormulaRefAuto{a = b,\, b = c \vdash a = c}{4,5}}
\end{tabproof}

\FormulaThmDelta
{A\subseteq B \dsep x\in A \vdash x=\iota_{A,B}(x)}%
{
  \DeltaRow{Mengen}{A\dsep B\dsep x}
}
\begin{tabproof}
  \proofstep{1}{A\subseteq B}{\rA}
  \proofstep{2}{x\in A}{\rA}

  \proofstep{1,2}{\iota_{A,B}(x)=x}{%
    \FormulaRefAuto{A\subseteq B \dsep x\in A \vdash \iota_{A,B}(x)=x}{1,2}}

  \proofstep{1,2}{x=\iota_{A,B}(x)}{%
    \FormulaRefAuto{a=b\vdash b=a}{3}}
\end{tabproof}


\subsubsection{Die Injektivität der Inklusionsabbildung}

\paragraph{Injektivität}
\FormulaThmDelta[Injektivität von \(\iota_{A,B}\)]{%
x\in A \dsep y\in A \dsep \iota_{A,B}(x)=\iota_{A,B}(y) \vdash x=y
}{
  \DeltaRow{Mengen}{A \dsep B \dsep x \dsep y}
}
\begin{tabproof}
  \proofstep{1}{\iota_{A,B}(x)=\iota_{A,B}(y)}{\rA}
  \proofstep{2}{x\in A}{\rA}
  \proofstep{3}{y\in A}{\rA}
  \proofstep{2}{\iota_{A,B}(x)=x}{%
    \FormulaRefAuto{A\subseteq B \dsep x\in A \vdash \iota_{A,B}(x)=x}{2}}
  \proofstep{3}{\iota_{A,B}(y)=y}{%
    \FormulaRefAuto{A\subseteq B \dsep x\in A \vdash \iota_{A,B}(x)=x}{3}}
  \proofstep{1,2,3}{x=y}{%
    \FormulaRefAuto{a=b\dsep a=c\dsep b=d\vdash c=d}{1,4,5}}
\end{tabproof}

\FormulaThmDelta[\(\iota_{A,B}\) als injektive Funktion]{%
A\subseteq B \vdash \iota_{A,B}\colon A\inj B
}{
  \DeltaRow{Mengen}{A\dsep B}
}
\begin{tabproof}
    \proofstep{1}{A\subseteq B}{\rA}
    \proofstep{1}{\iota_{A,B}\colon A\to B}{%
      \FormulaRefAuto{A\subseteq B\vdash \iota_{A,B}\colon A\to B}{1}}
    \proofstep{}{\forall x,y\in A\,(\iota_{A,B}(x)=\iota_{A,B}(y)\rightarrow x=y)}{%
      \FormulaRefAuto{x\in A \dsep y\in A \dsep \iota_{A,B}(x)=\iota_{A,B}(y) \vdash x=y}}
    \proofstep{1}{\iota_{A,B}\colon A\inj B}{%
      \FormulaRefAuto{Injektive Funktion}{2,3}}
\end{tabproof}


\subsection{Die Faserabbildung}

\subsubsection{Funktionsvorschrift}

\FormulaDefDelta[Funktionsvorschrift der Faser]{%
  y\in B \vdash t_{\Fib_F}(y) \coloneqq F^{-1}[\{y\}]%
}{
  \DeltaRow{Mengen}{A\dsep B\dsep y}
  \DeltaPrem{Funktionen}{F\colon A\to B}
  \DeltaRow{Funktionensymbol}{t_{\Fib_F}}
}

\FormulaThmDelta[Typisierungsaxiom für $t_{\Fib_F}$]{%
  y\in B \vdash t_{\Fib_F}(y)\in \powerset(A)%
}{
  \DeltaRow{Mengen}{A\dsep B\dsep y}
  \DeltaPrem{Funktionen}{F\colon A\to B}
  \DeltaRow{Funktionensymbol}{t_{\Fib_F}}
}
\begin{tabproof}
  \proofstep{1}{y\in B}{\rA}
  \proofstep{1}{t_{\Fib_F}(y)=F^{-1}[\{y\}]}{%
    \FormulaRefAuto{y\in B \vdash t_{\Fib_F}(y) \coloneqq F^{-1}[\{y\}]}{1}}
  \proofstep{1}{\{y\}\subseteq B}{%
    \FormulaRefAuto{x\in A\vdash \{x\}\subseteq A}{1}}
  \proofstep{1}{F^{-1}[\{y\}]=\{\,u\in A \mid F(u)\in \{y\}\,\}}{%
    \FormulaRefAuto{C\subseteq B \vdash F^{-1}[C] \coloneqq \{\,x\in A \mid F(x)\in C\,\}}{3}}
  \proofstep{1}{\{\,u\in A \mid F(u)\in \{y\}\,\}\subseteq A}{%
    \FormulaRefAuto{\{\,x\in A \mid P(x)\,\}\subseteq A}}
  \proofstep{1}{F^{-1}[\{y\}]\subseteq A}{\rIE{4,5}}
  \proofstep{1}{t_{\Fib_F}(y)\subseteq A}{\rIE{2,6}}
  \proofstep{1}{t_{\Fib_F}(y)\in \powerset(A)}{%
    \FormulaRefAuto{x \in \powerset(A)\;\eqvdash\; x \subseteq A}{7}}
\end{tabproof}

\subsubsection{Die Faser als Abbildung}

\FormulaThmDelta[Faser als Funktion]{%
  \exists! G\colon B\to \powerset(A)\,\forall y\in B\,G(y)=t_{\Fib_F}(y)%
}{
  \DeltaRow{Mengen}{A\dsep B\dsep y}
  \DeltaPrem{Funktionen}{F\colon A\to B}
  \DeltaRow{Funktionensymbol}{t_{\Fib_F}}
}
\begin{tabproof}
  \proofstep{}{\forall y\in B\,t_{\Fib_F}(y)\in \powerset(A)}{%
    \FormulaRefAuto{y\in B \vdash t_{\Fib_F}(y)\in \powerset(A)}}
  \proofstep{}{%
    \exists! G\colon B\to \powerset(A)\,\forall y\in B\,G(y)=t_{\Fib_F}(y)}{%
    \FormulaRefAuto{x\in A\vdash t(x)\coloneq y\dsep x\in A\vdash t(x)\in B\exists! F\colon A\to B\,\forall x\in A\,F(x) = t(x)}{1}}
\end{tabproof}

\FormulaDefDelta[Faser zu $F$]{%
  \Fib_F\coloneqq\iota G\left(G\colon B\to \powerset(A) \land \forall y\in B\, G(y)=t_{\Fib_F}(y)\right)%
}{
  \DeltaRow{Mengen}{A\dsep B}
  \DeltaPrem{Funktionen}{F\colon A\to B}
  \DeltaRow{Funktionensymbol}{t_{\Fib_F}}
}

\FormulaThmDelta[Faser ist eine Funktion]{%
  \Fib_F\colon B\to \powerset(A)%
}{
  \DeltaRow{Mengen}{A\dsep B}
  \DeltaPrem{Funktionen}{F\colon A\to B}
}
\begin{tabproof}
  \proofstep{}{\Fib_F\colon B\to \powerset(A)}{%
    \rAEa{\FormulaRefAuto{\Fib_F\coloneqq\iota G\left(G\colon B\to \powerset(A) \land \forall y\in B\, G(y)=t_{\Fib_F}(y)\right)}}}
\end{tabproof}

\FormulaThmDelta[Grundgleichung der Faser]{%
  y\in B \vdash \Fib_F(y)=F^{-1}[\{y\}]%
}{
  \DeltaRow{Mengen}{A\dsep B\dsep y}
  \DeltaPrem{Funktionen}{F\colon A\to B}
  \DeltaRow{Funktionensymbol}{t_{\Fib_F}}
}
\begin{tabproof}
  \proofstep{1}{y\in B}{\rA}
  \proofstep{}{%
    \forall y\in B\, \Fib_F(y)=t_{\Fib_F}(y)}{%
    \rAEb{\FormulaRefAuto{\Fib_F\coloneqq\iota G\left(G\colon B\to \powerset(A) \land \forall y\in B\, G(y)=t_{\Fib_F}(y)\right)}}}
  \proofstep{1}{\Fib_F(y)=t_{\Fib_F}(y)}{\rRE{1,\rUE{2}}}
  \proofstep{1}{t_{\Fib_F}(y)=F^{-1}[\{y\}]}{%
    \FormulaRefAuto{y\in B \vdash t_{\Fib_F}(y) \coloneqq F^{-1}[\{y\}]}{1}}
  \proofstep{1}{\Fib_F(y)=F^{-1}[\{y\}]}{%
    \FormulaRefAuto{a = b,\, b = c \vdash a = c}{3,4}}
\end{tabproof}

\subsubsection{Die Faser als Menge}

\FormulaThmDelta[Faser als Urbild eines Singletons]{%
  y\in B \vdash \Fib_F(y) = \{\,x\in A \mid F(x)=y\,\}%
}{
  \DeltaRow{Mengen}{A\dsep B\dsep x\dsep y}
  \DeltaPrem{Funktionen}{F\colon A\to B}
}
\begin{tabproofwide}
  \proofstepwidestar[1]{y\in B}{\rA}
  \proofstepwidestar[1]{\{y\}\subseteq B}{%
    \FormulaRefAuto{x\in A\vdash \{x\}\subseteq A}{1}}

  \proofstepwide[1]{\Fib_F(y)}{=}{F^{-1}[\{y\}]}{%
    \FormulaRefAuto{y\in B \vdash \Fib_F(y)=F^{-1}[\{y\}]}{1}}

  \proofstepwide[1]{}{=}{\{\,u\in A \mid F(u)\in \{y\}\,\}}{%
    \FormulaRefAuto{C\subseteq B \vdash F^{-1}[C] \coloneqq \{\,x\in A \mid F(x)\in C\,\}}{2}}

  \proofstepwide[1]{}{=}{\{\,u\in A \mid F(u)=y\,\}}{%
    \FormulaRefAuto{\{x\in A\mid F(x)\in \{a\}\}=\{x\in A\mid F(x)=a\}}}

  \proofstepwide[1]{\Fib_F(y)}{=}{\{\,u\in A \mid F(u)=y\,\}}{%
    \rChain{3,4,5}}
\end{tabproofwide}


\FormulaThmDelta{y\in B\dsep x\in \Fib_F(y) \vdash x\in A\land F(x)=y}
{
  \DeltaRow{Mengen}{A\dsep B\dsep x\dsep y}
  \DeltaPrem{Funktionen}{F\colon A\to B}
}
\begin{tabproof}
  \proofstep{1}{y\in B}{\rA}
  \proofstep{2}{x\in \Fib_F(y)}{\rA}
  \proofstep{1}{\Fib_F(y)=\{\,t\in A \mid F(t)=y\,\}}{%
    \FormulaRefAuto{y\in B \vdash \Fib_F(y) = \{\,x\in A \mid F(x)=y\,\}}{1}}
  \proofstep{1,2}{x\in \{\,t\in A \mid F(t)=y\,\}}{\rIE{3,2}}
  \proofstep{1,2}{x\in A \land F(x)=y}{%
    \FormulaRefAuto{x\in\{u\in A\mid P(u)\}\eqvdash x\in A\land P(x)}{4}}
\end{tabproof}

\FormulaThmDelta{y\in B\dsep x\in \Fib_F(y) \vdash x\in A}
{
  \DeltaRow{Mengen}{A\dsep B\dsep x\dsep y}
  \DeltaPrem{Funktionen}{F\colon A\to B}
}
\begin{tabproof}
  \proofstep{1}{y\in B}{\rA}
  \proofstep{2}{x\in \Fib_F(y)}{\rA}
  \proofstep{1,2}{x\in A\land F(x)=y}{\FormulaRefAuto{y\in B\dsep x\in \Fib_F(y) \vdash x\in A\land F(x)=y}{1,2}}
  \proofstep{1,2}{x\in A}{\rAEa{3}}
\end{tabproof}

\FormulaThmDelta{y\in B\dsep x\in \Fib_F(y) \vdash F(x)=y}
{
  \DeltaRow{Mengen}{A\dsep B\dsep x\dsep y}
  \DeltaPrem{Funktionen}{F\colon A\to B}
}
\begin{tabproof}
  \proofstep{1}{y\in B}{\rA}
  \proofstep{2}{x\in \Fib_F(y)}{\rA}
  \proofstep{1,2}{x\in A\land F(x)=y}{\FormulaRefAuto{y\in B\dsep x\in \Fib_F(y) \vdash x\in A\land F(x)=y}{1,2}}
  \proofstep{1,2}{F(x)=y}{\rAEb{3}}
\end{tabproof}

\FormulaThmDelta
{y\in B\dsep x\in A\dsep F(x)=y\vdash x\in \Fib_F(y)}
{
  \DeltaRow{Mengen}{A\dsep B\dsep x\dsep y}
  \DeltaPrem{Funktionen}{F\colon A\to B}
}
\begin{tabproof}
  \proofstep{1}{y\in B}{\rA}
  \proofstep{2}{x\in A}{\rA}
  \proofstep{3}{F(x)=y}{\rA}
  \proofstep{2,3}{x\in A\land F(x)=y}{\rAI{2,3}}
  \proofstep{2,3}{x\in \{\,t\in A \mid F(t)=y\,\}}{%
    \FormulaRefAuto{x\in\{u\in A\mid P(u)\}\eqvdash x\in A\land P(x)}{4}}
  \proofstep{1}{\Fib_F(y)=\{\,t\in A \mid F(t)=y\,\}}{%
    \FormulaRefAuto{y\in B \vdash \Fib_F(y) = \{\,x\in A \mid F(x)=y\,\}}{1}}
  \proofstep{1,2,3}{x\in \Fib_F(y)}{\rIE{6,5}}
\end{tabproof}


\FormulaThmDelta
{y\in B\vdash x\in\Fib_F(y)\leftrightarrow x\in A\land F(x)=y}
{
  \DeltaRow{Mengen}{A\dsep B\dsep x\dsep y}
  \DeltaPrem{Funktionen}{F\colon A\to B}
}
\begin{tabproof}
  \proofstep{1}{y\in B}{\rA}

  \proofstep{2}{x\in \Fib_F(y)}{\rA}
  \proofstep{1,2}{x\in A\land F(x)=y}{%
    \FormulaRefAuto{y\in B\dsep x\in \Fib_F(y) \vdash x\in A\land F(x)=y}{1,2}}
  \proofstep{1}{x\in \Fib_F(y)\rightarrow (x\in A\land F(x)=y)}{\rRI{2,3}}

  \proofstep{5}{x\in A\land F(x)=y}{\rA}
  \proofstep{5}{x\in A}{\rAEa{5}}
  \proofstep{5}{F(x)=y}{\rAEb{5}}
  \proofstep{1,5}{x\in \Fib_F(y)}{%
    \FormulaRefAuto{y\in B\dsep x\in A\dsep F(x)=y\vdash x\in \Fib_F(y)}{1,6,7}}
  \proofstep{1}{(x\in A\land F(x)=y)\rightarrow x\in \Fib_F(y)}{\rRI{5,8}}

  \proofstep{1}{x\in\Fib_F(y)\leftrightarrow (x\in A\land F(x)=y)}{\rLRI{4,9}}
\end{tabproof}

\FormulaThmDelta[Verschiedene Fasern sind disjunkt]{%
  y\in B \dsep z\in B \dsep y\neq z \vdash \Fib_F(y)\cap \Fib_F(z)=\varnothing%
}{
  \DeltaRow{Mengen}{A\dsep B\dsep y\dsep z\dsep a}
  \DeltaPrem{Funktionen}{F\colon A\to B}
}
\begin{tabproof}
  \proofstep{1}{y\in B}{\rA}
  \proofstep{2}{z\in B}{\rA}
  \proofstep{3}{y\neq z}{\rA}

  \proofstep{4}{a\in \Fib_F(y)\cap \Fib_F(z)}{\rA}
  \proofstep{4}{a\in \Fib_F(y)}{%
    \FormulaRefAuto{x \in A \cap B \vdash x \in A}{4}}
  \proofstep{4}{a\in \Fib_F(z)}{%
    \FormulaRefAuto{x \in A \cap B \vdash x \in B}{4}}

  \proofstep{1,4}{F(a)=y}{%
    \FormulaRefAuto{y\in B\dsep x\in \Fib_F(y) \vdash F(x)=y}{1,5}}
  \proofstep{2,4}{F(a)=z}{%
    \FormulaRefAuto{y\in B\dsep x\in \Fib_F(y) \vdash F(x)=y}{2,6}}

  \proofstep{1,2,4}{y=z}{%
    \FormulaRefAuto{a=b \dsep a=c \vdash b=c}{7,8}}
  \proofstep{1,2,3,4}{\bot}{\rBI{3,9}}

  \proofstep{1,2,3}{a\notin \Fib_F(y)\cap \Fib_F(z)}{\rCE{4,10}}
  \proofstep{1,2,3}{\forall a\,\bigl(a\notin \Fib_F(y)\cap \Fib_F(z)\bigr)}{\rUI{11}}
  \proofstep{1,2,3}{\Fib_F(y)\cap \Fib_F(z)=\varnothing}{%
    \FormulaRefAuto{\forall b\,b\notin X\vdash X=\varnothing}{12}}
\end{tabproof}

\FormulaThmDelta{%
  y\in B \dsep z\in B \dsep X=\Fib_F(y) \dsep Y=\Fib_F(z) \dsep X\neq Y
  \vdash X\cap Y=\varnothing%
}{
  \DeltaRow{Mengen}{A\dsep B\dsep X\dsep Y\dsep y\dsep z}
  \DeltaPrem{Funktionen}{F\colon A\to B}
}
\begin{tabproof}
  \proofstep{1}{y\in B}{\rA}
  \proofstep{2}{z\in B}{\rA}
  \proofstep{3}{X=\Fib_F(y)}{\rA}
  \proofstep{4}{Y=\Fib_F(z)}{\rA}
  \proofstep{5}{X\neq Y}{\rA}

  \proofstep{6}{y=z}{\rA}
  \proofstep{4,6}{Y=\Fib_F(y)}{%
    \FormulaRefAuto{x = y \dsep a = F(y) \vdash a = F(x)}{6,4}}

  \proofstep{3,4,6}{X=Y}{%
    \FormulaRefAuto{a = b,\, c = b \vdash a = c}{3,7}}

  \proofstep{3,4,5,6}{\bot}{\rBI{5,8}}
  \proofstep{1,2,3,4,5}{y\neq z}{\rCI{6,9}}

  \proofstep{1,2,3,4,5}{\Fib_F(y)\cap \Fib_F(z)=\varnothing}{%
    \FormulaRefAuto{y\in B \dsep z\in B \dsep y\neq z \vdash \Fib_F(y)\cap \Fib_F(z)=\varnothing}{1,2,10}}

  \proofstep{1,2,3,4,5}{X\cap Y=\varnothing}{%
    \FormulaRefAuto{X=U \dsep Y=V \dsep U\cap V=W \vdash X\cap Y=W}{3,4,11}}
\end{tabproof}



\FormulaThmDelta{%
  X\in \Fib_F[B] \dsep Y\in \Fib_F[B] \dsep X\neq Y \vdash X\cap Y=\varnothing%
}{
  \DeltaRow{Mengen}{A\dsep B\dsep X\dsep Y\dsep y\dsep z}
  \DeltaPrem{Funktionen}{F\colon A\to B}
}
\begin{tabproof}
  \proofstep{1}{X\in \Fib_F[B]}{\rA}
  \proofstep{2}{Y\in \Fib_F[B]}{\rA}
  \proofstep{3}{X\neq Y}{\rA}
  \proofstep{}{\Fib_F\colon B\to \powerset(A)}{%
    \FormulaRefAuto{\Fib_F\colon B\to \powerset(A)}}

  \proofstep{1}{\exists y\in B\,X=\Fib_F(y)}{%
    \FormulaRefAuto{F\colon A\to B\dsep y\in F[A] \vdash \exists x\in A\,y=F(x)}{4,1}}

  \proofstep{2}{\exists z\in B\,Y=\Fib_F(z)}{%
    \FormulaRefAuto{F\colon A\to B\dsep y\in F[A] \vdash \exists x\in A\,y=F(x)}{4,2}}

  \proofstep{7}{y\in B \land X=\Fib_F(y)}{\rA}
  \proofstep{7}{y\in B}{\rAEa{7}}
  \proofstep{7}{X=\Fib_F(y)}{\rAEb{7}}

  \proofstep{10}{z\in B \land Y=\Fib_F(z)}{\rA}
  \proofstep{10}{z\in B}{\rAEa{10}}
  \proofstep{10}{Y=\Fib_F(z)}{\rAEb{10}}

  \proofstep{3,7,10}{X\cap Y=\varnothing}{%
    \FormulaRefAuto{y\in B \dsep z\in B \dsep X=\Fib_F(y) \dsep Y=\Fib_F(z) \dsep X\neq Y \vdash X\cap Y=\varnothing}{8,11,9,12,3}}

  \proofstep{2,3,7}{X\cap Y=\varnothing}{\rEE{6,10,13}}
  \proofstep{1,2,3}{X\cap Y=\varnothing}{\rEE{5,7,14}}
\end{tabproof}

\subsubsection{Fasern surjektiver Abbildungen}

\FormulaThmDeltaR{y\in B \vdash \Fib_F(y)\neq\varnothing}
{F\colon A\sur B \dsep y\in B \vdash \Fib_F(y)\neq\varnothing}
{
  \DeltaRow{Mengen}{A\dsep B\dsep x\dsep y}
  \DeltaRow{Surjektive Funktion}{F\colon A\sur B}
}
\begin{tabproof}
  \proofstep{1}{y\in B}{\rA}
  \proofstep{1}{\exists x\in A\,F(x)=y}{\FormulaRefAuto{Surjektivität}{1}}
  \proofstep{3}{x\in A \land F(x)=y}{\rA}
  \proofstep{3}{x\in A}{\rAEa{3}}
  \proofstep{3}{F(x)=y}{\rAEb{3}}
  \proofstep{1,3}{x\in \Fib_F(y)}{%
    \FormulaRefAuto{y\in B\dsep x\in A\dsep F(x)=y \vdash x\in\Fib_F(y)}{1,4,5}}
  \proofstep{1,3}{\Fib_F(y)\neq\varnothing}{\FormulaRefAuto{a\in A\vdash A\neq\varnothing}{6}}
  \proofstep{1}{\Fib_F(y)\neq\varnothing}{\rEE{2,3,7}}
\end{tabproof}

\FormulaThmDelta
{X\in \Fib_F[B] \vdash X\neq\varnothing}
{
  \DeltaRow{Mengen}{A\dsep B\dsep X\dsep y}
  \DeltaRow{Surjektive Funktion}{F\colon A\sur B}
}
\begin{tabproof}
  \proofstep{1}{X\in \Fib_F[B]}{\rA}
  \proofstep{}{\Fib_F\colon B\to \powerset(A)}{%
    \FormulaRefAuto{\Fib_F\colon B\to \powerset(A)}}

  \proofstep{1}{\exists y\in B\,X=\Fib_F(y)}{%
    \FormulaRefAuto{F\colon A\to B\dsep y\in F[A] \vdash \exists x\in A\,y=F(x)}{2,1}}

  \proofstep{4}{y\in B \land X=\Fib_F(y)}{\rA}
  \proofstep{4}{y\in B}{\rAEa{4}}
  \proofstep{4}{X=\Fib_F(y)}{\rAEb{4}}

  \proofstep{5}{\Fib_F(y)\neq\varnothing}{%
    \FormulaRefAuto{F\colon A\sur B \dsep y\in B \vdash \Fib_F(y)\neq\varnothing}{5}}

  \proofstep{4}{X\neq\varnothing}{\rIE{6,7}}
  \proofstep{1}{X\neq\varnothing}{\rEE{3,4,8}}
\end{tabproof}

\FormulaThmDelta
{\DisjFam(\Fib_F[B])}
{
  \DeltaRow{Mengen}{A\dsep B\dsep X\dsep Y}
  \DeltaRow{Surjektive Funktion}{F\colon A\sur B}
}
\begin{tabproof}
  \proofstep{}{X\in \Fib_F[B] \vdash X\neq\varnothing}{%
    \FormulaRefAuto{X\in \Fib_F[B] \vdash X\neq\varnothing}}

  \proofstep{}{X\in \Fib_F[B]\dsep Y\in \Fib_F[B]\dsep X\neq Y \vdash X\cap Y=\varnothing}{%
    \FormulaRefAuto{X\in \Fib_F[B] \dsep Y\in \Fib_F[B] \dsep X\neq Y \vdash X\cap Y=\varnothing}}

  \proofstep{}{\DisjFam(\Fib_F[B])}{%
    \FormulaRefAuto{Disjunkte Familie}{1,2}}
\end{tabproof}

\subsection{Die Einschränkung einer Funktion}

\subsubsection{Definition der eingeschränkten Funktionsvorschrift}

\FormulaDefDelta[Funktionsvorschrift für die Einschränkung von \(F\) auf \(C\)]
{C\subseteq A \dsep x\in C \vdash t_{F\restriction_C}(x) \coloneqq F(x)}
{
  \DeltaRow{Mengen}{A\dsep B\dsep C\dsep x}
  \DeltaPrem{Funktionen}{F\colon A\to B}
  \DeltaRow{Funktionensymbol}{t_{F\restriction_C}}
}

% — Typisierungsaxiom / Wertebereich —
\FormulaThmDelta[Typisierungsaxiom für \(t_{F\restriction_C}\)]
{C\subseteq A \dsep x\in C \vdash t_{F\restriction_C}(x)\in B}
{
  \DeltaRow{Mengen}{A\dsep B\dsep C\dsep x}
  \DeltaPrem{Funktionen}{F\colon A\to B}
  \DeltaRow{Funktionensymbol}{t_{F\restriction_C}}
}
\begin{tabproof}
  \proofstep{1}{C\subseteq A}{\rA}
  \proofstep{2}{x\in C}{\rA}
  % aus C⊆A und x∈C folgt x∈A
  \proofstep{1,2}{x\in A}{\FormulaRefAuto{A\subseteq B, x\in A \vdash x\in B}{1,2}}
  % Typisierungsaxiom für F selbst: x∈A ⇒ F(x)∈ B
  \proofstep{3}{F(x)\in B}{\FormulaRefAuto{F\colon A\to B\dsep x\in A\vdash F(x)\in B}{3}}
  % Definition der Funktionsvorschrift der Einschränkung
  \proofstep{1,2,3}{t_{F\restriction_C}(x)=F(x)}{%
    \FormulaRefAuto{C\subseteq A \dsep x\in C \vdash t_{F\restriction_C}(x) \coloneqq F(x)}{1,2,3}}
  % Über Gleichung auf den Wertebereich schließen
  \proofstep{1,2,3}{t_{F\restriction_C}(x)\in B}{\rIE{6,5}}
\end{tabproof}

\subsubsection{Die Einschränkung als Funktion}

\FormulaThmDelta[Einschränkung als Funktion]{%
C\subseteq A \vdash
\exists! G\colon C\to B\,\forall x\in C\,G(x) = t_{F\restriction_C}(x)%
}{
  \DeltaRow{Mengen}{A\dsep B\dsep C\dsep x}
  \DeltaPrem{Funktionen}{F\colon A\to B}
  \DeltaRow{Funktionensymbol}{t_{F\restriction_C}}
}
\begin{tabproof}
  \proofstep{1}{C\subseteq A}{\rA}

  \proofstep{1}{\forall x\in C\, t_{F\restriction_C}(x) = F(x)}{\FormulaRefAuto{C\subseteq A \dsep x\in C \vdash t_{F\restriction_C}(x) \coloneqq F(x)}{1}}

  \proofstep{1}{\forall x\in C\, t_{F\restriction_C}(x)\in B}{\FormulaRefAuto{C\subseteq A \dsep x\in C \vdash t_{F\restriction_C}(x)\in B}{1}}
  % Existenz und Eindeutigkeit einer Funktion G: C→B mit dieser Vorschrift
  \proofstep{1}{\exists! G\colon C\to B\,\forall x\in C\,G(x) = t_{F\restriction_C}(x)}{%
    \FormulaRefAuto{x\in A\vdash t(x)\coloneq y\dsep x\in A\vdash t(x)\in B\exists! F\colon A\to B\,\forall x\in A\,F(x) = t(x)}{2,3}}
\end{tabproof}

\FormulaDefDelta[Einschränkung von \(F\) auf \(C\)]
{C\subseteq A \vdash F\restriction_C\coloneqq\iota G\left(G\colon C\to B \land \forall x\in C\, G(x)=t_{F\restriction_C}(x)\right)}
{
  \DeltaRow{Mengen}{A\dsep B\dsep C}
  \DeltaPrem{Funktionen}{F\colon A\to B}
  \DeltaRow{Funktionensymbol}{t_{F\restriction_C}}
}

\FormulaThmDelta[Einschränkung von \(F\) auf \(C\)]
{C\subseteq A \vdash F\restriction_C\colon C\to B}
{
  \DeltaRow{Mengen}{A\dsep B\dsep C}
  \DeltaPrem{Funktionen}{F\colon A\to B}
}
\begin{tabproof}
  \proofstep{}{\;F\restriction_C\colon C\to B}{%
    \rAEa{\FormulaRefAuto{C\subseteq A \vdash F\restriction_C\coloneqq\iota G\left(G\colon C\to B \land \forall x\in C\, G(x)=t_{F\restriction_C}(x)\right)}}}
\end{tabproof}

\FormulaThmDelta
{C\subseteq A \vdash \TotRel{F\restriction_C,C,B}}
{
  \DeltaRow{Mengen}{A\dsep B}
}
\begin{tabproof}
    \proofstep{1}{C\subseteq A }{\rA}
    \proofstep{1}{F\restriction_C\colon C\to B}{\FormulaRefAuto{C\subseteq A \vdash F\restriction_C\colon C\to B}{1}}
    \proofstep{1}{\TotRel{F\restriction_C,C,B}}{\FormulaRefAuto{F\colon A\to B\vdash\TotRel{F,A,B}}{2}}
\end{tabproof}

\FormulaThmDelta[Grundgleichung der Einschränkung]
{C\subseteq A \dsep x\in C \vdash F\restriction_C(x)=F(x)}
{
  \DeltaRow{Mengen}{A\dsep B\dsep C\dsep x}
  \DeltaPrem{Funktionen}{F\colon A\to B}
}
\begin{tabproof}
  \proofstep{1}{C\subseteq A}{\rA}
  \proofstep{2}{x\in C}{\rA}
  \proofstep{1}{%
    \forall x\in C\, F\restriction_C(x)=t_{F\restriction_C}(x)}{%
    \rAEb{\FormulaRefAuto{C\subseteq A \vdash F\restriction_C\coloneqq\iota G\left(G\colon C\to B \land \forall x\in C\, G(x)=t_{F\restriction_C}(x)\right)}{1}}}
  \proofstep{1,2}{F\restriction_C(x)=t_{F\restriction_C}(x)}{\rRE{2,\rUE{3}}}
  \proofstep{1,2}{t_{F\restriction_C}(x)=F(x)}{%
    \FormulaRefAuto{C\subseteq A \dsep x\in C \vdash t_{F\restriction_C}(x) \coloneqq F(x)}{1,2}}
  \proofstep{1,2}{F\restriction_C(x)=F(x)}{%
    \FormulaRefAuto{a = b,\, b = c \vdash a = c}{4,5}}
\end{tabproof}

\subsubsection{Bild der Einschränkung}

\FormulaThmDelta[Bild der Einschränkung]{%
  C\subseteq A \dsep F\colon A\to B \vdash (F\restriction_C)[C]=F[C]
}{%
  \DeltaRow{Mengen}{A\dsep B\dsep C\dsep x\dsep y}
  \DeltaRow{Funktionen}{F}
}
\begin{tabproofsplitwide}
  \proofpartwideind[Teilmengenrichtung \((F\restriction_C)[C]\subseteq F[C]\)]{%
    C\subseteq A \dsep F\colon A\to B \vdash (F\restriction_C)[C]\subseteq F[C]
  }[%
    C\subseteq A \dsep F\colon A\to B \vdash (F\restriction_C)[C]\subseteq F[C]
  ]
    \proofstepwidestar[1]{C\subseteq A}{\rA}
    \proofstepwidestar[2]{F\colon A\to B}{\rA}
    \proofstepwidestar[1,2]{F\restriction_C\colon C\to B}{%
      \FormulaRefAuto{C\subseteq A \vdash F\restriction_C\colon C\to B}{1}}
    \proofstepwidestar[4]{y\in (F\restriction_C)[C]}{\rA}
    \proofstepwidestar[1,2,4]{\exists x\in C\,y=F\restriction_C(x)}{%
      \FormulaRefAuto{F\colon A\to B\dsep y\in F[A] \vdash \exists x\in A\,y=F(x)}{3,4}}
    \proofstepwidestar[6]{x\in C\land y=F\restriction_C(x)}{\rA}
    \proofstepwidestar[6]{x\in C}{\rAEa{6}}
    \proofstepwidestar[6]{y=F\restriction_C(x)}{\rAEb{6}}
    \proofstepwidestar[1,6]{F\restriction_C(x)=F(x)}{%
      \FormulaRefAuto{C\subseteq A \dsep x\in C \vdash F\restriction_C(x)=F(x)}{1,7}}
    \proofstepwidestar[1,6]{y=F(x)}{\FormulaRefAuto{a=b,\, b=c \vdash a=c}{8,9}}
    \proofstepwidestar[1,6]{\exists x\in C\,y=F(x)}{\rEI{\rAI{7,10}}}
    \proofstepwidestar[1,2,6]{y\in F[C]}{%
      \FormulaRefAuto{C\subseteq A\dsep F\colon A\to B\dsep \exists x\in C\,y=F(x)\vdash y\in F[C]}{1,2,11}}
    \proofstepwidestar[1,2]{\forall y\,(y\in (F\restriction_C)[C]\rightarrow y\in F[C])}{\rUI{\rRI{4,12}}}
    \proofstepwidestar[1,2]{(F\restriction_C)[C]\subseteq F[C]}{%
      \FormulaRefAuto{A\subseteq B \coloneq \forall x\,(x\in A \rightarrow x\in B)}{13}}
  \closeproofpartwideind

  \proofpartwideind[Teilmengenrichtung \(F[C]\subseteq (F\restriction_C)[C]\)]{%
    C\subseteq A \dsep F\colon A\to B \vdash F[C]\subseteq (F\restriction_C)[C]
  }[%
    C\subseteq A \dsep F\colon A\to B \vdash F[C]\subseteq (F\restriction_C)[C]
  ]
    \proofstepwidestar[1]{C\subseteq A}{\rA}
    \proofstepwidestar[2]{F\colon A\to B}{\rA}
    \proofstepwidestar[1,2]{F\restriction_C\colon C\to B}{%
      \FormulaRefAuto{C\subseteq A \vdash F\restriction_C\colon C\to B}{1}}
    \proofstepwidestar[4]{y\in F[C]}{\rA}
    \proofstepwidestar[1,2,4]{\exists x\in C\,y=F(x)}{%
      \FormulaRefAuto{C\subseteq A\dsep F\colon A\to B\dsep y\in F[C]\vdash \exists x\in C\,y=F(x)}{1,2,4}}
    \proofstepwidestar[6]{x\in C\land y=F(x)}{\rA}
    \proofstepwidestar[6]{x\in C}{\rAEa{6}}
    \proofstepwidestar[6]{y=F(x)}{\rAEb{6}}
    \proofstepwidestar[1,6]{F\restriction_C(x)=F(x)}{%
      \FormulaRefAuto{C\subseteq A \dsep x\in C \vdash F\restriction_C(x)=F(x)}{1,7}}
    \proofstepwidestar[1,6]{F(x)=F\restriction_C(x)}{\FormulaRefAuto{a=b\vdash b=a}{9}}
    \proofstepwidestar[1,6]{y=F\restriction_C(x)}{\FormulaRefAuto{a=b,\, b=c \vdash a=c}{8,10}}
    \proofstepwidestar[1,6]{\exists x\in C\,y=F\restriction_C(x)}{\rEI{\rAI{7,11}}}
    \proofstepwidestar[1,2,6]{y\in (F\restriction_C)[C]}{%
      \FormulaRefAuto{F\colon A\to B\dsep \exists x\in A\,y=F(x)\vdash y\in F[A]}{3,12}}
    \proofstepwidestar[1,2]{\forall y\,(y\in F[C]\rightarrow y\in (F\restriction_C)[C])}{\rUI{\rRI{4,13}}}
    \proofstepwidestar[1,2]{F[C]\subseteq (F\restriction_C)[C]}{%
      \FormulaRefAuto{A\subseteq B \coloneq \forall x\,(x\in A \rightarrow x\in B)}{14}}
  \closeproofpartwideind

  \proofpartwide{Zusammenfassung}
    \proofstepwidestar[1]{C\subseteq A}{\rA}
    \proofstepwidestar[2]{F\colon A\to B}{\rA}
    \proofstepwidestar[1,2]{(F\restriction_C)[C]\subseteq F[C]}{%
      \FormulaRefAuto{C\subseteq A \dsep F\colon A\to B \vdash (F\restriction_C)[C]\subseteq F[C]}{1,2}}
    \proofstepwidestar[1,2]{F[C]\subseteq (F\restriction_C)[C]}{%
      \FormulaRefAuto{C\subseteq A \dsep F\colon A\to B \vdash F[C]\subseteq (F\restriction_C)[C]}{1,2}}
    \proofstepwidestar[1,2]{(F\restriction_C)[C]=F[C]}{%
      \FormulaRefAuto{A\subseteq B\dsep B\subseteq A \vdash A=B}{3,4}}
  \closeproofpartwide
\end{tabproofsplitwide}


\subsubsection{Die Auswertungseigenschaft}

\FormulaThmDelta{%
C\subseteq A \dsep x\in C \dsep y\in C \dsep F(x)=F(y)\vdash F\restriction_C(x)=F\restriction_C(y)
}{
  \DeltaRow{Mengen}{A \dsep B \dsep C \dsep x \dsep y}
  \DeltaPrem{Funktionen}{F\colon A\to B}
}
\begin{tabproof}
  \proofstep{1}{C\subseteq A}{\rA}
  \proofstep{2}{x\in C}{\rA}
  \proofstep{3}{y\in C}{\rA}
  \proofstep{4}{F(x)=F(y)}{\rA}
  \proofstep{1,2}{F\restriction_C(x)=F(x)}{%
    \FormulaRefAuto{C\subseteq A \dsep x\in C \vdash F\restriction_C(x)=F(x)}{1,2}}
  \proofstep{1,3}{F\restriction_C(y)=F(y)}{%
    \FormulaRefAuto{C\subseteq A \dsep x\in C \vdash F\restriction_C(x)=F(x)}{1,3}}
  \proofstep{1,3}{F(y)=F\restriction_C(y)}{%
    \FormulaRefAuto{a=b\vdash b=a}{6}}
  \proofstep{1,3,4}{F(x)=F\restriction_C(y)}{%
    \FormulaRefAuto{a = b,\, b = c \vdash a = c}{4,7}}
  \proofstep{1,2,3,4}{F\restriction_C(x)=F\restriction_C(y)}{%
    \FormulaRefAuto{a = b,\, b = c \vdash a = c}{5,8}}
\end{tabproof}

\FormulaThmDelta{%
C\subseteq A \dsep x\in C \dsep y\in C \dsep F\restriction_C(x)=F\restriction_C(y)\vdash F(x)=F(y)
}{
  \DeltaRow{Mengen}{A \dsep B \dsep C \dsep x \dsep y}
  \DeltaPrem{Funktionen}{F\colon A\to B}
}
\begin{tabproof}
  \proofstep{1}{C\subseteq A}{\rA}
  \proofstep{2}{x\in C}{\rA}
  \proofstep{3}{y\in C}{\rA}
  \proofstep{4}{F\restriction_C(x)=F\restriction_C(y)}{\rA}
  \proofstep{1,2}{F\restriction_C(x)=F(x)}{%
    \FormulaRefAuto{C\subseteq A \dsep x\in C \vdash F\restriction_C(x)=F(x)}{1,2}}
  \proofstep{1,3}{F\restriction_C(y)=F(y)}{%
    \FormulaRefAuto{C\subseteq A \dsep x\in C \vdash F\restriction_C(x)=F(x)}{1,3}}
  \proofstep{1,2}{F(x)=F\restriction_C(x)}{%
    \FormulaRefAuto{a=b\vdash b=a}{5}}
  \proofstep{1,2,4}{F(x)=F\restriction_C(y)}{%
    \FormulaRefAuto{a = b,\, b = c \vdash a = c}{7,4}}
  \proofstep{1,2,3,4}{F(x)=F(y)}{%
    \FormulaRefAuto{a = b,\, b = c \vdash a = c}{8,6}}
\end{tabproof}

\subsubsection{Erhalt von Eigenschaften}

\FormulaThmDelta[Injektivität der Einschränkung]{%
F\colon A\inj B\dsep C\subseteq A \dsep x\in C \dsep y\in C \dsep
F\restriction_C(x)=F\restriction_C(y) \vdash x=y
}{
  \DeltaRow{Mengen}{A\dsep B\dsep C\dsep x\dsep y}
  \DeltaRow{Funktionen}{F}
  \DeltaPrem{Funktionen}{F\restriction_C\colon C\to B}[\FormulaRefAuto{C\subseteq A \vdash F\restriction_C\colon C\to B}]
}
\begin{tabproof}
  \proofstep{1}{F\colon A\inj B}{\rA}
  \proofstep{2}{C\subseteq A}{\rA}
  \proofstep{3}{x\in C}{\rA}
  \proofstep{4}{y\in C}{\rA}
  \proofstep{5}{F\restriction_C(x)=F\restriction_C(y)}{\rA}
  \proofstep{2,3,4,5}{F(x)=F(y)}{%
    \FormulaRefAuto{C\subseteq A \dsep x\in C \dsep y\in C \dsep F\restriction_C(x)=F\restriction_C(y)\vdash F(x)=F(y)}{2,3,4,5}}
  \proofstep{2,3}{x\in A}{\FormulaRefAuto{A\subseteq B, x\in A\vdash x\in B}{2,3}}
  \proofstep{2,4}{y\in A}{\FormulaRefAuto{A\subseteq B, x\in A\vdash x\in B}{2,4}}
  \proofstep{1,2,3,4,5}{x=y}{%
    \FormulaRefAuto{F\colon A\inj B\dsep x\in A\dsep y\in A\dsep F(x)=F(y)\vdash x=y}{1,7,8,6}}
\end{tabproof}

\FormulaThmDelta[Einschränkung einer injektiven Funktion]{%
 F\colon A\inj B\dsep C\subseteq A \vdash F\restriction_C\colon C\inj B
}{
  \DeltaRow{Mengen}{A\dsep B\dsep C}
  \DeltaRow{Funktionen}{F}
}
\begin{tabproof}
  \proofstep{1}{F\colon A\inj B}{\rA}
  \proofstep{2}{C\subseteq A}{\rA}
  \proofstep{2}{F\restriction_C\colon C\to B}{%
    \FormulaRefAuto{C\subseteq A \vdash F\restriction_C\colon C\to B}{2}}
  \proofstep{1,2}{\forall x,y\in C\,(F\restriction_C(x)=F\restriction_C(y)\rightarrow x=y)}{%
    \FormulaRefAuto{F\colon A\inj B\dsep C\subseteq A \dsep x\in C \dsep y\in C \dsep
F\restriction_C(x)=F\restriction_C(y) \vdash x=y}{1,2}}
  \proofstep{1,2}{F\restriction_C\colon C\inj B}{%
    \FormulaRefAuto{Injektive Funktion}{3,4}}
\end{tabproof}

\FormulaThmDelta[Surjektivität der Einschränkung auf das Bild]{%
  C\subseteq A \dsep y\in F[C] \vdash \exists x\in C\, F\restriction_C(x)=y
}{
  \DeltaRow{Mengen}{A\dsep B\dsep C\dsep x\dsep y}
  \DeltaPrem{Funktionen}{F\colon A\to B}
}
\begin{tabproof}
  \proofstep{1}{C\subseteq A}{\rA}
  \proofstep{2}{y\in F[C]}{\rA}
  \proofstep{1,2}{\exists x\in C\, y=F(x)}{%
    \FormulaRefAuto{C\subseteq A\dsep F\colon A\to B\dsep y\in F[C]\vdash \exists x\in C\,y=F(x)}{1,2}}

  \proofstep{4}{x\in C\land y=F(x)}{\rA}
  \proofstep{4}{x\in C}{\rAEa{4}}
  \proofstep{4}{y=F(x)}{\rAEb{4}}
  \proofstep{1,4}{F\restriction_C(x)=F(x)}{%
    \FormulaRefAuto{C\subseteq A \dsep x\in C \vdash F\restriction_C(x)=F(x)}{1,5}}
  \proofstep{4}{F(x)=y}{\FormulaRefAuto{a=b\vdash b=a}{6}}
  \proofstep{1,4}{F\restriction_C(x)=y}{%
    \FormulaRefAuto{a=b,\, b=c \vdash a=c}{7,8}}
  \proofstep{1,4}{\exists x\in C\,F\restriction_C(x)=y}{\rEI{\rAI{5,9}}}

  \proofstep{1,2}{\exists x\in C\,F\restriction_C(x)=y}{\rEE{3,4,10}}
\end{tabproof}

\FormulaThmDelta[Einschränkung als surjektive Funktion auf das Bild]{%
C\subseteq A \vdash F\restriction_C\colon C\twoheadrightarrow F[C]
}{
  \DeltaRow{Mengen}{A\dsep B\dsep C}
  \DeltaPrem{Funktionen}{F\colon A\to B}
}
\begin{tabproof}
  \proofstep{1}{C\subseteq A}{\rA}
  \proofstep{1}{F\restriction_C\colon C\to B}{%
    \FormulaRefAuto{C\subseteq A \vdash F\restriction_C\colon C\to B}{1}}
  \proofstep{1}{\forall y\in F[C]\;\exists x\in C\,F\restriction_C(x)=y}{%
    \FormulaRefAuto{C\subseteq A \dsep y\in F[C] \vdash \exists x\in C\, F\restriction_C(x)=y}{1}}
  \proofstep{1}{F\restriction_C\colon C\twoheadrightarrow F[C]}{%
    \FormulaRefAuto{Surjektive Funktion}{2,3}}
\end{tabproof}

\subsubsection{Verkleinerung der Zielmenge}

\FormulaThmDelta[Verkleinerung der Zielmenge]{%
  F\colon A\to B\dsep \forall x\in A\,F(x)\in C
  \vdash
  F\colon A\to C
}{%
  \DeltaRow{Mengen}{A\dsep B\dsep C\dsep x\dsep y\dsep z}
  \DeltaRow{Funktionen}{F}
}
\begin{tabproof}
  \proofstep{1}{F\colon A\to B}{\rA}
  \proofstep{2}{\forall x\in A\,F(x)\in C}{\rA}

  \proofstep{3}{(x,y)\in F}{\rA}
  \proofstep{1,3}{x\in A}{\FormulaRefAuto{(x,y)\in F\vdash x\in A}{1,3}}
  \proofstep{1,2,3}{F(x)\in C}{\rRE{2,4}}
  \proofstep{1,3}{y=F(x)}{\FormulaRefAuto{(x,y)\in F\vdash y=F(x)}{1,3}}
  \proofstep{1,2,3}{y\in C}{\rIE{6,5}}
  \proofstep{1,2,3}{(x,y)\in A\times C}{%
    \FormulaRefAuto{x\in A\dsep y\in C\vdash (x,y)\in A\times C}{4,7}}
  \proofstep{1,2}{(x,y)\in F\rightarrow (x,y)\in A\times C}{\rRI{3,8}}
  \proofstep{1,2}{F\subseteq A\times C}{%
    \FormulaRefAuto{A \subseteq B \coloneq \forall x\,(x\in A \rightarrow x\in B)}{\rUI{9}}}

  \proofstep{1}{(x,y)\in F\land (x,z)\in F\rightarrow y=z}{%
    \FormulaRefAuto{(x,y)\in F\dsep (x,z)\in F \vdash y=z}{1}}
  \proofstep{1}{x\in A \rightarrow \exists y\,(x,y)\in F}{%
    \FormulaRefAuto{x\in A \vdash \exists y\,(x,y)\in F}{1}}

  \proofstep{1,2}{F\colon A\to C}{\FormulaRefAuto{Funktion}{10,11,12}}
\end{tabproof}

\FormulaThmDelta[Erweiterung der Zielmenge]{%
  F\colon A\to B\dsep B\subseteq C\vdash F\colon A\to C
}{%
  \DeltaRow{Mengen}{A\dsep B\dsep C\dsep x}
  \DeltaRow{Funktionen}{F}
}
\begin{tabproof}
  \proofstep{1}{F\colon A\to B}{\rA}
  \proofstep{2}{B\subseteq C}{\rA}
  \proofstep{3}{x\in A}{\rA}
  \proofstep{1,3}{F(x)\in B}{%
    \FormulaRefAuto{F\colon A\to B\dsep x\in A\vdash F(x)\in B}{1,3}}
  \proofstep{2,3}{F(x)\in C}{%
    \FormulaRefAuto{A\subseteq B,\, x\in A \vdash x\in B}{2,4}}
  \proofstep{1,2}{\forall x\in A\,F(x)\in C}{\rUI{\rRI{3,5}}}
  \proofstep{1,2}{F\colon A\to C}{%
    \FormulaRefAuto{F\colon A\to B\dsep \forall x\in A\,F(x)\in C\vdash F\colon A\to C}{1,6}}
\end{tabproof}

\FormulaThmDelta[Verkleinerung der Zielmenge erhält Injektivität]{%
  F\colon A\inj B\dsep \forall x\in A\,F(x)\in C
  \vdash
  F\colon A\inj C
}{%
  \DeltaDecl{Mengen}{A\dsep B\dsep C}%
  \DeltaDecl{Funktionen}{F}%
}
\begin{tabproof}
  \proofstep{1}{F\colon A\inj B}{\rA}
  \proofstep{2}{\forall x\in A\,F(x)\in C}{\rA}
  \proofstep{1}{F\colon A\to B}{\FormulaRefAuto{Injektive Funktion}{1}}
  \proofstep{1,2}{F\colon A\to C}{\FormulaRefAuto{F\colon A\to B\dsep \forall x\in A\,F(x)\in C\vdash F\colon A\to C}{3,2}}
  \proofstep{1}{x\in A\land y\in A\land F(x)=F(y)\rightarrow x=y}{%
    \FormulaRefAuto{F\colon A\inj B\dsep x\in A\dsep y\in A\dsep F(x)=F(y)\vdash x=y}{1}}
  \proofstep{1,2}{F\colon A\inj C}{\FormulaRefAuto{Injektive Funktion}{4,5}}
\end{tabproof}


\subsubsection{Einschränkung auf das Urbild einer Teilmenge}

\FormulaThmDelta[Einschränkung auf ein Urbild als Funktion]{%
  F\colon A\to B\dsep T\subseteq B \vdash F\restriction_{F^{-1}[T]}\colon F^{-1}[T]\to T
}{
  \DeltaRow{Mengen}{A\dsep B\dsep T\dsep x}
  \DeltaRow{Funktionen}{F}
}
\begin{tabproof}
  \proofstep{1}{F\colon A\to B}{\rA}
  \proofstep{2}{T\subseteq B}{\rA}
  \proofstep{1,2}{F^{-1}[T]\subseteq A}{%
    \FormulaRefAuto{F\colon A\to B\dsep C\subseteq B \vdash F^{-1}[C]\subseteq A}{1,2}}
  \proofstep{1,2}{F\restriction_{F^{-1}[T]}\colon F^{-1}[T]\to B}{%
    \FormulaRefAuto{C\subseteq A \vdash F\restriction_C\colon C\to B}{3}}
  \proofstep{5}{x\in F^{-1}[T]}{\rA}
  \proofstep{1,2,5}{x\in F^{-1}[T]\leftrightarrow (x\in A \land F(x)\in T)}{%
    \FormulaRefAuto{C\subseteq B \vdash \bigl(x\in F^{-1}[C] \leftrightarrow (x\in A \land F(x)\in C)\bigr)}{2}}
  \proofstep{1,2,5}{x\in A \land F(x)\in T}{%
    \FormulaRefAuto{P \leftrightarrow Q, P \vdash Q}{6,5}}
  \proofstep{1,2,5}{F(x)\in T}{\rAEb{7}}
  \proofstep{1,2,5}{F\restriction_{F^{-1}[T]}(x)=F(x)}{%
    \FormulaRefAuto{C\subseteq A \dsep x\in C \vdash F\restriction_C(x)=F(x)}{3,5}}
  \proofstep{1,2,5}{F\restriction_{F^{-1}[T]}(x)\in T}{\rIE{9,8}}
  \proofstep{1,2}{\forall x\in F^{-1}[T]\,F\restriction_{F^{-1}[T]}(x)\in T}{%
    \rUI{\rRI{5,10}}}
  \proofstep{1,2}{F\restriction_{F^{-1}[T]}\colon F^{-1}[T]\to T}{%
    \FormulaRefAuto{F\colon A\to B\dsep \forall x\in A\,F(x)\in C \vdash F\colon A\to C}{4,11}}
\end{tabproof}

\FormulaThmDelta[Einschränkung auf ein Urbild als injektive Funktion]{%
  F\colon A\inj B\dsep T\subseteq B \vdash F\restriction_{F^{-1}[T]}\colon F^{-1}[T]\inj T
}{
  \DeltaRow{Mengen}{A\dsep B\dsep T\dsep x}
  \DeltaRow{Funktionen}{F}
}
\begin{tabproof}
  \proofstep{1}{F\colon A\inj B}{\rA}
  \proofstep{2}{T\subseteq B}{\rA}
  \proofstep{1}{F\colon A\to B}{\FormulaRefAuto{Injektive Funktion}{1}}
  \proofstep{2,3}{F^{-1}[T]\subseteq A}{%
    \FormulaRefAuto{F\colon A\to B\dsep C\subseteq B \vdash F^{-1}[C]\subseteq A}{3,2}}
  \proofstep{1,4}{F\restriction_{F^{-1}[T]}\colon F^{-1}[T]\inj B}{%
    \FormulaRefAuto{F\colon A\inj B\dsep C\subseteq A \vdash F\restriction_C\colon C\inj B}{1,4}}
  \proofstep{2,3}{F\restriction_{F^{-1}[T]}\colon F^{-1}[T]\to T}{%
    \FormulaRefAuto{F\colon A\to B\dsep T\subseteq B \vdash F\restriction_{F^{-1}[T]}\colon F^{-1}[T]\to T}{3,2}}
  \proofstep{7}{x\in F^{-1}[T]}{\rA}
  \proofstep{2,6,7}{F\restriction_{F^{-1}[T]}(x)\in T}{%
    \FormulaRefAuto{F\colon A\to B\dsep x\in A\vdash F(x)\in B}{6,7}}
  \proofstep{2,6}{\forall x\in F^{-1}[T]\,F\restriction_{F^{-1}[T]}(x)\in T}{%
    \rUI{\rRI{7,8}}}
  \proofstep{1,2}{F\restriction_{F^{-1}[T]}\colon F^{-1}[T]\inj T}{%
    \FormulaRefAuto{F\colon A\inj B\dsep \forall x\in A\,F(x)\in C \vdash F\colon A\inj C}{5,9}}
\end{tabproof}

\FormulaThmDelta[Zeugen im Urbild einer Teilmenge]{%
  F\colon A\sur B \dsep T\subseteq B \dsep y\in T \vdash
  \exists x\in F^{-1}[T]\,F\restriction_{F^{-1}[T]}(x)=y
}{
  \DeltaRow{Mengen}{A\dsep B\dsep T\dsep x\dsep y}
  \DeltaRow{Funktionen}{F}
}
\begin{tabproof}
  \proofstep{1}{F\colon A\sur B}{\rA}
  \proofstep{2}{T\subseteq B}{\rA}
  \proofstep{3}{y\in T}{\rA}
  \proofstep{1}{F\colon A\to B}{\FormulaRefAuto{Surjektive Funktion}{1}}
  \proofstep{2,4}{F^{-1}[T]\subseteq A}{%
    \FormulaRefAuto{F\colon A\to B\dsep C\subseteq B \vdash F^{-1}[C]\subseteq A}{4,2}}
  \proofstep{2,3}{y\in B}{%
    \FormulaRefAuto{A\subseteq B, x\in A \vdash x\in B}{2,3}}
  \proofstep{1,2,3}{\exists x\in A\,F(x)=y}{%
    \FormulaRefAuto{y\in B\vdash \exists x\in A\; F(x)=y}{6}}
  \proofstep{8}{x\in A \land F(x)=y}{\rA}
  \proofstep{8}{x\in A}{\rAEa{8}}
  \proofstep{8}{F(x)=y}{\rAEb{8}}
  \proofstep{3,8}{F(x)\in T}{\rIE{10,3}}
  \proofstep{4,2}{x\in F^{-1}[T]\leftrightarrow (x\in A \land F(x)\in T)}{%
    \FormulaRefAuto{C\subseteq B \vdash \bigl(x\in F^{-1}[C] \leftrightarrow (x\in A \land F(x)\in C)\bigr)}{2}}
  \proofstep{8}{x\in A \land F(x)\in T}{\rAI{9,11}}
  \proofstep{2,8}{x\in F^{-1}[T]}{%
    \FormulaRefAuto{P \leftrightarrow Q, Q \vdash P}{12,13}}
  \proofstep{5,8}{F\restriction_{F^{-1}[T]}(x)=F(x)}{%
    \FormulaRefAuto{C\subseteq A \dsep x\in C \vdash F\restriction_C(x)=F(x)}{5,14}}
  \proofstep{5,8}{F\restriction_{F^{-1}[T]}(x)=y}{%
    \FormulaRefAuto{a = b,\, b = c \vdash a = c}{15,10}}
  \proofstep{5,8}{\exists x\in F^{-1}[T]\,F\restriction_{F^{-1}[T]}(x)=y}{%
    \rEI{\rAI{14,16}}}
  \proofstep{1,2,3}{\exists x\in F^{-1}[T]\,F\restriction_{F^{-1}[T]}(x)=y}{%
    \rEE{7,8,17}}
\end{tabproof}

\FormulaThmDelta[Einschränkung auf ein Urbild als surjektive Funktion]{%
  F\colon A\sur B \dsep T\subseteq B \vdash F\restriction_{F^{-1}[T]}\colon F^{-1}[T]\sur T
}{
  \DeltaRow{Mengen}{A\dsep B\dsep T\dsep y}
  \DeltaRow{Funktionen}{F}
}
\begin{tabproof}
  \proofstep{1}{F\colon A\sur B}{\rA}
  \proofstep{2}{T\subseteq B}{\rA}
  \proofstep{1}{F\colon A\to B}{\FormulaRefAuto{Surjektive Funktion}{1}}
  \proofstep{2,3}{F\restriction_{F^{-1}[T]}\colon F^{-1}[T]\to T}{%
    \FormulaRefAuto{F\colon A\to B\dsep T\subseteq B \vdash F\restriction_{F^{-1}[T]}\colon F^{-1}[T]\to T}{3,2}}
  \proofstep{1,2}{\forall y\in T\,\exists x\in F^{-1}[T]\,F\restriction_{F^{-1}[T]}(x)=y}{%
    \FormulaRefAuto{F\colon A\sur B \dsep T\subseteq B \dsep y\in T \vdash \exists x\in F^{-1}[T]\,F\restriction_{F^{-1}[T]}(x)=y}{1,2}}
  \proofstep{1,2}{F\restriction_{F^{-1}[T]}\colon F^{-1}[T]\sur T}{%
    \FormulaRefAuto{Surjektive Funktion}{4,5}}
\end{tabproof}

\FormulaThmDelta[Einschränkung auf ein Urbild als bijektive Funktion]{%
  F\colon A\bij B \dsep T\subseteq B \vdash F\restriction_{F^{-1}[T]}\colon F^{-1}[T]\bij T
}{
  \DeltaRow{Mengen}{A\dsep B\dsep T}
  \DeltaRow{Funktionen}{F}
}
\begin{tabproof}
  \proofstep{1}{F\colon A\bij B}{\rA}
  \proofstep{2}{T\subseteq B}{\rA}
  \proofstep{1}{F\colon A\inj B}{\FormulaRefAuto{Bijektive Funktion}{1}}
  \proofstep{1}{F\colon A\sur B}{\FormulaRefAuto{Bijektive Funktion}{1}}
  \proofstep{2,3}{F\restriction_{F^{-1}[T]}\colon F^{-1}[T]\inj T}{%
    \FormulaRefAuto{F\colon A\inj B\dsep T\subseteq B \vdash F\restriction_{F^{-1}[T]}\colon F^{-1}[T]\inj T}{3,2}}
  \proofstep{2,4}{F\restriction_{F^{-1}[T]}\colon F^{-1}[T]\sur T}{%
    \FormulaRefAuto{F\colon A\sur B \dsep T\subseteq B \vdash F\restriction_{F^{-1}[T]}\colon F^{-1}[T]\sur T}{4,2}}
  \proofstep{1,2}{F\restriction_{F^{-1}[T]}\colon F^{-1}[T]\bij T}{%
    \FormulaRefAuto{Bijektive Funktion}{5,6}}
\end{tabproof}

\subsubsection{Beispiele}

\paragraph{Projektionen}

\FormulaThmDelta[Surjektivität]
{x\in A\vdash \exists y\in R(\pi_1\restriction_R(y)=x)}
{
  \DeltaRow{Mengen}{A\dsep B\dsep x}
  \DeltaPrem{Totale Relationen}{\TotRel{R,A,B}}
}
\begin{tabproofwide}
  \proofstepwidestar[1]{x\in A}{\rA}
  \proofstepwidestar[]{R\subseteq A\times B}{\FormulaRefAuto{F \subseteq A \times B}}
  \proofstepwidestar[1]{\exists x' \, (x,x')\in R}{\FormulaRefAuto{x\in A \vdash \exists y\,(x,y)\in F}{1}}
  \proofstepwidestar[4]{(x,x')\in R}{\rA}
  \proofstepwidestar[4]{(x,x')\in A\times B}{\FormulaRefAuto{A\subseteq B, x\in A\vdash x\in B}{2,4}}
  \proofstepwide[4]{\pi_1\restriction_R(x,x')}{=}{\pi_1(x,x')}{\FormulaRefAuto{C\subseteq A \dsep x\in C \vdash F\restriction_C(x)=F(x)}{2,4}}
  \proofstepwide[4]{}{=}{x}{\FormulaRefAuto{(x,y)\in A\times B \vdash \pi_1(x,y)=x}{5}}
  \proofstepwide[4]{\pi_1\restriction_R(x,x')}{=}{x}{\rChain{6,7}}
  \proofstepwidestar[4]{\exists y\in R \, \pi_1\restriction_R(y)=x}{\rEI{\rAI{4,8}}}
  \proofstepwidestar[1]{\exists y\in R \, \pi_1\restriction_R(y)=x}{\rEE{3,4,9}}
\end{tabproofwide}

\FormulaThmDelta[Surjektive Funktion]{%
\pi_1\restriction_R\colon R\sur A
}{
  \DeltaRow{Mengen}{A\dsep B}
  \DeltaPrem{Totale Relationen}{\TotRel{R,A,B}}
}
\begin{tabproof}
    \proofstep{}{\pi_1\restriction_R\colon R\to A}{\FormulaRefAuto{C\subseteq A \vdash F\restriction_C\colon C\to B}{\FormulaRefAuto{R\subseteq A\times B}}}
    \proofstep{}{\forall x\in A \exists y\in R\,(\pi_1\restriction_R(y)=x)}{\FormulaRefAuto{x\in A\vdash \exists y\in R(\pi_1\restriction_R(y)=x)}}
    \proofstep{}{\pi_1\restriction_R\colon R\sur A}{\FormulaRefAuto{Surjektive Funktion}{1,2}}
\end{tabproof}

\subsection{Korestriktionen}



\FormulaThmDelta[Korestriktion auf das Bild ist surjektiv]{%
  F\colon A\to B \vdash F\colon A\sur F[A]
}{%
  \DeltaRow{Mengen}{A\dsep B\dsep x\dsep y\dsep z}
  \DeltaRow{Funktionen}{F}
}
\begin{tabproof}
  \proofstep{1}{F\colon A\to B}{\rA}

  \proofstep{2}{x\in A}{\rA}
  \proofstep{}{F(x)=F(x)}{\rII}
  \proofstep{2}{\exists z\in A\,F(x)=F(z)}{\rEI{\rAI{2,3}}}
  \proofstep{1,2}{F(x)\in F[A]}{%
    \FormulaRefAuto{F\colon A\to B\dsep \exists x\in A\,y=F(x)\vdash y\in F[A]}{1,4}}
  \proofstep{1}{\forall x\in A\,F(x)\in F[A]}{\rUI{\rRI{2,5}}}
  \proofstep{1}{F\colon A\to F[A]}{%
    \FormulaRefAuto{F\colon A\to B\dsep \forall x\in A\,F(x)\in C \vdash F\colon A\to C}{1,6}}

  \proofstep{8}{y\in F[A]}{\rA}
  \proofstep{1,8}{\exists x\in A\,y=F(x)}{%
    \FormulaRefAuto{F\colon A\to B\dsep y\in F[A] \vdash \exists x\in A\,y=F(x)}{1,8}}
  \proofstep{10}{x\in A\land y=F(x)}{\rA}
  \proofstep{10}{x\in A}{\rAEa{10}}
  \proofstep{10}{y=F(x)}{\rAEb{10}}
  \proofstep{10}{F(x)=y}{\FormulaRefAuto{a = b \vdash b = a}{12}}
  \proofstep{10}{\exists x\in A\,F(x)=y}{\rEI{\rAI{11,13}}}
  \proofstep{1,8}{\exists x\in A\,F(x)=y}{\rEE{9,10,14}}
  \proofstep{1}{\forall y\in F[A]\,\exists x\in A\,F(x)=y}{\rUI{\rRI{8,15}}}

  \proofstep{1}{F\colon A\sur F[A]}{\FormulaRefAuto{Surjektive Funktion}{7,16}}
\end{tabproof}

\FormulaThmDelta[Korestriktion auf das Bild einer injektiven Funktion ist bijektiv]{%
  F\colon A\inj B \vdash F\colon A\bij F[A]
}{%
  \DeltaRow{Mengen}{A\dsep B\dsep x\dsep z}
  \DeltaRow{Funktionen}{F}
}
\begin{tabproof}
  \proofstep{1}{F\colon A\inj B}{\rA}
  \proofstep{1}{F\colon A\to B}{\FormulaRefAuto{Injektive Funktion}{1}}
  \proofstep{3}{x\in A}{\rA}
  \proofstep{3}{F(x)=F(x)}{\rII}
  \proofstep{1,3}{\exists z\in A\,F(x)=F(z)}{\rEI{\rAI{3,4}}}
  \proofstep{1,3}{F(x)\in F[A]}{%
    \FormulaRefAuto{F\colon A\to B\dsep \exists x\in A\,y=F(x)\vdash y\in F[A]}{2,5}}
  \proofstep{1}{\forall x\in A\,F(x)\in F[A]}{\rUI{\rRI{3,6}}}
  \proofstep{1}{F\colon A\inj F[A]}{%
    \FormulaRefAuto{F\colon A\inj B\dsep \forall x\in A\,F(x)\in C \vdash F\colon A\inj C}{1,7}}
  \proofstep{1}{F\colon A\sur F[A]}{%
    \FormulaRefAuto{F\colon A\to B \vdash F\colon A\sur F[A]}{2}}
  \proofstep{1}{F\colon A\bij F[A]}{%
    \FormulaRefAuto{F\colon A\inj B\dsep F\colon A\sur B \vdash F\colon A\bij B}{8,9}}
\end{tabproof}




\subsection{Die Umkehrfunktion}

\subsubsection{Definition der Umkehrfunktionsvorschrift}

\FormulaDefDelta[Funktionsvorschrift für die Umkehrfunktion von \(F\)]
{y\in B \vdash t_{F^{-1}}(y) \coloneqq \iota x\bigl(x\in A \land F(x)=y\bigr)}
{
  \DeltaRow{Mengen}{A\dsep B\dsep x\dsep y}
  \DeltaPrem{Funktionen}{F\colon A\bij B}
  \DeltaRow{Funktionensymbol}{t_{F^{-1}}}
}

% — Typisierungsaxiom / Wertebereich —
\FormulaThmDelta[Typisierungsaxiom für \(t_{F^{-1}}\)]
{y\in B \vdash t_{F^{-1}}(y)\in A}
{
  \DeltaRow{Mengen}{A\dsep B\dsep x\dsep y}
  \DeltaPrem{Funktionen}{F\colon A\bij B}
  \DeltaRow{Funktionensymbol}{t_{F^{-1}}}
}
\begin{tabproof}
  \proofstep{1}{y\in B}{\rA}
  \proofstep{1}{t_{F^{-1}}(y)\in A}{\rAEa{\FormulaRefAuto{y\in B \vdash t_{F^{-1}}(y) \coloneqq \iota x\bigl(x\in A \land F(x)=y\bigr)}{1}}}
\end{tabproof}

\FormulaThmDelta
{y\in B \vdash F(t_{F^{-1}}(y))=y}
{
  \DeltaRow{Mengen}{A\dsep B\dsep x\dsep y}
  \DeltaPrem{Funktionen}{F\colon A\bij B}
  \DeltaRow{Funktionensymbol}{t_{F^{-1}}}
}
\begin{tabproof}
  \proofstep{1}{y\in B}{\rA}
  \proofstep{1}{F(t_{F^{-1}}(y))=y}{\rAEb{\FormulaRefAuto{y\in B \vdash t_{F^{-1}}(y) \coloneqq \iota x\bigl(x\in A \land F(x)=y\bigr)}{1}}}
\end{tabproof}

\subsubsection{Die Umkehrfunktion als Funktion}

\FormulaThmDelta[Umkehrfunktion als Funktion]{%
\exists! G\colon B\to A\,\forall y\in B\,G(y) = t_{F^{-1}}(y)%
}{
  \DeltaRow{Mengen}{A\dsep B\dsep x\dsep y}
  \DeltaPrem{Funktionen}{F\colon A\bij B}
  \DeltaRow{Funktionensymbol}{t_{F^{-1}}}
}
\begin{tabproof}
  \proofstep{}{\forall y\in B\,t_{F^{-1}}(y)=\iota x\bigl(x\in A\land F(x)=y\bigr)}{\FormulaRefAuto{y\in B \vdash t_{F^{-1}}(y) \coloneqq \iota x\bigl(x\in A \land F(x)=y\bigr)}}
  \proofstep{}{\forall y\in B\,t_{F^{-1}}(y)\in A}{\FormulaRefAuto{y\in B \vdash t_{F^{-1}}(y)\in A}}
  \proofstep{}{\exists! G\colon B\to A\,\forall y\in B\,G(y) = t_{F^{-1}}(y)}{\FormulaRefAuto{x\in A\vdash t(x)\coloneq y\dsep x\in A\vdash t(x)\in B\exists! F\colon A\to B\,\forall x\in A\,F(x) = t(x)}{1,2}}
\end{tabproof}

\FormulaDefDelta[Umkehrfunktion von \(F\)]
{F^{-1}\coloneqq\iota G\left(G\colon B\to A \land 
  \forall y\in B\, G(y)=t_{F^{-1}}(y)\right)}
{
  \DeltaRow{Mengen}{A\dsep B}
  \DeltaPrem{Funktionen}{F\colon A\bij B}
  \DeltaRow{Funktionensymbol}{t_{F^{-1}}}
}

\FormulaThmDeltaR[Umkehrfunktion von \(F\)]
{F^{-1}\colon B\to A}{F\colon A\bij B\vdash F^{-1}\colon B\to A}
{
  \DeltaRow{Mengen}{A\dsep B}
  \DeltaPrem{Funktionen}{F\colon A\bij B}
}
\begin{tabproof}
  \proofstep{}{F^{-1}\colon B\to A}{%
    \rAEa{\FormulaRefAuto{\exists! G\colon B\to A\,\forall y\in B\,G(y) = t_{F^{-1}}(y)}}}
\end{tabproof}

\FormulaThmDelta
{\TotRel{F^{-1},B,A}}
{
  \DeltaRow{Mengen}{A\dsep B}
  \DeltaPrem{Funktionen}{F\colon A\bij B}
}
\begin{tabproof}
    \proofstep{}{F^{-1}\colon B\to A}{\FormulaRefAuto{F\colon A\bij B\vdash F^{-1}\colon B\to A}}
    \proofstep{1}{\TotRel{F^{-1},B,A}}{\FormulaRefAuto{F\colon A\to B\vdash\TotRel{F,A,B}}{1}}
\end{tabproof}

\FormulaThmDelta[Umkehrfunktion von \(F\)]
{y\in B\vdash F^{-1}(y)=t_{F^{-1}}(y)}
{
  \DeltaRow{Mengen}{A\dsep B}
  \DeltaPrem{Funktionen}{F\colon A\bij B}
}
\begin{tabproof}
  \proofstep{1}{y\in B}{\rA}
  \proofstep{1}{F^{-1}(y)=t_{F^{-1}}(y)}{%
    \rRE{1,\rAEb{\FormulaRefAuto{\exists! G\colon B\to A\,\forall y\in B\,G(y) = t_{F^{-1}}(y)}}}}
\end{tabproof}

\FormulaThmDelta[Typisierungsaxiom für \(t_{F^{-1}}\)]
{y\in B \vdash F^{-1}(y)\in A}
{
  \DeltaRow{Mengen}{A\dsep B\dsep x\dsep y}
  \DeltaPrem{Funktionen}{F\colon A\bij B}
}

\subsubsection{Grundgleichungen der Umkehrfunktion}

\FormulaThmDelta[Links-Umkehrung]
{x\in A \vdash F^{-1}(F(x)) = x}
{
  \DeltaRow{Mengen}{A\dsep B\dsep x}
  \DeltaPrem{Funktionen}{F\colon A\bij B}
}
\begin{tabproof}
  \proofstep{1}{x\in A}{\rA}
  \proofstep{1}{F(x)\in B}{%
    \FormulaRefAuto{F\colon A\to B\dsep x\in A\vdash F(x)\in B}{1}}
  \proofstep{1}{t_{F^{-1}}(F(x))\in A}{%
    \FormulaRefAuto{y\in B \vdash t_{F^{-1}}(y)\in A}{1}}
  \proofstep{1}{F(t_{F^{-1}}(F(x))) = F(x)}{\FormulaRefAuto{y\in B \vdash F(t_{F^{-1}}(y))=y}{2}}
  \proofstep{1}{F^{-1}(F(x))=t_{F^{-1}}(F(x))}{\FormulaRefAuto{y\in B\vdash F^{-1}(y)=t_{F^{-1}}(y)}{2}}
  \proofstep{1}{t_{F^{-1}}(F(x))=x}{%
    \FormulaRefAuto{F\colon A\inj B\dsep x\in A\dsep y\in A\dsep F(x)=F(y)\vdash x=y}{%
      \FormulaRefAuto{F\colon A\bij B \vdash F\colon A\inj B},3,2,4}}
  \proofstep{1}{F^{-1}(F(x))=x}{\FormulaRefAuto{a=b,b=c\vdash a=c}{5,6}}
\end{tabproof}

\FormulaThmDeltaR[Rechts-Umkehrung]
{y\in B \vdash F(F^{-1}(y)) = y}{F\colon A\bij B \dsep y\in B \vdash F(F^{-1}(y)) = y}
{
  \DeltaRow{Mengen}{A\dsep B\dsep y}
  \DeltaPrem{Funktionen}{F\colon A\bij B}
}
\begin{tabproof}
  \proofstep{1}{y\in B}{\rA}
  \proofstep{1}{t_{F^{-1}}(y)=F^{-1}(y)}{%
    \FormulaRefAuto{a=b\vdash b=a}{\FormulaRefAuto{y\in B\vdash F^{-1}(y)=t_{F^{-1}}(y)}{1}}
  }
  \proofstep{1}{F(t_{F^{-1}}(y))=y}{%
    \FormulaRefAuto{y\in B \vdash F(t_{F^{-1}}(y))=y}{2}}
  \proofstep{1}{F(F^{-1}(y))=y}{%
    \rIE{2,3}}
\end{tabproof}

\subsubsection{Mengenbezogene Grundgleichungen der Umkehrfunktion}

\FormulaThmDelta[Bild des Urbilds einer Teilmenge unter einer Bijektion]{%
  Y\subseteq B \dsep F\colon A\bij B \vdash F[F^{-1}[Y]]=Y
}{%
  \DeltaRow{Mengen}{A\dsep B\dsep Y}
  \DeltaRow{Funktionen}{F}
}
\begin{tabproof}
  \proofstep{1}{Y\subseteq B}{\rA}
  \proofstep{2}{F\colon A\bij B}{\rA}
  \proofstep{1,2}{F^{-1}[Y]\subseteq A}{%
    \FormulaRefAuto{C\subseteq B\dsep F\colon A\bij B \vdash F^{-1}[C]\subseteq A}{1,2}}
  \proofstep{2}{F\colon A\to B}{\FormulaRefAuto{Bijektive Funktion}{2}}
  \proofstep{1,2}{F\restriction_{F^{-1}[Y]}\colon F^{-1}[Y]\bij Y}{%
    \FormulaRefAuto{F\colon A\bij B \dsep T\subseteq B \vdash F\restriction_{F^{-1}[T]}\colon F^{-1}[T]\bij T}{2,1}}
  \proofstep{1,2}{(F\restriction_{F^{-1}[Y]})[F^{-1}[Y]]=Y}{%
    \FormulaRefAuto{F\colon A\bij B \vdash F[A]=B}{5}}
  \proofstep{1,2}{(F\restriction_{F^{-1}[Y]})[F^{-1}[Y]]=F[F^{-1}[Y]]}{%
    \FormulaRefAuto{C\subseteq A \dsep F\colon A\to B \vdash (F\restriction_C)[C]=F[C]}{3,4}}
  \proofstep{1,2}{F[F^{-1}[Y]]=(F\restriction_{F^{-1}[Y]})[F^{-1}[Y]]}{%
    \FormulaRefAuto{a=b\vdash b=a}{7}}
  \proofstep{1,2}{F[F^{-1}[Y]]=Y}{%
    \FormulaRefAuto{a=b,\, b=c \vdash a=c}{8,6}}
\end{tabproof}

\subsubsection{Die Bijektionseigenschaften der Umkehrfunktion}

\FormulaThmDelta[Injektivität von \(F^{-1}\)]%
{x\in B\dsep y\in B\dsep F^{-1}(x)=F^{-1}(y)\vdash x=y}%
{
  \DeltaRow{Mengen}{x\dsep y\dsep A\dsep B}
  \DeltaPrem{Bijektive Funktionen}{F\colon A\bij B}
}
\begin{tabproof}
  \proofstep{1}{x\in B}{\rA}
  \proofstep{2}{y\in B}{\rA}
  \proofstep{3}{F^{-1}(x)=F^{-1}(y)}{\rA}
  \proofstep{}{F^{-1}\colon B\to A}{\FormulaRefAuto{F\colon A\bij B\vdash F^{-1}\colon B\to A}}
  \proofstep{1}{F^{-1}(x)\in A}{\FormulaRefAuto{F\colon A\to B\dsep x\in A\vdash F(x)\in B}{1,4}}
  \proofstep{2}{F^{-1}(y)\in A}{\FormulaRefAuto{F\colon A\to B\dsep x\in A\vdash F(x)\in B}{2,4}}
  \proofstep{1}{F(F^{-1}(x))=x}{%
    \FormulaRefAuto{F\colon A\bij B \dsep y\in B \vdash F(F^{-1}(y)) = y}{1}}
  \proofstep{2}{F(F^{-1}(y))=y}{%
    \FormulaRefAuto{F\colon A\bij B \dsep y\in B \vdash F(F^{-1}(y)) = y}{2}}
  \proofstep{1,2,3}{F(F^{-1}(x))=F(F^{-1}(y))}{%
    \FormulaRefAuto{x\in A\dsep y\in A\dsep x=y\vdash F(x)=F(y)}{3}}
  \proofstep{1,2,3}{x=y}{%
    \FormulaRefAuto{a = b\dsep a = c\dsep b=d\vdash c = d}{9,7,8}}
\end{tabproof}

\FormulaThmDelta[Surjektivität von \(F^{-1}\)]%
{x\in A\vdash \exists y\in B\; F^{-1}(y)=x}%
{
  \DeltaRow{Mengen}{x\dsep y\dsep A\dsep B}
  \DeltaPrem{Bijektive Funktionen}{F\colon A\bij B}
}
\begin{tabproof}
  \proofstep{1}{x\in A}{\rA}
  % Typisierung von F:
  \proofstep{1}{F(x)\in B}{%
    \FormulaRefAuto{F\colon A\to B\dsep x\in A\vdash F(x)\in B}{1}}
  % Links-Umkehrung: F^{-1}(F(x))=x
  \proofstep{1}{F^{-1}(F(x))=x}{%
    \FormulaRefAuto{x\in A \vdash F^{-1}(F(x)) = x}{1}}
  % Zeugenbildung y \coloneq F(x)
  \proofstep{1}{\exists y\in B\; F^{-1}(y)=x}{%
    \rEI{\rAI{2,3}}}
\end{tabproof}

\FormulaThmDeltaR[\(F^{-1}\) als bijektive Funktion]%
{F^{-1}\colon B\bij A}{F\colon A\bij B\vdash F^{-1}\colon B\bij A}%
{
  \DeltaRow{Mengen}{A\dsep B}
  \DeltaPrem{Bijektive Funktionen}{F\colon A\bij B}
}
\begin{tabproof}
  \proofstep{}{F^{-1}\colon B\to A}{\FormulaRefAuto{F\colon A\bij B\vdash F^{-1}\colon B\to A}}
  \proofstep{}{\forall x,y\in B\, (F^{-1}(x)=F^{-1}(y)\rightarrow x=y)}{%
    \FormulaRefAuto{x\in B\dsep y\in B\dsep F^{-1}(x)=F^{-1}(y)\vdash x=y}}
  \proofstep{}{\forall x\in A\exists y\in B\, F^{-1}(y)=x}{%
    \FormulaRefAuto{x\in A\vdash \exists y\in B\; F^{-1}(y)=x}}
  \proofstep{}{F^{-1}\colon B\bij A}{%
    \FormulaRefAuto{Bijektive Funktion}{1,2,3}}
\end{tabproof}


\FormulaThmDelta[Eigeninversität der Identität]{%
\Id_A^{-1}=\Id_A
}{
  \DeltaRow{Mengen}{A}
}
\begin{tabproof}
  \proofstep{}{%
    \Id_A\colon A\bij A}{%
    \FormulaRefAuto{\Id_A\colon A\bij A}}

  \proofstep{2}{%
    x\in A}{%
    \rA}

  \proofstep{2}{%
    \Id_A^{-1}(\Id_A(x)) = x}{%
    \FormulaRefAuto{x\in A \vdash F^{-1}(F(x)) = x}{1,2}}

  \proofstep{2}{%
    \Id_A(x)=x}{%
    \FormulaRefAuto{x\in A\vdash \Id_A(x)=x}{2}}

  \proofstep{2}{%
    \Id_A^{-1}(x)=x}{%
    \rIE{4,3}}

  \proofstep{2}{%
    \Id_A(x)=\Id_A(x)}{%
    \rII}

  \proofstep{2}{%
    x=\Id_A(x)}{%
    \rIE{4,6}}

  \proofstep{2}{%
    \Id_A^{-1}(x)=\Id_A(x)}{%
    \rIE{7,5}}

  \proofstep{}{%
    x\in A \rightarrow \Id_A^{-1}(x)=\Id_A(x)}{%
    \rRI{2,8}}

  \proofstep{}{%
    \forall x\in A\, \Id_A^{-1}(x)=\Id_A(x)}{%
    \rUI{9}}

  \proofstep{}{%
    \forall x\in A\,(\Id_A^{-1}(x)=\Id_A(x)) \leftrightarrow \Id_A^{-1}=\Id_A}{%
    \FormulaRefAuto{\forall x\in A (F(x)=G(x)) \eqvdash F=G}}

  \proofstep{}{%
    \Id_A^{-1}=\Id_A}{%
    \FormulaRefAuto{P\leftrightarrow Q, P\vdash Q}{11,10}}
\end{tabproof}


\subsubsection{Weitere Eigenschaften}

\FormulaThmDelta{%
\forall y\in B\,\bigl(F(G(y))=y\bigr) \vdash G = F^{-1}
}{
  \DeltaRow{Mengen}{x \dsep y \dsep A \dsep B}
  \DeltaPrem{Bijektive Funktionen}{F\colon A\bij B}
  \DeltaPrem{Funktionen}{G\colon B\to A}
}
\begin{tabproof}
  \proofstep{1}{\forall y\in B\,\bigl(F(G(y))=y\bigr)}{\rA}
  \proofstep{2}{y\in B}{\rA}
  \proofstep{1,2}{F(G(y))=y}{\rRE{\rUE{1},2}}
  \proofstep{1,2}{F(F^{-1}(y))=y}{\FormulaRefAuto{F\colon A\bij B \dsep y\in B \vdash F(F^{-1}(y)) = y}{2}}
  \proofstep{2}{G(y)\in A}{\FormulaRefAuto{F\colon A\to B\dsep x\in A\vdash F(x)\in B}{2}}
  \proofstep{2}{F^{-1}(y)\in A}{\FormulaRefAuto{F\colon A\to B\dsep x\in A\vdash F(x)\in B}{2}}
  \proofstep{1,2}{G(y)=F^{-1}(y)}{%
    \FormulaRefAuto{F\colon A\inj B\dsep x\in A\dsep y\in A\dsep F(x)=F(y)\vdash x=y}{%
      \FormulaRefAuto{F\colon A\bij B \vdash F\colon A\inj B},5,6,3,4}}
  \proofstep{1}{\forall y\,\Bigl(G(y)=F^{-1}(y)\Bigr)}{\rUI{\rRI{2,7}}}
  \proofstep{1}{G=F^{-1}}{\FormulaRefAuto{\forall x\in A (F(x)=G(x)) \eqvdash F=G}{8}}
\end{tabproof}

\FormulaThmDelta{%
\forall x\in A\,\bigl(F^{-1}(G(x))=x\bigr) \vdash G = F
}{
  \DeltaRow{Mengen}{x \dsep A \dsep B}
  \DeltaPrem{Bijektive Funktionen}{F\colon A\bij B}
  \DeltaPrem{Funktionen}{G\colon A\to B}
}
\begin{tabproof}
  \proofstep{1}{\forall x\in A\,\bigl(F^{-1}(G(x))=x\bigr)}{\rA}
  \proofstep{2}{x\in A}{\rA}
  \proofstep{1,2}{F^{-1}(G(x))=x}{\rRE{\rUE{1},2}}
  \proofstep{2}{F^{-1}\bigl(F(x)\bigr)=x}{\FormulaRefAuto{x\in A \vdash F^{-1}(F(x)) = x}{2}}
  \proofstep{2}{G(x)\in B}{\FormulaRefAuto{F\colon A\to B\dsep x\in A\vdash F(x)\in B}{2}}
  \proofstep{2}{F(x)\in B}{\FormulaRefAuto{F\colon A\to B\dsep x\in A\vdash F(x)\in B}{2}}
  \proofstep{1,2}{F^{-1}(G(x))=F^{-1}(F(x))}{\FormulaRefAuto{a=b, c=b\vdash a=c}{3,4}}
  \proofstep{1,2}{G(x)=F(x)}{\FormulaRefAuto{x\in B\dsep y\in B\dsep F^{-1}(x)=F^{-1}(y)\vdash x=y}{5,6,7}}
  \proofstep{1}{\forall x\in A\,\Bigl(G(x)=F(x)\Bigr)}{\rUI{\rRI{2,8}}}
  \proofstep{1}{G=F}{\FormulaRefAuto{\forall x\in A (F(x)=G(x)) \eqvdash F=G}{9}}
\end{tabproof}

\subsection{Die Komposition von Funktionen}

\subsubsection{Funktionsvorschrift und Typisierung}

\FormulaDefDelta[Funktionsvorschrift für die Komposition \(G\circ F\)]
{x\in A \vdash t_{G\circ F}(x) \coloneqq G(F(x))}
{
  \DeltaRow{Mengen}{A\dsep B\dsep C\dsep x}
  \DeltaPrem{Funktionen}{F\colon A\to B\dsep G\colon B\to C}
  \DeltaRow{Funktionensymbol}{t_{G\circ F}}
}

% — Typisierungsaxiom / Wertebereich —
\FormulaThmDelta[Typisierungsaxiom für \(t_{G\circ F}\)]
{x\in A \vdash t_{G\circ F}(x)\in C}
{
  \DeltaRow{Mengen}{A\dsep B\dsep C\dsep x}
  \DeltaPrem{Funktionen}{F\colon A\to B\dsep G\colon B\to C}
  \DeltaRow{Funktionensymbol}{t_{G\circ F}}
}
\begin{tabproof}
  \proofstep{1}{x\in A}{\rA}
  % Typ von F
  \proofstep{1}{F(x)\in B}{%
    \FormulaRefAuto{F\colon A\to B\dsep x\in A\vdash F(x)\in B}{1}}
  % Typ von G
  \proofstep{1}{G(F(x))\in C}{%
    \FormulaRefAuto{F\colon A\to B\dsep x\in A\vdash F(x)\in B}{2}}
  % Definition der Vorschrift
  \proofstep{1}{t_{G\circ F}(x)=G(F(x))}{%
    \FormulaRefAuto{x\in A \vdash t_{G\circ F}(x) \coloneqq G(F(x))}{1}}
  % Wertebereich per Gleichheitselimination
  \proofstep{1}{t_{G\circ F}(x)\in C}{\rIE{4,3}}
\end{tabproof}

\subsubsection{Die Komposition als Funktion}

\FormulaThmDelta[Komposition als Funktion]{%
\exists! H\colon A\to C\,\forall x\in A\,H(x) = t_{G\circ F}(x)%
}{
  \DeltaRow{Mengen}{A\dsep B\dsep C\dsep x}
  \DeltaPrem{Funktionen}{F\colon A\to B\dsep G\colon B\to C}
  \DeltaRow{Funktionensymbol}{t_{G\circ F}}
}
\begin{tabproof}
  % Vorschrift global
  \proofstep{}{\forall x\in A\,t_{G\circ F}(x)=G(F(x))}{%
    \FormulaRefAuto{x\in A \vdash t_{G\circ F}(x) \coloneqq G(F(x))}}
  % Typ global
  \proofstep{}{\forall x\in A\,t_{G\circ F}(x)\in C}{%
    \FormulaRefAuto{x\in A \vdash t_{G\circ F}(x)\in C}}
  % Existenz und Eindeutigkeit der Funktion H: A→C mit dieser Vorschrift
  \proofstep{}{\exists! H\colon A\to C\,\forall x\in A\,H(x) = t_{G\circ F}(x)}{%
    \FormulaRefAuto{x\in A\vdash t(x)\coloneq y\dsep x\in A\vdash t(x)\in B\exists! F\colon A\to B\,\forall x\in A\,F(x) = t(x)}{1,2}}
\end{tabproof}

\subsubsection{Definition und Grundgleichung der Komposition}

\FormulaDefDelta[Komposition $G\circ F$]{%
  G\circ F \coloneqq
  \iota H\Bigl(H\colon A\to C \land \forall x\in A\, H(x)=t_{G\circ F}(x)\Bigr)%
}{
  \DeltaRow{Mengen}{A\dsep B\dsep C\dsep x}
  \DeltaPrem{Funktionen}{F\colon A\to B\dsep G\colon B\to C}
  \DeltaRow{Funktionensymbol}{t_{G\circ F}}
}

\FormulaThmDelta[Komposition als Abbildung]{%
  G\circ F\colon A\to C%
}{
  \DeltaRow{Mengen}{A\dsep B\dsep C}
  \DeltaPrem{Funktionen}{F\colon A\to B\dsep G\colon B\to C}
}
\begin{tabproof}
  \proofstep{}{G\circ F\colon A\to C}{%
    \rAEa{\FormulaRefAuto{G\circ F\coloneqq\iota H\left(H\colon A\to C \land \forall x\in A\, H(x)=t_{G\circ F}(x)\right)}}}
\end{tabproof}

% =========================================================
% Hilfssatz: Codomain-Wechsel (über Inklusion + Komposition)
% Empfohlener Ort in B03: Abschnitt 3.15.11 direkt nach
%   - Def. 3.15.11.8  (\iota_{A,B})
%   - Def./Th. zur Komposition (Def. 3.15.11.16, Th. 3.15.11.79)
% =========================================================

\FormulaThmDelta{%
  B\subseteq C \dsep F\colon A\to B \vdash \iota_{B,C}\circ F\colon A\to C%
}{%
  \DeltaDecl{Mengen}{A\dsep B\dsep C}%
  \DeltaDecl{Funktionen}{F}%
}
\begin{tabproof}
  \proofstep{1}{B\subseteq C}{\rA}
  \proofstep{2}{F\colon A\to B}{\rA}
  \proofstep{1}{\iota_{B,C}\colon B\to C}{%
    \FormulaRefAuto{A\subseteq B\vdash \iota_{A,B}\colon A\to B}{1}}
  \proofstep{1,2}{\iota_{B,C}\circ F\colon A\to C}{%
    \FormulaRefAuto{G\circ F\colon A\to C}{2,3}}
\end{tabproof}

% ------------------------------------------------------------
% (1) Grundgleichung der Komposition (vertauschte Namen)
% ------------------------------------------------------------
\FormulaThmDelta[Grundgleichung der Komposition (für $F\circ G$)]%
{x\in A \vdash (F\circ G)(x)=F(G(x))}%
{
  \DeltaRow{Mengen}{A\dsep B\dsep C\dsep x}
  \DeltaPrem{Funktionen}{G\colon A\to B\dsep F\colon B\to C}
}
\begin{tabproofwide}
  \proofstepwidestar[1]{x\in A}{\rA}

  \proofstepwidestar[]{%
    \forall x\in A\, (F\circ G)(x)=t_{F\circ G}(x)}{%
    \rAEb{\FormulaRefAuto{G\circ F \coloneqq
  \iota H\Bigl(H\colon A\to C \land \forall x\in A\, H(x)=t_{G\circ F}(x)\Bigr)}}}

  \proofstepwide[1]{(F\circ G)(x)}{=}{t_{F\circ G}(x)}{\rRE{1,\rUE{2}}}

  \proofstepwide[1]{t_{F\circ G}(x)}{=}{F(G(x))}{%
    \FormulaRefAuto{x\in A \vdash t_{G\circ F}(x) \coloneqq G(F(x))}{1}}

  \proofstepwide[1]{(F\circ G)(x)}{=}{F(G(x))}{%
    \FormulaRefAuto{a=b,\, b=c \vdash a=c}{3,4}}
\end{tabproofwide}


% ------------------------------------------------------------
% (2) Rückwärts-Variante (saubere Gleichungskette / Symmetrie)
% ------------------------------------------------------------
\FormulaThmDelta%
{x\in A \vdash F(G(x))=(F\circ G)(x)}%
{
  \DeltaRow{Mengen}{A\dsep B\dsep C\dsep x}
  \DeltaPrem{Funktionen}{G\colon A\to B\dsep F\colon B\to C}
}
\begin{tabproof}
  \proofstep{1}{x\in A}{\rA}
  \proofstep{1}{(F\circ G)(x)=F(G(x))}{%
    \FormulaRefAuto{x\in A \vdash (F\circ G)(x)=F(G(x))}{1}}
  \proofstep{1}{F(G(x))=(F\circ G)(x)}{%
    \FormulaRefAuto{a=b\vdash b=a}{2}}
\end{tabproof}

% ------------------------------------------------------------
% (3) Grundgleichung der Komposition für drei Komponenten
% ------------------------------------------------------------
\FormulaThmDelta[Grundgleichung der Komposition für drei Komponenten]%
{x\in A \vdash ((F\circ G)\circ H)(x)=F(G(H(x)))}%
{
  \DeltaRow{Mengen}{A\dsep B\dsep C\dsep D\dsep x}
  \DeltaPrem{Funktionen}{H\colon A\to B\dsep G\colon B\to C\dsep F\colon C\to D}
}
\begin{tabproofwide}
  \proofstepwidestar[1]{x\in A}{\rA}

  \proofstepwidestar[1]{H(x)\in B}{%
    \FormulaRefAuto{H\colon A\to B\dsep x\in A\vdash H(x)\in B}{1}}

  \proofstepwide[1]{((F\circ G)\circ H)(x)}{=}{(F\circ G)(H(x))}{%
    \FormulaRefAuto{x\in A \vdash (F\circ G)(x)=F(G(x))}{1}}

  \proofstepwide[1]{}{=}{F(G(H(x)))}{%
    \FormulaRefAuto{x\in A \vdash (F\circ G)(x)=F(G(x))}{2}}

  \proofstepwide[1]{((F\circ G)\circ H)(x)}{=}{F(G(H(x)))}{\rChain{3,4}}
\end{tabproofwide}


% ------------------------------------------------------------
% (4) Rückwärts-Variante für drei Komponenten
% ------------------------------------------------------------
\FormulaThmDelta%
{x\in A \vdash F(G(H(x)))=((F\circ G)\circ H)(x)}%
{
  \DeltaRow{Mengen}{A\dsep B\dsep C\dsep D\dsep x}
  \DeltaPrem{Funktionen}{H\colon A\to B\dsep G\colon B\to C\dsep F\colon C\to D}
}
\begin{tabproof}
  \proofstep{1}{x\in A}{\rA}
  \proofstep{1}{((F\circ G)\circ H)(x)=F(G(H(x)))}{%
    \FormulaRefAuto{x\in A \vdash ((F\circ G)\circ H)(x)=F(G(H(x)))}{1}}
  \proofstep{1}{F(G(H(x)))=((F\circ G)\circ H)(x)}{%
    \FormulaRefAuto{a=b\vdash b=a}{2}}
\end{tabproof}

\subsubsection{Die Auswertungseigenschaft}

\FormulaThmDelta{%
x\in A\dsep y\in A\dsep G(F(x))=G(F(y))\vdash (G\circ F)(x)=(G\circ F)(y)
}{
  \DeltaRow{Mengen}{x \dsep A \dsep B \dsep C}
  \DeltaPrem{Funktionen}{F\colon A\to B \dsep G\colon B\to C}
}
\begin{tabproof}
\proofstep{1}{x\in A}{\rA}
\proofstep{2}{y\in A}{\rA}
\proofstep{3}{G(F(x))=G(F(y))}{\rA}
\proofstep{1}{(G\circ F)(x)=G(F(x))}{\FormulaRefAuto{x\in A \vdash (G\circ F)(x)=G(F(x))}{1}}
\proofstep{2}{(G\circ F)(y)=G(F(y))}{\FormulaRefAuto{x\in A \vdash (G\circ F)(x)=G(F(x))}{2}}
\proofstep{3}{(G\circ F)(x)=G(F(y))}{\rIE{4,3}}
\proofstep{3}{(G\circ F)(x)=(G\circ F)(y)}{\rIE{5,6}}
\end{tabproof}

\FormulaThmDelta{%
x\in A\dsep y\in A\dsep (G\circ F)(x)=(G\circ F)(y) \vdash G(F(x))=G(F(y))
}{
  \DeltaRow{Mengen}{x \dsep A \dsep B \dsep C}
  \DeltaPrem{Funktionen}{F\colon A\to B \dsep G\colon B\to C}
}
\begin{tabproof}
  \proofstep{1}{x\in A}{\rA}
  \proofstep{2}{y\in A}{\rA}
  \proofstep{3}{(G\circ F)(x)=(G\circ F)(y)}{\rA}

  \proofstep{1}{(G\circ F)(x)=G(F(x))}{\FormulaRefAuto{x\in A \vdash (G\circ F)(x)=G(F(x))}{1}}
  \proofstep{2}{(G\circ F)(y)=G(F(y))}{\FormulaRefAuto{x\in A \vdash (G\circ F)(x)=G(F(x))}{2}}

  \proofstep{3}{G(F(x))=(G\circ F)(y)}{\rIE{4,3}}
  \proofstep{3}{G(F(x))=G(F(y))}{\rIE{5,6}}
\end{tabproof}

\FormulaThmDeltaR[Ungleichheit auf Kompositionswerte heben]{%
  \begin{aligned}[t]
    &x\in A\dsep y\in A\dsep G(F(x))\neq G(F(y))
    \vdash{}\\
    &(G\circ F)(x)\neq (G\circ F)(y)
  \end{aligned}
}{%
  x\in A\dsep y\in A\dsep G(F(x))\neq G(F(y))
  \vdash (G\circ F)(x)\neq (G\circ F)(y)
}{%
  \DeltaDecl{Mengen}{A\dsep x\dsep y}%
  \DeltaDecl{Funktionen}{F\dsep G}%
}
\begin{tabproof}
  \proofstep{1}{x\in A}{\rA}
  \proofstep{2}{y\in A}{\rA}
  \proofstep{3}{G(F(x))\neq G(F(y))}{\rA}
  \proofstep{4}{(G\circ F)(x)=(G\circ F)(y)}{\rA}
  \proofstep{1,2,4}{G(F(x))=G(F(y))}{%
    \FormulaRefAuto{x\in A\dsep y\in A\dsep (G\circ F)(x)=(G\circ F)(y) \vdash G(F(x))=G(F(y))}{1,2,4}}
  \proofstep{1,2,3,4}{\bot}{\rBI{3,5}}
  \proofstep{1,2,3}{(G\circ F)(x)\neq (G\circ F)(y)}{\rCI{4,6}}
\end{tabproof}

\subsubsection{Erhalt der Injektivität}

% — Erhalt von Injektivität / Surjektivität / Bijektivität —
\FormulaThmDelta[Erhalt der Injektivität]{%
x\in A \dsep y\in A \dsep (G\circ F)(x)=(G\circ F)(y) \vdash x=y
}{
  \DeltaRow{Mengen}{x \dsep y \dsep A \dsep B \dsep C}
  \DeltaRow{Injektive Funktionen}{F\colon A\inj B \dsep G\colon B\inj C}
}
\begin{tabproof}
  \proofstep{1}{(G\circ F)(x)=(G\circ F)(y)}{\rA}
  \proofstep{2}{x\in A}{\rA}
  \proofstep{3}{y\in A}{\rA}
  \proofstep{2}{F(x)\in B}{\FormulaRefAuto{F\colon A\to B\dsep x\in A\vdash F(x)\in B}{2}}
  \proofstep{3}{F(y)\in B}{\FormulaRefAuto{F\colon A\to B\dsep x\in A\vdash F(x)\in B}{3}}
  \proofstep{1}{G(F(x))=G(F(y))}{\rIE{\FormulaRefAuto{x\in A \vdash (G\circ F)(x)=G(F(x))},1}}
  \proofstep{1,2,3}{F(x)=F(y)}{\FormulaRefAuto{Injektivität}{4,5,6}}
  \proofstep{1,2,3}{x=y}{\FormulaRefAuto{Injektivität}{2,3,7}}
\end{tabproof}


\FormulaThmDeltaR[Die Komposition als injektive Funktion]{%
G\circ F\colon A\inj C
}{F\colon A\inj B \dsep G\colon B\inj C\vdash G\circ F\colon A\inj C}{
  \DeltaRow{Mengen}{A\dsep B\dsep C}
  \DeltaPrem{Injektive Funktionen}{F\colon A\inj B \dsep G\colon B\inj C}
}
\begin{tabproof}
    \proofstep{}{G\circ F\colon A\to C}{\FormulaRefAuto{G\circ F\colon A\to C}}
    \proofstep{}{\forall x,y\in A\,((G\circ F)(x)=(G\circ F)(y) \rightarrow x=y)}{\FormulaRefAuto{x\in A \dsep y\in A \dsep (G\circ F)(x)=(G\circ F)(y) \vdash x=y}}
    \proofstep{}{G\circ F\colon A\inj C}{\FormulaRefAuto{Injektive Funktion}{1,2}}
\end{tabproof}

\FormulaThmDelta{%
  B\subseteq C \dsep F\colon A\inj B \vdash \iota_{B,C}\circ F\colon A\inj C%
}{%
  \DeltaDecl{Mengen}{A\dsep B\dsep C}%
  \DeltaDecl{Funktionen}{F}%
}
\begin{tabproof}
  \proofstep{1}{B\subseteq C}{\rA}
  \proofstep{2}{F\colon A\inj B}{\rA}
  \proofstep{1}{\iota_{B,C}\colon B\inj C}{%
    \FormulaRefAuto{A\subseteq B\vdash \iota_{A,B}\colon A\inj B}{1}}
  \proofstep{1,2}{\iota_{B,C}\circ F\colon A\inj C}{%
    \FormulaRefAuto{F\colon A\inj B \dsep G\colon B\inj C\vdash G\circ F\colon A\inj C}{2,3}}
\end{tabproof}

\subsubsection{Erhalt der Surjektivität}

\FormulaThmDelta[Erhalt der Surjektivität]{%
z\in C \vdash \exists x\in A\, (G\circ F)(x)=z
}{
  \DeltaRow{Mengen}{x \dsep y \dsep z \dsep A \dsep B \dsep C}
  \DeltaPrem{Surjektive Funktionen}{F\colon A\sur B \dsep G\colon B\sur C}
}
\begin{tabproof}
  \proofstep{1}{z\in C}{\rA}
  \proofstep{1}{\exists y\in B\, G(y)=z}{\FormulaRefAuto{Surjektivität}{1}} % (Schema: Surj. von G)
  \proofstep{3}{y\in B \land G(y)=z}{\rA}
  \proofstep{3}{y\in B}{\rAEa{3}}
  \proofstep{3}{G(y)=z}{\rAEb{3}}
  \proofstep{3}{\exists x\in A\, F(x)=y}{\FormulaRefAuto{y\in B \vdash \exists x\in A\, F(x)=y}{4}}
  \proofstep{7}{x\in A \land F(x)=y}{\rA}
  \proofstep{7}{x\in A}{\rAEa{7}}
  \proofstep{7}{F(x)=y}{\rAEb{7}}
  \proofstep{3,7}{G(F(x))=z}{\rIE{9,5}}
  \proofstep{3,7}{(G\circ F)(x)=G(F(x))}{\FormulaRefAuto{x\in A \vdash (G\circ F)(x)=G(F(x))}{8}}
  \proofstep{3,7}{(G\circ F)(x)=z}{\rIE{11,10}}
  \proofstep{3,7}{\exists x\in A\,(G\circ F)(x)=z}{\rEI{\rAI{8,12}}}
  \proofstep{3}{\exists x\in A\,(G\circ F)(x)=z}{\rEE{6,7,13}}
  \proofstep{1}{\exists x\in A\,(G\circ F)(x)=z}{\rEE{2,3,14}}
\end{tabproof}

\FormulaThmDelta[Die Komposition als surjektive Funktion]{%
G\circ F\colon A\sur C
}{
  \DeltaRow{Mengen}{A\dsep B\dsep C}
  \DeltaPrem{Surjektive Funktionen}{F\colon A\sur B \dsep G\colon B\sur C}
}
\begin{tabproof}
    \proofstep{}{G\circ F\colon A\to C}{\FormulaRefAuto{G\circ F\colon A\to C}}
    \proofstep{}{\forall z\in C\exists x\in A\, (G\circ F)(x)=z}{\FormulaRefAuto{z\in C \vdash \exists x\in A\, (G\circ F)(x)=z}}
    \proofstep{}{G\circ F\colon A\sur C}{\FormulaRefAuto{Surjektive Funktion}{1,2}}
\end{tabproof}

\subsubsection{Erhalt der Bijektivität}

\FormulaThmDelta[Die Komposition als bijektive Funktion]{%
F\colon A\bij B \dsep G\colon B\bij C\vdash G\circ F\colon A\bij C
}{
  \DeltaRow{Mengen}{A\dsep B\dsep C}
  \DeltaPrem{Funktionen}{F\colon A\to B \dsep G\colon B\to C}
}
\begin{tabproof}
    \proofstep{1}{F\colon A\bij B}{\rA}
    \proofstep{2}{G\colon B\bij C}{\rA}
    \proofstep{1}{F\colon A\sur B}{\FormulaRefAuto{F\colon A\bij B\vdash F\colon A\sur B}{1}}
    \proofstep{1}{F\colon A\inj B}{\FormulaRefAuto{F\colon A\bij B\vdash F\colon A\inj B}{1}}
    \proofstep{2}{G\colon B\sur C}{\FormulaRefAuto{F\colon A\bij B\vdash F\colon A\sur B}{2}}
    \proofstep{2}{G\colon B\inj C}{\FormulaRefAuto{F\colon A\bij B\vdash F\colon A\inj B}{2}}
    \proofstep{1,2}{ G\circ F\colon A\inj C}{\FormulaRefAuto{F\colon A\inj B \dsep G\colon B\inj C\vdash G\circ F\colon A\inj C}{4,6}}
    \proofstep{1,2}{ G\circ F\colon A\sur C}{\FormulaRefAuto{G\circ F\colon A\sur C}{3,5}}
    \proofstep{1,2}{ G\circ F\colon A\bij C}{\FormulaRefAuto{F\colon A\inj B\dsep F\colon A\sur B\vdash F\colon A\bij B}{7,8}}
\end{tabproof}


\subsubsection{Komposition mit Inklusionen und Einschränkung}

\FormulaThmDelta{%
A\subseteq B\vdash F\restriction_A = F\circ \iota_{A,B}%
}{
  \DeltaRow{Mengen}{A \dsep B \dsep C \dsep x}
  \DeltaPrem{Funktionen}{F\colon B\to C}
  \DeltaPrem{Inklusionsabbildung}{\iota_{A,B}\colon A\to B}[\FormulaRefAuto{A\subseteq B\vdash \iota_{A,B}\colon A\to B}]
  \DeltaPrem{Komposition}{F\circ \iota_{A,B}\colon A\to C}[\FormulaRefAuto{G\circ F\colon A\to C}]
  \DeltaPrem{Einschränkung}{F\restriction_A\colon A\to C}[\FormulaRefAuto{C\subseteq A \vdash F\restriction_C\colon C\to B}]
}
\begin{tabproofwide}
  \proofstepwidestar[1]{A\subseteq B}{\rA}

  \proofstepwidestar[2]{x\in A}{\rA}

  \proofstepwidestar[1,2]{x=\iota_{A,B}(x)}{%
    \FormulaRefAuto{A\subseteq B \dsep x\in A \vdash x=\iota_{A,B}(x)}{1,2}}

  \proofstepwide[1,2]{F\restriction_A(x)}{=}{F(x)}{
    \FormulaRefAuto{C\subseteq A \dsep x\in C \vdash F\restriction_C(x)=F(x)}{1,2}}

  \proofstepwide[1,2]{}{=}{F(\iota_{A,B}(x))}{
    \rIE{3,4}}    

  \proofstepwide[1,2]{}{=}{(F\circ\iota_{A,B})(x)}{
    \FormulaRefAuto{x\in A \vdash F(G(x))=(F\circ G)(x)}{2}}   
    
  \proofstepwide[1,2]{F\restriction_A(x)}{=}{(F\circ\iota_{A,B})(x)}{
    \rChain{4,6}}

  \proofstepwidestar[1]{\forall x\in A\,\bigl(F\restriction_A(x)=(F\circ \iota_{A,B})(x)\bigr)}{%
    \rUI{\rRI{2,7}}}

  \proofstepwidestar[1]{F\restriction_A = F\circ \iota_{A,B}}{%
    \FormulaRefAuto{\forall x\in A (F(x)=G(x)) \eqvdash F=G}{8}}
\end{tabproofwide}

\subsubsection{Rechenregeln}

\FormulaThmDelta{%
G = H \vdash F\circ G = F\circ H
}{
  \DeltaRow{Mengen}{A \dsep B \dsep C \dsep D}
  \DeltaPrem{Funktionen}{G,H\colon A\to B \dsep F\colon B\to C}
}
\begin{tabproof}
  \proofstep{1}{G=H}{\rA}
  \proofstep{}{F\circ G=F\circ G}{\rII}
  \proofstep{}{F\circ G=F\circ H}{\rIE{1,2}}
\end{tabproof}

\FormulaThmDelta{%
G = H \vdash G\circ F = H\circ F
}{
  \DeltaRow{Mengen}{x \dsep A \dsep B \dsep C}
  \DeltaPrem{Funktionen}{F\colon A\to B \dsep G,H\colon B\to C}
}
\begin{tabproof}
  \proofstep{1}{G=H}{\rA}
  \proofstep{}{G\circ F=G\circ F}{\rII}
  \proofstep{}{G\circ F=H\circ F}{\rIE{1,2}}
\end{tabproof}

% — Assoziativität der Komposition —
\FormulaThmDelta[Assoziativität]{%
H\circ (G\circ F) = (H\circ G)\circ F
}{
  \DeltaRow{Mengen}{A \dsep B \dsep C \dsep D}
  \DeltaPrem{Funktionen}{F\colon A\to B \dsep G\colon B\to C \dsep H\colon C\to D}
}
\begin{tabproofwide}
  \proofstepwidestar[1]{x\in A}{\rA}
  \proofstepwidestar[1]{F(x)\in B}{\FormulaRefAuto{F\colon A\to B\dsep x\in A\vdash F(x)\in B}{1}}
  \proofstepwidestar[1]{(G\circ F)(x)=G(F(x))}{\FormulaRefAuto{x\in A \vdash (G\circ F)(x)=G(F(x))}{1}}
  \proofstepwide[1]{(H\circ(G\circ F))(x)}{=}{H((G\circ F)(x))}{\FormulaRefAuto{x\in A \vdash (G\circ F)(x)=G(F(x))}{1}}
  \proofstepwide[1]{}{=}{H(G(F(x)))}{\rIE{3,4}}
  \proofstepwide[1]{}{=}{(H\circ G)(F(x))}{\FormulaRefAuto{x\in A \vdash (G\circ F)(x)=G(F(x))}{2}}
  \proofstepwide[1]{}{=}{((H\circ G)\circ F)(x)}{\FormulaRefAuto{x\in A \vdash (G\circ F)(x)=G(F(x))}{1}}
  \proofstepwide[1]{(H\circ(G\circ F))(x)}{=}{((H\circ G)\circ F)(x)}{\rChain{4,7}}
  \proofstepwide[]{\forall x\in A\,((H\circ(G\circ F))(x)}{=}{((H\circ G)\circ F)(x))}{\rUI{\rRI{1,8}}}
  \proofstepwide[]{H\circ (G\circ F)}{=}{(H\circ G)\circ F}{\FormulaRefAuto{\forall x\in A (F(x)=G(x)) \eqvdash F=G}{9}}
\end{tabproofwide}

% — Linksneutralität —
\FormulaThmDelta[Linksneutralität]{%
\Id_B \circ F = F
}{
  \DeltaRow{Mengen}{A \dsep B}
  \DeltaPrem{Funktionen}{F\colon A\to B}
}
\begin{tabproofwide}
  \proofstepwidestar[1]{x\in A}{\rA}
  \proofstepwidestar[1]{F(x)\in B}{\FormulaRefAuto{F\colon A\to B\dsep x\in A\vdash F(x)\in B}{1}}

  \proofstepwide[1]{(\Id_B\circ F)(x)}{=}{\Id_B(F(x))}{\FormulaRefAuto{x\in A \vdash (G\circ F)(x)=G(F(x))}{1}}
  \proofstepwide[1]{}{=}{F(x)}{\FormulaRefAuto{x\in A \vdash \Id_A(x)=x}{3}}
  \proofstepwide[1]{(\Id_B\circ F)(x)}{=}{F(x)}{\rChain{3,4}}
  
  \proofstepwide[]{\forall x\in A\,\bigl((\Id_B\circ F)(x)}{=}{F(x)\bigr)}{\rUI{\rRI{1,5}}}
  \proofstepwide[]{\Id_B\circ F}{=}{F}{\FormulaRefAuto{\forall x\in A\,(F(x)=G(x)) \eqvdash F=G}{6}}
\end{tabproofwide}

% — Rechtsneutralität —
\FormulaThmDelta[Rechtsneutralität]{%
F \circ \Id_A = F
}{
  \DeltaRow{Mengen}{A \dsep B}
  \DeltaPrem{Funktionen}{F\colon A\to B}
}
\begin{tabproofwide}
  \proofstepwidestar[1]{x\in A}{\rA}

  % Fixpunkteigenschaft der Identität auf A
  \proofstepwidestar[1]{\Id_A(x)=x}{\FormulaRefAuto{x\in A \vdash \Id_A(x)=x}{1}}

  % Komposition auswerten
  \proofstepwide[1]{(F\circ \Id_A)(x)}{=}{F(\Id_A(x))}{\FormulaRefAuto{x\in A \vdash (G\circ F)(x)=G(F(x))}{1}}

  % Ersetze \Id_A(x) durch x
  \proofstepwide[1]{}{=}{F(x)}{\rIE{2,3}}

  \proofstepwide[1]{(F\circ \Id_A)(x)}{=}{F(x)}{\rChain{3,4}}
  
  % Verallgemeinern und Extensionalität
  \proofstepwide[]{\forall x\in A\,\bigl((F\circ \Id_A)(x)}{=}{F(x)\bigr)}{\rUI{\rRI{1,5}}}
  \proofstepwide[]{F\circ \Id_A}{=}{F}{\FormulaRefAuto{\forall x\in A\,(F(x)=G(x)) \eqvdash F=G}{6}}
\end{tabproofwide}

\FormulaThmDelta{%
F\circ F^{-1} = \Id_B
}{
  \DeltaRow{Mengen}{A \dsep B}
  \DeltaPrem{Bijektive Funktionen}{F\colon A\bij B}
}
\begin{tabproof}
    \proofstep{1}{y\in B}{\rA}
    \proofstep{1}{F(F^{-1}(y))=y}{\FormulaRefAuto{F\colon A\bij B\dsep y\in B \vdash F\bigl(F^{-1}(y)\bigr)=y}{1}}
    \proofstep{1}{(F\circ F^{-1})(y)=F(F^{-1}(y))}{\FormulaRefAuto{x\in A \vdash (G\circ F)(x)=G(F(x))}{2}}
    \proofstep{1}{(F\circ F^{-1})(y)=y}{\rIE{2,3}}
    \proofstep{1}{\Id_B(y)=y}{\FormulaRefAuto{x\in A \vdash \Id_A(x)=x}{1}}
    \proofstep{1}{(F\circ F^{-1})(y)=\Id_B(y)}{\rIE{4,5}}
    \proofstep{1}{\forall y\in B\,(F\circ F^{-1})(y)=\Id_B(y)}{\rUI{\rRI{1,6}}}
    \proofstep{1}{F\circ F^{-1}=\Id_B}{\FormulaRefAuto{\forall x\in A (F(x)=G(x)) \eqvdash F=G}{7}}    
\end{tabproof}

\FormulaThmDelta{%
F^{-1}\circ F = \Id_A
}{
  \DeltaRow{Mengen}{x \dsep A \dsep B}
  \DeltaPrem{Bijektive Funktionen}{F\colon A\bij B}
}
\begin{tabproof}
  \proofstep{1}{x\in A}{\rA}
  \proofstep{1}{F^{-1}(F(x))=x}{\FormulaRefAuto{x\in A \vdash F^{-1}(F(x))=x}{1}}
  \proofstep{1}{(F^{-1}\circ F)(x)=F^{-1}(F(x))}{\FormulaRefAuto{x\in A \vdash (G\circ F)(x)=G(F(x))}{1}}
  \proofstep{1}{(F^{-1}\circ F)(x)=x}{\rIE{2,3}}
  \proofstep{1}{\Id_A(x)=x}{\FormulaRefAuto{x\in A \vdash \Id_A(x)=x}{1}}
  \proofstep{1}{(F^{-1}\circ F)(x)=\Id_A(x)}{\rIE{4,5}}
  \proofstep{ }{\forall x\in A\,\bigl((F^{-1}\circ F)(x)=\Id_A(x)\bigr)}{\rUI{\rRI{1,6}}}
  \proofstep{ }{F^{-1}\circ F=\Id_A}{\FormulaRefAuto{\forall x\in A (F(x)=G(x)) \eqvdash F=G}{7}}
\end{tabproof}

% ------------------------------------------------------------
% Hilfslemma: die breite Kettenrechnung auslagern
% ------------------------------------------------------------
\FormulaThmDelta{%
  x\in A \vdash (G^{-1}\circ F^{-1})\bigl((F\circ G)(x)\bigr)=x%
}{
  \DeltaRow{Mengen}{x \dsep A \dsep B \dsep C}
  \DeltaPrem{Bijektive Funktionen}{G\colon A\bij B \dsep F\colon B\bij C}
}
\begin{tabproofwide}
  \proofstepwidestar[1]{x\in A}{\rA}
  \proofstepwidestar[1]{G(x)\in B}{\FormulaRefAuto{F\colon A\to B\dsep x\in A\vdash F(x)\in B}{1}}
  \proofstepwidestar[1]{F(G(x))\in C}{\FormulaRefAuto{F\colon A\to B\dsep x\in A\vdash F(x)\in B}{2}}

  \proofstepwidestar[1]{(F\circ G)(x)=F(G(x))}{\FormulaRefAuto{x\in A \vdash (G\circ F)(x)=G(F(x))}{1}}

  \proofstepwidestar[1]{F^{-1}\bigl(F(G(x))\bigr)=G(x)}{\FormulaRefAuto{x\in A \vdash F^{-1}\bigl(F(x)\bigr)=x}{2}}
  \proofstepwidestar[1]{G^{-1}\bigl(G(x)\bigr)=x}{\FormulaRefAuto{x\in A \vdash F^{-1}\bigl(F(x)\bigr)=x}{1}}

  \proofstepwide[1]{%
    \begin{aligned}[t]
      (G^{-1}\circ F^{-1})\\
      \bigl((F\circ G)(x)\bigr)
    \end{aligned}
  }{=}{%
    \begin{aligned}[t]
      (G^{-1}\circ F^{-1})\\
      \bigl(F(G(x))\bigr)
    \end{aligned}
  }{\rIE{4,7}}

  \proofstepwide[1]{}{=}{G^{-1}\Bigl(F^{-1}\bigl(F(G(x))\bigr)\Bigr)}{\FormulaRefAuto{x\in A \vdash (G\circ F)(x)=G(F(x))}{3}}
  \proofstepwide[1]{}{=}{G^{-1}\bigl(G(x)\bigr)}{\rIE{5,8}}
  \proofstepwide[1]{}{=}{x}{\rIE{6,9}}

  \proofstepwide[1]{(G^{-1}\circ F^{-1})\bigl((F\circ G)(x)\bigr)}{=}{x}{\rChain{7,10}}
\end{tabproofwide}


% ------------------------------------------------------------
% Inversenformel für Kompositionen (kurzer Beweis: Block ersetzt)
% ------------------------------------------------------------
\FormulaThmDelta[Inverse der Komposition]{%
  G^{-1}\circ F^{-1} = (F\circ G)^{-1}%
}{
  \DeltaRow{Mengen}{x \dsep A \dsep B \dsep C}
  \DeltaPrem{Bijektive Funktionen}{G\colon A\bij B \dsep F\colon B\bij C}
}
\begin{tabproof}
  \proofstep{1}{x\in A}{\rA}

  \proofstep{1}{(G^{-1}\circ F^{-1})\bigl((F\circ G)(x)\bigr)=x}{%
    \FormulaRefAuto{x\in A \vdash (G^{-1}\circ F^{-1})\bigl((F\circ G)(x)\bigr)=x}{1}%
  }

  \proofstep{}{%
    \forall x\in A\,\Bigl((G^{-1}\circ F^{-1})\bigl((F\circ G)(x)\bigr)=x\Bigr)%
  }{%
    \rUI{\rRI{1,2}}%
  }

  \proofstep{}{%
    G^{-1}\circ F^{-1} = (F\circ G)^{-1}%
  }{%
    \FormulaRefAuto{\forall y\in B\,\bigl(F(G(y))=y\bigr) \vdash G = F^{-1}}{3}%
  }
\end{tabproof}




\subsubsection{Eigenschaften der Komponenten}

\FormulaThmDeltaR{%
x\in A\dsep y\in A\dsep F(x)=F(y)\vdash x=y
}{G\circ F\colon A\to C bijektiv\dsep x\in A\dsep y\in A\dsep F(x)=F(y)\vdash x=y}{
  \DeltaRow{Mengen}{A\dsep B\dsep C}
  \DeltaPrem{Funktionen}{F\colon A\to B\dsep G\colon B\to C}
  \DeltaPrem{Bijektive Funktionen}{G\circ F\colon A\bij C }
}
\begin{tabproof}
    \proofstep{1}{x\in A}{\rA}
    \proofstep{2}{y\in A}{\rA}
    \proofstep{3}{F(x)=F(y)}{\rA}
    \proofstep{1}{F(x)\in B}{\FormulaRefAuto{F\colon A\to B\dsep x\in A\vdash F(x)\in B}{1}}
    \proofstep{2}{F(y)\in B}{\FormulaRefAuto{F\colon A\to B\dsep x\in A\vdash F(x)\in B}{2}}
    \proofstep{1,2,3}{G(F(x))=G(F(y))}{\FormulaRefAuto{x\in A\dsep y\in A\dsep x=y\vdash F(x)=F(y)}{4,5,3}}
    \proofstep{1,2,3}{(G\circ F)(x)=(G\circ F)(y)}{\FormulaRefAuto{x\in A\dsep y\in A\dsep G(F(x))=G(F(y))\vdash (G\circ F)(x)=(G\circ F)(y)}{6}}
    \proofstep{1,2,3}{x=y}{\FormulaRefAuto{Injektivität}{1,2,7}}
\end{tabproof}

\FormulaThmDeltaR{%
y\in C\vdash \exists x\in B\; G(x)=y
}{G\circ F\colon A\to C bijektiv\dsep y\in C\vdash \exists x\in B\; G(x)=y}{
  \DeltaRow{Mengen}{A\dsep B\dsep C}
  \DeltaPrem{Funktionen}{F\colon A\to B\dsep G\colon B\to C}
  \DeltaPrem{Bijektive Funktionen}{G\circ F\colon A\bij C }
}
\begin{tabproof}
    \proofstep{1}{y\in C}{\rA}
    \proofstep{1}{\exists x\in A\, (G\circ F)(x)=y}{\FormulaRefAuto{y\in B\vdash \exists x\in A\; F(x)=y}{1}}
    \proofstep{3}{x\in A\land (G\circ F)(x)=y}{\rA}
    \proofstep{3}{x\in A}{\rAEa{3}}
    \proofstep{3}{(G\circ F)(x)=y}{\rAEb{3}}
    \proofstep{3}{(G\circ F)(x)=G(F(x))}{\FormulaRefAuto{x\in A \vdash (G\circ F)(x)=G(F(x))}{4}}
    \proofstep{3}{G(F(x))=y}{\FormulaRefAuto{a=b,\, a=c \vdash b=c}{6,5}}
    \proofstep{3}{F(x)\in B}{\FormulaRefAuto{F\colon A\to B\dsep x\in A\vdash F(x)\in B}{4}}
    \proofstep{3}{F(x)\in B\land G(F(x))=y}{\rAI{8,7}}
    \proofstep{3}{\exists x\in B\, G(x)=y}{\rEI{9}}
    \proofstep{1}{\exists x\in B\, G(x)=y}{\rEE{2,3,10}}
\end{tabproof}

\FormulaThmDeltaR{%
x\in B\dsep y\in B\dsep G(x)=G(y)\vdash x=y
}{G\circ F\colon A\to C bijektiv\dsep F\colon A\to B surjektiv\dsep x\in B\dsep y\in B\dsep G(x)=G(y)\vdash x=y}{
  \DeltaRow{Mengen}{A\dsep B\dsep C}
  \DeltaPrem{Funktionen}{G\colon B\to C}
  \DeltaPrem{Surjektive Funktionen}{F\colon A\sur B}
  \DeltaPrem{Bijektive Funktionen}{G\circ F\colon A\bij C }
}
\begin{tabproof}
    \proofstep{1}{x\in B}{\rA}
    \proofstep{2}{y\in B}{\rA}
    \proofstep{3}{G(x)=G(y)}{\rA}
    \proofstep{1}{\exists u\in A\,F(u)=x}{\FormulaRefAuto{y\in B\vdash \exists x\in A\; F(x)=y}{1}}
    \proofstep{2}{\exists v\in A\,F(v)=y}{\FormulaRefAuto{y\in B\vdash \exists x\in A\; F(x)=y}{2}}
    \proofstep{6}{u\in A\land F(u)=x}{\rA}
    \proofstep{6}{u\in A}{\rAEa{6}}
    \proofstep{6}{F(u)=x}{\rAEb{6}}
    \proofstep{9}{v\in A\land F(v)=y}{\rA}
    \proofstep{9}{v\in A}{\rAEa{9}}
    \proofstep{9}{F(v)=y}{\rAEb{9}}
    \proofstep{3,6,9}{G(F(u))=G(F(v))}{\FormulaRefAuto{a=b\dsep c=d \dsep P(b,d)\vdash P(a,c)}{8,11,3}}
    \proofstep{3,6,9}{(G\circ F)(u)=(G\circ F)(v)}{\FormulaRefAuto{x\in A\dsep y\in A\dsep G(F(x))=G(F(y))\vdash (G\circ F)(x)=(G\circ F)(y)}{7,10,12}}
    \proofstep{3,6,9}{u=v}{\FormulaRefAuto{x\in A \dsep y\in A \dsep (G\circ F)(x)=(G\circ F)(y) \vdash x=y}{7,10,13}}
    \proofstep{3,6,9}{F(u)=F(v)}{\FormulaRefAuto{x\in A\dsep y\in A\dsep x=y\vdash F(x)=F(y)}{7,10,14}}
    \proofstep{3,6,9}{x=y}{\rIE{11,\rIE{8,15}}}
    \proofstep{1,3,9}{x=y}{\rEE{4,6,16}}
    \proofstep{1,2,3}{x=y}{\rEE{5,9,17}}
\end{tabproof}

\FormulaThmDelta{%
 G\circ F\colon A\bij C\vdash F\colon A\inj B
}{
  \DeltaRow{Mengen}{A \dsep B \dsep C}
  \DeltaPrem{Funktionen}{F\colon A\to B \dsep G\colon B\to C}
}
\begin{tabproof}
    \proofstep{1}{ G\circ F\colon A\bij C}{\rA}
    \proofstep{1}{ \forall x,y\in A\, (F(x)=F(y)\rightarrow x=y)}{\FormulaRefAuto{G\circ F\colon A\to C bijektiv\dsep x\in A\dsep y\in A\dsep F(x)=F(y)\vdash x=y}{1}}
    \proofstep{1}{F\colon A\inj B}{\FormulaRefAuto{Injektive Funktion}{2}}
\end{tabproof}

\FormulaThmDelta{%
 G\circ F\colon A\bij C\vdash G\colon B\sur C
}{
  \DeltaRow{Mengen}{A \dsep B \dsep C}
  \DeltaPrem{Funktionen}{F\colon A\to B \dsep G\colon B\to C}
}
\begin{tabproof}
    \proofstep{1}{ G\circ F\colon A\bij C}{\rA}
    \proofstep{1}{\forall y\in C\exists x\in B\; G(x)=y}{\FormulaRefAuto{G\circ F\colon A\to C bijektiv\dsep y\in C\vdash \exists x\in B\; G(x)=y}{1}}
    \proofstep{1}{G\colon B\sur C}{\FormulaRefAuto{Surjektive Funktion}{2}}
\end{tabproof}


\FormulaThmDelta{%
 F\colon A\sur B\dsep G\circ F\colon A\bij C\vdash G\colon B\bij C
}{
  \DeltaRow{Mengen}{A \dsep B \dsep C}
  \DeltaPrem{Funktionen}{F\colon A\to B \dsep G\colon B\to C}
}
\begin{tabproof}
    \proofstep{1}{ F\colon A\sur B}{\rA}
    \proofstep{2}{ G\circ F\colon A\bij C}{\rA}
    \proofstep{2}{G\colon B\sur C}{\FormulaRefAuto{ G\circ F\colon A\bij C\vdash G\colon B\sur C}{2}}
    \proofstep{1,2}{\forall x,y\in B\, (G(x)=G(y)\rightarrow x=y)}{\FormulaRefAuto{G\circ F\colon A\to C bijektiv\dsep F\colon A\to B surjektiv\dsep x\in B\dsep y\in B\dsep G(x)=G(y)\vdash x=y}{1,2}}
    \proofstep{1,2}{G\colon B\inj C}{\FormulaRefAuto{Injektive Funktion}{4}}
    \proofstep{1,2}{G\colon B\bij C}{\FormulaRefAuto{F\colon A\inj B\dsep F\colon A\sur B\vdash F\colon A\bij B}{5,3}}
\end{tabproof}

\FormulaThmDelta{%
F\circ G = \Id_A\dsep G\circ F =\Id_B \vdash F\colon B\bij A
}{
  \DeltaRow{Mengen}{A}
  \DeltaPrem{Funktionen}{G\colon A\to B \dsep F\colon B\to A}
}
\begin{tabproof}
    \proofstep{1}{ F\circ G = \Id_A}{\rA}
    \proofstep{2}{ G\circ F =\Id_B}{\rA}
    \proofstep{}{\Id_A\colon A\bij A}{\FormulaRefAuto{\Id_A\colon A\bij A}}
    \proofstep{}{\Id_B\colon B\bij B}{\FormulaRefAuto{\Id_A\colon A\bij A}}
    \proofstep{1}{F\circ G\colon A\bij A}{\rIE{1,3}}
    \proofstep{2}{G\circ F\colon B\bij B}{\rIE{2,4}}
    \proofstep{1}{F\colon B\sur A}{\FormulaRefAuto{G\circ F\colon A\bij C\vdash G\colon B\sur C}{5}}
    \proofstep{2}{F\colon B\inj A}{\FormulaRefAuto{G\circ F\colon A\bij C\vdash F\colon A\inj B}{6}}
    \proofstep{1,2}{F\colon B\bij A}{\FormulaRefAuto{F\colon A\inj B\dsep F\colon A\sur B\vdash F\colon A\bij B}{8,7}}
\end{tabproof}

\FormulaThmDelta{%
F\circ G = \Id_A\dsep G\circ F =\Id_B \vdash G\colon A\bij B
}{
  \DeltaRow{Mengen}{A}
  \DeltaPrem{Funktionen}{G\colon A\to B \dsep F\colon B\to A}
}
\begin{tabproof}
    \proofstep{1}{ F\circ G = \Id_A}{\rA}
    \proofstep{2}{ G\circ F =\Id_B}{\rA}
    \proofstep{}{\Id_A\colon A\bij A}{\FormulaRefAuto{\Id_A\colon A\bij A}}
    \proofstep{}{\Id_B\colon B\bij B}{\FormulaRefAuto{\Id_A\colon A\bij A}}
    \proofstep{1}{F\circ G\colon A\bij A}{\rIE{1,3}}
    \proofstep{2}{G\circ F\colon B\bij B}{\rIE{2,4}}
    \proofstep{1}{G\colon A\sur B}{\FormulaRefAuto{G\circ F\colon A\bij C\vdash G\colon B\sur C}{6}}
    \proofstep{2}{G\colon A\inj B}{\FormulaRefAuto{G\circ F\colon A\bij C\vdash F\colon A\inj B}{5}}
    \proofstep{1,2}{F\colon B\bij A}{\FormulaRefAuto{F\colon A\inj B\dsep F\colon A\sur B\vdash F\colon A\bij B}{8,7}}
\end{tabproof}

\FormulaThmDelta{%
F\circ G = \Id_A\dsep G\circ F =\Id_B \vdash G=F^{-1}
}{
  \DeltaRow{Mengen}{A}
  \DeltaPrem{Funktionen}{G\colon A\to B \dsep F\colon B\to A}
}
\begin{tabproofwide}
    \proofstepwidestar[1]{ F\circ G = \Id_A}{\rA}
    \proofstepwidestar[2]{ G\circ F =\Id_B}{\rA}
    \proofstepwidestar[1,2]{F\colon B\bij A}{\FormulaRefAuto{F\circ G = \Id_A\dsep G\circ F =\Id_B \vdash F\colon B\bij A}{1,2}}
    \proofstepwide[]{G}{=}{G\circ \Id_A}{\FormulaRefAuto{F \circ \Id_A = F}}
    \proofstepwide[1,2]{}{=}{G\circ (F\circ F^{-1})}{\rIE{\FormulaRefAuto{F\circ F^{-1} = \Id_B}{3},4}}
    \proofstepwide[1,2]{}{=}{(G\circ F)\circ F^{-1}}{\FormulaRefAuto{H\circ (G\circ F) = (H\circ G)\circ F}{5}}
    \proofstepwide[1,2]{}{=}{\Id_B\circ F^{-1}}{\rIE{2,6}}
    \proofstepwide[1,2]{}{=}{F^{-1}}{\FormulaRefAuto{\Id_B\circ F = F}}
    \proofstepwide[1,2]{G}{=}{F^{-1}}{\rChain{4,8}}
\end{tabproofwide}

% ------------------------------------------------------------
% Hilfssatz: Auswertung der Komposition (rückwärts)
% ------------------------------------------------------------


% ------------------------------------------------------------
% Hauptsatz: Rechtsinverse ist injektiv
% ------------------------------------------------------------
\FormulaThmDelta[Rechtsinverse ist injektiv]{%
  F\circ G=\Id_B\vdash G\colon B\inj A%
}{
  \DeltaRow{Mengen}{A\dsep B\dsep x\dsep y}
  \DeltaPrem{Funktionen}{F\colon A\to B \dsep G\colon B\to A}
}
\begin{tabproofwide}
  \proofstepwidestar[1]{F\circ G=\Id_B}{\rA}

  \proofstepwidestar[2]{x\in B}{\rA}
  \proofstepwidestar[3]{y\in B}{\rA}
  \proofstepwidestar[4]{G(x)=G(y)}{\rA}

  \proofstepwidestar[1]{\Id_B=F\circ G}{%
    \FormulaRefAuto{X=Y\vdash Y=X}}

  \proofstepwide[2]{x}{=}{\Id_B(x)}{%
    \FormulaRefAuto{x\in A\vdash x=\Id_A(x)}}

  \proofstepwide[1,2]{}{=}{(F\circ G)(x)}{%
    \rIE{5,2}}

  \proofstepwide[2]{}{=}{F(G(x))}{%
    \FormulaRefAuto{x\in B \vdash (F\circ G)(x)=F(G(x))}}

  \proofstepwide[2,3,4]{}{=}{F(G(y))}{%
    \rIE{4,8}}

  \proofstepwide[3]{}{=}{(F\circ G)(y)}{%
    \FormulaRefAuto{x\in B \vdash F(G(x))=(F\circ G)(x)}}

  \proofstepwide[1,3]{}{=}{\Id_B(y)}{%
    \rIE{1,3}}

  \proofstepwide[3]{}{=}{y}{%
    \FormulaRefAuto{x\in B\vdash \Id_B(x)=x}}

  \proofstepwide[1,2,3,4]{x}{=}{y}{%
    \rChain{6,12}}

  \proofstepwidestar[1]{G\colon B\inj A}{%
    \FormulaRefAuto{Injektive Funktion}{13}}
\end{tabproofwide}

\FormulaThmDelta{%
  G\colon A\sur B\dsep F\colon B\to A\dsep
  \forall y\in B\,\bigl(G(F(y))=y\bigr)\vdash
  F\colon B\inj A
}{%
  \DeltaDecl{Mengen}{A\dsep B\dsep y}%
  \DeltaDecl{Funktionen}{F\dsep G}%
}
\begin{tabproof}
  \proofstep{1}{G\colon A\sur B}{\rA}
  \proofstep{2}{F\colon B\to A}{\rA}
  \proofstep{3}{\forall y\in B\,\bigl(G(F(y))=y\bigr)}{\rA}

  \proofstep{4}{y\in B}{\rA}

  \proofstep{2,4}{(G\circ F)(y)=G(F(y))}{%
    \FormulaRefAuto{x\in A \vdash (G\circ F)(x)=G(F(x))}{4}}

  \proofstep{3,4}{G(F(y))=y}{\rRE{4,\rUE{3}}}

  \proofstep{2,3,4}{(G\circ F)(y)=y}{%
    \FormulaRefAuto{a = b,\, b = c \vdash a = c}{5,6}}

  \proofstep{4}{\Id_B(y)=y}{%
    \FormulaRefAuto{x\in A\vdash \Id_A(x)=x}{4}}

  \proofstep{4}{y=\Id_B(y)}{%
    \FormulaRefAuto{a=b\vdash b=a}{8}}

  \proofstep{2,3,4}{(G\circ F)(y)=\Id_B(y)}{%
    \FormulaRefAuto{a = b,\, b = c \vdash a = c}{7,9}}

  \proofstep{2,3}{\forall y\in B\,\bigl((G\circ F)(y)=\Id_B(y)\bigr)}{%
    \rUI{\rRI{4,10}}}

  \proofstep{1,2,3}{G\circ F=\Id_B}{%
    \FormulaRefAuto{\forall x\in A\,(F(x)=G(x)) \eqvdash F=G}{11}}

  \proofstep{1,2,3}{F\colon B\inj A}{%
    \FormulaRefAuto{F\circ G=\Id_B\vdash G\colon B\inj A}{1,2,12}}
\end{tabproof}

\subsection{Die induzierte Potenzmengenabbildung}

\FormulaDefDelta[Funktionsvorschrift der induzierten Potenzmengenabbildung]{%
  F\colon A\to B \dsep X\in\powerset(A)\vdash \tPowMap{F}(X)\coloneqq F[X]
}{%
  \DeltaRow{Mengen}{A\dsep B\dsep X}
  \DeltaRow{Funktionen}{F}
}

\FormulaThmDelta[Typisierung der Funktionsvorschrift der induzierten Potenzmengenabbildung]{%
  F\colon A\to B \dsep X\in\powerset(A)\vdash \tPowMap{F}(X)\in\powerset(B)
}{%
  \DeltaRow{Mengen}{A\dsep B\dsep X}
  \DeltaRow{Funktionen}{F}
}
\begin{tabproof}
  \proofstep{1}{F\colon A\to B}{\rA}
  \proofstep{2}{X\in\powerset(A)}{\rA}
  \proofstep{2}{X\subseteq A}{%
    \FormulaRefAuto{x \in \powerset(A)\eqvdash x \subseteq A}{2}}
  \proofstep{1,2}{F[X]\subseteq B}{%
    \FormulaRefAuto{C\subseteq A \dsep F\colon A\to B \vdash F[C]\subseteq B}{3,1}}
  \proofstep{1,2}{\tPowMap{F}(X)\coloneqq F[X]}{%
    \FormulaRefAuto{F\colon A\to B \dsep X\in\powerset(A)\vdash \tPowMap{F}(X)\coloneqq F[X]}{1,2}}
  \proofstep{1,2}{\tPowMap{F}(X)\subseteq B}{\rIE{5,4}}
  \proofstep{1,2}{\tPowMap{F}(X)\in\powerset(B)}{%
    \FormulaRefAuto{x \in \powerset(A)\eqvdash x \subseteq A}{6}}
\end{tabproof}

\FormulaThmDelta[Existenz der induzierten Potenzmengenabbildung]{%
  F\colon A\to B \vdash
  \exists ! G\colon \powerset(A)\to \powerset(B)\,
  \forall X\in\powerset(A)\,G(X)=\tPowMap{F}(X)
}{%
  \DeltaRow{Mengen}{A\dsep B\dsep X}
  \DeltaRow{Funktionen}{F\dsep G}
}
\begin{tabproof}
  \proofstep{1}{F\colon A\to B}{\rA}
  \proofstep{1}{%
    \begin{aligned}[t]
      &\exists ! G\colon \powerset(A)\to \powerset(B)\,{}\\
      &\forall X\in\powerset(A)\,G(X)=\tPowMap{F}(X)
    \end{aligned}
  }{%
    \FormulaRefAuto{x\in A\vdash t(x)\coloneq y\dsep x\in A\vdash t(x)\in B\exists! F\colon A\to B\,\forall x\in A\,F(x) = t(x)}{%
      \FormulaRefAuto{F\colon A\to B \dsep X\in\powerset(A)\vdash \tPowMap{F}(X)\coloneqq F[X]}{1},
      \FormulaRefAuto{F\colon A\to B \dsep X\in\powerset(A)\vdash \tPowMap{F}(X)\in\powerset(B)}{1}%
    }%
  }
\end{tabproof}

\FormulaDefDelta[Induzierte Potenzmengenabbildung]{%
  F\colon A\to B \vdash
  \PowMap{F}\coloneqq \iota G\Bigl(
    G\colon \powerset(A)\to \powerset(B)
    \land
    \forall X\in\powerset(A)\,G(X)=\tPowMap{F}(X)
  \Bigr)
}{%
  \DeltaRow{Mengen}{A\dsep B\dsep X}
  \DeltaRow{Funktionen}{F\dsep G\dsep \PowMap{F}}
}

\FormulaThmDelta[Typisierung der induzierten Potenzmengenabbildung]{%
  F\colon A\to B \vdash \PowMap{F}\colon \powerset(A)\to \powerset(B)
}{%
  \DeltaRow{Mengen}{A\dsep B\dsep X}
  \DeltaRow{Funktionen}{F\dsep \PowMap{F}}
}
\begin{tabproof}
  \proofstep{1}{F\colon A\to B}{\rA}
  \proofstep{1}{\PowMap{F}\colon \powerset(A)\to \powerset(B)}{%
    \rAEa{\FormulaRefAuto{%
      F\colon A\to B \vdash
      \PowMap{F}\coloneqq \iota G\Bigl(
        G\colon \powerset(A)\to \powerset(B)
        \land
        \forall X\in\powerset(A)\,G(X)=\tPowMap{F}(X)
      \Bigr)
    }{1}}}
\end{tabproof}

\FormulaThmDelta[Explizite Werte der induzierten Potenzmengenabbildung]{%
  F\colon A\to B \dsep X\in\powerset(A)\vdash \PowMap{F}(X)=F[X]
}{%
  \DeltaRow{Mengen}{A\dsep B\dsep X}
  \DeltaRow{Funktionen}{F\dsep \PowMap{F}}
}
\begin{tabproof}
  \proofstep{1}{F\colon A\to B}{\rA}
  \proofstep{2}{X\in\powerset(A)}{\rA}
  \proofstep{1}{\forall X\in\powerset(A)\,\PowMap{F}(X)=\tPowMap{F}(X)}{%
    \rAEb{\FormulaRefAuto{%
      F\colon A\to B \vdash
      \PowMap{F}\coloneqq \iota G\Bigl(
        G\colon \powerset(A)\to \powerset(B)
        \land
        \forall X\in\powerset(A)\,G(X)=\tPowMap{F}(X)
      \Bigr)
    }{1}}}
  \proofstep{1,2}{\PowMap{F}(X)=\tPowMap{F}(X)}{\rUE{3}}
  \proofstep{1,2}{\tPowMap{F}(X)\coloneqq F[X]}{%
    \FormulaRefAuto{F\colon A\to B \dsep X\in\powerset(A)\vdash \tPowMap{F}(X)\coloneqq F[X]}{1,2}}
  \proofstep{1,2}{\PowMap{F}(X)=F[X]}{\rIE{4,5}}
\end{tabproof}

\FormulaThmDelta[Injektivität der induzierten Potenzmengenabbildung]{%
  F\colon A\bij B \dsep X\in\powerset(A)\dsep Y\in\powerset(A)\dsep
  \PowMap{F}(X)=\PowMap{F}(Y) \vdash X=Y
}{%
  \DeltaRow{Mengen}{A\dsep B\dsep X\dsep Y}
  \DeltaRow{Funktionen}{F\dsep \PowMap{F}}
}
\begin{tabproofsplitwide}
  \proofpartwideindR[Hilfsschritt \(F[X]=F[Y]\)]{%
    F\colon A\bij B \dsep X\in\powerset(A)\dsep Y\in\powerset(A)\dsep
    \PowMap{F}(X)=\PowMap{F}(Y) \vdash F[X]=F[Y]
  }{%
    F\colon A\bij B \dsep X\in\powerset(A)\dsep Y\in\powerset(A)\dsep
    \PowMap{F}(X)=\PowMap{F}(Y) \vdash F[X]=F[Y]
  }
    \proofstepwidestar[1]{F\colon A\bij B}{\rA}
    \proofstepwidestar[2]{X\in\powerset(A)}{\rA}
    \proofstepwidestar[3]{Y\in\powerset(A)}{\rA}
    \proofstepwidestar[4]{\PowMap{F}(X)=\PowMap{F}(Y)}{\rA}
    \proofstepwidestar[1]{F\colon A\to B}{\FormulaRefAuto{Bijektive Funktion}{1}}
    \proofstepwidestar[1,2]{\PowMap{F}(X)=F[X]}{%
      \FormulaRefAuto{F\colon A\to B \dsep X\in\powerset(A)\vdash \PowMap{F}(X)=F[X]}{5,2}}
    \proofstepwidestar[1,3]{\PowMap{F}(Y)=F[Y]}{%
      \FormulaRefAuto{F\colon A\to B \dsep X\in\powerset(A)\vdash \PowMap{F}(X)=F[X]}{5,3}}
    \proofstepwidestar[1,2]{F[X]=\PowMap{F}(X)}{\FormulaRefAuto{a=b\vdash b=a}{6}}
    \proofstepwidestar[1,3,4]{\PowMap{F}(X)=F[Y]}{\rIE{4,7}}
    \proofstepwidestar[1,2,3,4]{F[X]=F[Y]}{\rIE{8,9}}
  \closeproofpartwideindR

  \proofpartwide{Schluss}
    \proofstepwidestar[1]{F\colon A\bij B}{\rA}
    \proofstepwidestar[2]{X\in\powerset(A)}{\rA}
    \proofstepwidestar[3]{Y\in\powerset(A)}{\rA}
    \proofstepwidestar[4]{\PowMap{F}(X)=\PowMap{F}(Y)}{\rA}
    \proofstepwidestar[1]{F\colon A\inj B}{\FormulaRefAuto{F\colon A\bij B \vdash F\colon A\inj B}{1}}
    \proofstepwidestar[1,2,3,4]{F[X]=F[Y]}{%
      \FormulaRefAuto{F\colon A\bij B \dsep X\in\powerset(A)\dsep Y\in\powerset(A)\dsep \PowMap{F}(X)=\PowMap{F}(Y) \vdash F[X]=F[Y]}{1,2,3,4}}
    \proofstepwidestar[1,2,3,4]{X=Y}{%
      \FormulaRefAuto{F\colon A\inj B \dsep X\in\powerset(A)\dsep Y\in\powerset(A)\dsep F[X]=F[Y] \vdash X=Y}{5,2,3,6}}
  \closeproofpartwide
\end{tabproofsplitwide}

\FormulaThmDelta[Surjektivität der induzierten Potenzmengenabbildung]{%
  F\colon A\bij B \dsep Y\in\powerset(B)\vdash \exists X\in\powerset(A)\,\PowMap{F}(X)=Y
}{%
  \DeltaRow{Mengen}{A\dsep B\dsep X\dsep Y}
  \DeltaRow{Funktionen}{F\dsep \PowMap{F}}
}
\begin{tabproof}
  \proofstep{1}{F\colon A\bij B}{\rA}
  \proofstep{2}{Y\in\powerset(B)}{\rA}
  \proofstep{2}{Y\subseteq B}{\FormulaRefAuto{x \in \powerset(A)\eqvdash x \subseteq A}{2}}
  \proofstep{1,2}{F^{-1}[Y]\subseteq A}{%
    \FormulaRefAuto{C\subseteq B\dsep F\colon A\bij B \vdash F^{-1}[C]\subseteq A}{3,1}}
  \proofstep{1,2}{F^{-1}[Y]\in\powerset(A)}{%
    \FormulaRefAuto{x \in \powerset(A)\eqvdash x \subseteq A}{4}}
  \proofstep{1}{F\colon A\to B}{\FormulaRefAuto{Bijektive Funktion}{1}}
  \proofstep{1,2}{F[F^{-1}[Y]]=Y}{%
    \FormulaRefAuto{Y\subseteq B \dsep F\colon A\bij B \vdash F[F^{-1}[Y]]=Y}{3,1}}
  \proofstep{1,2}{\PowMap{F}(F^{-1}[Y])=F[F^{-1}[Y]]}{%
    \FormulaRefAuto{F\colon A\to B \dsep X\in\powerset(A)\vdash \PowMap{F}(X)=F[X]}{6,5}}
  \proofstep{1,2}{\PowMap{F}(F^{-1}[Y])=Y}{%
    \FormulaRefAuto{a=b,\, b=c \vdash a=c}{8,7}}
  \proofstep{1,2}{\exists X\in\powerset(A)\,\PowMap{F}(X)=Y}{\rEI{\rAI{5,9}}}
\end{tabproof}

\FormulaThmDelta[Induzierte Potenzmengenabbildung einer Bijektion]{%
  F\colon A\bij B \vdash \PowMap{F}\colon \powerset(A)\bij \powerset(B)
}{%
  \DeltaRow{Mengen}{A\dsep B\dsep X\dsep Y}
  \DeltaRow{Funktionen}{F\dsep \PowMap{F}}
}
\begin{tabproof}
  \proofstep{1}{F\colon A\bij B}{\rA}
  \proofstep{1}{F\colon A\to B}{\FormulaRefAuto{Bijektive Funktion}{1}}
  \proofstep{1}{\PowMap{F}\colon \powerset(A)\to \powerset(B)}{%
    \FormulaRefAuto{F\colon A\to B \vdash \PowMap{F}\colon \powerset(A)\to \powerset(B)}{2}}
  \proofstep{1}{\PowMap{F}\colon \powerset(A)\inj \powerset(B)}{%
    \FormulaRefAuto{Injektive Funktion}{3,\FormulaRefAuto{F\colon A\bij B \dsep X\in\powerset(A)\dsep Y\in\powerset(A)\dsep \PowMap{F}(X)=\PowMap{F}(Y) \vdash X=Y}{1}}}
  \proofstep{1}{\PowMap{F}\colon \powerset(A)\sur \powerset(B)}{%
    \FormulaRefAuto{Surjektive Funktion}{3,\FormulaRefAuto{F\colon A\bij B \dsep Y\in\powerset(B)\vdash \exists X\in\powerset(A)\,\PowMap{F}(X)=Y}{1}}}
  \proofstep{1}{\PowMap{F}\colon \powerset(A)\bij \powerset(B)}{%
    \FormulaRefAuto{F\colon A\inj B\dsep F\colon A\sur B \vdash F\colon A\bij B}{4,5}}
\end{tabproof}

\subsection{Erweiterung einer Abbildung um einen Punkt}


\FormulaDefDelta[Funktionsvorschrift der Erweiterung]{%
  x\in A\cup\{a\}\vdash
  \tExtMap{f}{a}{b}(x)\coloneqq \IfThenElse{x=a}{b}{f(x)}%
}{%
  \DeltaDecl{Mengen}{A\dsep a\dsep M\dsep b\dsep x}%
  \DeltaDecl{Funktionen}{f}%
  \DeltaDecl{Funktionensymbol}{\tExtMap{f}{a}{b}}%
}

\FormulaThmDelta[Typisierungsaxiom für $\tExtMap{f}{a}{b}$]{%
  a\notin A \dsep f\colon A\to M \dsep b\in M \dsep x\in A\cup\{a\}
  \vdash \tExtMap{f}{a}{b}(x)\in M%
}{%
  \DeltaDecl{Mengen}{A\dsep a\dsep M\dsep b\dsep x}%
  \DeltaDecl{Funktionen}{f}%
  \DeltaDecl{Funktionensymbol}{\tExtMap{f}{a}{b}}%
}
\begin{tabproof}
  \proofstep{1}{a\notin A}{\rA}
  \proofstep{2}{f\colon A\to M}{\rA}
  \proofstep{3}{b\in M}{\rA}
  \proofstep{4}{x\in A\cup\{a\}}{\rA}
  \proofstep{4}{x\in A \lor x\in\{a\}}{%
    \FormulaRefAuto{z \in A \cup B \eqvdash z \in A \lor z \in B}{4}}

  \proofstep{6}{x\in A}{\rA}
  \proofstep{2,6}{f(x)\in M}{%
    \FormulaRefAuto{F\colon A\to B\dsep x\in A\vdash F(x)\in B}{2,6}}
  \proofstep{1,6}{x\neq a}{%
    \FormulaRefAuto{x\in A \dsep a\notin A \vdash x\neq a}{6,1}}
  \proofstep{1,2,6}{\IfThenElse{x=a}{b}{f(x)}=f(x)}{%
    \FormulaRefAuto{\neg P \vdash \IfThenElse{P}{a}{b}=b}{8}}
  \proofstep{1,2,6}{f(x)=\IfThenElse{x=a}{b}{f(x)}}{%
    \FormulaRefAuto{a=b\vdash b=a}{9}}
  \proofstep{1,2,6}{\IfThenElse{x=a}{b}{f(x)}\in M}{\rIE{10,7}}
  \proofstep{1,2,3,6}{\tExtMap{f}{a}{b}(x)\in M}{%
    \rIE{%
      \FormulaRefAuto{x\in A\cup\{a\}\vdash \tExtMap{f}{a}{b}(x)\coloneqq \IfThenElse{x=a}{b}{f(x)}}{4},%
      11}}

  \proofstep{13}{x\in\{a\}}{\rA}
  \proofstep{13}{x=a}{%
    \FormulaRefAuto{z \in \{a\} \eqvdash z = a}{13}}
  \proofstep{13}{\IfThenElse{x=a}{b}{f(x)}=b}{%
    \FormulaRefAuto{P \vdash \IfThenElse{P}{a}{b}=a}{14}}
  \proofstep{13}{b=\IfThenElse{x=a}{b}{f(x)}}{%
    \FormulaRefAuto{a=b\vdash b=a}{15}}
  \proofstep{3,13}{\IfThenElse{x=a}{b}{f(x)}\in M}{\rIE{16,3}}
  \proofstep{3,13}{\tExtMap{f}{a}{b}(x)\in M}{%
    \rIE{%
      \FormulaRefAuto{x\in A\cup\{a\}\vdash \tExtMap{f}{a}{b}(x)\coloneqq \IfThenElse{x=a}{b}{f(x)}}{4},%
      17}}

  \proofstep{1,2,3,4}{\tExtMap{f}{a}{b}(x)\in M}{\rOE{5,6,12,13,18}}
\end{tabproof}

\FormulaThmDeltaR[Erweiterung als Funktion zur Vorschrift]{%
  \begin{aligned}[t]
    &a\notin A \dsep f\colon A\to M \dsep b\in M \vdash{}\\
    &\exists! g\colon (A\cup\{a\})\to M\,
    \forall x\in A\cup\{a\}\,g(x)=\tExtMap{f}{a}{b}(x)
  \end{aligned}
}{%
  a\notin A \dsep f\colon A\to M \dsep b\in M \vdash
  \exists! g\colon (A\cup\{a\})\to M\,
  \forall x\in A\cup\{a\}\,g(x)=\tExtMap{f}{a}{b}(x)
}{%
  \DeltaDecl{Mengen}{A\dsep a\dsep M\dsep b\dsep x}%
  \DeltaDecl{Funktionen}{f\dsep g}%
  \DeltaDecl{Funktionensymbol}{\tExtMap{f}{a}{b}}%
}
\begin{tabproofwide}
  \proofstepwidestar[1]{a\notin A}{\rA}
  \proofstepwidestar[2]{f\colon A\to M}{\rA}
  \proofstepwidestar[3]{b\in M}{\rA}

  \proofstepwidestar[4]{x\in A\cup\{a\}}{\rA}
  \proofstepwidestar[1,2,3,4]{\tExtMap{f}{a}{b}(x)\in M}{%
    \FormulaRefAuto{a\notin A \dsep f\colon A\to M \dsep b\in M \dsep x\in A\cup\{a\}\vdash \tExtMap{f}{a}{b}(x)\in M}{1,2,3,4}}
  \proofstepwidestar[1,2,3]{\forall x\in A\cup\{a\}\,\tExtMap{f}{a}{b}(x)\in M}{\rUI{\rRI{4,5}}}

  \proofstepwidestar[1,2,3]{%
    \begin{aligned}[t]
      &\exists! g\colon (A\cup\{a\})\to M\,\\
      &\forall x\in A\cup\{a\}\,g(x)=\tExtMap{f}{a}{b}(x)
    \end{aligned}
  }{%
    \FormulaRefAuto{%
      x\in A\vdash t(x)\coloneq y\dsep x\in A\vdash t(x)\in B
      \exists! F\colon A\to B\,\forall x\in A\,F(x)=t(x)
    }{6}}
\end{tabproofwide}

\FormulaDefDeltaR[Erweiterungsabbildung]{%
  \begin{aligned}[t]
    &a\notin A \dsep f\colon A\to M \dsep b\in M \vdash{}\\
    &\ExtMap{f}{a}{b}\coloneqq \iota g\Bigl(
      g\colon (A\cup\{a\})\to M \land{}\\
    &\forall x\in A\cup\{a\}\,g(x)=\tExtMap{f}{a}{b}(x)
    \Bigr)
  \end{aligned}
}{%
  a\notin A \dsep f\colon A\to M \dsep b\in M \vdash
  \ExtMap{f}{a}{b}\coloneqq \iota g\Bigl(
    g\colon (A\cup\{a\})\to M \land
    \forall x\in A\cup\{a\}\,g(x)=\tExtMap{f}{a}{b}(x)
  \Bigr)
}{%
  \DeltaDecl{Mengen}{A\dsep a\dsep M\dsep b\dsep x}%
  \DeltaDecl{Funktionen}{f\dsep g}%
  \DeltaDecl{Funktionensymbol}{\tExtMap{f}{a}{b}\dsep \ExtMap{f}{a}{b}}%
}

\FormulaThmDelta[Erweiterungsabbildung ist eine Funktion]{%
  a\notin A \dsep f\colon A\to M \dsep b\in M \vdash
  \ExtMap{f}{a}{b}\colon (A\cup\{a\})\to M%
}{%
  \DeltaDecl{Mengen}{A\dsep a\dsep M\dsep b}%
  \DeltaDecl{Funktionen}{f}%
  \DeltaDecl{Funktionensymbol}{\ExtMap{f}{a}{b}}%
}
\begin{tabproof}
  \proofstep{1}{a\notin A}{\rA}
  \proofstep{2}{f\colon A\to M}{\rA}
  \proofstep{3}{b\in M}{\rA}

  \proofstep{1,2,3}{\ExtMap{f}{a}{b}\colon (A\cup\{a\})\to M}{%
    \rAEa{\FormulaRefAuto{%
      a\notin A \dsep f\colon A\to M \dsep b\in M \vdash
      \ExtMap{f}{a}{b}\coloneqq \iota g\Bigl(
        g\colon (A\cup\{a\})\to M \land
        \forall x\in A\cup\{a\}\,g(x)=\tExtMap{f}{a}{b}(x)
      \Bigr)}{1,2,3}}}
\end{tabproof}

\FormulaThmDelta[Grundgleichung der Erweiterung]{%
  a\notin A \dsep f\colon A\to M \dsep b\in M \dsep x\in A\cup\{a\}
  \vdash \ExtMap{f}{a}{b}(x)=\tExtMap{f}{a}{b}(x)%
}{%
  \DeltaDecl{Mengen}{A\dsep a\dsep M\dsep b\dsep x}%
  \DeltaDecl{Funktionen}{f}%
  \DeltaDecl{Funktionensymbol}{\tExtMap{f}{a}{b}\dsep \ExtMap{f}{a}{b}}%
}
\begin{tabproof}
  \proofstep{1}{a\notin A}{\rA}
  \proofstep{2}{f\colon A\to M}{\rA}
  \proofstep{3}{b\in M}{\rA}
  \proofstep{4}{x\in A\cup\{a\}}{\rA}

  \proofstep{1,2,3}{\forall x\in A\cup\{a\}\,\ExtMap{f}{a}{b}(x)=\tExtMap{f}{a}{b}(x)}{%
    \rAEb{\FormulaRefAuto{%
      a\notin A \dsep f\colon A\to M \dsep b\in M \vdash
      \ExtMap{f}{a}{b}\coloneqq \iota g\Bigl(
        g\colon (A\cup\{a\})\to M \land
        \forall x\in A\cup\{a\}\,g(x)=\tExtMap{f}{a}{b}(x)
      \Bigr)}{1,2,3}}}
  \proofstep{1,2,3,4}{\ExtMap{f}{a}{b}(x)=\tExtMap{f}{a}{b}(x)}{\rRE{4,\rUE{5}}}
\end{tabproof}

\FormulaThmDelta[Erweiterung trifft den neuen Punkt]{%
  a\notin A \dsep f\colon A\to M \dsep b\in M \vdash \ExtMap{f}{a}{b}(a)=b%
}{%
  \DeltaDecl{Mengen}{A\dsep a\dsep M\dsep b}%
  \DeltaDecl{Funktionen}{f}%
  \DeltaDecl{Funktionensymbol}{\tExtMap{f}{a}{b}\dsep \ExtMap{f}{a}{b}}%
}
\begin{tabproof}
  \proofstep{1}{a\notin A}{\rA}
  \proofstep{2}{f\colon A\to M}{\rA}
  \proofstep{3}{b\in M}{\rA}

  \proofstep{4}{a\in A\cup\{a\}}{%
    \FormulaRefAuto{a\in A\cup\{a\}}}
  \proofstep{1,2,3}{\ExtMap{f}{a}{b}(a)=\tExtMap{f}{a}{b}(a)}{%
    \FormulaRefAuto{a\notin A \dsep f\colon A\to M \dsep b\in M \dsep x\in A\cup\{a\}\vdash
      \ExtMap{f}{a}{b}(x)=\tExtMap{f}{a}{b}(x)}{1,2,3,4}}
  \proofstep{4}{\tExtMap{f}{a}{b}(a)=\IfThenElse{a=a}{b}{f(a)}}{%
    \FormulaRefAuto{x\in A\cup\{a\}\vdash \tExtMap{f}{a}{b}(x)\coloneqq \IfThenElse{x=a}{b}{f(x)}}{4}}
  \proofstep{}{ \IfThenElse{a=a}{b}{f(a)}=b }{%
    \FormulaRefAuto{P \vdash \IfThenElse{P}{a}{b}=a}{\rII}}
  \proofstep{1,2,3}{\tExtMap{f}{a}{b}(a)=b}{%
    \FormulaRefAuto{a=b,\,b=c\vdash a=c}{6,8}}
  \proofstep{1,2,3}{\ExtMap{f}{a}{b}(a)=b}{%
    \FormulaRefAuto{a=b,\,b=c\vdash a=c}{5,9}}
\end{tabproof}

\FormulaThmDelta[Erweiterung stimmt auf $A$ mit $f$ überein]{%
  a\notin A \dsep f\colon A\to M \dsep b\in M \vdash
  \forall x\in A\,\ExtMap{f}{a}{b}(x)=f(x)%
}{%
  \DeltaDecl{Mengen}{A\dsep a\dsep M\dsep b\dsep x}%
  \DeltaDecl{Funktionen}{f}%
  \DeltaDecl{Funktionensymbol}{\tExtMap{f}{a}{b}\dsep \ExtMap{f}{a}{b}}%
}
\begin{tabproof}
  \proofstep{1}{a\notin A}{\rA}
  \proofstep{2}{f\colon A\to M}{\rA}
  \proofstep{3}{b\in M}{\rA}

  \proofstep{4}{x\in A}{\rA}
  \proofstep{4}{x\in A\cup\{a\}}{%
    \FormulaRefAuto{z\in A\vdash z\in A\cup B}{4}}
  \proofstep{1,2,3,4}{\ExtMap{f}{a}{b}(x)=\tExtMap{f}{a}{b}(x)}{%
    \FormulaRefAuto{a\notin A \dsep f\colon A\to M \dsep b\in M \dsep x\in A\cup\{a\}\vdash
      \ExtMap{f}{a}{b}(x)=\tExtMap{f}{a}{b}(x)}{1,2,3,5}}
  \proofstep{1,4}{x\neq a}{%
    \FormulaRefAuto{x\in A \dsep a\notin A \vdash x\neq a}{4,1}}
  \proofstep{1,2,4}{\IfThenElse{x=a}{b}{f(x)}=f(x)}{%
    \FormulaRefAuto{\neg P \vdash \IfThenElse{P}{a}{b}=b}{7}}
  \proofstep{1,2,4}{\tExtMap{f}{a}{b}(x)=f(x)}{%
    \FormulaRefAuto{x\in A\cup\{a\}\vdash \tExtMap{f}{a}{b}(x)\coloneqq \IfThenElse{x=a}{b}{f(x)}}{5}}
  \proofstep{1,2,4}{\tExtMap{f}{a}{b}(x)=f(x)}{%
    \FormulaRefAuto{a=b,\,b=c\vdash a=c}{11,8}}
  \proofstep{1,2,3,4}{\ExtMap{f}{a}{b}(x)=f(x)}{%
    \FormulaRefAuto{a=b,\,b=c\vdash a=c}{6,12}}
  \proofstep{1,2,3}{\forall x\in A\,\ExtMap{f}{a}{b}(x)=f(x)}{\rUI{\rRI{4,13}}}
\end{tabproof}

\FormulaThmDelta[Auswertungsformel der Erweiterung]{%
  a\notin A \dsep f\colon A\to M \dsep b\in M \dsep x\in A\cup\{a\}
  \vdash \ExtMap{f}{a}{b}(x)=\IfThenElse{x=a}{b}{f(x)}%
}{%
  \DeltaDecl{Mengen}{A\dsep a\dsep M\dsep b\dsep x}%
  \DeltaDecl{Funktionen}{f}%
  \DeltaDecl{Funktionensymbol}{\tExtMap{f}{a}{b}\dsep \ExtMap{f}{a}{b}}%
}
\begin{tabproof}
  \proofstep{1}{a\notin A}{\rA}
  \proofstep{2}{f\colon A\to M}{\rA}
  \proofstep{3}{b\in M}{\rA}
  \proofstep{4}{x\in A\cup\{a\}}{\rA}

  \proofstep{1,2,3,4}{\ExtMap{f}{a}{b}(x)=\tExtMap{f}{a}{b}(x)}{%
    \FormulaRefAuto{a\notin A \dsep f\colon A\to M \dsep b\in M \dsep x\in A\cup\{a\}\vdash
      \ExtMap{f}{a}{b}(x)=\tExtMap{f}{a}{b}(x)}{1,2,3,4}}

  \proofstep{4}{\tExtMap{f}{a}{b}(x)=\IfThenElse{x=a}{b}{f(x)}}{%
    \FormulaRefAuto{x\in A\cup\{a\}\vdash
      \tExtMap{f}{a}{b}(x)\coloneqq \IfThenElse{x=a}{b}{f(x)}}{4}}

  \proofstep{1,2,3,4}{\ExtMap{f}{a}{b}(x)=\IfThenElse{x=a}{b}{f(x)}}{%
    \FormulaRefAuto{a=b,\,b=c\vdash a=c}{5,6}}
\end{tabproof}

\FormulaThmDelta[Punkt-Erweiterung]{%
  \begin{aligned}[t]
    &a\notin A \dsep f\colon A\to M \dsep b\in M\\
    &\vdash\ \exists g\colon (A\cup\{a\})\to M\,\Bigl(
      g(a)=b \land \forall x\in A\,g(x)=f(x)
    \Bigr)
  \end{aligned}%
}{%
  \DeltaDecl{Mengen}{A\dsep a\dsep M\dsep b\dsep x}%
  \DeltaDecl{Funktionen}{f\dsep g}%
  \DeltaDecl{Funktionensymbol}{\ExtMap{f}{a}{b}}%
}
\begin{tabproof}
  \proofstep{1}{a\notin A}{\rA}
  \proofstep{2}{f\colon A\to M}{\rA}
  \proofstep{3}{b\in M}{\rA}

  \proofstep{1,2,3}{\ExtMap{f}{a}{b}\colon (A\cup\{a\})\to M}{%
    \FormulaRefAuto{a\notin A \dsep f\colon A\to M \dsep b\in M \vdash
      \ExtMap{f}{a}{b}\colon (A\cup\{a\})\to M}{1,2,3}}

  \proofstep{1,2,3}{\ExtMap{f}{a}{b}(a)=b}{%
    \FormulaRefAuto{a\notin A \dsep f\colon A\to M \dsep b\in M \vdash
      \ExtMap{f}{a}{b}(a)=b}{1,2,3}}

  \proofstep{1,2,3}{\forall x\in A\,\ExtMap{f}{a}{b}(x)=f(x)}{%
    \FormulaRefAuto{a\notin A \dsep f\colon A\to M \dsep b\in M \vdash
      \forall x\in A\,\ExtMap{f}{a}{b}(x)=f(x)}{1,2,3}}

  \proofstep{1,2,3}{%
    \ExtMap{f}{a}{b}(a)=b \land \forall x\in A\,\ExtMap{f}{a}{b}(x)=f(x)}{%
    \rAI{5,6}}

\proofstep{1,2,3}{%
  \begin{aligned}[t]
    \exists g\colon (A\cup\{a\})\to M\,\Bigl(
      &g(a)=b \\
      &\land \forall x\in A\,g(x)=f(x)
    \Bigr)
  \end{aligned}%
}{\rEI{\rAI{4,7}}}
\end{tabproof}



\subsection{Abbildungen mit leerem Definitions- oder Zielbereich}

\FormulaThmDelta[Jede Funktion aus $\varnothing$ ist leer]{%
  f\colon \varnothing\to M \vdash f=\varnothing%
}{%
  \DeltaDecl{Mengen}{M\dsep x\dsep y\dsep u}%
  \DeltaDecl{Funktionen}{f}%
}
\begin{tabproof}
  \proofstep{1}{f\colon \varnothing\to M}{\rA}
  \proofstep{1}{f\subseteq \varnothing\times M}{\FormulaRefAuto{Funktion}{1}}
  \proofstep{3}{u\in f}{\rA}
  \proofstep{1,3}{u\in \varnothing\times M}{\FormulaRefAuto{A \subseteq B,\, x \in A \vdash x \in B}{2,3}}
  \proofstep{}{\varnothing\times M=\varnothing}{\FormulaRefAuto{\varnothing\times A=\varnothing}}
  \proofstep{1,3}{u\in\varnothing}{\rIE{5,4}}
  \proofstep{}{u\notin\varnothing}{\FormulaRefAuto{x \not\in \varnothing}}
  \proofstep{1,3}{\bot}{\rBI{6,7}}
  \proofstep{1}{u\notin f}{\rCI{2,8}}
  \proofstep{1}{\forall u\,u\notin f}{\rUI{9}}
  \proofstep{1}{f=\varnothing}{\FormulaRefAuto{\forall x\, x\notin A\vdash A=\varnothing}{10}}
\end{tabproof}

% =========================================================
% Abbildungen aus der leeren Menge
% =========================================================

\FormulaThmDelta[Leere Relation ist eine Funktion]{%
  \varnothing\colon \varnothing\to M%
}{%
  \DeltaDecl{Mengen}{M}%
}
\begin{tabproof}
  \proofstep{}{ \Id_M\colon M\to M }{%
    \FormulaRefAuto{IdA}{}} % (hier: A := M)

  \proofstep{}{ \varnothing \subseteq M }{%
    \FormulaRefAuto{\varnothing\subseteq A}{}}

  \proofstep{1,2}{ \Id_M\restriction_{\varnothing}\colon \varnothing\to M }{%
    \FormulaRefAuto{C\subseteq A \vdash F\restriction _C\colon C\to B}{}}

  \proofstep{1,2}{ \Id_M\restriction_{\varnothing}=\varnothing }{%
    \FormulaRefAuto{f\colon \varnothing \to M \vdash f=\varnothing}{3}}

  \proofstep{3,4}{ \varnothing\colon \varnothing\to M }{%
    \rIE{4,3}}
\end{tabproof}

\FormulaThmDelta[Eindeutige Existenz einer Funktion aus $\varnothing$]{%
  \exists! f\,\bigl(f\colon \varnothing\to M\bigr)%
}{%
  \DeltaDecl{Mengen}{M}%
  \DeltaDecl{Funktionen}{f\dsep g}%
}
\begin{tabproof}
  \proofstep{}{ \varnothing\colon \varnothing\to M }{%
    \FormulaRefAuto{\varnothing\colon \varnothing\to M}}

  \proofstep{}{ \exists f\,\bigl(f\colon \varnothing\to M\bigr) }{%
    \rEI{1}}

  \proofstep{3}{ f\colon \varnothing\to M }{\rA}
  \proofstep{4}{ g\colon \varnothing\to M }{\rA}

  \proofstep{3}{ f=\varnothing }{%
    \FormulaRefAuto{f\colon \varnothing\to M \vdash f=\varnothing}{3}}

  \proofstep{4}{ g=\varnothing }{%
    \FormulaRefAuto{f\colon \varnothing\to M \vdash f=\varnothing}{4}}

  \proofstep{3,4}{ f=g }{%
    \FormulaRefAuto{a = b,\, c = b \vdash a = c}{5,6}}

  \proofstep{}{ \exists! f\,\bigl(f\colon \varnothing\to M\bigr) }{%
    \rUEI{2,3,4,7}}
\end{tabproof}

\FormulaThmDelta[Leere Relation ist injektiv]{%
  \varnothing\colon \varnothing\inj M%
}{%
  \DeltaDecl{Mengen}{M\dsep x\dsep y}%
}
\begin{tabproof}
  \proofstep{}{ \varnothing\colon \varnothing\to M }{%
    \FormulaRefAuto{\varnothing\colon \varnothing\to M}}

  \proofstep{2}{x\in \varnothing \land y\in \varnothing \land \varnothing(x)=\varnothing(y)}{\rA}
  \proofstep{2}{x\in \varnothing}{\rAEa{2}}
  \proofstep{}{x\notin \varnothing}{\FormulaRefAuto{x \not\in \varnothing}}
  \proofstep{2}{\bot}{\rBI{3,4}}
  \proofstep{2}{x=y}{\FormulaRefAuto{P\dsep \neg P \vdash Q}{3,4}}
  \proofstep{}{(x\in \varnothing \land y\in \varnothing \land \varnothing(x)=\varnothing(y))\rightarrow x=y}{\rRI{2,6}}

  \proofstep{}{ \varnothing\colon \varnothing\inj M }{%
    \FormulaRefAuto{Injektive Funktion}{1,7}}
\end{tabproof}

\FormulaThmDelta[Leere Relation ist surjektiv auf $\varnothing$]{%
  \varnothing\colon \varnothing\sur \varnothing%
}{%
  \DeltaDecl{Mengen}{y\dsep x}%
}
\begin{tabproof}
  \proofstep{}{ \varnothing\colon \varnothing\to \varnothing }{%
    \FormulaRefAuto{\varnothing\colon \varnothing\to M}}

  \proofstep{2}{y\in \varnothing}{\rA}
  \proofstep{}{y\notin \varnothing}{\FormulaRefAuto{x \not\in \varnothing}}
  \proofstep{2}{\bot}{\rBI{2,3}}
  \proofstep{2}{\exists x\in \varnothing\,(\varnothing(x)=y)}{%
    \FormulaRefAuto{P\dsep \neg P \vdash Q}{2,3}}
  \proofstep{}{y\in \varnothing\rightarrow \exists x\in \varnothing\,(\varnothing(x)=y)}{\rRI{2,5}}

  \proofstep{}{ \varnothing\colon \varnothing\sur \varnothing }{%
    \FormulaRefAuto{Surjektive Funktion}{1,6}}
\end{tabproof}

\FormulaThmDelta[Leere Relation ist bijektiv auf $\varnothing$]{%
  \varnothing\colon \varnothing\bij \varnothing%
}{%
  \DeltaDecl{Mengen}{}
}
\begin{tabproof}
  \proofstep{}{ \varnothing\colon \varnothing\inj \varnothing }{%
    \FormulaRefAuto{\varnothing\colon \varnothing\inj M}}

  \proofstep{}{ \varnothing\colon \varnothing\sur \varnothing }{%
    \FormulaRefAuto{\varnothing\colon \varnothing\sur \varnothing}}

  \proofstep{}{ \varnothing\colon \varnothing\bij \varnothing }{%
    \FormulaRefAuto{F\colon A\inj B\dsep F\colon A\sur B \vdash F\colon A\bij B}{1,2}}
\end{tabproof}

\FormulaThmDelta{%
  A\neq \varnothing \vdash \neg \exists f\colon A\to \varnothing%
}{%
  \DeltaDecl{Mengen}{A\dsep x}%
  \DeltaDecl{Funktionen}{f}%
}
\begin{tabproof}
  \proofstep{1}{A\neq\varnothing}{\rA}
  \proofstep{1}{\exists x\,(x\in A)}{%
    \FormulaRefAuto{A\neq\varnothing \eqvdash \exists x\,(x \in A)}{1}}

  \proofstep{3}{\exists f\colon A\to\varnothing}{\rA}

  \proofstep{4}{x\in A}{\rA}
  \proofstep{5}{f\colon A\to\varnothing}{\rA}

  \proofstep{4,5}{(x,f(x))\in f}{%
    \FormulaRefAuto{x\in A\vdash (x,F(x))\in F}{4,5}}

  \proofstep{4,5}{f(x)\in\varnothing}{%
    \FormulaRefAuto{(x,y)\in F\vdash y\in B}{6}}

  \proofstep{}{f(x)\notin\varnothing}{%
    \FormulaRefAuto{x \not\in \varnothing}}

  \proofstep{4,5}{\bot}{\rBI{7,8}}

  \proofstep{1,5}{\bot}{\rEE{2,4,9}}
  \proofstep{1,3}{\bot}{\rEE{3,5,10}}

  \proofstep{1}{\neg \exists f\colon A\to\varnothing}{\rCI{3,11}}
\end{tabproof}

\section{Gleichmächtigkeit}

\FormulaDefDeltaK[Begriff der gleichmächtigen Mengen]{A \EqCard B \coloneqq \exists F\colon A\bij B}{Gleichmächtigkeit}{
  \DeltaRow{Mengen}{A\dsep B}
  \DeltaRow{\textbf{Neue Symbole}}{}
  \DeltaRow{Gleichmächtigkeit}{}
}

\FormulaThmDelta[Reflexivität der Gleichmächtigkeit]{%
  A \EqCard A%
}{
  \DeltaRow{Mengen}{A}
}
\begin{tabproof}
  \proofstep{}{\Id_A\colon A\bij A}{\FormulaRefAuto{\Id_A\colon A\bij A}}
  \proofstep{}{\exists F\colon A\bij A}{\rEI{1}}
  \proofstep{}{A \EqCard A}{\FormulaRefAuto{A \EqCard B \coloneqq \exists F\colon A\bij B}{2}}
\end{tabproof}

\FormulaThmDelta[Symmetrie der Gleichmächtigkeit]{%
  A \EqCard B \vdash B \EqCard A%
}{
  \DeltaRow{Mengen}{A\dsep B}
  \DeltaRow{Funktionen}{F\dsep G}
}
\begin{tabproof}
  \proofstep{1}{A \EqCard B}{\rA}
  \proofstep{1}{\exists F\colon A\bij B}{\FormulaRefAuto{A \EqCard B \coloneqq \exists F\colon A\bij B}{1}}

  \proofstep{3}{F\colon A\bij B}{\rA}
  \proofstep{3}{F^{-1}\colon B\bij A}{\FormulaRefAuto{F\colon A\bij B\vdash F^{-1}\colon B\bij A}{3}}
  \proofstep{3}{\exists G\colon B\bij A}{\rEI{4}}

  \proofstep{1}{\exists G\colon B\bij A}{\rEE{2,3,5}}
  \proofstep{1}{B \EqCard A}{\FormulaRefAuto{A \EqCard B \coloneqq \exists F\colon A\bij B}{6}}
\end{tabproof}

\FormulaThmDelta[Transitivität der Gleichmächtigkeit]{%
  A \EqCard B \dsep B \EqCard C \vdash A \EqCard C%
}{
  \DeltaRow{Mengen}{A\dsep B\dsep C}
  \DeltaRow{Funktionen}{F\dsep G\dsep H}
}
\begin{tabproof}
  \proofstep{1}{A \EqCard B}{\rA}
  \proofstep{2}{B \EqCard C}{\rA}

  \proofstep{1}{\exists F\colon A\bij B}{\FormulaRefAuto{A \EqCard B \coloneqq \exists F\colon A\bij B}{1}}
  \proofstep{2}{\exists G\colon B\bij C}{\FormulaRefAuto{A \EqCard B \coloneqq \exists F\colon A\bij B}{2}}

  \proofstep{5}{F\colon A\bij B}{\rA}
  \proofstep{6}{G\colon B\bij C}{\rA}
  \proofstep{5,6}{G\circ F\colon A\bij C}{\FormulaRefAuto{F\colon A\bij B \dsep G\colon B\bij C\vdash G\circ F\colon A\bij C}{5,6}}
  \proofstep{5,6}{\exists H\colon A\bij C}{\rEI{7}}
  \proofstep{5}{\exists H\colon A\bij C}{\rEE{4,6,8}}
  \proofstep{1,2}{\exists H\colon A\bij C}{\rEE{3,5,9}}

  \proofstep{1,2}{A \EqCard C}{\FormulaRefAuto{A \EqCard B \coloneqq \exists F\colon A\bij B}{10}}
\end{tabproof}


\FormulaThmDelta[Potenzmengen gleichmächtiger Mengen sind gleichmächtig]{%
  A \EqCard B \vdash \powerset(A)\EqCard \powerset(B)
}{%
  \DeltaRow{Mengen}{A\dsep B}
  \DeltaRow{Funktionen}{F\dsep G\dsep \widehat{F}}
}
\begin{tabproof}
  \proofstep{1}{A \EqCard B}{\rA}
  \proofstep{1}{\exists F\colon A\bij B}{%
    \FormulaRefAuto{A \EqCard B \coloneqq \exists F\colon A\bij B}{1}}

  \proofstep{3}{F\colon A\bij B}{\rA}
    \proofstep{3}{\PowMap{F}\colon \powerset(A)\bij \powerset(B)}{%
      \FormulaRefAuto{F\colon A\bij B \vdash \PowMap{F}\colon \powerset(A)\bij \powerset(B)}{3}}
  \proofstep{3}{\exists G\colon \powerset(A)\bij \powerset(B)}{\rEI{4}}

  \proofstep{1}{\exists G\colon \powerset(A)\bij \powerset(B)}{\rEE{2,3,5}}
  \proofstep{1}{\powerset(A)\EqCard \powerset(B)}{%
    \FormulaRefAuto{A \EqCard B \coloneqq \exists F\colon A\bij B}{6}}
\end{tabproof}

